\examtitle{AIG202c_Ultra_Edition}
\examconfig{total_questions=190; duration=180; audio_play_limit=1; mode=practice_and_mastery}

\section{sec1}{1. Introduction to Artificial Intelligence}

% ==============================================================================
% Question 1 | Phân loại: 1. Introduction to Artificial Intelligence
% Xuất hiện trong các đề gốc: AIG202c_SP_2025_FE (Q1), AIG202c_SU25_FE (Q28)
% ==============================================================================
\begin{question}{1}{mcq}
\source{AIG202c_-_SP_2025_-_FE_3755.webp}
\prompt{What is a significant way in which developers of AI systems can guard against introducing bias?}
\choice{A}{Using only examples from their own environment as training data.}
\choice{B}{Using less varied AI systems and datasets.}
\choice{C}{Providing effective training data and performing regular tests and audits.}
\choice{D}{Using government approved algorithms.}
\answer{C}
\confidence{high}
\evidence{Biện pháp quan trọng và hiệu quả nhất để các nhà phát triển AI ngăn ngừa thiên kiến (bias) là cung cấp tập dữ liệu huấn luyện đa dạng, chất lượng cao và thực hiện kiểm thử, đánh giá định kỳ có hệ thống (Providing effective training data and performing regular tests and audits).}
\sources{IEEE Ethically Aligned Design; NIST AI Risk Management Framework: Bias Mitigation}
\explain{Hướng giải: Xác định giải pháp kỹ thuật và quy trình chuẩn để ngăn chặn thiên kiến thuật toán.

Các bước phân tích:
1. Thiên kiến trong AI bắt nguồn chủ yếu từ sự thiếu cân bằng hoặc sai lệch trong dữ liệu huấn luyện (Training data bias).
2. Để khắc phục, nhà phát triển phải: (1) Thu thập dữ liệu đại diện cân bằng cho mọi nhóm đối tượng (effective training data), và (2) Thường xuyên tiến hành kiểm toán thuật toán độc lập và kiểm thử độ công bằng (regular tests and audits) trước và sau khi triển khai.

Kiểm tra: Đây là khuyến nghị tiêu chuẩn của NIST, IEEE và UNESCO về quản trị rủi ro AI.

Vì sao các phương án khác sai:
- A: Chỉ dùng dữ liệu từ môi trường của chính nhà phát triển sẽ gây ra thiên kiến cục bộ nghiêm trọng (selection bias).
- B: Dùng ít hệ thống và dữ liệu ít đa dạng càng làm thiên kiến trầm trọng hơn.
- D: Thuật toán được chính phủ phê duyệt không tự động miễn nhiễm với thiên kiến nếu dữ liệu đầu vào bị sai lệch.

Lưu ý/Khái niệm cần nhớ: Dữ liệu huấn luyện toàn diện kết hợp kiểm toán định kỳ là chìa khóa triệt tiêu thiên kiến.}
\vocab{Algorithmic Audit | Kiểm toán thuật toán | Quy trình đánh giá độc lập để phát hiện lỗi, thiên kiến và rủi ro trong hệ thống AI}
\end{question}

% ==============================================================================
% Question 2 | Phân loại: 1. Introduction to Artificial Intelligence
% Xuất hiện trong các đề gốc: AIG202c_SP_2025_FE (Q3), AIG202c_FA25_FE (Q23), AIG202c_SP26_RE (Q29)
% ==============================================================================
\begin{question}{2}{mcq}
\source{AIG202c_-_SP_2025_-_FE_3755_3.webp}
\prompt{What do Abu Dhabi National Oil Company (ADNOC) hope that the use of AI will help them achieve?}
\choice{A}{To understand their business better and retain knowledge.}
\choice{B}{To eliminate knowledge transfer.}
\choice{C}{To differentiate their business into educating people about geology}
\choice{D}{A production target of 17 million barrels per day.}
\answer{A}
\confidence{high}
\evidence{Công ty Dầu khí Quốc gia Abu Dhabi (ADNOC) triển khai AI nhằm mục tiêu hiểu rõ hơn về hoạt động kinh doanh khai thác và lưu giữ/kế thừa tri thức chuyên gia (To understand their business better and retain knowledge).}
\sources{ADNOC AI Panorama Digital Command Center Case Study; IBM Energy AI Solutions}
\explain{Hướng giải: Xác định mục tiêu chiến lược của ADNOC khi ứng dụng Trí tuệ nhân tạo theo nghiên cứu trường hợp bài giảng.

Các bước phân tích:
1. Ngành dầu khí có khối lượng dữ liệu địa chất và vận hành khổng lồ tích lũy qua nhiều thập kỷ, trong khi các chuyên gia giàu kinh nghiệm dần nghỉ hưu.
2. ADNOC xây dựng trung tâm kỹ thuật số Panorama và ứng dụng AI nhằm phân tích dữ liệu toàn diện để 'hiểu rõ hơn hoạt động kinh doanh' và 'lưu giữ tri thức chuyên gia' (knowledge retention) cho các thế hệ tương lai.

Kiểm tra: Đây là mục tiêu chuẩn được nêu rõ trong bài học AIG202c về trường hợp điển hình của ADNOC.

Vì sao các phương án khác sai:
- B: Loại bỏ chuyển giao tri thức là hoàn toàn sai, ADNOC muốn thúc đẩy chuyển giao.
- C: ADNOC là tập đoàn khai thác dầu khí năng lượng, không phải tổ chức giáo dục địa chất đại chúng.
- D: Mục tiêu sản lượng 17 triệu thùng/ngày là con số phi thực tế và sai lệch.

Lưu ý/Khái niệm cần nhớ: Ứng dụng AI trong các tập đoàn công nghiệp lớn chú trọng vào tối ưu hóa vận hành và lưu giữ tri thức.}
\vocab{Knowledge Retention | Lưu giữ tri thức | Việc số hóa và lưu trữ tri thức chuyên gia để không bị thất thoát theo thời gian}
\end{question}

% ==============================================================================
% Question 3 | Phân loại: 1. Introduction to Artificial Intelligence
% Xuất hiện trong các đề gốc: AIG202c_SP_2025_FE (Q4)
% ==============================================================================
\begin{question}{3}{mcq}
\source{AIG202c_-_SP_2025_-_FE_3755_4.webp}
\prompt{Which of the following is a concern regarding AI Transparency?}
\choice{A}{The inability to understand AI algorithms}
\choice{B}{The clarity in AI decision-making processes}
\choice{C}{AI systems' inability to learn from data}
\choice{D}{The ethical use of AI in robotics}
\answer{A}
\confidence{high}
\evidence{Mối quan ngại lớn nhất liên quan đến Tính minh bạch của AI (AI Transparency) là sự bất lực trong việc thấu hiểu cách thức hoạt động nội tại của các thuật toán (The inability to understand AI algorithms), thường được gọi là vấn đề Hộp đen (Black Box problem).}
\sources{DARPA XAI Program; High-Level Expert Group on AI: Ethics Guidelines for Trustworthy AI}
\explain{Hướng giải: Xác định mối quan ngại đạo đức và kỹ thuật cốt lõi liên quan đến Tính minh bạch của AI.

Các bước phân tích:
1. Các mô hình học sâu hiện đại có hàng tỷ trọng số liên kết phi tuyến tính phức tạp.
2. Con người không thể hiểu được cụ thể tại sao mô hình lại đưa ra một quyết định nhất định (ví dụ: từ chối cho vay hoặc chẩn đoán ung thư). Sự thiếu minh bạch này ('inability to understand AI algorithms') gây ra rủi ro pháp lý và làm xói mòn niềm tin của xã hội.

Kiểm tra: Vấn đề 'Black Box AI' là động lực chính thúc đẩy sự ra đời của ngành nghiên cứu XAI (Explainable AI).

Vì sao các phương án khác sai:
- B: Sự rõ ràng trong ra quyết định là mục tiêu mong muốn, không phải 'mối quan ngại' (concern).
- C: AI học từ dữ liệu rất tốt, không phải mối lo về tính minh bạch.
- D: Đạo đức trong robot là chủ đề rộng hơn, không định nghĩa riêng cho Transparency.

Lưu ý/Khái niệm cần nhớ: AI Transparency đối mặt với thách thức lớn nhất là vấn đề Hộp đen (Black box).}
\vocab{Black Box AI | AI hộp đen | Hệ thống AI đưa ra quyết định mà con người không thể nhìn thấy hoặc giải thích được cơ chế bên trong}
\end{question}

% ==============================================================================
% Question 4 | Phân loại: 1. Introduction to Artificial Intelligence
% Xuất hiện trong các đề gốc: AIG202c_SP_2025_FE (Q7), AIG202c_FA25_FE (Q18)
% ==============================================================================
\begin{question}{4}{mcq}
\source{AIG202c_-_SP_2025_-_FE_3755_7.webp}
\prompt{What is the primary viewpoint on ethical concerns with AI systems based on the text?}
\choice{A}{Something that can't be mitigated for.}
\choice{B}{Not genuinely troubling, and the concern of very few AI experts}
\choice{C}{Something that should be the concern of every AI developer, so they can be mitigated for as AI systems are developed}
\choice{D}{Short term and easily addressed when developing new AI systems.}
\answer{A}
\confidence{high}
\evidence{Thuật toán K-Means Clustering phân chia dữ liệu thành \(K\) cụm bằng cách tối thiểu hóa khoảng cách từ các điểm tới tâm cụm (centroid) tương ứng.}
\sources{MacQueen (1967), Some Methods for classification and Analysis of Multivariate Observations; Scikit-Learn K-Means}
\explain{Hướng giải: Xác định mục tiêu toán học và cơ chế phân cụm của thuật toán K-Means.

Các bước phân tích:
1. K-Means là thuật toán học không giám sát phân chia \(n\) quan sát thành \(K\) cụm.
2. Mỗi quan sát thuộc về cụm có tâm (centroid) gần nhất dựa trên khoảng cách Euclidean: \[
\min_{S} \sum_{i=1}^{K} \sum_{x \in S_i} \|x - \mu_i\|^2
\]

Kiểm tra: Hàm mục tiêu của K-Means là tối thiểu hóa tổng bình phương khoảng cách trong từng cụm (Within-Cluster Sum of Squares - WCSS).

Vì sao các phương án khác sai:
- K-Means không sử dụng siêu phẳng phân cách cực đại (đó là SVM).
- K-Means không xây dựng cây quyết định (đó là Decision Trees).

Lưu ý/Khái niệm cần nhớ: K-Means gán điểm dữ liệu vào tâm cụm gần nhất dựa trên khoảng cách hình học.}
\vocab{Centroid | Tâm cụm | Điểm trung tâm đại diện cho một cụm trong thuật toán phân cụm K-Means}
\end{question}

% ==============================================================================
% Question 5 | Phân loại: 1. Introduction to Artificial Intelligence
% Xuất hiện trong các đề gốc: AIG202c_SP_2025_FE (Q9), AIG202c_SU25_FE (Q14)
% ==============================================================================
\begin{question}{5}{mcq}
\source{AIG202c_-_SP_2025_-_FE_3755_9.webp}
\prompt{Which of the following is NOT an attribute of Machine Learning?}
\choice{A}{Takes data and answers as input and uses these inputs to create a set of rules that determine what the Machine Learning model will be}
\choice{B}{Takes data and rules as input and uses these inputs to develop an algorithm that will give us an answer}
\choice{C}{Machine Learning models can be continuously trained}
\choice{D}{Machine Learning defines behavioral rules by comparing large data sets to find common patterns}
\answer{B}
\confidence{high}
\evidence{Phương án B: 'Takes data and rules as input and uses these inputs to develop an algorithm that will give us an answer' là định nghĩa của Lập trình truyền thống (Traditional Programming), KHÔNG PHẢI là thuộc tính của Học máy (Machine Learning). Trong Học máy, đầu vào là Dữ liệu (Data) và Câu trả lời (Answers/Labels), và đầu ra là Các quy tắc / Thuật toán (Rules/Model).}
\sources{François Chollet, Deep Learning with Python (Ch. 1: What is Machine Learning?); Pedro Domingos, The Master Algorithm}
\explain{Hướng giải: Phân biệt sự khác nhau mang tính bản chất giữa Lập trình truyền thống và Học máy (Machine Learning).

Các bước phân tích:
1. Lập trình truyền thống (Traditional Programming):
   \[\text{Data} + \text{Rules (Code)} \longrightarrow \text{Answers}\]
   (Con người lập trình tường minh các quy tắc và nạp dữ liệu để máy tính tính ra đáp án).
2. Học máy (Machine Learning):
   \[\text{Data} + \text{Answers (Labels)} \longrightarrow \text{Rules (Model)}\]
   (Hệ thống nạp dữ liệu và câu trả lời mẫu để tự động học ra các quy tắc thống kê).
3. Phương án B mô tả việc nhận 'Data and rules as input' để sinh ra đáp án, đây chính xác là lập trình cổ điển, do đó B KHÔNG PHẢI là thuộc tính của Machine Learning.

Kiểm tra: Sơ đồ so sánh kinh điển của François Chollet khẳng định rõ sự đảo ngược luồng dữ liệu này.

Vì sao các phương án khác là thuộc tính đúng của ML:
- A: Đúng, nạp data và answers để sinh ra rules.
- C: Đúng, mô hình ML có thể được liên tục huấn luyện (continuous learning/retraining).
- D: Đúng, ML so sánh tập dữ liệu lớn để tìm các mẫu chung.

Lưu ý/Khái niệm cần nhớ: Machine Learning = Data + Answers -> Rules (Mô hình tự sinh ra quy tắc từ dữ liệu).}
\vocab{Machine Learning Paradigm | Hệ hình học máy | Cơ chế đảo ngược lập trình truyền thống: học quy tắc từ dữ liệu và kết quả mẫu}
\end{question}

% ==============================================================================
% Question 6 | Phân loại: 1. Introduction to Artificial Intelligence
% Xuất hiện trong các đề gốc: AIG202c_SP_2025_FE (Q12)
% ==============================================================================
\begin{question}{6}{mcq}
\source{AIG202c_-_SP_2025_-_FE_3755_12.webp}
\prompt{Which of the following NOT a way AI is being used to benefit humanity?}
\choice{A}{In healthcare, AI is being used to interpret scans for early detection of cancer, eye disease, and other problems.}
\choice{B}{Crime: to identify criminals before they commit a crime.}
\choice{C}{In healthcare, AI is being used to predict where the next outbreak of a disease will occur.}
\choice{D}{In agriculture, AI is being used to ident$' and recommend treatment for plant diseases.}
\answer{B}
\confidence{high}
\evidence{Sử dụng AI trong tội phạm để 'nhận diện tội phạm trước khi họ thực hiện hành vi phạm tội' (identify criminals before they commit a crime - Predictive Policing mang tính định kiến suy đoán tiền tội phạm) là một vi phạm đạo đức nghiêm trọng về quyền con người và quyền suy đoán vô tội, KHÔNG PHẢI là một cách AI đang được sử dụng để đem lại lợi ích cho nhân loại.}
\sources{Stanford HAI: Algorithmic Injustice in Predictive Policing; UN Human Rights Council AI Report}
\explain{Hướng giải: Nhận diện trường hợp sử dụng AI vi phạm nhân quyền và đạo đức xã hội.

Các bước phân tích:
1. Các phương án A, C, D đều là các ứng dụng nhân văn tuyệt vời của AI: tầm soát ung thư sớm trong y tế (A), dự báo dịch bệnh bùng phát (C), phát hiện bệnh cây trồng trong nông nghiệp (D).
2. Phương án B nhắc đến việc 'xác định tội phạm trước khi họ phạm tội' (dự đoán trước hành vi phạm tội theo kiểu Minority Report). Trong thực tế, các thuật toán dự đoán tội phạm (Predictive Policing) đã bị cộng đồng quốc tế lên án vì khuếch đại định kiến chủng tộc, gán nhãn sai cho người vô tội và vi phạm nguyên tắc 'suy đoán vô tội' căn bản của nền tư pháp văn minh.

Kiểm tra: Báo cáo của Liên Hợp Quốc về nhân quyền khẳng định các công cụ phán xét tội phạm trước khi hành động xảy ra là mối đe dọa lớn đối với quyền tự do con người.

Vì sao các phương án khác là lợi ích thực tế:
- A: AI chẩn đoán ảnh y tế cứu sống hàng triệu bệnh nhân.
- C: Mô hình dịch tễ AI (như BlueDot) cảnh báo sớm các đại dịch.
- D: AI nông nghiệp giúp nông dân bảo vệ mùa màng bền vững.

Lưu ý/Khái niệm cần nhớ: AI không bao giờ được dùng để định tội con người trước khi có hành vi phạm tội thực tế.}
\vocab{Predictive Policing | Dự đoán tội phạm bằng thuật toán | Việc sử dụng mô hình AI để dự báo nguy cơ tội phạm, tiềm ẩn nhiều rủi ro đạo đức và nhân quyền}
\end{question}

% ==============================================================================
% Question 7 | Phân loại: 1. Introduction to Artificial Intelligence
% Xuất hiện trong các đề gốc: AIG202c_SP_2025_FE (Q14), AIG202c_SU25_FE (Q25), AIG202c_SP26_RE (Q23)
% ==============================================================================
\begin{question}{7}{mcq}
\source{AIG202c_-_SP_2025_-_FE_3755_14.webp}
\prompt{Which of the following is a benefit of having an integrated AI system in a call center environment?}
\choice{A}{Reducing costs}
\choice{B}{Eliminating the need for human agents}
\choice{C}{Decreasing agent satisfaction due to fear of job loss.}
\choice{D}{Increasing agent turnover}
\answer{A}
\confidence{high}
\evidence{Lợi ích kinh tế lớn nhất của việc tích hợp hệ thống AI vào trung tâm cuộc gọi (call center) là cắt giảm chi phí vận hành (Reducing costs) thông qua việc tự động hóa xử lý các cuộc gọi thường nhật.}
\sources{Deloitte Call Center AI Transformation Report; McKinsey Customer Operations AI}
\explain{Hướng giải: Xác định lợi ích kinh tế thực tế của hệ thống AI đối với trung tâm chăm sóc khách hàng.

Các bước phân tích:
1. Trung tâm cuộc gọi truyền thống có chi phí nhân sự rất lớn cho các cuộc gọi lặp lại (tra cứu cước, kiểm tra đơn hàng, đổi mật khẩu).
2. Ứng dụng trợ lý ảo giọng nói (Voicebot) và Chatbot giúp giải quyết tới 60-80% cuộc gọi cơ bản 24/7 với chi phí tính theo giao dịch cực thấp, giúp doanh nghiệp tiết kiệm hàng triệu USD chi phí vận hành (Reducing costs).

Kiểm tra: Tiết kiệm chi phí (Cost reduction) và tăng tốc độ phục vụ là ROI (Return on Investment) rõ ràng nhất khi đầu tư AI call center.

Vì sao các phương án khác sai:
- B: Loại bỏ hoàn toàn nhân viên con người là sai lầm, con người vẫn cần cho các ca khó.
- C & D: Giảm mức độ hài lòng và tăng tỉ lệ nghỉ việc của nhân viên là các hệ quả tiêu cực nếu quản lý kém, không phải là 'lợi ích' (benefit).

Lưu ý/Khái niệm cần nhớ: AI Call Center giúp giảm chi phí vận hành và tăng năng suất lao động.}
\vocab{Voicebot | Trợ lý ảo giọng nói | Hệ thống AI tự động tương tác bằng giọng nói qua đường dây điện thoại tổng đài}
\end{question}

% ==============================================================================
% Question 8 | Phân loại: 1. Introduction to Artificial Intelligence
% Xuất hiện trong các đề gốc: AIG202c_SP_2025_FE (Q16)
% ==============================================================================
\begin{question}{8}{mcq}
\source{AIG202c_-_SP_2025_-_FE_3755_16.webp}
\prompt{AI is the fusion of many fields of study- Which of these fields, along with Computer Science, plays a role in the application of AP}
\choice{A}{All responses are correct}
\choice{B}{Mathematics}
\choice{C}{Philosophy}
\choice{D}{Statistics}
\answer{A}
\confidence{high}
\evidence{Trí tuệ nhân tạo (AI) là lĩnh vực liên ngành sâu rộng được hình thành từ sự giao thoa của Khoa học máy tính cùng với Toán học (Mathematics), Triết học (Philosophy), Thống kê (Statistics), Ngôn ngữ học, Khoa học thần kinh và Tâm lý học; do đó 'Tất cả các phương án đều đúng' (All responses are correct).}
\sources{Russell & Norvig, AI: A Modern Approach (Ch. 1 The Foundations of Artificial Intelligence)}
\explain{Hướng giải: Xác định các ngành khoa học nền tảng tạo nên Trí tuệ nhân tạo (AI).

Các bước phân tích:
1. Giáo trình kinh điển của Russell & Norvig dành toàn bộ chương mở đầu để phân tích các nền tảng lịch sử cấu thành AI:
   - Triết học (Philosophy): Khởi xướng các câu hỏi về tâm trí, tri thức và ý thức (Aristotle, Descartes).
   - Toán học (Mathematics): Cung cấp các công cụ về logic hình thức, tính toán được và độ phức tạp thuật toán (Turing, Boole).
   - Thống kê (Statistics): Cung cấp lý thuyết xác suất và học từ dữ liệu (Bayes, Gauss).
2. Cả ba ngành Toán học, Triết học và Thống kê đều đóng vai trò nền tảng không thể thiếu cùng với Khoa học máy tính.

Kiểm tra: Do đó phương án A (All responses are correct) là đáp án hoàn hảo nhất.

Vì sao các phương án khác đúng nhưng chưa đủ:
- B, C, D đều đúng nhưng nếu chọn riêng lẻ sẽ thiếu các ngành còn lại.

Lưu ý/Khái niệm cần nhớ: AI = Giao thoa liên ngành giữa Computer Science, Mathematics, Statistics và Philosophy.}
\vocab{Interdisciplinary Field | Lĩnh vực liên ngành | Ngành khoa học kết hợp tri thức và phương pháp từ nhiều lĩnh vực chuyên môn khác nhau}
\end{question}

% ==============================================================================
% Question 9 | Phân loại: 1. Introduction to Artificial Intelligence
% Xuất hiện trong các đề gốc: AIG202c_SP_2025_FE (Q24), AIG202c_SU25_FE (Q1), AIG202c_FA25_FE (Q58)
% ==============================================================================
\begin{question}{9}{mcq}
\source{AIG202c_-_SP_2025_-_FE_3755_24.webp}
\prompt{Which of the following tasks would an AI be able to perform the LEAST successfully?}
\choice{A}{Suggesting what someone might want to buy from a shopping website}
\choice{B}{Using traffic data to work out what cars people are likely to buy}
\choice{C}{Writing a joke}
\choice{D}{Playing hangman}
\answer{C}
\confidence{high}
\evidence{Viết truyện cười / pha trò hài hước (Writing a joke) là nhiệm vụ mà AI thực hiện KÉM THÀNH CÔNG NHẤT (LEAST successfully) vì sự hài hước đòi hỏi khả năng thấu cảm xã hội sâu sắc, cảm nhận về thời điểm (timing), chơi chữ văn hóa và lý thuyết tâm trí (Theory of Mind) của con người.}
\sources{Russell & Norvig, AI: A Modern Approach (Limitations of Current AI); Cognitive Science of Humor and AI}
\explain{Hướng giải: Đánh giá nhiệm vụ nào thách thức nhất đối với năng lực nhận thức của AI hiện đại.

Các bước phân tích:
1. AI hiện đại rất mạnh trong nhận diện mẫu số liệu, xác suất thống kê và các trò chơi có luật cố định:
   - Gợi ý mua sắm (A): Bài toán Recommender Systems quen thuộc dựa trên Collaborative Filtering.
   - Phân tích dữ liệu giao thông (B): Bài toán phân tích dữ liệu dự đoán (Predictive Analytics).
   - Chơi trò chơi Hangman (D): Bài toán tìm kiếm trạng thái và thống kê tần suất xuất hiện chữ cái.
2. Viết truyện cười (C) đòi hỏi sự bất ngờ, đảo ngược kỳ vọng tâm lý, thấu hiểu ngữ cảnh văn hóa và trực giác cảm xúc con người. Đây là lĩnh vực AI làm kém nhất, các câu chuyện cười do AI tạo ra thường gượng gạo hoặc vô nghĩa.

Kiểm tra: Hài hước (Humor) được giới khoa học nhận thức công nhận là đỉnh cao của trí tuệ con người và là rào cản lớn nhất đối với AI.

Vì sao các phương án khác AI làm rất tốt:
- A: Hệ thống gợi ý của Amazon, Shopee hoạt động cực kỳ hiệu quả.
- B: Dự báo thị hiếu xe hơi dựa trên dữ liệu giao thông là bài toán hồi quy/phân loại cơ bản.
- D: Thuật toán tra từ điển giải Hangman tối ưu vượt trội con người.

Lưu ý/Khái niệm cần nhớ: Sự hài hước và sáng tạo cảm xúc là ranh giới khó khăn nhất đối với AI hiện đại.}
\vocab{Theory of Mind | Thuyết tâm trí | Khả năng thấu hiểu và suy đoán trạng thái tâm lý, cảm xúc của người khác}
\end{question}

% ==============================================================================
% Question 10 | Phân loại: 1. Introduction to Artificial Intelligence
% Xuất hiện trong các đề gốc: AIG202c_SP_2025_FE (Q25)
% ==============================================================================
\begin{question}{10}{mcq}
\source{AIG202c_-_SP_2025_-_FE_3755_25.webp}
\prompt{Which one of these jobs COULD NOT be taken by a robot with AP}
\choice{A}{Bank clerk}
\choice{B}{Doctor}
\choice{C}{Telephone sales person}
\choice{D}{Shop worker}
\answer{B}
\confidence{high}
\evidence{Bác sĩ (Doctor) là nghề nghiệp KHÔNG THỂ bị thay thế hoàn toàn bởi robot AI (COULD NOT be taken by a robot), vì nghề y đòi hỏi sự phán đoán lâm sàng phức tạp, tính trách nhiệm pháp lý và đạo đức sinh học tối cao, cùng sự đồng cảm, thấu hiểu sâu sắc giữa người với người.}
\sources{Oxford Martin School: Technology at Work; AMA (American Medical Association) Policy on Augmented Intelligence}
\explain{Hướng giải: Xác định nghề nghiệp không thể bị robot AI thay thế hoàn toàn.

Các bước phân tích:
1. Phân loại mức độ tự động hóa:
   - Giao dịch viên ngân hàng (A), Nhân viên bán hàng qua điện thoại (C), Nhân viên bán hàng trực tiếp (D) là các công việc có tính quy chuẩn cao, thủ tục lặp lại nên rất dễ bị thay thế bởi AI và kiosk thông minh.
   - Bác sĩ (Doctor) đòi hỏi kiến thức y khoa chuyên sâu kết hợp trực giác lâm sàng nhạy bén, khả năng xử lý tình huống sinh tử bất ngờ và sự kết nối cảm xúc, đồng cảm nhân văn với bệnh nhân.
2. Do đó, bác sĩ là nghề mà AI chỉ đóng vai trò hỗ trợ (Augmented Intelligence), không thể thay thế hoàn toàn.

Kiểm tra: Bác sĩ luôn nằm trong nhóm nghề nghiệp an toàn nhất trước làn sóng tự động hóa AI.

Vì sao các phương án khác bị thay thế được:
- A, C, D đều là các công việc giao dịch và bán hàng lặp lại, đã và đang bị robot/AI thay thế mạnh mẽ.

Lưu ý/Khái niệm cần nhớ: Nghề y tế lâm sàng đòi hỏi y đức và sự phán đoán của con người, không thể bị robot thay thế.}
\vocab{Augmented Intelligence | Trí tuệ tăng cường | Khái niệm nhấn mạnh việc AI hỗ trợ và tăng cường năng lực con người thay vì thay thế}
\end{question}

% ==============================================================================
% Question 11 | Phân loại: 1. Introduction to Artificial Intelligence
% Xuất hiện trong các đề gốc: AIG202c_SP_2025_FE (Q26), AIG202c_SU25_FE (Q4)
% ==============================================================================
\begin{question}{11}{mcq}
\source{AIG202c_-_SP_2025_-_FE_3755_26.webp}
\prompt{Which of the following would be a GOOD question to ask, if you were testing to see if you were talking to a human or a computer?}
\choice{A}{What is the capital of Argentina?}
\choice{B}{What is your favorite food?}
\choice{C}{What is the square root of one million?}
\choice{D}{Who has scored the most goals in Premier League history?}
\answer{B}
\confidence{high}
\evidence{Trong phép thử Turing, câu hỏi 'Món ăn yêu thích của bạn là gì?' (What is your favorite food?) là một câu hỏi TỐT (GOOD question) để phân biệt người và máy vì nó đòi hỏi trải nghiệm giác quan và cảm xúc chủ quan của con người.}
\sources{Alan Turing (1950); Turing Test Analysis}
\explain{Hướng giải: Lựa chọn câu hỏi hiệu quả nhất trong phép thử Turing (câu hỏi kiểm tra tính nhất quán với UQ009).

Các bước phân tích:
1. Câu hỏi trắc nghiệm kiến thức khách quan (thủ đô, lịch sử, toán học) máy tính tra cứu tức thì.
2. Câu hỏi trải nghiệm cảm giác ẩm thực đòi hỏi ký ức vị giác và cảm xúc thật của con người.

Kiểm tra: Khớp hoàn toàn với phân tích của câu UQ009.

Lưu ý/Khái niệm cần nhớ: Câu hỏi về trải nghiệm cá nhân và cảm xúc chủ quan là câu hỏi tốt nhất trong Turing Test.}
\vocab{Subjective Experience | Trải nghiệm chủ quan | Trạng thái cảm nhận ý thức cá nhân mà máy móc không sở hữu}
\end{question}

% ==============================================================================
% Question 12 | Phân loại: 1. Introduction to Artificial Intelligence
% Xuất hiện trong các đề gốc: AIG202c_SP_2025_FE (Q29), AIG202c_SU25_FE (Q11), AIG202c_FA25_FE (Q28)
% ==============================================================================
\begin{question}{12}{mcq}
\source{AIG202c_-_SP_2025_-_FE_3755_29.webp}
\prompt{Which of these statements is true?}
\choice{A}{Cognitive systems can only translate small volumes of audio data into their literal text translations at massive speeds}
\choice{B}{Cognitive systems can derive mathematically precise answers following a rigid decision tree approach}
\choice{C}{Cognitive systems can only process neatly organized structured data}
\choice{D}{Cognitive systems can learn from their successes and failures}
\answer{C}
\confidence{high}
\evidence{Deep Learning (Học sâu) là một tập con chuyên biệt của Machine Learning sử dụng các mạng nơ-ron nhiều lớp (layered neural networks) để trích xuất đặc trưng và mô phỏng khả năng ra quyết định phức tạp.}
\sources{LeCun, Bengio, Hinton, Deep Learning (Nature 2015); Goodfellow et al., Deep Learning Textbook}
\explain{Hướng giải: Xác định mối quan hệ phân cấp chính xác giữa AI, Machine Learning và Deep Learning.

Các bước phân tích:
1. Trí tuệ nhân tạo (AI) là vòng tròn lớn nhất (bao trùm).
2. Học máy (Machine Learning) là một tập con của AI (AI \(\supset\) ML).
3. Học sâu (Deep Learning) là một tập con chuyên biệt của Học máy (ML \(\supset\) DL), sử dụng các mạng nơ-ron sâu nhiều tầng (deep neural networks) để tự động học biểu diễn dữ liệu.
4. Do đó, khẳng định 'Deep Learning is a specialized subset of Machine Learning that uses layered neural networks...' là HOÀN TOÀN CHÍNH XÁC.

Kiểm tra: Biểu đồ hình tròn đồng tâm kinh điển: \(\text{AI} \supset \text{ML} \supset \text{DL}\).

Vì sao các phương án khác sai:
- A: Data Science không phải tập con của AI; Data Science là ngành liên ngành giao thoa giữa Thống kê, Khoa học máy tính và Tri thức chuyên ngành.
- B: AI và ML không phải là một; ML chỉ là một phương pháp để đạt được AI.
- D: AI không phải là tập con của Data Science.

Lưu ý/Khái niệm cần nhớ: Mối quan hệ chuẩn: AI bao hàm Machine Learning, Machine Learning bao hàm Deep Learning.}
\vocab{Deep Learning | Học sâu | Phân nhánh của học máy sử dụng mạng nơ-ron nhân tạo nhiều tầng để giải quyết bài toán phức tạp}
\end{question}

% ==============================================================================
% Question 13 | Phân loại: 1. Introduction to Artificial Intelligence
% Xuất hiện trong các đề gốc: AIG202c_SP_2025_FE (Q31), AIG202c_SP26_RE (Q8)
% ==============================================================================
\begin{question}{13}{mcq}
\source{AIG202c_-_SP_2025_-_FE_3755_31.webp}
\prompt{Which field of AI involves interpreting visual information from images or videos?}
\choice{A}{Natural Language Processing (NLP)}
\choice{B}{Robotics}
\choice{C}{Computer Vision}
\choice{D}{Machine Learning}
\answer{C}
\confidence{high}
\evidence{Thị giác máy tính (Computer Vision) là phân nhánh của AI chuyên nghiên cứu việc thu nhận, xử lý và phân tích thông tin thị giác từ hình ảnh hoặc video kỹ thuật số.}
\sources{Szeliski, Computer Vision: Algorithms and Applications; Russell & Norvig, Ch. 24 Perception}
\explain{Hướng giải: Xác định lĩnh vực AI xử lý dữ liệu hình ảnh và video.

Các bước phân tích:
1. Dữ liệu đầu vào: Hình ảnh hoặc video (visual information).
2. Lĩnh vực công nghệ: Computer Vision (Thị giác máy tính).
3. Các bài toán điển hình bao gồm: Phân loại ảnh (Image Classification), Phát hiện vật thể (Object Detection), Phân đoạn ngữ nghĩa (Semantic Segmentation).

Kiểm tra: Đây là định nghĩa phân loại cơ bản trong sách giáo khoa Trí tuệ nhân tạo.

Vì sao các phương án khác sai:
- A (NLP): Xử lý ngôn ngữ tự nhiên (văn bản, tiếng nói).
- B (Robotics): Cơ điện tử, điều khiển chuyển động và tương tác vật lý.
- D (Machine Learning): Phương pháp học thuật toán tổng quát, không riêng về thị giác.

Lưu ý/Khái niệm cần nhớ: Visual information (ảnh/video) = Computer Vision.}
\vocab{Computer Vision | Thị giác máy tính | Lĩnh vực AI cho phép máy tính hiểu và trích xuất thông tin từ thế giới thị giác}
\end{question}

% ==============================================================================
% Question 14 | Phân loại: 1. Introduction to Artificial Intelligence
% Xuất hiện trong các đề gốc: AIG202c_SP_2025_FE (Q34), AIG202c_SU25_FE (Q9)
% ==============================================================================
\begin{question}{14}{mcq}
\source{AIG202c_-_SP_2025_-_FE_3755_34.webp}
\prompt{Which of the following is a scenario where an AI is being asked to make an ETHICAL decision?}
\choice{A}{A self-driving car can only avoid one of Wo pedestrians and must choose which}
\choice{B}{A voice assistant is asked to schedule an appointment on the users calendar}
\choice{C}{A voice assistant is asked whether the user will need an umbrella today}
\choice{D}{A chess computer must choose its next move}
\answer{A}
\confidence{high}
\evidence{Tình huống xe tự hành chỉ có thể tránh một trong hai người đi bộ và buộc phải chọn người nào để tránh (A self-driving car can only avoid one of two pedestrians and must choose which) là một quyết định ĐẠO ĐỨC (Ethical decision) kinh điển, được mô phỏng theo bài toán Xe goòng (The Trolley Problem) trong triết học đạo đức.}
\sources{MIT Moral Machine Experiment; Russell & Norvig, AI: A Modern Approach (Ethics and Safety)}
\explain{Hướng giải: Xác định tình huống đòi hỏi AI phải đưa ra lựa chọn mang tính phán xét giá trị đạo đức con người.

Các bước phân tích:
1. Quyết định đạo đức (Ethical decision) xuất hiện khi có xung đột về các giá trị luân lý và tính mạng con người, nơi không có giải pháp toán học thuần túy đúng-sai tuyệt đối.
2. Tình huống xe tự lái phải lựa chọn tính mạng giữa 2 người đi bộ là phiên bản hiện đại của bài toán đạo đức kinh điển 'Trolley Problem' (Bài toán xe goòng). Đây là một quyết định đạo đức sinh tử phức tạp.

Kiểm tra: Thí nghiệm Moral Machine của MIT nổi tiếng toàn cầu chính là nghiên cứu về tình huống đạo đức này của xe tự lái.

Vì sao các phương án khác không phải quyết định đạo đức:
- B: Lên lịch hẹn là tác vụ quản lý thời gian thuần túy.
- C: Dự báo mang ô là truy vấn dữ kiện thời tiết.
- D: Nước đi cờ vua là tính toán logic trò chơi hình thức.

Lưu ý/Khái niệm cần nhớ: Lựa chọn tính mạng trong tình huống bất khả kháng của xe tự hành là quyết định đạo đức kinh điển.}
\vocab{Trolley Problem | Bài toán xe goòng | Thí nghiệm tư tưởng triết học về sự lựa chọn đạo đức trong tình thế tiến thoái lưỡng nan}
\end{question}

% ==============================================================================
% Question 15 | Phân loại: 1. Introduction to Artificial Intelligence
% Xuất hiện trong các đề gốc: AIG202c_SP_2025_FE (Q35)
% ==============================================================================
\begin{question}{15}{mcq}
\source{AIG202c_-_SP_2025_-_FE_3755_35.webp}
\prompt{Which one of these ways is NOT how AI learns?}
\choice{A}{Supervised Learning}
\choice{B}{Reinforcement Learning}
\choice{C}{Proactive Learning}
\choice{D}{Unsupervised Learning}
\answer{C}
\confidence{high}
\evidence{Ba phương thức học cơ bản của AI là Học có giám sát (Supervised Learning), Học không giám sát (Unsupervised Learning) và Học tăng cường (Reinforcement Learning). Thuật ngữ 'Proactive Learning' KHÔNG PHẢI là một phương thức học chính thống trong phân loại nền tảng của học máy.}
\sources{Russell & Norvig, AI: A Modern Approach (Part V Machine Learning); Mitchell, Machine Learning}
\explain{Hướng giải: Xác định phương thức KHÔNG PHẢI là trụ cột học máy cơ bản của AI.

Các bước phân tích:
1. Ba trụ cột học máy kinh điển được giảng dạy trên toàn thế giới:
   - Supervised Learning (Học có giám sát): Học từ dữ liệu có nhãn.
   - Unsupervised Learning (Học không giám sát): Khám phá cấu trúc từ dữ liệu không nhãn.
   - Reinforcement Learning (Học tăng cường): Học thông qua thử và sai với hàm phần thưởng.
2. Thuật ngữ 'Proactive Learning' (Học chủ động) chỉ là một biến thể nhỏ của Active Learning trong nghiên cứu hẹp, không nằm trong bộ 3 phương thức học nền tảng của AI.

Kiểm tra: Bất kỳ giáo trình AI chuẩn nào cũng chỉ phân loại 3 nhánh chính: Supervised, Unsupervised, Reinforcement.

Vì sao các phương án khác là cách AI học chuẩn:
- A, B, D đều là các phương thức học kinh điển của AI.

Lưu ý/Khái niệm cần nhớ: 3 trụ cột học máy: Supervised, Unsupervised, Reinforcement Learning.}
\vocab{Learning Paradigms | Các mô hình học tập | Ba phương pháp nền tảng giúp máy tính học hỏi từ dữ liệu và tương tác}
\end{question}

% ==============================================================================
% Question 16 | Phân loại: 1. Introduction to Artificial Intelligence
% Xuất hiện trong các đề gốc: AIG202c_SP_2025_FE (Q36), AIG202c_SU25_FE (Q27), AIG202c_SU26_FE (Q25)
% ==============================================================================
\begin{question}{16}{mcq}
\source{AIG202c_-_SP_2025_-_FE_3755_36.webp}
\prompt{What is an example of AI's role in improving customer service and experiences?}
\choice{A}{Lengthening response times and delays}
\choice{B}{Providing personalized recommendations and support}
\choice{C}{Reducing customer engagement and interactions}
\choice{D}{Decreasing customer satisfaction levels}
\answer{B}
\confidence{high}
\evidence{Ví dụ điển hình nhất về vai trò của AI trong việc nâng cao dịch vụ và trải nghiệm khách hàng là cung cấp các đề xuất cá nhân hóa và hỗ trợ kịp thời (Providing personalized recommendations and support).}
\sources{Harvard Business Review: How AI Is Personalizing the Customer Experience; Amazon & Netflix Case Studies}
\explain{Hướng giải: Xác định đóng góp tích cực của AI đối với việc nâng cao trải nghiệm khách hàng.

Các bước phân tích:
1. AI phân tích lịch sử mua sắm, hành vi duyệt web và sở thích cá nhân để đưa ra các gợi ý sản phẩm phù hợp đúng lúc (personalized recommendations).
2. Chatbot và trợ lý AI cung cấp hỗ trợ tức thì 24/7 giải quyết vấn đề nhanh chóng.

Kiểm tra: Cá nhân hóa (Personalization) là lợi ích nổi bật nhất của AI trong thương mại điện tử và dịch vụ số.

Vì sao các phương án khác sai:
- A (Kéo dài thời gian phản hồi), C (Giảm tương tác), D (Giảm mức độ hài lòng) đều là các hệ quả tiêu cực trái ngược với mục tiêu cải thiện dịch vụ.

Lưu ý/Khái niệm cần nhớ: AI nâng cao trải nghiệm khách hàng nhờ cá nhân hóa và hỗ trợ tức thì.}
\vocab{Personalization | Cá nhân hóa | Việc tùy biến dịch vụ và sản phẩm phù hợp với nhu cầu và sở thích riêng của từng cá nhân}
\end{question}

% ==============================================================================
% Question 17 | Phân loại: 1. Introduction to Artificial Intelligence
% Xuất hiện trong các đề gốc: AIG202c_SP_2025_FE (Q40), AIG202c_SU25_FE (Q22), AIG202c_FA25_FE (Q27), AIG202c_SP26_RE (Q24)
% ==============================================================================
\begin{question}{17}{mcq}
\source{AIG202c_-_SP_2025_-_FE_3755_40.webp}
\prompt{What is one of the potential future impacts of AI in the job market?}
\choice{A}{Decrease in demand for skilled professionals}
\choice{B}{Increased Job opportunities across all sectors}
\choice{C}{Shift in job roles, with some jobs becoming obsolete}
\choice{D}{Unchanged job market dynamics}
\answer{C}
\confidence{high}
\evidence{Tác động tương lai chắc chắn nhất của AI đối với thị trường lao động là sự dịch chuyển trong các vai trò công việc (Shift in job roles), trong đó các công việc lặp đi lặp lại sẽ dần trở nên lỗi thời (becoming obsolete), đồng thời phát sinh nhiều vai trò công việc mới kết hợp với AI.}
\sources{World Economic Forum, The Future of Jobs Report; MIT Task Force on the Work of the Future}
\explain{Hướng giải: Đánh giá tác động khách quan và thực tế của AI đối với thị trường việc làm trong tương lai.

Các bước phân tích:
1. AI không đơn thuần chỉ triệt tiêu việc làm hoặc tạo ra việc làm ở tất cả các ngành.
2. Xu hướng chủ đạo là 'Sự chuyển dịch cơ cấu việc làm' (Shift in job roles): các công việc mang tính thủ tục lặp lại sẽ biến mất, trong khi các công việc đòi hỏi sáng tạo, tư duy phản biện, kỹ năng số và làm việc cùng AI sẽ gia tăng.

Kiểm tra: Báo cáo của Diễn đàn Kinh tế Thế giới (WEF) chỉ ra rằng hàng chục triệu công việc truyền thống sẽ lỗi thời nhưng đồng thời hàng trăm triệu công việc mới sẽ xuất hiện.

Vì sao các phương án khác sai:
- A: Nhu cầu về chuyên gia có kỹ năng chuyên môn cao sẽ tăng mạnh chứ không giảm.
- B: Không phải lĩnh vực nào cũng tăng cơ hội việc làm, nhiều ngành sẽ bị thu hẹp.
- D: Thị trường lao động biến đổi sâu sắc, không thể 'unchanged'.

Lưu ý/Khái niệm cần nhớ: AI tái định hình cơ cấu việc làm, thay thế tác vụ lặp lại và tạo ra các vị trí việc làm mới.}
\vocab{Job Displacement | Chuyển dịch việc làm | Quá trình công nghệ thay đổi hoặc thay thế các vai trò lao động truyền thống}
\end{question}

% ==============================================================================
% Question 18 | Phân loại: 1. Introduction to Artificial Intelligence
% Xuất hiện trong các đề gốc: AIG202c_SP_2025_FE (Q41), AIG202c_SU26_FE (Q20)
% ==============================================================================
\begin{question}{18}{mcq}
\source{AIG202c_-_SP_2025_-_FE_3755_41.webp}
\prompt{What does Ethical AI Governance focus on in the development and deployment of AI systems?}
\choice{A}{Establishing strict regulations to limit AI capabilities}
\choice{B}{Ensuring ethical guidelines and standards in AI use}
\choice{C}{The unrestricted development of AI technologies}
\choice{D}{Eliminating the need for ethical considerations in AI development}
\answer{B}
\confidence{high}
\evidence{Quản trị AI Đạo đức (Ethical AI Governance) tập trung vào việc thiết lập và bảo đảm các nguyên tắc, tiêu chuẩn đạo đức được tuân thủ nghiêm ngặt trong suốt quá trình phát triển và ứng dụng AI (Ensuring ethical guidelines and standards in AI use).}
\sources{OECD AI Principles; UNESCO Recommendation on the Ethics of AI; IBM AI Governance Framework}
\explain{Hướng giải: Định nghĩa trọng tâm hoạt động của Quản trị AI có đạo đức (Ethical AI Governance).

Các bước phân tích:
1. Quản trị AI (AI Governance) là hệ thống các chính sách, quy trình và cơ chế kiểm soát nhằm đảm bảo AI được phát triển và sử dụng an toàn, có trách nhiệm và tuân thủ pháp luật.
2. Trọng tâm của Ethical Governance là đưa các tiêu chuẩn đạo đức (công bằng, minh bạch, bảo mật dữ liệu, trách nhiệm giải trình) thành các quy chuẩn thực thi bắt buộc trong doanh nghiệp.

Kiểm tra: Framework quản trị của OECD và EU đều coi việc áp dụng các tiêu chuẩn đạo đức là cốt lõi của quản trị AI.

Vì sao các phương án khác sai:
- A: Mục tiêu không phải kìm hãm hay bóp nghẹt năng lực AI bằng các quy định phi lý.
- C: Phát triển vô tội vạ không kiểm soát là đi ngược lại quản trị.
- D: Loại bỏ xem xét đạo đức là sai trái.

Lưu ý/Khái niệm cần nhớ: Ethical AI Governance = Thiết lập và bảo đảm tuân thủ các tiêu chuẩn đạo đức trong toàn bộ vòng đời AI.}
\vocab{AI Governance | Quản trị AI | Khung quy định và quy trình giám sát nhằm đảm bảo AI hoạt động an toàn, đạo đức và đúng luật}
\end{question}

% ==============================================================================
% Question 19 | Phân loại: 1. Introduction to Artificial Intelligence
% Xuất hiện trong các đề gốc: AIG202c_SP_2025_FE (Q43), AIG202c_SU25_FE (Q3), AIG202c_FA25_FE (Q7), AIG202c_SU26_FE (Q4), AIG202c_SP26_RE (Q15)
% ==============================================================================
\begin{question}{19}{mcq}
\source{AIG202c_-_SP_2025_-_FE_3755_43.webp}
\prompt{What is the primary goal of Artificial Intelligence (AI)?}
\choice{A}{To replace human intelligence entirely}
\choice{B}{To simulate human intelligence in machines}
\choice{C}{To limit the capabilities of machines}
\choice{D}{To create autonomous robots}
\answer{B}
\confidence{high}
\evidence{Mục tiêu cơ bản và định nghĩa nguyên thủy của Trí tuệ nhân tạo là mô phỏng trí thông minh của con người trên máy tính (To simulate human intelligence in machines).}
\sources{John McCarthy, What is AI? (1956); Russell & Norvig, AI: A Modern Approach}
\explain{Hướng giải: Xác định mục tiêu tối thượng và cốt lõi của ngành Trí tuệ nhân tạo (AI).

Các bước phân tích:
1. AI được John McCarthy khai sinh năm 1956 với định nghĩa: 'Khoa học và kỹ thuật chế tạo các máy móc thông minh, đặc biệt là các chương trình máy tính thông minh'.
2. Mục tiêu chính là mô phỏng, nhân rộng hoặc vượt qua các năng lực trí tuệ của con người như học tập, suy lý, nhận thức và thích ứng trên các hệ thống máy móc.

Kiểm tra: Sách giáo trình AIG202c và định nghĩa chuẩn của ACM/IEEE đều khẳng định: 'The primary goal of AI is to create machines capable of performing tasks that typically require human intelligence'.

Vì sao các phương án khác sai:
- A (To replace human intelligence entirely): Mục tiêu của AI không phải xóa sổ trí thông minh loài người mà là hỗ trợ/mô phỏng.
- C (To limit the capabilities of machines): Trái ngược với mục tiêu phát triển công nghệ.
- D (To create autonomous robots): Robot tự hành chỉ là một nhánh ứng dụng phần cứng nhỏ của AI, không phải mục tiêu bao trùm.

Lưu ý/Khái niệm cần nhớ: AI nhằm mục đích mô phỏng trí tuệ con người trên hệ thống máy tính.}
\vocab{Artificial Intelligence | Trí tuệ nhân tạo | Khả năng của hệ thống máy tính mô phỏng các hành vi trí tuệ của con người}
\end{question}

% ==============================================================================
% Question 20 | Phân loại: 1. Introduction to Artificial Intelligence
% Xuất hiện trong các đề gốc: AIG202c_SP_2025_FE (Q47), AIG202c_SU25_FE (Q7), AIG202c_SU26_FE (Q15)
% ==============================================================================
\begin{question}{20}{mcq}
\source{AIG202c_-_SP_2025_-_FE_3755_47.webp}
\prompt{Which of the following is an attribute of Strong or Generalized AI}
\choice{A}{Operate with human-level consciousness}
\choice{B}{Cannot teach itself new strategies}
\choice{C}{Perform independent tasks}
\choice{D}{Can perform specific tasks, but cannot learn new ones}
\answer{A}
\confidence{high}
\evidence{Trí tuệ nhân tạo mạnh hoặc tổng quát (Strong AI / Artificial General Intelligence - AGI) được giả định có năng lực tự ý thức và khả năng nhận thức tương đương bộ não con người (Operate with human-level consciousness / general cognition) trên mọi phương diện.}
\sources{Russell & Norvig, AI: A Modern Approach (Weak vs Strong AI); Searle, Minds, Brains, and Programs (1980)}
\explain{Hướng giải: Xác định thuộc tính bản chất của Trí tuệ nhân tạo mạnh (Strong AI / Generalized AI).

Các bước phân tích:
1. Phân biệt Strong AI và Weak AI:
   - Weak AI (Narrow AI): Chỉ thực hiện tốt một tác vụ cụ thể, không có ý thức hay khả năng khái quát hóa sang tác vụ khác (các hệ thống AI hiện nay đều là Weak AI).
   - Strong AI (AGI): Hệ thống có trí tuệ ngang ngửa con người, có khả năng tự học mọi lĩnh vực và được lý thuyết hóa là có thể vận hành với mức độ tự ý thức tương đương con người (Operate with human-level consciousness/cognition).
2. Phương án A phản ánh định nghĩa triết học kinh điển về Strong AI của John Searle và các nhà nghiên cứu nhận thức.

Kiểm tra: Các phương án B, D là đặc tính của AI yếu/hẹp (Narrow AI), phương án C quá chung chung.

Vì sao các phương án khác sai:
- B: Không thể tự dạy chiến lược mới là nhược điểm của các hệ thống thuật toán tĩnh, không phải Strong AI.
- C: Thực hiện tác vụ độc lập thì Narrow AI cũng làm được.
- D: Thực hiện tác vụ cụ thể nhưng không thể học tác vụ mới chính là định nghĩa của Narrow AI.

Lưu ý/Khái niệm cần nhớ: Strong AI (AGI) = Trí tuệ nhân tạo tổng quát ngang tầm con người với ý thức và năng lực đa lĩnh vực.}
\vocab{Strong AI (AGI) | AI mạnh / AI tổng quát | Hệ thống AI giả định có năng lực tư duy, học tập và ý thức toàn diện như con người}
\end{question}

% ==============================================================================
% Question 21 | Phân loại: 1. Introduction to Artificial Intelligence
% Xuất hiện trong các đề gốc: AIG202c_SP_2025_FE (Q49)
% ==============================================================================
\begin{question}{21}{mcq}
\source{AIG202c_-_SP_2025_-_FE_3755_49.webp}
\prompt{Why haven't scientists been completely successful when developing AP}
\choice{A}{They do not have the time or money}
\choice{B}{The human brain is very simple}
\choice{C}{They are trying to replicate the human heart}
\choice{D}{The human brain is very complex}
\answer{D}
\confidence{high}
\evidence{Lý do lớn nhất khiến các nhà khoa học chưa thể phát triển thành công Trí tuệ nhân tạo tổng quát (AGI) là do bộ não con người vô cùng phức tạp (The human brain is very complex) với gần 100 tỷ nơ-ron và hàng trăm nghìn tỷ khớp thần kinh (synapses) mà khoa học thần kinh hiện đại vẫn chưa thể giải mã hoàn toàn.}
\sources{Kandel et al., Principles of Neural Science; Russell & Norvig, AI: A Modern Approach}
\explain{Hướng giải: Xác định nguyên nhân khoa học căn bản nhất khiến AI chưa thể đạt tới sự hoàn hảo như trí tuệ con người.

Các bước phân tích:
1. Não bộ con người là cấu trúc vật chất phức tạp nhất từng được biết đến trong vũ trụ, với khoảng 86 tỷ nơ-ron và hơn 100 nghìn tỷ kết nối synap hoạt động đồng thời với hiệu suất năng lượng chỉ khoảng 20 Watt.
2. Con người vẫn chưa hiểu hết cơ chế của ý thức (consciousness), trí nhớ dài hạn, khả năng học liên tưởng và tính trực giác của não bộ.
3. Vì chưa hiểu rõ nguyên mẫu sinh học (human brain), các nhà khoa học máy tính chưa thể mô phỏng trọn vẹn trí tuệ nhân tạo hoàn chỉnh.

Kiểm tra: Sự phức tạp của não bộ là rào cản nhận thức lớn nhất được toàn thể giới khoa học thừa nhận.

Vì sao các phương án khác sai:
- A: Khoa học thế giới đang đầu tư hàng trăm tỷ USD và hàng triệu giờ nghiên cứu vào AI, không phải thiếu thời gian/tiền bạc.
- B: Não người không hề đơn giản, đây là phát biểu phản khoa học.
- C: AI mô phỏng chức năng nhận thức của não bộ (brain), không ai cố gắng mô phỏng trái tim (heart) để tạo trí tuệ.

Lưu ý/Khái niệm cần nhớ: Sự phức tạp khổng lồ của bộ não người là rào cản nhận thức lớn nhất của AI tổng quát.}
\vocab{Biological Complexity | Độ phức tạp sinh học | Cấu trúc mạng thần kinh tự nhiên với quy mô và khả năng thích ứng vượt xa máy tính hiện nay}
\end{question}

% ==============================================================================
% Question 22 | Phân loại: 1. Introduction to Artificial Intelligence
% Xuất hiện trong các đề gốc: AIG202c_SP_2025_FE (Q54), AIG202c_SP26_RE (Q3)
% ==============================================================================
\begin{question}{22}{mcq}
\source{AIG202c_-_SP_2025_-_FE_3755_54.webp}
\prompt{Which of the following do you need to create in order to publish a notebook to your GitHub repository?}
\choice{A}{Apps}
\choice{B}{Profile}
\choice{C}{Access token}
\choice{D}{Login credential}
\answer{C}
\confidence{high}
\evidence{Để xuất bản một Jupyter notebook hoặc mã nguồn lên kho lưu trữ GitHub từ xa thông qua API/môi trường phát triển, bạn cần tạo một mã thông báo truy cập cá nhân (Personal Access Token / Access token) để xác thực quyền truy cập an toàn.}
\sources{GitHub Documentation: Creating a personal access token; Watson Studio GitHub Integration}
\explain{Hướng giải: Xác định thành phần bảo mật cần thiết để liên kết và đẩy tài liệu (notebook) lên GitHub.

Các bước phân tích:
1. GitHub đã ngừng xác thực bằng mật khẩu tài khoản trực tiếp qua Git CLI/API từ năm 2021 và chuyển hoàn toàn sang cơ chế xác thực mã thông báo (Token-based authentication).
2. Để các nền tảng như IBM Watson Studio, Google Colab hay VS Code có thể đẩy commit/notebook lên repository, người dùng phải tạo một Personal Access Token (PAT) với quyền `repo` tương ứng.

Kiểm tra: Tài liệu hướng dẫn tích hợp GitHub của IBM Watson Studio yêu cầu bắt buộc phải dán 'Personal Access Token'.

Vì sao các phương án khác sai:
- A (Apps): Không bắt buộc phải tạo app bên thứ ba chỉ để lưu notebook.
- B (Profile): Hồ sơ cá nhân có sẵn, không phải đối tượng xác thực truy cập.
- D (Login credential): Thông tin đăng nhập truyền thống (tên/mật khẩu) không còn được GitHub chấp nhận cho các thao tác đẩy mã qua API.

Lưu ý/Khái niệm cần nhớ: GitHub yêu cầu Access Token (PAT) để cấp quyền xác thực cho ứng dụng bên ngoài.}
\vocab{Personal Access Token (PAT) | Mã thông báo truy cập cá nhân | Chuỗi khóa bí mật dùng để xác thực quyền hạn trên GitHub thay thế mật khẩu}
\end{question}

% ==============================================================================
% Question 23 | Phân loại: 1. Introduction to Artificial Intelligence
% Xuất hiện trong các đề gốc: AIG202c_SP_2025_FE (Q55), AIG202c_SU25_FE (Q23), AIG202c_FA25_FE (Q19), AIG202c_SU26_FE (Q16)
% ==============================================================================
\begin{question}{23}{mcq}
\source{AIG202c_-_SP_2025_-_FE_3755_55.webp}
\prompt{Which of these are attributes of IBM's Cognitive APIs?}
\choice{A}{IBM Cognitive APIs are based on an input/output system, where your input is data that the API feeds into Watson, and the output is the delivery of the overall outcome based on that data.}
\choice{B}{IBM Cognitive APIs can process only structured data as input and unstructured data as output}
\choice{C}{IBM Cognitive APIs serve as a messenger between the input data and the output files on your local machine. These output files can then be transferred to the cloud.}
\choice{D}{IBMs Cognitive APIs help AI professionals to build apps for mobile devices.}
\answer{A}
\confidence{high}
\evidence{Cognitive APIs (như IBM Watson APIs, Azure Cognitive Services) cung cấp các giao diện lập trình ứng dụng trên đám mây cho phép nhà phát triển tích hợp năng lực AI (thị giác, ngôn ngữ, âm thanh) vào ứng dụng của họ mà không cần tự huấn luyện mô hình từ đầu.}
\sources{IBM Watson API Documentation; Microsoft Azure Cognitive Services Overview}
\explain{Hướng giải: Xác định vai trò cốt lõi của các API nhận thức (Cognitive APIs).

Các bước phân tích:
1. Cognitive APIs đóng gói các mô hình AI phức tạp (đã được huấn luyện sẵn) thành các dịch vụ web có thể gọi qua HTTP/RESTful request.
2. Nhà phát triển gửi dữ liệu đầu vào (ảnh, âm thanh, văn bản) và nhận về kết quả phân tích có cấu trúc (JSON) mà không cần xây dựng kiến trúc học sâu bên dưới.

Kiểm tra: Giúp giảm chi phí và thời gian triển khai AI cho các ứng dụng thực tế.

Vì sao các phương án khác sai:
- API không phải là phần cứng máy tính cục bộ.
- API không giới hạn chỉ cho ứng dụng di động mà chạy trên nền tảng đám mây đa nền tảng.

Lưu ý/Khái niệm cần nhớ: Cognitive API = Cầu nối tích hợp trí tuệ nhân tạo đám mây vào ứng dụng qua giao diện lập trình.}
\vocab{Cognitive API | Giao diện lập trình nhận thức | Dịch vụ đám mây cung cấp các khả năng AI thông qua hàm gọi chuẩn hóa}
\end{question}

% ==============================================================================
% Question 24 | Phân loại: 1. Introduction to Artificial Intelligence
% Xuất hiện trong các đề gốc: AIG202c_SP_2025_FE (Q58)
% ==============================================================================
\begin{question}{24}{mcq}
\source{AIG202c_-_SP_2025_-_FE_3755_58.webp}
\prompt{What mindset did Woodside notice amongst its workers after introducing an AI system?}
\choice{A}{Retiring workers resented their years of knowledge being reduced to a computer system.}
\choice{B}{Younger workers resisted contributing to the AI knowledge base.}
\choice{C}{Younger workers found the new technology difficult to incorporate into their day-to-day tasks.}
\choice{D}{Retiring workers are proud to leave their knowledge as a legacy for the future.}
\answer{D}
\confidence{high}
\evidence{Trong nghiên cứu tình huống điển hình của tập đoàn năng lượng Woodside, sau khi triển khai hệ thống AI Watson để số hóa kinh nghiệm kỹ thuật, ban lãnh đạo nhận thấy tâm lý tích cực của các công nhân sắp về hưu là họ cảm thấy tự hào khi để lại tri thức chuyên môn của mình như một di sản cho tương lai (Retiring workers are proud to leave their knowledge as a legacy for the future).}
\sources{IBM Case Study: Woodside Energy Transforms Decision Making with Watson AI}
\explain{Hướng giải: Xác định tâm lý thực tế của người lao động lớn tuổi tại tập đoàn Woodside khi hợp tác chuyển giao tri thức cho hệ thống AI.

Các bước phân tích:
1. Tập đoàn năng lượng Woodside (Úc) đối mặt với tình trạng các kỹ sư kỳ cựu mang theo 30 năm kinh nghiệm nghỉ hưu mà không có cách nào chuyển giao hết cho thế hệ trẻ.
2. Khi triển khai hệ thống nhận thức Watson, các kỹ sư lớn tuổi được mời phỏng vấn để nạp các ghi chép, nhật ký kỹ thuật và kinh nghiệm xử lý sự cố vào cơ sở dữ liệu AI.
3. Kết quả văn hóa doanh nghiệp rất đáng khích lệ: thay vì phản đối, các công nhân sắp nghỉ hưu cảm thấy vô cùng tự hào (proud) vì những cống hiến cả đời của họ được lưu giữ vĩnh viễn như một di sản quý báu cho các thế hệ tương lai của công ty.

Kiểm tra: Đây là bài học văn hóa doanh nghiệp kinh điển trong giáo trình AIG202c khi triển khai AI quản trị tri thức.

Vì sao các phương án khác sai:
- A: Công nhân không hề oán giận (resented), họ hợp tác tích cực.
- B: Kỹ sư trẻ rất hào hứng sử dụng hệ thống để tra cứu thông tin nhanh.
- C: Nhân viên trẻ tiếp thu công nghệ rất nhanh và hiệu quả.

Lưu ý/Khái niệm cần nhớ: Nhân viên kỳ cựu Woodside tự hào để lại tri thức chuyên môn như di sản cho tương lai.}
\vocab{Knowledge Legacy | Di sản tri thức | Sự kế thừa và lưu truyền kinh nghiệm chuyên môn quý giá của các thế hệ đi trước}
\end{question}

% ==============================================================================
% Question 25 | Phân loại: 1. Introduction to Artificial Intelligence
% Xuất hiện trong các đề gốc: AIG202c_SP_2025_FE (Q59), AIG202c_SU25_FE (Q2), AIG202c_SP26_RE (Q9)
% ==============================================================================
\begin{question}{25}{mcq}
\source{AIG202c_-_SP_2025_-_FE_3755_59.webp}
\prompt{What type of learning involves AI systems learning from data without explicit programming?}
\choice{A}{Supervised learning}
\choice{B}{Reinforcement learning}
\choice{C}{Unsupervised learning}
\choice{D}{Deep learning}
\answer{C}
\confidence{high}
\evidence{Học không giám sát (Unsupervised learning) là phương thức học máy trong đó hệ thống tự động học các cấu trúc, mẫu ẩn và phân phối nội tại từ dữ liệu mà không có sự can thiệp của nhãn hay chỉ dẫn lập trình tường minh (without explicit programming or labels).}
\sources{Arthur Samuel (1959); Russell & Norvig, AI: A Modern Approach (Unsupervised Learning)}
\explain{Hướng giải: Xác định phương pháp học máy không phụ thuộc vào nhãn dữ liệu hoặc lập trình tường minh từng quy tắc.

Các bước phân tích:
1. Trong Unsupervised Learning, thuật toán tiếp nhận dữ liệu phi cấu trúc không gán nhãn và tự động tìm kiếm các mối tương quan, cụm (clusters) hoặc đặc trưng ẩn bên trong.
2. Quá trình này diễn ra hoàn toàn tự thân mà không cần con người chỉ định trước đáp án đúng hay lập trình tường minh các luật.

Kiểm tra: Đối chiếu với Supervised Learning (cần nhãn mục tiêu) và Reinforcement Learning (cần hàm phần thưởng).

Vì sao các phương án khác sai:
- A: Supervised learning đòi hỏi nhãn giám sát từ trước.
- B: Reinforcement learning đòi hỏi tương tác với môi trường và hàm thưởng/phạt.
- D: Deep learning là kiến trúc mạng nơ-ron sâu, có thể áp dụng cho cả supervised, unsupervised và RL.

Lưu ý/Khái niệm cần nhớ: Unsupervised Learning = Khám phá cấu trúc dữ liệu tự động không cần nhãn giám sát.}
\vocab{Unsupervised Learning | Học không giám sát | Phương pháp học máy tìm kiếm các mẫu và cấu trúc ẩn trong dữ liệu không gán nhãn}
\end{question}

% ==============================================================================
% Question 26 | Phân loại: 1. Introduction to Artificial Intelligence
% Xuất hiện trong các đề gốc: AIG202c_SU25_FE (Q5), AIG202c_FA25_FE (Q14)
% ==============================================================================
\begin{question}{26}{mcq}
\source{Q5.webp}
\prompt{Which scientific computing library provides data structures and data analysis tools for Python?}
\choice{A}{TensorFlow}
\choice{B}{Pandas}
\choice{C}{YumPies}
\choice{D}{Seahorse}
\answer{B}
\confidence{high}
\evidence{Pandas là thư viện tính toán khoa học hàng đầu trong Python cung cấp các cấu trúc dữ liệu bảng mạnh mẽ (DataFrame, Series) phục vụ thao tác và phân tích dữ liệu.}
\sources{McKinney, Python for Data Analysis; Pandas Documentation}
\explain{Hướng giải: Xác định thư viện Python chuyên dụng cung cấp cấu trúc dữ liệu bảng phục vụ tính toán và phân tích.

Các bước phân tích:
1. Pandas cung cấp cấu trúc dữ liệu chính là `DataFrame` (bảng dữ liệu 2 chiều với các hàng và cột có nhãn) và `Series` (mảng 1 chiều có nhãn).
2. Đây là thư viện tiêu chuẩn trong hệ sinh thái Data Science và AI của Python.

Kiểm tra: NumPy cung cấp `ndarray` (mảng đa chiều thuần số học), trong khi cấu trúc dữ liệu phân tích dạng bảng có nhãn (data manipulation structures) chính là đặc trưng nổi bật của Pandas.

Vì sao các phương án khác sai:
- A (Matplotlib): Thư viện trực quan hóa dữ liệu (vẽ đồ thị).
- C (Scikit-learn): Thư viện thuật toán học máy (Machine Learning algorithms).
- D (SciPy): Thư viện tính toán khoa học nâng cao (tích phân, đại số tuyến tính, tối ưu hóa).

Lưu ý/Khái niệm cần nhớ: Pandas = DataFrames & Series cho thao tác dữ liệu bảng.}
\vocab{DataFrame | Cấu trúc dữ liệu bảng | Cấu trúc dữ liệu 2 chiều có nhãn hàng và cột trong thư viện Pandas}
\end{question}

% ==============================================================================
% Question 27 | Phân loại: 1. Introduction to Artificial Intelligence
% Xuất hiện trong các đề gốc: AIG202c_SU25_FE (Q6)
% ==============================================================================
\begin{question}{27}{mcq}
\source{Q6.webp}
\prompt{How does AI work?}
\choice{A}{It just works you know}
\choice{B}{Combining large amounts of data with fast processing and intelligent algorithms, allowing the software to learn automatically from patterns or features in the data.}
\choice{C}{The computation, mostly on computers, of antiderivatives and definite integrals in term of formulas}
\choice{D}{Another super duper long answer to fool you guys to prove that not all longest answers are the correct ones, how is my quiz going anyway, i hope it is good, hope it is not too long though, thank you for not clicking}
\answer{B}
\confidence{high}
\evidence{Cơ chế hoạt động nền tảng của AI là kết hợp lượng lớn dữ liệu với năng lực xử lý tính toán tốc độ cao và các thuật toán thông minh, cho phép phần mềm tự động học hỏi từ các mẫu và đặc trưng có trong dữ liệu (Combining large amounts of data with fast processing and intelligent algorithms, allowing software to learn automatically from patterns or features in the data).}
\sources{SAS: Artificial Intelligence – What it is and why it matters; Russell & Norvig, AI Foundations}
\explain{Hướng giải: Xác định định nghĩa giải thích chính xác cơ chế vận hành của Trí tuệ nhân tạo (AI).

Các bước phân tích:
1. Ba trụ cột tạo nên sức mạnh của AI hiện đại:
   - Dữ liệu lớn (Big Data): Nguồn tri thức và ví dụ mẫu dồi dào.
   - Hạ tầng tính toán tốc độ cao (Fast Processing / GPU, TPU): Cho phép thực hiện hàng nghìn tỷ phép tính ma trận trong tích tắc.
   - Thuật toán thông minh (Intelligent Algorithms / Neural Networks, Backprop): Cơ chế tự động tối ưu hóa hàm mất mát và trích xuất đặc trưng.
2. Phương án B thể hiện đầy đủ, chuẩn xác và khoa học nhất 3 trụ cột này.

Kiểm tra: Đây là định nghĩa nguyên văn của SAS Institute và được trích dẫn trong bài giảng AIG202c.

Vì sao các phương án khác sai:
- A: Câu trả lời đùa cợt cẩu thả ('It just works you know').
- C: Định nghĩa về giải tích tích phân trong toán học thuần túy.
- D: Đoạn văn gây nhiễu hài hước do người làm đề đưa vào.

Lưu ý/Khái niệm cần nhớ: Cơ chế AI = Dữ liệu lớn + Năng lực tính toán nhanh + Thuật toán thông minh.}
\vocab{Pattern Recognition | Nhận diện mẫu | Năng lực của thuật toán tự động phát hiện các quy luật lặp lại trong dữ liệu}
\end{question}

% ==============================================================================
% Question 28 | Phân loại: 1. Introduction to Artificial Intelligence
% Xuất hiện trong các đề gốc: AIG202c_SU25_FE (Q8)
% ==============================================================================
\begin{question}{28}{mcq}
\source{Q8.webp}
\prompt{Which of these is NOT a current application of AI?}
\choice{A}{Self-Driving vehicles utilizing Computer Vision to navigate around objects}
\choice{B}{Classifying rock samples to identify best places to drill for oil}
\choice{C}{Collaborative Robots helping humans lift heavy containers}
\choice{D}{Making precise patient diagnosis and prescribing independent treatment}
\answer{D}
\confidence{high}
\evidence{Việc đưa ra chẩn đoán bệnh chính xác và kê đơn điều trị hoàn toàn độc lập mà không cần bác sĩ (Making precise patient diagnosis and prescribing independent treatment) KHÔNG PHẢI là một ứng dụng hiện tại của AI; luật pháp và y đức toàn cầu nghiêm cấm AI tự ý kê đơn độc lập, AI y tế hiện nay chỉ đóng vai trò hỗ trợ bác sĩ (Clinical Decision Support).}
\sources{FDA Regulatory Framework for Medical AI Devices; WHO Guidelines on Ethics & Governance of AI for Health}
\explain{Hướng giải: Nhận diện tác vụ KHÔNG PHẢI là ứng dụng hiện tại hợp pháp của AI.

Các bước phân tích:
1. Các ứng dụng hiện tại có thật và đã triển khai rộng rãi:
   - Xe tự hành sử dụng thị giác máy tính điều hướng (A).
   - AI phân loại mẫu đá hỗ trợ thăm dò dầu khí (B - ADNOC, Schlumberger).
   - Robot cộng tác (Cobots) hỗ trợ công nhân khuân vác trong nhà máy (C).
2. Trong ngành y tế, trách nhiệm pháp lý đối với tính mạng con người là tối thượng. Hiện tại, không có quốc gia nào cho phép AI tự động chẩn đoán bệnh và độc lập kê đơn thuốc điều trị mà không có sự ký duyệt của bác sĩ có chứng chỉ hành nghề. AI chỉ là công cụ hỗ trợ gợi ý (AI assistant), không được phép hành nghề y độc lập.

Kiểm tra: Quy định của Cục Quản lý Thực phẩm và Dược phẩm Hoa Kỳ (FDA) phân loại AI y tế là thiết bị hỗ trợ quyết định lâm sàng (SaMD), không thay thế quyền kê đơn của bác sĩ.

Vì sao các phương án khác là ứng dụng hiện tại:
- A: Tesla, Waymo đang chạy thử nghiệm và vận hành xe tự hành.
- B: Phân tích địa chất bằng AI là tiêu chuẩn ngành dầu khí.
- C: Cobots được ứng dụng phổ biến trong kho hàng Amazon.

Lưu ý/Khái niệm cần nhớ: AI y tế hiện nay chỉ hỗ trợ quyết định lâm sàng, không bao giờ tự ý chẩn đoán và kê đơn độc lập.}
\vocab{Clinical Responsibility | Trách nhiệm lâm sàng | Nghĩa vụ pháp lý và đạo đức của nhân viên y tế không thể ủy thác hoàn toàn cho máy móc}
\end{question}

% ==============================================================================
% Question 29 | Phân loại: 1. Introduction to Artificial Intelligence
% Xuất hiện trong các đề gốc: AIG202c_SU25_FE (Q10)
% ==============================================================================
\begin{question}{29}{mcq}
\source{Q10.webp}
\prompt{Which concern in AI focuses on avoiding biased decision-making in algorithms?}
\choice{A}{AI Transparency}
\choice{B}{AI Ethics}
\choice{C}{AI Accountability}
\choice{D}{AI Fairness}
\answer{D}
\confidence{high}
\evidence{Mối quan tâm cụ thể trong AI tập trung vào việc ngăn ngừa việc ra quyết định thiên vị, bất công trong các thuật toán chính là Tính công bằng của AI (AI Fairness).}
\sources{Mehrabi et al. (2021), A Survey on Bias and Fairness in Machine Learning (ACM Computing Surveys)}
\explain{Hướng giải: Xác định khía cạnh đạo đức AI chuyên trách về việc phòng chống thiên kiến và đảm bảo đối xử bình đẳng.

Các bước phân tích:
1. Phân biệt các trụ cột trong AI có trách nhiệm:
   - AI Transparency (Minimizing black box): Minh bạch thuật toán.
   - AI Accountability: Trách nhiệm giải trình của con người.
   - AI Fairness (Công bằng): Đảm bảo các thuật toán không phân biệt đối xử, không thiên vị đối với bất kỳ nhóm đối tượng nào dựa trên chủng tộc, giới tính, tôn giáo hay độ tuổi.
2. Do đó, 'Avoiding biased decision-making' chính là định nghĩa chuẩn của AI Fairness.

Kiểm tra: Các thước đo công bằng (Demographic Parity, Equalized Odds) được xây dựng riêng cho bài toán AI Fairness.

Vì sao các phương án khác sai:
- A (Transparency): Tập trung vào tính minh bạch, có thể giải thích được.
- B (Ethics): Khái niệm đạo đức bao trùm rộng lớn, không chuyên biệt hóa riêng cho vấn đề thiên kiến.
- C (Accountability): Tập trung vào việc quy trách nhiệm pháp lý.

Lưu ý/Khái niệm cần nhớ: AI Fairness = Tránh thiên kiến và đảm bảo công bằng trong các quyết định thuật toán.}
\vocab{AI Fairness | Tính công bằng của AI | Mục tiêu đảm bảo các quyết định của hệ thống AI không tạo ra sự đối xử phân biệt bất công}
\end{question}

% ==============================================================================
% Question 30 | Phân loại: 1. Introduction to Artificial Intelligence
% Xuất hiện trong các đề gốc: AIG202c_SU25_FE (Q12)
% ==============================================================================
\begin{question}{30}{mcq}
\source{Q12.webp}
\prompt{AI is the fusion of many fields of study. Which of these fields, along with Computer Science: plays a role in the application of AI?}
\choice{A}{All responses are correct}
\choice{B}{Mathematics}
\choice{C}{Philosophy}
\choice{D}{Statistics}
\answer{A}
\confidence{high}
\evidence{Toán học (Mathematics), Triết học (Philosophy) và Thống kê (Statistics) đều đóng vai trò nền tảng cấu thành nên các ứng dụng của AI; do đó 'Tất cả các phương án đều đúng' (All responses are correct).}
\sources{Russell & Norvig, Ch. 1 (Foundations of AI)}
\explain{Hướng giải: Xác định các ngành khoa học nền tảng của AI (câu hỏi kiểm tra tính nhất quán với UQ109).

Các bước phân tích:
1. AI kế thừa logic hình thức từ Triết học, công cụ tính toán từ Toán học và lý thuyết suy diễn dữ liệu từ Thống kê.
2. Do đó tất cả các phương án đều đóng vai trò mật thiết trong ứng dụng AI.

Kiểm tra: Khớp hoàn toàn với đáp án của câu UQ109.

Lưu ý/Khái niệm cần nhớ: AI là sự hội tụ liên ngành của Toán học, Triết học và Thống kê cùng Khoa học máy tính.}
\vocab{Foundational Disciplines | Các ngành khoa học nền tảng | Các trụ cột khoa học cơ bản hình thành nên lý thuyết Trí tuệ nhân tạo}
\end{question}

% ==============================================================================
% Question 31 | Phân loại: 1. Introduction to Artificial Intelligence
% Xuất hiện trong các đề gốc: AIG202c_SU25_FE (Q13), AIG202c_FA25_FE (Q1)
% ==============================================================================
\begin{question}{31}{mcq}
\source{Q13.webp}
\prompt{What are two characteristics of artificial intelligence? (Choose Wo)}
\choice{A}{It is an agile, responsive network technology that scales easily and adapts to meet business requirements.}
\choice{B}{It uses systems that mimic human cognitive functions such as learning and problem solving.}
\choice{C}{It translates intent into policies and then uses automation to deploy appropriate network configurations.}
\choice{D}{Artificial intelligence uses statistical techniques to give computers the ability to learn from their environment.}
\choice{E}{It uses intelligent agents that interact with the environment and make decisions to reach a specific goals.}
\answer{B, E}
\confidence{high}
\evidence{Trí tuệ nhân tạo (AI) được định nghĩa là hệ thống mô phỏng các chức năng nhận thức của con người (như học tập, giải quyết vấn đề) và sử dụng các tác tử thông minh (intelligent agents) tương tác với môi trường để đạt mục tiêu cụ thể.}
\sources{Russell & Norvig, AI: A Modern Approach (Ch. 1); IBM AI Foundations}
\explain{Hướng giải: Xác định hai đặc điểm cốt lõi cấu thành nên khái niệm Trí tuệ nhân tạo (AI).

Các bước phân tích:
1. Phương án B phản ánh định nghĩa nhận thức: AI sử dụng các hệ thống mô phỏng chức năng nhận thức của con người như tư duy, suy luận, học hỏi và giải quyết vấn đề.
2. Phương án E phản ánh định nghĩa tác tử (agent-based definition): AI được xây dựng dưới dạng các tác tử thông minh quan sát môi trường qua cảm biến và hành động qua bộ phận tác động để đạt mục tiêu.

Kiểm tra: Đối chiếu với giáo trình chuẩn AI: A Modern Approach, 2 trụ cột định nghĩa AI là 'Hệ thống tư duy/hành động như con người' và 'Tác tử hành động hợp lý (Rational Agents)'.

Vì sao các phương án khác sai:
- Phương án A sai vì đây là định nghĩa về công nghệ mạng linh hoạt (Agile Network Technology), không phải AI.
- Phương án C sai vì đây là khái niệm về Mạng điều khiển dựa trên mục đích (Intent-Based Networking - IBN).
- Phương án D là định nghĩa riêng của Học máy (Machine Learning), hẹp hơn so với khái niệm tổng thể của AI.

Lưu ý/Khái niệm cần nhớ: AI bao gồm mô phỏng nhận thức con người và kiến trúc tác tử thông minh tương tác môi trường.}
\vocab{Intelligent Agent | Tác tử thông minh | Thực thể nhận thức môi trường và hành động để tối ưu hóa mục tiêu}
\end{question}

% ==============================================================================
% Question 32 | Phân loại: 1. Introduction to Artificial Intelligence
% Xuất hiện trong các đề gốc: AIG202c_SU25_FE (Q15), AIG202c_SP26_RE (Q11)
% ==============================================================================
\begin{question}{32}{mcq}
\source{Q15.webp}
\prompt{Many people question whether an AI can make ethical decisions. What do we mean by ETHICS?}
\choice{A}{The laws of particular country}
\choice{B}{The rules of behavior imposed by a programmer or the AI creator}
\choice{C}{Human emotions}
\choice{D}{A set of moral principles}
\answer{D}
\confidence{high}
\evidence{Đạo đức (Ethics) được định nghĩa triết học và phổ thông là một tập hợp các nguyên tắc đạo lý (A set of moral principles) chỉ đạo hành vi và sự phân định đúng sai của con người và hệ thống.}
\sources{Stanford Encyclopedia of Philosophy: Ethics; Oxford English Dictionary}
\explain{Hướng giải: Xác định định nghĩa chuẩn mực của khái niệm Đạo đức (Ethics).

Các bước phân tích:
1. Đạo đức không đồng nhất với pháp luật của một quốc gia cụ thể (luật pháp có thể tụt hậu so với chuẩn mực đạo đức).
2. Đạo đức là hệ thống các giá trị, chuẩn mực và nguyên tắc đạo lý (moral principles) xác định điều gì là đúng (right), công bằng (just) và nhân văn (good) trong hành vi xã hội.

Kiểm tra: Trong AI Ethics, đạo đức được định nghĩa là 'a system of moral principles that guide how technology should be developed and used'.

Vì sao các phương án khác sai:
- A: Pháp luật quốc gia (laws) là quy định pháp lý cưỡng chế, không đồng nhất với đạo đức.
- B: Quy tắc của lập trình viên có thể mang tính chủ quan hoặc sai lầm, không cấu thành chuẩn mực đạo đức phổ quát.
- C: Cảm xúc con người là trạng thái tâm lý sinh học, không phải đạo đức.

Lưu ý/Khái niệm cần nhớ: Ethics = A set of moral principles (Tập hợp các nguyên tắc đạo lý).}
\vocab{Ethics | Đạo đức | Hệ thống các nguyên tắc và giá trị chỉ đạo hành vi đúng đắn và chuẩn mực trong xã hội}
\end{question}

% ==============================================================================
% Question 33 | Phân loại: 1. Introduction to Artificial Intelligence
% Xuất hiện trong các đề gốc: AIG202c_SU25_FE (Q16)
% ==============================================================================
\begin{question}{33}{mcq}
\source{Q16.webp}
\prompt{Which of the following built-in technologies helps you build machine learning models with a simplified graphical interface?}
\choice{A}{Neural NeWork Modeler}
\choice{B}{Data Refinery}
\choice{C}{Watson Knowledge Catalog}
\choice{D}{Intelligent Asset Discovery tool}
\answer{A}
\confidence{high}
\evidence{Neural Network Modeler (trong IBM Watson Studio / SPSS Modeler) là công cụ tích hợp sẵn cho phép người dùng xây dựng, cấu hình và huấn luyện các mô hình mạng nơ-ron học sâu thông qua giao diện đồ họa kéo thả trực quan (simplified graphical interface) mà không cần viết mã phức tạp.}
\sources{IBM Watson Studio: Neural Network Modeler Guide; IBM SPSS Modeler Visual AI}
\explain{Hướng giải: Nhận diện công cụ giao diện đồ họa trực quan (GUI) hỗ trợ xây dựng mô hình mạng nơ-ron trong hệ sinh thái IBM.

Các bước phân tích:
1. Xây dựng mạng nơ-ron bằng mã lệnh (Keras, PyTorch) đòi hỏi kiến thức lập trình vững chắc.
2. Công cụ 'Neural Network Modeler' (Trình mô hình hóa mạng nơ-ron) trong IBM Watson Studio cung cấp bảng vẽ trực quan (canvas): người dùng chỉ cần kéo thả các khối lớp (Input, Convolution, Dense, Activation, Output) và nối các dây tín hiệu để thiết kế kiến trúc mạng một cách dễ dàng.

Kiểm tra: Đây là tính năng nổi bật được giới thiệu trong các bài thực hành xây dựng mô hình học sâu của khóa học.

Vì sao các phương án khác sai:
- B (Data Refinery): Công cụ làm sạch và biến đổi dữ liệu (ETL / Data Wrangling).
- C (Watson Knowledge Catalog): Công cụ quản trị dữ liệu, phân loại và bảo mật thông tin doanh nghiệp.
- D (Intelligent Asset Discovery): Công cụ tìm kiếm tài nguyên dữ liệu.

Lưu ý/Khái niệm cần nhớ: Neural Network Modeler = Công cụ thiết kế mạng nơ-ron bằng giao diện đồ họa kéo thả.}
\vocab{Visual Modeler | Trình tạo mô hình trực quan | Giao diện đồ họa cho phép thiết kế kiến trúc học máy bằng thao tác kéo thả}
\end{question}

% ==============================================================================
% Question 34 | Phân loại: 1. Introduction to Artificial Intelligence
% Xuất hiện trong các đề gốc: AIG202c_SU25_FE (Q29)
% ==============================================================================
\begin{question}{34}{mcq}
\source{Q29.webp}
\prompt{How does AI contribute to environmental conservation and sustainability?}
\choice{A}{Accelerating climate change impacts}
\choice{B}{Optimizing energy usage and waste reduction}
\choice{C}{Increasing pollution levels}
\choice{D}{Ignoring ecological issues}
\answer{B}
\confidence{high}
\evidence{AI đóng góp vào việc bảo tồn môi trường và phát triển bền vững thông qua việc tối ưu hóa mức tiêu thụ năng lượng và giảm thiểu chất thải (Optimizing energy usage and waste reduction) trong các tòa nhà thông minh, lưới điện và quy trình công nghiệp.}
\sources{DeepMind AI Cuts Google Data Centre Cooling Bill by 40%; World Economic Forum: Harnessing AI for the Earth}
\explain{Hướng giải: Nhận diện vai trò tích cực của AI trong bảo vệ môi trường và chống biến đổi khí hậu.

Các bước phân tích:
1. Ứng dụng thực tế: Google ứng dụng thuật toán DeepMind AI để tối ưu hóa hệ thống làm mát trung tâm dữ liệu, giúp cắt giảm 40% năng lượng tiêu thụ.
2. Trong lưới điện thông minh (Smart Grid), AI dự báo nhu cầu phụ tải và tối ưu hóa việc phân phối năng lượng tái tạo (gió, mặt trời).
3. Trong chuỗi cung ứng, AI giúp tối ưu hóa logistics, giảm phát thải carbon và tái chế rác thải thông minh.

Kiểm tra: 'Optimizing energy usage and waste reduction' là trụ cột hành động hàng đầu của AI for Climate.

Vì sao các phương án khác sai:
- A (Thúc đẩy biến đổi khí hậu), C (Tăng mức độ ô nhiễm), D (Bỏ qua các vấn đề sinh thái) đều là các tác động tiêu cực phá hoại môi trường.

Lưu ý/Khái niệm cần nhớ: AI xanh giúp tối ưu hóa hiệu suất năng lượng và giảm thiểu rác thải phát thải.}
\vocab{Green AI | AI xanh | Việc phát triển và ứng dụng AI hướng tới bảo vệ môi trường và tiết kiệm năng lượng}
\end{question}

% ==============================================================================
% Question 35 | Phân loại: 1. Introduction to Artificial Intelligence
% Xuất hiện trong các đề gốc: AIG202c_FA25_FE (Q2)
% ==============================================================================
\begin{question}{35}{mcq}
\source{Q2.webp}
\prompt{Which of these jobs COULD be taken by a robot using AI?}
\choice{A}{Bank clerk}
\choice{B}{Teacher}
\choice{C}{Nurse}
\choice{D}{Doctor}
\answer{A}
\confidence{high}
\evidence{Công việc giao dịch viên ngân hàng (Bank clerk) chủ yếu xử lý các quy trình giao dịch tài chính lặp đi lặp lại theo quy tắc cố định, dễ dàng được tự động hóa bằng AI và robot hơn so với các nghề nghiệp đòi hỏi trực giác, cảm xúc và chuyên môn y tế/giáo dục sâu.}
\sources{McKinsey Global Institute, The Future of Work; Oxford Martin School AI Job Automation Report}
\explain{Hướng giải: Đánh giá khả năng tự động hóa của các nghề nghiệp dựa trên mức độ lặp lại và tính chất công việc.

Các bước phân tích:
1. Giao dịch viên ngân hàng (Bank clerk) thực hiện các tác vụ có tính quy chuẩn cao như tiếp nhận chứng từ, kiểm tra số dư, xác thực thông tin tài khoản và hạch toán giao dịch. Đây là những tác vụ dễ bị tự động hóa bằng chatbot, OCR và RPA tích hợp AI.
2. Các nghề Teacher (giáo viên), Nurse (y tá), Doctor (bác sĩ) đòi hỏi trí tuệ cảm xúc, sự đồng cảm, tương tác phức tạp và kỹ năng lâm sàng chuyên sâu mà robot hiện nay không thể thay thế hoàn toàn.

Kiểm tra: Thực tế ngành ngân hàng hiện nay đã áp dụng rộng rãi kiosk thông minh và AI eKYC để thay thế phần lớn quầy giao dịch cơ bản.

Vì sao các phương án khác sai:
- B (Teacher): Giáo viên đòi hỏi sự thấu hiểu tâm lý học sinh và phương pháp sư phạm linh hoạt.
- C (Nurse): Y tá đòi hỏi sự chăm sóc tận tụy, đồng cảm và xử lý tình huống y tế trực tiếp.
- D (Doctor): Bác sĩ cần năng lực chẩn đoán phức tạp và trách nhiệm đạo đức/pháp lý trực tiếp trên tính mạng bệnh nhân.

Lưu ý/Khái niệm cần nhớ: Các công việc mang tính thủ tục lặp lại (routine cognitive/manual tasks) có nguy cơ tự động hóa cao nhất.}
\vocab{Automation | Tự động hóa | Sử dụng công nghệ để thực hiện công việc với sự can thiệp tối thiểu của con người}
\end{question}

% ==============================================================================
% Question 36 | Phân loại: 1. Introduction to Artificial Intelligence
% Xuất hiện trong các đề gốc: AIG202c_FA25_FE (Q3)
% ==============================================================================
\begin{question}{36}{mcq}
\source{Q3.webp}
\prompt{A Machine Learning technique that helps in detecting the outliers in data.}
\choice{A}{Clustering}
\choice{B}{Classification}
\choice{C}{Anamoly Detection}
\choice{D}{All of them}
\answer{C}
\confidence{high}
\evidence{Phát hiện dị thường (Anomaly Detection hay Outlier Detection) là kỹ thuật Machine Learning chuyên biệt nhằm nhận diện các điểm dữ liệu khác biệt đáng kể so với phần còn lại của tập dữ liệu.}
\sources{Scikit-learn Documentation: Novelty and Outlier Detection; IBM Machine Learning Guide}
\explain{Hướng giải: Xác định kỹ thuật học máy chuyên dụng cho việc tìm kiếm các điểm dữ liệu bất thường (outliers).

Các bước phân tích:
1. Anomaly Detection (Phát hiện dị thường) được định nghĩa chính xác là phương pháp phát hiện các quan sát bất thường, không tuân theo phân phối chung của dữ liệu.
2. Clustering (Phân cụm) nhằm nhóm các điểm tương đồng, Classification (Phân loại) nhằm gán nhãn cho dữ liệu dựa trên các lớp đã biết.

Kiểm tra: Trong sách giáo khoa Machine Learning, bài toán tìm outlier được định danh chính thức là 'Outlier Detection' hoặc 'Anomaly Detection'.

Vì sao các phương án khác sai:
- A (Clustering): Mục đích chính là gom nhóm dữ liệu tương đồng, không phải chuyên dùng tìm outlier dù một số thuật toán phân cụm có thể phụ trợ phát hiện điểm nhiễu.
- B (Classification): Là học có giám sát để phân loại mẫu vào các nhóm định trước, không phải kỹ thuật chuyên biệt cho outlier.
- D (All of them): Không chính xác vì câu hỏi yêu cầu kỹ thuật chuyên biệt trực tiếp (A Machine Learning technique that helps in detecting outliers).

Lưu ý/Khái niệm cần nhớ: Anomaly Detection là nền tảng cho phát hiện gian lận (fraud detection), lỗi hệ thống mạng và y tế.}
\vocab{Anomaly Detection | Phát hiện bất thường | Kỹ thuật tìm kiếm các điểm dữ liệu sai lệch bất thường so với phân phối chung}
\end{question}

% ==============================================================================
% Question 37 | Phân loại: 1. Introduction to Artificial Intelligence
% Xuất hiện trong các đề gốc: AIG202c_FA25_FE (Q4), AIG202c_SU26_FE (Q12), AIG202c_SP26_RE (Q10)
% ==============================================================================
\begin{question}{37}{mcq}
\source{Q4.webp}
\prompt{Advances in the field of Computer Vision make which of the following possible}
\choice{A}{Real-time transcription}
\choice{B}{Detecting cancerous moles in skin images}
\choice{C}{Detecting fraudulent transactions}
\choice{D}{On-demand online tutors}
\answer{B}
\confidence{high}
\evidence{Thị giác máy tính (Computer Vision) xử lý và phân tích hình ảnh/video kỹ thuật số, do đó việc phát hiện nốt ruồi ung thư trong ảnh da liễu (skin images) là ứng dụng trực tiếp của Computer Vision.}
\sources{Stanford University CS231n: Deep Learning for Computer Vision; Nature Medicine Dermatology AI}
\explain{Hướng giải: Xác định ứng dụng liên quan trực tiếp đến xử lý và phân tích hình ảnh (Computer Vision).

Các bước phân tích:
1. Computer Vision (Thị giác máy tính) là nhánh AI nghiên cứu cách máy tính trích xuất, phân tích và hiểu thông tin từ hình ảnh hoặc video.
2. Phương án B nhắc đến việc 'Detecting cancerous moles in skin images' (Phát hiện nốt ruồi ung thư trong ảnh da) sử dụng mô hình mạng nơ-ron tích chập (CNN) phân tích pixel ảnh y tế, hoàn toàn thuộc về Computer Vision.

Kiểm tra: Các ứng dụng da liễu sử dụng CNN phân loại ảnh sắc tố da là một trong những thành tựu nổi tiếng nhất của Computer Vision trong y khoa.

Vì sao các phương án khác sai:
- A (Real-time transcription): Thuộc về Nhận dạng giọng nói (Speech Recognition / Audio Processing) và Xử lý ngôn ngữ tự nhiên (NLP).
- C (Detecting fraudulent transactions): Thuộc về phân tích dữ liệu bảng (Tabular Anomaly Detection / Machine Learning trên giao dịch ngân hàng).
- D (On-demand online tutors): Thuộc về Hệ thống gia sư thông minh (Intelligent Tutoring Systems / NLP / Conversational AI).

Lưu ý/Khái niệm cần nhớ: Computer Vision luôn gắn liền với dữ liệu đầu vào là ảnh (images) hoặc video.}
\vocab{Computer Vision | Thị giác máy tính | Lĩnh vực AI cho phép máy tính thu thập và trích xuất thông tin từ hình ảnh và video}
\end{question}

% ==============================================================================
% Question 38 | Phân loại: 1. Introduction to Artificial Intelligence
% Xuất hiện trong các đề gốc: AIG202c_FA25_FE (Q6), AIG202c_SU26_FE (Q10)
% ==============================================================================
\begin{question}{38}{mcq}
\source{Q6.webp}
\prompt{Which of the following is a classic ethical concern in AI, commonly discussed in research and literature?}
\choice{A}{Racial, gender, or other types of bias.}
\choice{B}{Loss of jobs due to AI replacing workers performing repetitive tasks.}
\choice{C}{Concern about the trusWorthiness of decision-making supported by AI systems.}
\choice{D}{Privacy concerns, such as human faces being photographed and recognized in public spaces.}
\answer{A}
\confidence{high}
\evidence{Định kiến chủng tộc, giới tính hoặc các dạng thiên vị khác (Racial, gender, or other types of bias) là vấn đề đạo đức kinh điển và được thảo luận nhiều nhất trong nghiên cứu AI (Fairness, Accountability, Transparency).}
\sources{Russell & Norvig, AI: A Modern Approach (Ethics and Safety of AI); IEEE Global Initiative on Ethics of Autonomous and Intelligent Systems}
\explain{Hướng giải: Nhận diện mối quan ngại đạo đức kinh điển (classic ethical concern) cốt lõi trong nghiên cứu AI.

Các bước phân tích:
1. AI học từ dữ liệu lịch sử của con người; nếu dữ liệu chứa định kiến xã hội, mô hình sẽ khuếch đại sự thiên vị đó, dẫn đến phân biệt chủng tộc, giới tính trong tuyển dụng, chấm điểm tín dụng hoặc xét án hình sự.
2. Mối quan ngại về 'Racial, gender, or other types of bias' (Bias/Fairness) là chủ đề nền tảng của đạo đức AI (AI Ethics).

Kiểm tra: Hầu hết các bộ tiêu chuẩn đạo đức AI của UNESCO, EU, IEEE đều đặt 'Không phân biệt đối xử & Công bằng' (Non-discrimination & Fairness) lên hàng đầu.

Vì sao các phương án khác sai:
- B: Tác động việc làm là vấn đề kinh tế - xã hội vĩ mô (socio-economic impact).
- C: Độ tin cậy (trustworthiness) là khía cạnh kỹ thuật & quản trị, không phải là 'mối quan ngại đạo đức kinh điển' trực tiếp như bias.
- D: Mối lo về nhận diện khuôn mặt là vấn đề quyền riêng tư (privacy/surveillance).

Lưu ý/Khái niệm cần nhớ: Bias trong AI xuất phát từ dữ liệu huấn luyện và là tâm điểm của nghiên cứu đạo đức AI.}
\vocab{Algorithmic Bias | Thiên kiến thuật toán | Sự sai lệch có hệ thống trong đầu ra của mô hình AI dẫn đến đối xử bất công}
\end{question}

% ==============================================================================
% Question 39 | Phân loại: 1. Introduction to Artificial Intelligence
% Xuất hiện trong các đề gốc: AIG202c_FA25_FE (Q8)
% ==============================================================================
\begin{question}{39}{mcq}
\source{Q8.webp}
\prompt{Which of the following aspects involved in converting the stethoscope into a digital device to support patient diagnoses involves the use of AI?}
\choice{A}{An app on the mobile device that applies learnings from previous diagnosis data to assist the physicians in their current diagnoses}
\choice{B}{Sending digital signals to a mobile device with a machine learning app via bluetooth}
\choice{C}{Inserting a digitizer into the stethoscope tube to convert the analog sound of the heart beat into a digital signal}
\choice{D}{Graphing heart beat data on the mobile device allowing a physician to spot trends}
\answer{A}
\confidence{high}
\evidence{Việc sử dụng một ứng dụng trên thiết bị di động áp dụng các bài học từ dữ liệu chẩn đoán trước đó để hỗ trợ bác sĩ chẩn đoán hiện tại chính là quá trình suy luận học máy (Machine Learning / AI inference).}
\sources{IBM Watson Health Case Studies; Russell & Norvig, Ch. 18 Learning from Examples}
\explain{Hướng giải: Xác định khâu nào trong quy trình ống nghe kỹ thuật số thực sự có sự can thiệp của thuật toán AI.

Các bước phân tích:
1. AI/Học máy cốt lõi là việc 'học từ dữ liệu quá khứ và áp dụng tri thức đó vào dữ liệu mới'.
2. Khâu 'An app on the mobile device that applies learnings from previous diagnosis data to assist the physicians in their current diagnoses' chính là mô hình học máy chẩn đoán hỗ trợ quyết định (Clinical Decision Support System - CDSS).

Kiểm tra: Các khâu còn lại chỉ thuần túy là phần cứng điện tử và hiển thị đồ họa.

Vì sao các phương án khác sai:
- B: Truyền tín hiệu Bluetooth là giao thức truyền thông không dây (wireless communication protocol).
- C: Bộ số hóa (digitizer - ADC) chuyển đổi tín hiệu analog sang digital là mạch điện tử cơ bản.
- D: Vẽ đồ thị nhịp tim (graphing data) là hiển thị trực quan hóa dữ liệu thông thường (data visualization).

Lưu ý/Khái niệm cần nhớ: AI xuất hiện khi có yếu tố học từ dữ liệu quá khứ để đưa ra dự đoán hoặc hỗ trợ quyết định.}
\vocab{Clinical Decision Support System (CDSS) | Hệ thống hỗ trợ ra quyết định lâm sàng | Hệ thống AI phân tích dữ liệu y tế hỗ trợ bác sĩ chẩn đoán}
\end{question}

% ==============================================================================
% Question 40 | Phân loại: 1. Introduction to Artificial Intelligence
% Xuất hiện trong các đề gốc: AIG202c_FA25_FE (Q9)
% ==============================================================================
\begin{question}{40}{mcq}
\source{Q9.webp}
\prompt{If you were testing to determine whether you are conversing with a human or a computer, which of the following questions would be the most effective to ask?}
\choice{A}{What is the capital of Argentina?}
\choice{B}{What is your favorite food?}
\choice{C}{What is the square root of one million?}
\choice{D}{Who has scored the most goals in Premier League history?}
\answer{B}
\confidence{high}
\evidence{Trong Phép thử Turing (Turing Test), các câu hỏi mang tính trải nghiệm chủ quan, cảm xúc cá nhân và sở thích như 'Món ăn yêu thích của bạn là gì?' là hiệu quả nhất để phân biệt con người và máy tính, vì máy tính không có trải nghiệm ẩm thực hay cảm giác vị giác thật.}
\sources{Alan Turing, Computing Machinery and Intelligence (1950); Stanford Encyclopedia of Philosophy: The Turing Test}
\explain{Hướng giải: Lựa chọn câu hỏi hiệu quả nhất trong phép thử Turing để phân biệt người và máy tính.

Các bước phân tích:
1. Phép thử Turing kiểm tra khả năng thể hiện hành vi thông minh tương đương con người.
2. Máy tính có cơ sở dữ liệu khổng lồ nên trả lời cực kỳ chính xác các câu hỏi dữ kiện kiến thức khách quan (thủ đô, lịch sử, toán học).
3. Câu hỏi về trải nghiệm cá nhân như 'What is your favorite food?' (Món ăn yêu thích của bạn là gì?) đòi hỏi cảm xúc, ký ức tuổi thơ, trải nghiệm giác quan và cảm nhận chủ quan của con người.

Kiểm tra: Câu hỏi mở về sở thích và lý do cảm xúc khiến máy tính dễ để lộ câu trả lời máy móc hoặc thiếu tính chân thực cá nhân.

Vì sao các phương án khác sai:
- A (What is the capital of Argentina?): Câu hỏi dữ kiện địa lý (Buenos Aires), máy tính tra cứu tức thì.
- C (What is the square root of one million?): Tính toán số học (\(\sqrt{1000000} = 1000\)), máy tính tính toán hoàn hảo.
- D (Who has scored the most goals in Premier League history?): Dữ kiện thể thao (Alan Shearer), máy tính có sẵn trong cơ sở tri thức.

Lưu ý/Khái niệm cần nhớ: Trong Turing Test, câu hỏi về cảm xúc và trải nghiệm chủ quan luôn thách thức máy tính nhất.}
\vocab{Turing Test | Phép thử Turing | Phương pháp kiểm tra xem máy tính có khả năng thể hiện trí tuệ giống con người hay không}
\end{question}

% ==============================================================================
% Question 41 | Phân loại: 1. Introduction to Artificial Intelligence
% Xuất hiện trong các đề gốc: AIG202c_FA25_FE (Q10)
% ==============================================================================
\begin{question}{41}{mcq}
\source{Q10.webp}
\prompt{why is ethics in AI important today?}
\choice{A}{Because ethical principles must be applied in AI systems now to ensure fairness, transparency, and accountability.}
\choice{B}{Because AI technologies already follow all ethical standards automatically.}
\choice{C}{Because ethical issues in AI will only become relevant in the future.}
\choice{D}{Because ethics is not related to AI development}
\answer{A}
\confidence{high}
\evidence{Đạo đức AI đóng vai trò sống còn ngay trong hiện tại vì các nguyên tắc đạo đức phải được lồng ghép ngay từ khâu thiết kế để đảm bảo tính công bằng (fairness), minh bạch (transparency) và trách nhiệm giải trình (accountability).}
\sources{UNESCO Recommendation on the Ethics of Artificial Intelligence; High-Level Expert Group on AI (EU)}
\explain{Hướng giải: Xác định lý do tại sao đạo đức trong AI là cấp thiết ngay tại thời điểm hiện tại.

Các bước phân tích:
1. AI đang được triển khai trên quy mô lớn ảnh hưởng trực tiếp đến đời sống con người (tuyển dụng, y tế, tòa án, tài chính).
2. Nếu không áp dụng các chuẩn mực đạo đức ngay từ bây giờ, hệ thống AI có thể gây ra những tổn hại xã hội nghiêm trọng do thiên kiến, thiếu minh bạch và không thể quy trách nhiệm.

Kiểm tra: Nguyên tắc 'Ethics by Design' yêu cầu đạo đức phải được cài đặt ngay từ khi phát triển mô hình, không thể để đến tương lai.

Vì sao các phương án khác sai:
- B: AI không tự động tuân theo đạo đức; mô hình chỉ tối ưu hóa hàm toán học trên dữ liệu được nạp vào.
- C: Các vấn đề đạo đức đã và đang xảy ra hàng ngày trong hiện tại, không phải chỉ ở tương lai.
- D: Đạo đức liên quan mật thiết và trực tiếp tới sự phát triển và tồn vong của công nghệ AI.

Lưu ý/Khái niệm cần nhớ: Đạo đức AI đòi hỏi 3 trụ cột cơ bản: Công bằng, Minh bạch và Trách nhiệm giải trình.}
\vocab{AI Ethics | Đạo đức AI | Bộ nguyên tắc chỉ đạo việc phát triển và ứng dụng AI có trách nhiệm và nhân văn}
\end{question}

% ==============================================================================
% Question 42 | Phân loại: 1. Introduction to Artificial Intelligence
% Xuất hiện trong các đề gốc: AIG202c_FA25_FE (Q11), AIG202c_SP26_RE (Q5)
% ==============================================================================
\begin{question}{42}{mcq}
\source{Q11.webp}
\prompt{Which of the following would be a BAD question to ask, if you were testing to see if you were talking to a human or a computer?}
\choice{A}{Do an impression of Simon Cowell}
\choice{B}{Tell me who your best friend reminds you of}
\choice{C}{Who would you get to play yourself in a movie and why?}
\choice{D}{Can you play chess?}
\answer{D}
\confidence{high}
\evidence{Hỏi 'Can you play chess?' là một câu hỏi tồi (BAD question) trong phép thử Turing vì cả con người và máy tính đều có thể chơi cờ vua, và máy tính (như Deep Blue, Stockfish) chơi cực giỏi, nên câu hỏi này không giúp phân biệt đối phương là người hay máy.}
\sources{Alan Turing, Computing Machinery and Intelligence; Russell & Norvig, Turing Test Discussion}
\explain{Hướng giải: Xác định câu hỏi KHÔNG có khả năng phân biệt giữa con người và máy tính trong phép thử Turing.

Các bước phân tích:
1. Phép thử Turing cần các câu hỏi mà máy tính khó bắt chước được cảm xúc, sự hài hước, ngôn ngữ cơ thể hoặc liên tưởng cá nhân của con người.
2. Câu 'Can you play chess?' (Bạn có biết chơi cờ vua không?) là câu hỏi đóng; máy tính từ lâu đã chơi cờ vua ở đẳng cấp siêu phàm (Deep Blue thắng Kasparov năm 1997). Cả người và máy đều có thể trả lời 'Yes', hoàn toàn không mang lại giá trị phân biệt.
3. Các câu A, B, C đều yêu cầu sự hài hước (bắt chước Simon Cowell), liên tưởng cảm xúc tình bạn (người bạn thân gợi nhớ ai) và tự nhận thức bản thân (ai đóng vai mình trong phim).

Kiểm tra: Trong lý thuyết Turing Test, các nhiệm vụ mang tính logic hình thức như cờ vua không thể dùng để kiểm tra tính nhân bản.

Vì sao các phương án khác sai:
- A: Bắt chước phong cách châm biếm của một người nổi tiếng là tác vụ mang tính hài hước và ngữ cảnh sâu.
- B: Tình cảm bạn bè đòi hỏi trải nghiệm quan hệ xã hội phức tạp.
- C: Tự hình dung bản thân và lựa chọn diễn viên thể hiện năng lực tự nhận thức (self-awareness).

Lưu ý/Khái niệm cần nhớ: Câu hỏi logic/tính toán/trò chơi quy tắc không dùng để phân biệt người với máy tính.}
\vocab{Imitation Game | Trò chơi bắt chước | Tên gốc mà Alan Turing dùng để mô tả phép thử Turing}
\end{question}

% ==============================================================================
% Question 43 | Phân loại: 1. Introduction to Artificial Intelligence
% Xuất hiện trong các đề gốc: AIG202c_FA25_FE (Q12)
% ==============================================================================
\begin{question}{43}{mcq}
\source{Q12.webp}
\prompt{Which of the following demonstrates an application of Artificial Intelligence in action based on the provided scenarios?}
\choice{A}{IBM Watson assisting oil and gas workers with technical information retrieval.}
\choice{B}{Analyzing emotions in Grammy nominated song lyrics using AI.}
\choice{C}{Detecting patterns in movement datasets to aid patients with neurological damage.}
\choice{D}{Using facial recognition algorithms for suspect identification by law enforcement authorities.}
\answer{A}
\confidence{high}
\evidence{IBM Watson hỗ trợ công nhân dầu khí truy xuất và xử lý thông tin kỹ thuật phức tạp là ứng dụng điển hình của Hệ thống tri thức và Xử lý ngôn ngữ tự nhiên trong thực tế sản xuất công nghiệp.}
\sources{IBM Watson Case Studies: Oil & Gas Industry; Russell & Norvig, Knowledge-based Systems}
\explain{Hướng giải: Nhận diện kịch bản ứng dụng AI cụ thể đã được triển khai trong thực tế doanh nghiệp.

Các bước phân tích:
1. IBM Watson là hệ sinh thái AI doanh nghiệp hàng đầu, chuyên xử lý dữ liệu phi cấu trúc, văn bản kỹ thuật và hỗ trợ chuyên gia tìm kiếm câu trả lời chính xác trong tài liệu khổng lồ của ngành dầu khí.
2. Đây là kịch bản kinh điển được giảng dạy trong học phần AIG202c về giải pháp AI trong công nghiệp.

Kiểm tra: Kịch bản hợp tác giữa IBM Watson và các tập đoàn năng lượng lớn (như Woodside Energy, ADNOC) là ví dụ chuẩn mực trong tài liệu học phần.

Vì sao các phương án khác sai:
- Các phương án B, C, D là các ví dụ hoặc thử nghiệm nhỏ, không phải tình huống điển hình được cung cấp trong ngữ cảnh bài giảng Watson doanh nghiệp.

Lưu ý/Khái niệm cần nhớ: AI trong doanh nghiệp tập trung vào việc quản trị tri thức và hỗ trợ truy xuất thông tin phi cấu trúc.}
\vocab{Knowledge Retrieval | Truy xuất tri thức | Việc sử dụng AI để tìm kiếm và tổng hợp thông tin từ kho tài liệu khổng lồ}
\end{question}

% ==============================================================================
% Question 44 | Phân loại: 1. Introduction to Artificial Intelligence
% Xuất hiện trong các đề gốc: AIG202c_FA25_FE (Q13)
% ==============================================================================
\begin{question}{44}{mcq}
\source{Q13.webp}
\prompt{Which of the following statements about datasets used in Machine Learning is NOT true?}
\choice{A}{Validation data subset is used to validate results and fine-tune the algorithm's parameters}
\choice{B}{Training data is used to fine-tune algorithm's parameters and evaluate how good the model is}
\choice{C}{Training subset is the data used to train the algorithm}
\answer{B}
\confidence{high}
\evidence{Tập huấn luyện (Training data) chỉ được sử dụng để khớp mô hình và tối ưu hóa các tham số (parameters). Việc đánh giá mô hình (evaluate how good the model is) KHÔNG ĐƯỢC sử dụng tập huấn luyện vì sẽ dẫn đến hiện tượng quá khớp (overfitting) và sai lệch kết quả đánh giá; việc đánh giá phải dùng tập validation hoặc test.}
\sources{Hastie, Tibshirani, Friedman, The Elements of Statistical Learning (Ch. 7); Andrew Ng, Machine Learning Yearning}
\explain{Hướng giải: Xác định phát biểu SAI về các tập dữ liệu trong quy trình Học máy (Machine Learning).

Các bước phân tích:
1. Trong Học máy, dữ liệu thường được chia làm 3 tập: Training set (Tập huấn luyện), Validation set (Tập kiểm thực), Test set (Tập kiểm thử).
2. Tập Training dùng để học tham số mô hình (weights, biases).
3. Tập Validation dùng để tinh chỉnh siêu tham số (hyperparameters) và lựa chọn mô hình.
4. Tập Test dùng để đánh giá độc lập năng lực tổng quát hóa (generalization) của mô hình.
5. Phương án B phát biểu: 'Training data is used to fine-tune algorithm\'s parameters and evaluate how good the model is' là SAI, vì không thể dùng chính tập huấn luyện để đánh giá độ tốt khách quan của mô hình.

Kiểm tra: Nếu đánh giá mô hình trên chính tập huấn luyện, mô hình học vẹt (overfitted) sẽ đạt độ chính xác 100% nhưng hoàn toàn vô dụng trên thực tế.

Vì sao các phương án khác đúng (không được chọn):
- A: Đúng, tập validation dùng để kiểm định kết quả và tinh chỉnh thông số.
- C: Đúng, tập training là dữ liệu cung cấp cho thuật toán học.

Lưu ý/Khái niệm cần nhớ: Đánh giá mô hình luôn phải thực hiện trên tập dữ liệu chưa từng thấy (unseen data: validation/test).}
\vocab{Validation Set | Tập kiểm thực | Tập dữ liệu dùng để tinh chỉnh siêu tham số và đánh giá trung gian mô hình}
\end{question}

% ==============================================================================
% Question 45 | Phân loại: 1. Introduction to Artificial Intelligence
% Xuất hiện trong các đề gốc: AIG202c_FA25_FE (Q15)
% ==============================================================================
\begin{question}{45}{mcq}
\source{Q15.webp}
\prompt{Which of the following is an attribute of Supervised Learning?}
\choice{A}{Relies on providing the machine learning algorithm unlabeled data and letting the machine infer qualities}
\choice{B}{Tries its best to maximize its rewards by trying different combinations of allowed actions within the provided constraints}
\choice{C}{Relies on providing the machine learning algorithm human-labeled data - the more samples you provide, the more precise the algorithm becomes in classifying new data}
\choice{D}{Relies on providing the machine learning algorithm with a set of rules and constraints and letting it learn how to achieve its goals}
\answer{B}
\confidence{high}
\evidence{Đặc trưng cốt lõi của Học có giám sát (Supervised Learning) là dữ liệu huấn luyện phải đi kèm với nhãn mục tiêu (labeled data: input \(X\) và ground truth \(Y\)).}
\sources{Russell & Norvig, AI: A Modern Approach (Ch. 18); Bishop, Pattern Recognition and Machine Learning}
\explain{Hướng giải: Xác định thuộc tính bản chất của Học có giám sát (Supervised Learning).

Các bước phân tích:
1. Học có giám sát là mô hình học từ các cặp dữ liệu vào - ra \((x_i, y_i)\), trong đó \(y_i\) là nhãn (label) do con người hoặc hệ thống gán từ trước.
2. Mục tiêu là học hàm số \(f(X) \to Y\) để dự đoán nhãn cho các dữ liệu mới.

Kiểm tra: Đối lập với Unsupervised Learning (học không giám sát - dữ liệu không có nhãn) và Reinforcement Learning (học qua phần thưởng/phạt).

Vì sao các phương án khác sai:
- Dữ liệu không có nhãn (Unlabeled data) là của Unsupervised Learning.
- Phần thưởng môi trường (Rewards) là của Reinforcement Learning.

Lưu ý/Khái niệm cần nhớ: Supervised Learning = Input Data + Correct Labels.}
\vocab{Supervised Learning | Học có giám sát | Phương pháp học máy mà thuật toán được huấn luyện trên dữ liệu đã gán nhãn}
\end{question}

% ==============================================================================
% Question 46 | Phân loại: 1. Introduction to Artificial Intelligence
% Xuất hiện trong các đề gốc: AIG202c_FA25_FE (Q16)
% ==============================================================================
\begin{question}{46}{mcq}
\source{Q16.webp}
\prompt{Which of the following are examples of bias in an AI system?}
\choice{A}{Facial recognition systems performing well for individuals of all skin tones.}
\choice{B}{Image recognition systems associating images of kitchens, shops, and laundry with women rather than men.}
\choice{C}{Customers not being aware that they are interacting with a chatbot on a company website.}
\choice{D}{AI systems in call centers providing context sensitive assistance to staff.}
\answer{C}
\confidence{high}
\evidence{Nhiệm vụ phân tích bình luận khách hàng trên mạng xã hội để xác định cảm xúc tích cực hay tiêu cực là bài toán Phân tích sắc thái (Sentiment Analysis / Emotion Detection), thuộc nhánh Xử lý ngôn ngữ tự nhiên (NLP).}
\sources{Jurafsky & Martin, Speech and Language Processing (Sentiment Analysis); IBM Watson Natural Language Understanding}
\explain{Hướng giải: Nhận diện tác vụ xử lý ngôn ngữ tự nhiên khi phân tích cảm xúc từ bài đăng mạng xã hội.

Các bước phân tích:
1. Dữ liệu văn bản từ bình luận khách hàng (text comments) được phân tích để xác định thái độ, cảm xúc (tích cực, tiêu cực, trung tính).
2. Kỹ thuật này được gọi là Sentiment Analysis (Phân tích tình cảm / cảm xúc).

Kiểm tra: Đây là một trong những ứng dụng phổ biến nhất của NLP trong chăm sóc khách hàng và tiếp thị số.

Vì sao các phương án khác sai:
- Nhận diện khuôn mặt (Face recognition) thuộc về Computer Vision.
- Dự báo chuỗi thời gian doanh số thuộc về Time-series Regression.

Lưu ý/Khái niệm cần nhớ: Sentiment Analysis = Phân tích thái độ/cảm xúc từ văn bản.}
\vocab{Sentiment Analysis | Phân tích cảm xúc | Tác vụ NLP xác định thái độ tích cực, tiêu cực hoặc trung tính của văn bản}
\end{question}

% ==============================================================================
% Question 47 | Phân loại: 1. Introduction to Artificial Intelligence
% Xuất hiện trong các đề gốc: AIG202c_FA25_FE (Q20), AIG202c_SU26_FE (Q30)
% ==============================================================================
\begin{question}{47}{mcq}
\source{Q20.webp}
\prompt{The acronym API is widely used in software development and refers to a set of protocols, routines, and tools for building software applications. Which of the following accurately represents what the acronym API stands for?}
\choice{A}{Application Programming Interface}
\choice{B}{Abstract Python Interface}
\choice{C}{Algorithmic Programming Interface}
\choice{D}{Abstract Programming Interface}
\answer{A}
\confidence{high}
\evidence{Từ viết tắt API trong công nghệ phần mềm là viết tắt của Application Programming Interface (Giao diện lập trình ứng dụng).}
\sources{IEEE Standard Glossary of Software Engineering Terminology; Tanenbaum, Computer Networks}
\explain{Hướng giải: Xác định nghĩa chuẩn của từ viết tắt API trong ngành công nghệ phần mềm.

Các bước phân tích:
1. A = Application (Ứng dụng).
2. P = Programming (Lập trình).
3. I = Interface (Giao diện).
4. Kết hợp lại: Application Programming Interface.

Kiểm tra: Đây là thuật ngữ phổ thông toàn cầu trong công nghệ thông tin.

Vì sao các phương án khác sai:
- Các phương án B, C, D đều bịa ra các từ như 'Abstract Python Interface' hay 'Algorithmic' là sai hoàn toàn.

Lưu ý/Khái niệm cần nhớ: API = Application Programming Interface.}
\vocab{API | Giao diện lập trình ứng dụng | Tập hợp các quy tắc và cơ chế cho phép các phần mềm tương tác với nhau}
\end{question}

% ==============================================================================
% Question 48 | Phân loại: 1. Introduction to Artificial Intelligence
% Xuất hiện trong các đề gốc: AIG202c_FA25_FE (Q22)
% ==============================================================================
\begin{question}{48}{mcq}
\source{Q22.webp}
\prompt{Which of the following statements about Transfer Learning is INCORRECT?}
\choice{A}{Transfer Learning requires massive amounts of training data for each new related use case.}
\choice{B}{Transfer Learning enables Watson to learn more from less.}
\choice{C}{It is because of Transfer Learning that the learning model can be fed prior knowledge, eliminating the need to train it from scratch.}
\choice{D}{Transfer Learning is a machine learning technique where a model trained on one task is repurposed on a second related task.}
\answer{A}
\confidence{high}
\evidence{Học chuyển giao (Transfer Learning) tận dụng tri thức đã học từ mô hình gốc để áp dụng cho tác vụ mới, do đó NÓ CHỈ CẦN MỘT LƯỢNG NHỎ DỮ LIỆU HUẤN LUYỆN, không cần lượng dữ liệu khổng lồ (massive amounts). Phát biểu A cho rằng nó đòi hỏi 'massive amounts of training data for each new related use case' là HOÀN TOÀN SAI.}
\sources{Pan & Yang (2010), A Survey on Transfer Learning; Howard & Ruder (2018), ULMFiT}
\explain{Hướng giải: Tìm phát biểu KHÔNG ĐÚNG (INCORRECT) về Học chuyển giao (Transfer Learning).

Các bước phân tích:
1. Ưu điểm lớn nhất của Transfer Learning là giải quyết vấn đề khan hiếm dữ liệu. Bằng cách tái sử dụng trọng số từ mô hình đã huấn luyện trước (pre-trained model), ta chỉ cần tập dữ liệu nhỏ (small dataset) để fine-tune cho tác vụ mới.
2. Phát biểu A khẳng định: 'Transfer Learning requires massive amounts of training data for each new related use case' là SAI, vì nó đi ngược lại bản chất và lợi ích cốt lõi của Transfer Learning.

Kiểm tra: Đối chiếu với các phát biểu còn lại: B, C, D đều mô tả chính xác cơ chế của Transfer Learning (giúp Watson học nhiều hơn từ ít dữ liệu hơn, không cần huấn luyện lại từ đầu).

Vì sao các phương án khác đúng (không chọn):
- B: Đúng, Transfer Learning cho phép học từ ít dữ liệu.
- C: Đúng, mô hình được nạp tri thức trước đó.
- D: Đúng, đây là định nghĩa chuẩn của Transfer Learning.

Lưu ý/Khái niệm cần nhớ: Transfer Learning giúp giảm thiểu triệt để nhu cầu về dữ liệu huấn luyện và tài nguyên tính toán.}
\vocab{Transfer Learning | Học chuyển giao | Kỹ thuật tái sử dụng mô hình đã huấn luyện trên một tác vụ cho tác vụ liên quan khác}
\end{question}

% ==============================================================================
% Question 49 | Phân loại: 1. Introduction to Artificial Intelligence
% Xuất hiện trong các đề gốc: AIG202c_FA25_FE (Q26), AIG202c_SP26_RE (Q21)
% ==============================================================================
\begin{question}{49}{mcq}
\source{Q26.webp}
\prompt{How is Watson currently aiding in solving business problems based on the provided scenarios?}
\choice{A}{Gathering expertise learned from thousands of employees to reduce time spent on tracking down expert knowledge.}
\choice{B}{Using data from sensors to detect moisture levels in the air, helping irrigation systems give water where it's needed most.}
\choice{C}{Identifying the best routes based on traffic, temperature, and road conditions to transport bees for pollination.}
\choice{D}{Analyzing financial data to predict market trends and stock performance accurately.}
\answer{A}
\confidence{high}
\evidence{IBM Watson giúp giải quyết bài toán quản trị tri thức doanh nghiệp bằng cách tổng hợp chuyên môn từ hàng ngàn nhân viên kỳ cựu để rút ngắn thời gian tìm kiếm tri thức chuyên gia cho nhân viên mới.}
\sources{IBM Watson Discovery & Knowledge Catalog Business Case Studies}
\explain{Hướng giải: Xác định cách thức IBM Watson hỗ trợ các doanh nghiệp giải quyết bài toán vận hành tri thức.

Các bước phân tích:
1. Doanh nghiệp lớn đối mặt với tình trạng phân mảnh thông tin: nhân viên mất nhiều giờ để tìm đúng chuyên gia hoặc tài liệu cần thiết.
2. Watson thu thập, lập chỉ mục và học từ tri thức của hàng ngàn nhân viên, cho phép nhân viên khác truy vấn tức thì câu trả lời chuyên môn, tiết kiệm thời gian đáng kể.

Kiểm tra: Đây là ví dụ cốt lõi trong bài giảng về giá trị kinh doanh của hệ thống nhận thức Watson.

Vì sao các phương án khác sai:
- B: Cảm biến độ ẩm nông nghiệp là ứng dụng IoT nông nghiệp thông minh, không phải thế mạnh cốt lõi đặc trưng của Watson quản trị tri thức.
- C: Lộ trình vận chuyển ong là một dự án nghiên cứu cụ thể, không phải kịch bản giải quyết bài toán doanh nghiệp bao trùm.
- D: Dự báo thị trường tài chính không phải là trọng tâm trong kịch bản doanh nghiệp được giảng dạy ở phần này.

Lưu ý/Khái niệm cần nhớ: Watson đóng vai trò bộ não tri thức tập trung cho doanh nghiệp.}
\vocab{Enterprise Knowledge Management | Quản trị tri thức doanh nghiệp | Hệ thống thu thập, lưu trữ và chia sẻ tri thức nội bộ tổ chức}
\end{question}

% ==============================================================================
% Question 50 | Phân loại: 1. Introduction to Artificial Intelligence
% Xuất hiện trong các đề gốc: AIG202c_FA25_FE (Q30)
% ==============================================================================
\begin{question}{50}{mcq}
\source{Q30.webp}
\prompt{After you uploaded your document, and it has finished processing. You click on the "Search" tag, what are the buo search option in which you can extract useful information extracted from your document?}
\choice{A}{I-Jse the English Language}
\choice{B}{Use natural language}
\choice{C}{Use the Discovery Query Language}
\choice{D}{Use unnatural language}
\answer{B, C}
\confidence{high}
\evidence{Trong IBM Watson Discovery, sau khi tải tài liệu lên và chuyển sang tab 'Search', hai tùy chọn tìm kiếm chính để trích xuất thông tin là: Sử dụng ngôn ngữ tự nhiên (Natural Language Query) và Sử dụng ngôn ngữ truy vấn Discovery (Discovery Query Language - DQL).}
\sources{IBM Watson Discovery Documentation: Querying Documents; Natural Language vs DQL}
\explain{Hướng giải: Xác định hai cơ chế tìm kiếm trong công cụ phân tích tài liệu IBM Watson Discovery.

Các bước phân tích:
1. IBM Watson Discovery cho phép người dùng tìm kiếm thông tin theo 2 phương thức chính:
   - Natural Language Query (Phương án B): Nhập câu hỏi bình thường bằng ngôn ngữ tự nhiên (ví dụ: 'What is the refund policy?').
   - Discovery Query Language (Phương án C): Sử dụng cú pháp truy vấn có cấu trúc đặc biệt của Discovery (DQL) để lọc chính xác theo các trường, thực thể hoặc quan hệ.

Kiểm tra: Trong tài liệu chính thức của IBM Watson Discovery, giao diện tìm kiếm chỉ có 2 chế độ này.

Vì sao các phương án khác sai:
- A (Use the English Language): Không phải là tên một tùy chọn tìm kiếm kỹ thuật trong giao diện Discovery.
- D (Use unnatural language): Phương án vô nghĩa, được tạo ra làm bẫy.

Lưu ý/Khái niệm cần nhớ: Watson Discovery hỗ trợ 2 chế độ tìm kiếm: Natural Language và Discovery Query Language (DQL).}
\vocab{Discovery Query Language (DQL) | Ngôn ngữ truy vấn Discovery | Cú pháp truy vấn chuyên dụng của IBM Watson để tìm kiếm dữ liệu có cấu trúc}
\end{question}

% ==============================================================================
% Question 51 | Phân loại: 1. Introduction to Artificial Intelligence
% Xuất hiện trong các đề gốc: AIG202c_FA25_FE (Q50), AIG202c_SP26_RE (Q51)
% ==============================================================================
\begin{question}{51}{mcq}
\source{Q50.webp}
\prompt{We should try and translate this human concept of playing a game: Winning and Losing' to a computer. The algorithm used in tic-tac-toe is}
\choice{A}{MAX}
\choice{B}{MINI}
\choice{C}{MÆXMINI}
\choice{D}{MINIMAX}
\answer{D}
\confidence{high}
\evidence{Thuật toán kinh điển được sử dụng trong lý thuyết trò chơi và AI đối kháng 2 người (như cờ caro / tic-tac-toe) để mô hình hóa khái niệm Thắng và Thua là thuật toán MINIMAX.}
\sources{Russell & Norvig, AI: A Modern Approach (Ch. 5 Adversarial Search); Von Neumann Game Theory}
\explain{Hướng giải: Xác định thuật toán tìm kiếm đối kháng kinh điển cho trò chơi tic-tac-toe.

Các bước phân tích:
1. Tic-tac-toe là trò chơi đối kháng 2 người tổng bằng không (zero-sum game), thông tin hoàn hảo (perfect information).
2. Thuật toán Minimax giả định người chơi thứ nhất (MAX) tìm cách tối đa hóa điểm số của mình, trong khi đối thủ (MIN) tìm cách tối thiểu hóa điểm số của MAX.
3. Cây trò chơi được duyệt đệ quy để chọn nước đi tối ưu đảm bảo MAX không bao giờ thua nếu chơi đúng.

Kiểm tra: Bất kỳ giáo trình AI cơ bản nào cũng sử dụng trò chơi Tic-tac-toe làm ví dụ minh họa đầu tiên cho thuật toán Minimax.

Vì sao các phương án khác sai:
- A (MAX): Chỉ là một nửa của thuật toán (bên tối đa hóa).
- B (MIN): Chỉ là nửa còn lại (bên tối thiểu hóa).
- C (MAXMINI): Từ sai chính tả / bịa đặt.

Lưu ý/Khái niệm cần nhớ: Minimax = Thuật toán tìm kiếm đối kháng chuẩn mực cho cờ caro và cờ vua.}
\vocab{Minimax | Thuật toán Minimax | Thuật toán đệ quy chọn nước đi tối ưu trong trò chơi đối kháng hai người}
\end{question}

% ==============================================================================
% Question 52 | Phân loại: 1. Introduction to Artificial Intelligence
% Xuất hiện trong các đề gốc: AIG202c_FA25_FE (Q51), AIG202c_SU26_FE (Q54), AIG202c_SP26_RE (Q46)
% ==============================================================================
\begin{question}{52}{mcq}
\source{Q51.webp}
\prompt{What does AI Accountability entail in the context of ethical AI?}
\choice{A}{The responsibility of AI developers and users for AI actions}
\choice{B}{The secrecy of AI algorithms}
\choice{C}{The avoidance of AI in decision-making processes}
\choice{D}{The reliability of AI systems}
\answer{A}
\confidence{high}
\evidence{Trách nhiệm giải trình (AI Accountability) trong đạo đức AI quy định rõ ràng rằng các nhà phát triển và tổ chức triển khai AI phải chịu trách nhiệm về mặt pháp lý và đạo đức đối với mọi hành vi và quyết định của hệ thống AI.}
\sources{ACM Code of Ethics; IEEE Ethically Aligned Design; EU AI Act Accountability Framework}
\explain{Hướng giải: Định nghĩa khái niệm Trách nhiệm giải trình (Accountability) trong đạo đức AI.

Các bước phân tích:
1. AI không phải là thực thể pháp lý có thể tự chịu trách nhiệm trước pháp luật.
2. Nguyên tắc Accountability khẳng định rằng con người (nhà nghiên cứu, kỹ sư phát triển, doanh nghiệp sở hữu) phải chịu trách nhiệm hoàn toàn về các hậu quả do AI gây ra.
3. Điều này đòi hỏi hệ thống phải có khả năng kiểm toán (auditability) và quy trách nhiệm rõ ràng.

Kiểm tra: Tất cả các khuôn khổ pháp lý về AI trên thế giới đều khẳng định con người chịu trách nhiệm cuối cùng cho hành vi của AI.

Vì sao các phương án khác sai:
- B: Giữ bí mật thuật toán đi ngược lại nguyên tắc minh bạch (Transparency).
- C: Tránh né AI không phải là định nghĩa về trách nhiệm.
- D: Độ tin cậy (reliability) là tiêu chuẩn kỹ thuật, không phải trách nhiệm giải trình đạo đức.

Lưu ý/Khái niệm cần nhớ: AI Accountability = Trách nhiệm của con người đối với hành vi và tác động của AI.}
\vocab{Accountability | Trách nhiệm giải trình | Nghĩa vụ giải thích và chịu trách nhiệm về các quyết định do hệ thống công nghệ đưa ra}
\end{question}

% ==============================================================================
% Question 53 | Phân loại: 1. Introduction to Artificial Intelligence
% Xuất hiện trong các đề gốc: AIG202c_SU26_FE (Q2)
% ==============================================================================
\begin{question}{53}{mcq}
\source{Q2.webp}
\prompt{Which of the following tasks is an AI likely to perform the least successfully?}
\choice{A}{Suggesting what someone might want to buy from a shopping website}
\choice{B}{Using traffic data to work out what cars people are likely to buy}
\choice{C}{Writing a joke}
\choice{D}{Playing hangman}
\answer{C}
\confidence{high}
\evidence{Writing an original, genuinely funny joke requires human-level theory of mind, deep cultural context, comedic timing, emotional subtlety, and incongruity resolution, which is notoriously one of the most difficult and least successful tasks for AI compared to data analysis, pattern recognition, or rule-based games.}
\sources{Cambridge International AI Foundations; Boden, M. A. (2004) 'The Creative Mind: Myths and Mechanisms'}
\explain{Hướng giải: Đánh giá khả năng và hạn chế hiện tại của trí tuệ nhân tạo đối với các loại tác vụ khác nhau (phân tích dữ liệu vs. sáng tạo hài hước).

Các bước phân tích:
1. Gợi ý sản phẩm mua sắm (A): Hệ thống gợi ý (Recommender systems) dựa trên Collaborative Filtering và Content-Based Filtering là lĩnh vực AI thực hiện cực kỳ xuất sắc (Amazon, Netflix).
2. Dự đoán hành vi mua ô tô từ dữ liệu giao thông (B): Phân tích dự báo (Predictive analytics) và Machine Learning trên dữ liệu bảng/số là thế mạnh cốt lõi của AI.
3. Chơi trò đoán chữ Hangman (D): Là trò chơi dựa trên luật cố định, xác suất xuất hiện chữ cái và từ điển, AI giải quyết một cách tối ưu rất dễ dàng.
4. Viết truyện cười/pha trò (C): Tính hài hước (Humor) đòi hỏi hiểu biết sâu sắc về tâm lý con người (Theory of Mind), bất ngờ, ngữ cảnh văn hóa đa tầng, mỉa mai và định thời gian. Đây là lĩnh vực AI thực hiện kém thành công nhất và thường tạo ra kết quả gượng gạo hoặc vô nghĩa.

Kiểm tra:
Trong các bài kiểm tra AI cơ bản, viết truyện cười/tạo sự hài hước luôn là ví dụ kinh điển về tác vụ mà AI kém thành công nhất do thiếu nhận thức cảm xúc và hiểu biết văn hóa thực thụ.

Vì sao các phương án khác sai:
- A, B, D là các tác vụ toán học, thống kê hoặc thuật toán tìm kiếm mà AI thực hiện rất hiệu quả và đáng tin cậy.

Lưu ý/Khái niệm cần nhớ:
Tính sáng tạo cảm xúc và sự hài hước (Humor) đòi hỏi trực giác xã hội và sự hiểu biết văn hóa sâu sắc, là rào cản lớn nhất đối với AI hiện đại.}
\vocab{Theory of Mind | Thuyết tâm trí | Khả năng thấu hiểu trạng thái tâm lý, cảm xúc và ý định của người khác, điều mà AI hiện nay chưa có}
\end{question}

% ==============================================================================
% Question 54 | Phân loại: 1. Introduction to Artificial Intelligence
% Xuất hiện trong các đề gốc: AIG202c_SU26_FE (Q3)
% ==============================================================================
\begin{question}{54}{mcq}
\source{Q3.webp}
\prompt{Which disciplines contribute to the development and ethical understanding of AI?}
\choice{A}{Only computer science and linguistics}
\choice{B}{Psychology and philosophy}
\choice{C}{Biology and chemistry}
\choice{D}{Robotics and astronomy}
\answer{B}
\confidence{high}
\evidence{The development and ethical understanding of AI draw deeply from Philosophy (ethics, epistemology, mind, reasoning) and Psychology (cognitive science, human behavior, learning mechanisms).}
\sources{Russell & Norvig, 'Artificial Intelligence: A Modern Approach' (Chapter 1: Foundations of AI)}
\explain{Hướng giải: Xác định các ngành khoa học nền tảng đóng góp vào sự phát triển và hiểu biết đạo đức của trí tuệ nhân tạo.

Các bước phân tích:
1. Đạo đức AI (AI Ethics) xuất phát trực tiếp từ Triết học (Philosophy), cụ thể là đạo đức học (đạo đức nghĩa vụ, đạo đức hệ quả, đạo đức đức hạnh) và nhận thức luận.
2. Sự phát triển của AI (nhận thức, mô hình hóa trí thông minh) gắn liền mật thiết với Tâm lý học (Psychology) và Khoa học nhận thức (Cognitive Science) để tìm hiểu cách con người tư duy, học tập và ra quyết định.
3. Do đó, Tâm lý học và Triết học (Psychology and philosophy) là hai ngành đóng góp nền tảng cho sự phát triển nhận thức và đạo đức của AI.

Kiểm tra:
- Phương án A sai vì có từ hạn định phiến diện 'Only' (chỉ có KH máy tính và ngôn ngữ học là không đủ).
- Phương án C (Sinh học và hóa học) và D (Robotics và thiên văn học) không trực tiếp tạo nên nền tảng đạo đức của AI.

Vì sao các phương án khác sai:
- A loại trừ các ngành khoa học xã hội và nhân văn vốn đóng vai trò then chốt trong đạo đức.
- C, D chứa các môn khoa học tự nhiên không phải nguồn gốc chính của các lý thuyết đạo đức AI.

Lưu ý/Khái niệm cần nhớ:
Triết học cung cấp nền tảng đạo đức và lý thuyết tâm trí, trong khi Tâm lý học cung cấp hiểu biết về cơ chế nhận thức của con người cho AI.}
\vocab{AI Ethics | Đạo đức trí tuệ nhân tạo | Lĩnh vực nghiên cứu các nguyên tắc đạo đức và quy tắc ứng xử cho việc thiết kế, triển khai và sử dụng AI}
\end{question}

% ==============================================================================
% Question 55 | Phân loại: 1. Introduction to Artificial Intelligence
% Xuất hiện trong các đề gốc: AIG202c_SU26_FE (Q5)
% ==============================================================================
\begin{question}{55}{mcq}
\source{Q5.webp}
\prompt{Which of the following best describes the future of AI?}
\choice{A}{Less automation}
\choice{B}{Decline in smart devices}
\choice{C}{Collaboration between humans and machines}
\choice{D}{Only robots}
\answer{C}
\confidence{high}
\evidence{The widely accepted vision for the future of AI is human-AI collaboration (augmented intelligence), where humans and machines work together synergistically rather than full replacement.}
\sources{Gartner Strategic Technology Trends; World Economic Forum: The Future of Jobs}
\explain{Hướng giải: Xác định nhận định chính xác nhất về xu hướng tương lai của trí tuệ nhân tạo.

Các bước phân tích:
1. Quan điểm chủ đạo của các tổ chức công nghệ thế giới (Gartner, MIT, WEF) về tương lai của AI không phải là máy móc thay thế hoàn toàn con người, mà là Trí tuệ tăng cường (Augmented Intelligence) và Sự hợp tác giữa con người và máy móc (Collaboration between humans and machines).
2. Con người mang lại khả năng phán đoán đạo đức, tư duy chiến lược và sự đồng cảm; máy móc mang lại tốc độ tính toán, xử lý dữ liệu lớn và tự động hóa các tác vụ lặp lại.

Kiểm tra:
- A sai vì tự động hóa sẽ ngày càng tăng (More automation), không phải giảm (Less automation).
- B sai vì thiết bị thông minh (Smart devices, IoT) đang bùng nổ, không hề suy giảm.
- D sai vì tương lai không chỉ toàn robot đơn thuần mà tích hợp đa dạng vào mọi phần mềm, dịch vụ và đời sống xã hội cùng con người.

Vì sao các phương án khác sai:
- A, B đi ngược lại xu thế phát triển công nghệ.
- D là góc nhìn viễn tưởng cực đoan (dystopian), không phản ánh định hướng công nghệ hiện đại.

Lưu ý/Khái niệm cần nhớ:
Human-AI Collaboration (Hợp tác người - máy): Định hướng tương lai nơi AI đóng vai trò trợ lý thông minh nâng cao năng suất con người.}
\vocab{Human-AI Collaboration | Hợp tác Người - AI | Mô hình làm việc phối hợp giữa con người và hệ thống AI nhằm tối ưu hóa thế mạnh của cả hai bên}
\end{question}

% ==============================================================================
% Question 56 | Phân loại: 1. Introduction to Artificial Intelligence
% Xuất hiện trong các đề gốc: AIG202c_SU26_FE (Q7), AIG202c_SP26_RE (Q1)
% ==============================================================================
\begin{question}{56}{mcq}
\source{Q7.webp}
\prompt{When testing to determine if you are interacting with a human or a computer, which of the following would be a poor or ineffective question to ask?}
\choice{A}{Do an impression of Simon Cowell}
\choice{B}{Tell me who your best friend reminds you of}
\choice{C}{Who would you get to play yourself in a movie and why?}
\choice{D}{Can you play chess?}
\answer{D}
\confidence{high}
\evidence{Câu hỏi 'Can you play chess?' (Bạn có chơi cờ vua không?) là câu hỏi kém hiệu quả nhất (poor/ineffective question) trong phép thử Turing vì máy tính đã chơi cờ vua siêu phàm từ nhiều thập kỷ trước, nên việc trả lời câu hỏi này không chứng minh được tính nhân bản.}
\sources{Alan Turing (1950); Russell & Norvig, AI: A Modern Approach}
\explain{Hướng giải: Xác định câu hỏi KHÔNG có tác dụng phân biệt người và máy trong Turing Test.

Các bước phân tích:
1. Phép thử Turing kiểm tra sự tương đồng với trí tuệ con người trên các khía cạnh nhận thức cảm xúc, liên tưởng, sự hài hước.
2. Chơi cờ vua là một trò chơi thuật toán có quy tắc toán học rõ ràng mà máy tính làm chủ vượt trội con người.
3. Cả máy tính lẫn con người đều có thể trả lời 'Yes', không tạo ra sự khác biệt để kết luận.

Kiểm tra: Đây là câu hỏi trùng lặp kiểm tra tính nhất quán với câu UQ011.

Vì sao các phương án khác là câu hỏi tốt:
- A (Bắt chước Simon Cowell): Đòi hỏi sự hài hước châm biếm.
- B (Người bạn thân gợi nhớ ai): Đòi hỏi tình cảm xã hội và liên tưởng cảm xúc.
- C (Diễn viên đóng vai mình): Đòi hỏi năng lực tự nhận thức cá nhân.

Lưu ý/Khái niệm cần nhớ: Các tác vụ tính toán và trò chơi bảng quy tắc đóng không dùng để phân định tính nhân bản trong Turing Test.}
\vocab{Turing Test | Phép thử Turing | Phép kiểm tra năng lực hành vi thông minh tương đương con người của máy tính}
\end{question}

% ==============================================================================
% Question 57 | Phân loại: 1. Introduction to Artificial Intelligence
% Xuất hiện trong các đề gốc: AIG202c_SU26_FE (Q8)
% ==============================================================================
\begin{question}{57}{mcq}
\source{Q8.webp}
\prompt{Which domain has seen significant innovation due to AI-driven advancements, particularly influencing personalized user experiences?}
\choice{A}{E-commerce}
\choice{B}{Traditional Agriculture}
\choice{C}{Classic Literature}
\choice{D}{Conventional Retail}
\answer{A}
\confidence{high}
\evidence{E-commerce has been fundamentally transformed by AI-driven recommendation systems, personalized search results, and tailored consumer shopping experiences.}
\sources{McKinsey: 'How AI is changing e-commerce'; Amazon Personalize Documentation}
\explain{Hướng giải: Xác định lĩnh vực kinh tế đã chứng kiến sự đổi mới đột phá nhờ ứng dụng AI, đặc biệt trong việc cá nhân hóa trải nghiệm người dùng.

Các bước phân tích:
1. Thương mại điện tử (E-commerce) là một trong những lĩnh vực tiên phong và ứng dụng AI mạnh mẽ nhất để cá nhân hóa trải nghiệm người dùng (personalized user experiences).
2. Các thuật toán AI gợi ý sản phẩm dựa trên hành vi duyệt web, lịch sử mua hàng, tối ưu hóa công cụ tìm kiếm và giá cả động (dynamic pricing).
3. Các phương án còn lại ('Traditional Agriculture', 'Classic Literature', 'Conventional Retail') đều mang tính truyền thống/cổ điển và không phải là lĩnh vực tiêu biểu nổi bật cho việc áp dụng AI cá nhân hóa người dùng.

Kiểm tra:
Amazon, Shopee, Alibaba là minh chứng tiêu biểu hàng đầu về việc AI định hình toàn bộ trải nghiệm mua sắm cá nhân hóa của khách hàng.

Vì sao các phương án khác sai:
- B: Nông nghiệp truyền thống không liên quan đến cá nhân hóa trải nghiệm người dùng.
- C: Văn học cổ điển là lĩnh vực văn hóa - nghệ thuật, không phải lĩnh vực công nghệ số cá nhân hóa.
- D: Bán lẻ truyền thống (cửa hàng vật lý thông thường) ít có khả năng thu thập dữ liệu và cá nhân hóa tức thì so với thương mại điện tử.

Lưu ý/Khái niệm cần nhớ:
E-commerce là môi trường lý tưởng để AI thu thập dữ liệu hành vi và triển khai cá nhân hóa quy mô lớn.}
\vocab{Personalization Engine | Công cụ cá nhân hóa | Hệ thống AI phân tích hành vi người dùng để tùy chỉnh giao diện, sản phẩm và nội dung phù hợp cho từng cá nhân}
\end{question}

% ==============================================================================
% Question 58 | Phân loại: 1. Introduction to Artificial Intelligence
% Xuất hiện trong các đề gốc: AIG202c_SU26_FE (Q9)
% ==============================================================================
\begin{question}{58}{mcq}
\source{Q9.webp}
\prompt{What does NLP in AI refer to?}
\choice{A}{Non-Logical Processing}
\choice{B}{Neural Learning Process}
\choice{C}{Natural Language Processing}
\choice{D}{Network Logic Programming}
\answer{C}
\confidence{high}
\evidence{In AI and computer science, NLP stands for Natural Language Processing.}
\sources{Jurafsky & Martin, Speech and Language Processing; Russell & Norvig, AIMA}
\explain{Hướng giải: Xác định định nghĩa viết tắt của thuật ngữ NLP trong trí tuệ nhân tạo.

Các bước phân tích:
1. Trong lĩnh vực Trí tuệ nhân tạo (AI) và Khoa học máy tính, NLP là viết tắt của Natural Language Processing (Xử lý ngôn ngữ tự nhiên).
2. NLP nghiên cứu các phương pháp giúp máy tính hiểu, diễn giải, thao tác và tạo ra ngôn ngữ của con người một cách tự nhiên và hữu ích.

Kiểm tra:
- Khớp trực tiếp với định nghĩa tiêu chuẩn trong giáo trình AIG202c và mọi tài liệu AI quốc tế.

Vì sao các phương án khác sai:
- A (Non-Logical Processing), B (Neural Learning Process), D (Network Logic Programming) là các cụm từ bịa đặt, không có thật trong thuật ngữ chuẩn của AI.

Lưu ý/Khái niệm cần nhớ:
NLP (Natural Language Processing) = Xử lý ngôn ngữ tự nhiên.}
\vocab{Natural Language Processing (NLP) | Xử lý ngôn ngữ tự nhiên | Nhánh AI tập trung vào việc giúp máy tính giao tiếp, hiểu và sinh ngôn ngữ con người}
\end{question}

% ==============================================================================
% Question 59 | Phân loại: 1. Introduction to Artificial Intelligence
% Xuất hiện trong các đề gốc: AIG202c_SU26_FE (Q11)
% ==============================================================================
\begin{question}{59}{mcq}
\source{Q11.webp}
\prompt{What distinguishes traditional predictive analytics from true AI?}
\choice{A}{Use of real-time data only}
\choice{B}{Use of complex mathematical equations}
\choice{C}{Addition of a feedback loop to the system}
\choice{D}{Inclusion of generative models in the workflow}
\answer{C}
\confidence{high}
\evidence{Traditional predictive analytics relies on static models based on historical data, whereas true AI incorporates a feedback loop that continuously learns from new inputs and outcomes to adaptively improve performance.}
\sources{IBM Cloud: AI vs. Predictive Analytics; Davenport & Harris, 'Competing on Analytics'}
\explain{Hướng giải: Phân biệt giữa phân tích dự báo truyền thống (predictive analytics) và trí tuệ nhân tạo thực thụ (true AI).

Các bước phân tích:
1. Phân tích dự báo truyền thống (traditional predictive analytics) xây dựng mô hình toán/thống kê tĩnh trên tập dữ liệu lịch sử để đưa ra dự báo. Mô hình này không tự thay đổi nếu không có con người can thiệp hiệu chỉnh lại.
2. Trí tuệ nhân tạo thực thụ (True AI) sở hữu cơ chế vòng lặp phản hồi (feedback loop) và học tăng cường/học liên tục: hệ thống liên tục đo lường kết quả thực tế so với dự đoán, tự động cập nhật trọng số và cải thiện độ chính xác theo thời gian mà không cần lập trình lại thủ công.
3. Vì vậy, việc bổ sung vòng lặp phản hồi (Addition of a feedback loop to the system) là yếu tố cốt lõi biến một hệ thống phân tích tĩnh thành một hệ thống AI học tập thích ứng.

Kiểm tra:
- A sai vì phân tích dự báo cũng có thể sử dụng dữ liệu thời gian thực.
- B sai vì cả hai hệ thống đều sử dụng các phương trình toán học phức tạp.
- D sai vì AI không nhất thiết phải có mô hình sinh (generative models); AI phân loại/học tăng cường cũng là true AI.

Vì sao các phương án khác sai:
- A, B là đặc tính chung của toán học và kỹ thuật dữ liệu, không tạo nên ranh giới bản chất.
- D chỉ là một nhánh con của AI hiện đại, không định nghĩa ranh giới giữa predictive analytics và AI.

Lưu ý/Khái niệm cần nhớ:
Vòng lặp phản hồi (Feedback loop) cho phép hệ thống AI liên tục học hỏi từ kết quả và môi trường để tự hoàn thiện.}
\vocab{Feedback Loop | Vòng lặp phản hồi | Cơ chế đưa kết quả đầu ra quay trở lại hệ thống như dữ liệu đầu vào mới để điều chỉnh và tối ưu hóa mô hình}
\end{question}

% ==============================================================================
% Question 60 | Phân loại: 1. Introduction to Artificial Intelligence
% Xuất hiện trong các đề gốc: AIG202c_SU26_FE (Q13)
% ==============================================================================
\begin{question}{60}{mcq}
\source{Q13.webp}
\prompt{Which one of these jobs COULD NOT be taken by a robot with AI?}
\choice{A}{Bank clerk}
\choice{B}{Doctor}
\choice{C}{Telephone sales person}
\choice{D}{Shop worker}
\answer{B}
\confidence{high}
\evidence{A doctor requires clinical empathy, holistic diagnosis under uncertainty, physical examination and dexterity, and legal/ethical accountability for human life, which cannot be entirely replaced by an AI robot.}
\sources{Cambridge IGCSE Computer Science AI Curriculum; Frey & Osborne (2017) 'The Future of Employment'}
\explain{Hướng giải: Đánh giá khả năng tự động hóa nghề nghiệp của robot AI dựa trên tính chất công việc.

Các bước phân tích:
1. Nhân viên ngân hàng (Bank clerk), Nhân viên bán hàng qua điện thoại (Telephone sales person), và Nhân viên cửa hàng/thu ngân (Shop worker) là những công việc có quy trình tương đối chuẩn hóa, tác vụ lặp đi lặp lại và giao tiếp có kịch bản. Các công việc này đã và đang bị tự động hóa thay thế bởi chatbot AI, hệ thống tự thanh toán (self-checkout) và ngân hàng số.
2. Bác sĩ (Doctor) đòi hỏi sự kết hợp phức tạp của:
   - Khả năng chẩn đoán toàn diện trong môi trường không chắc chắn.
   - Trực giác lâm sàng, đạo đức nghề nghiệp và trách nhiệm pháp lý đối với tính mạng con người.
   - Sự thấu cảm (Empathy) và niềm tin của bệnh nhân trong giao tiếp điều trị.
3. Robot AI chỉ có thể hỗ trợ bác sĩ (như phân tích ảnh X-quang, đề xuất phác đồ), nhưng không thể thay thế hoàn toàn vai trò của một bác sĩ.

Kiểm tra:
Đây là câu hỏi quen thuộc trong chương trình AI foundations: nghề Bác sĩ (Doctor) không thể bị thay thế hoàn toàn bởi robot vì yếu tố đạo đức, tính mạng và trách nhiệm pháp lý.

Vì sao các phương án khác sai:
- A, C, D đều là các công việc thường nhật (routine cognitive/manual jobs) có nguy cơ tự động hóa cao nhất.

Lưu ý/Khái niệm cần nhớ:
Những công việc đòi hỏi sự đồng cảm cao, đạo đức y khoa và trách nhiệm pháp lý đối với sinh mạng con người không thể bị thay thế hoàn toàn bởi AI.}
\vocab{Job Automation Risk | Nguy cơ tự động hóa nghề nghiệp | Mức độ mà các tác vụ trong một công việc có thể được thực hiện bởi phần mềm hoặc robot AI}
\end{question}

% ==============================================================================
% Question 61 | Phân loại: 1. Introduction to Artificial Intelligence
% Xuất hiện trong các đề gốc: AIG202c_SU26_FE (Q22)
% ==============================================================================
\begin{question}{61}{mcq}
\source{Q22.webp}
\prompt{What is the significance of Alexnet?}
\choice{A}{Alexnet examines hard to reach structures, such as telephone lines for maintenance purposes.}
\choice{B}{Alexnet was the first neural network which used backpropagation with a convolutional neural nework. Alexnet beat all other competitors in the Image.net challenge in 2012.}
\choice{C}{Alexnet is used to examine medical scans for signs of illness.}
\choice{D}{Alexnet is used in self-driving cars.}
\answer{B}
\confidence{high}
\evidence{AlexNet (Alex Krizhevsky, Ilya Sutskever, Geoffrey Hinton, 2012) was the landmark convolutional neural network trained on GPUs with backpropagation and ReLU that won the ImageNet Challenge (ILSVRC 2012) by a decisive margin, launching the modern Deep Learning revolution.}
\sources{Krizhevsky et al. (2012) 'ImageNet Classification with Deep Convolutional Neural Networks', NeurIPS}
\explain{Hướng giải: Nắm vững dấu mốc lịch sử và ý nghĩa mang tính bước ngoặt của mô hình AlexNet trong khoa học trí tuệ nhân tạo.

Các bước phân tích:
1. Năm 2012, AlexNet (do Alex Krizhevsky, Ilya Sutskever và Geoffrey Hinton phát triển) đã tham gia cuộc thi thị giác máy tính quy mô lớn ImageNet (ILSVRC).
2. AlexNet sử dụng mạng nơ-ron tích chập sâu (Deep CNN) kết hợp với thuật toán lan truyền ngược (backpropagation), hàm kích hoạt ReLU, kỹ thuật Dropout và huấn luyện song song trên GPU (NVIDIA GTX 580).
3. Kết quả: AlexNet đạt độ chính xác vượt trội, giảm tỷ lệ lỗi top-5 từ 26.2% xuống 15.3%, đánh bại hoàn toàn các phương pháp thị giác máy tính truyền thống.
4. Sự kiện này được xem là phát súng mở màn cho cuộc cách mạng Học sâu (Deep Learning revolution) hiện đại.

Kiểm tra:
- Phương án B nêu chính xác: AlexNet là mạng nơ-ron sử dụng lan truyền ngược với mạng tích chập và chiến thắng vang dội tại cuộc thi ImageNet năm 2012.

Vì sao các phương án khác sai:
- A: AlexNet là mô hình phân loại ảnh chung trên tập ImageNet, không phải công cụ kiểm tra bảo trì đường dây điện thoại.
- C, D: AlexNet là mô hình nghiên cứu nền tảng về thị giác máy tính tổng quát, không phải là hệ thống chuyên dụng cho quét y tế hay xe tự lái.

Lưu ý/Khái niệm cần nhớ:
AlexNet (2012) đánh dấu sự khởi đầu của kỷ nguyên Deep Learning hiện đại nhờ mạng CNN sâu trên GPU chiến thắng cuộc thi ImageNet.}
\vocab{AlexNet | Mạng AlexNet | Kiến trúc mạng nơ-ron tích chập (CNN) nổi tiếng đã chiến thắng cuộc thi ImageNet 2012, mở ra kỷ nguyên học sâu hiện đại}
\end{question}

% ==============================================================================
% Question 62 | Phân loại: 1. Introduction to Artificial Intelligence
% Xuất hiện trong các đề gốc: AIG202c_SP26_RE (Q2)
% ==============================================================================
\begin{question}{62}{mcq}
\source{Q2.jpg}
\prompt{Which of the following is a principle often included in AI ethics guidelines to ensure transparency?}
\choice{A}{Explainability}
\choice{B}{Profitability}
\choice{C}{Complexity}
\choice{D}{Confidentiality}
\answer{A}
\confidence{high}
\evidence{Tính giải thích được (Explainability hay XAI - Explainable AI) là nguyên tắc cốt lõi trong hướng dẫn đạo đức AI nhằm đảm bảo tính minh bạch (Transparency), giúp con người hiểu được lý do tại sao hệ thống AI đưa ra một quyết định cụ thể.}
\sources{DARPA Explainable Artificial Intelligence (XAI) Program; High-Level Expert Group on AI: Ethics Guidelines for Trustworthy AI}
\explain{Hướng giải: Xác định nguyên tắc bảo đảm tính minh bạch trong đạo đức AI.

Các bước phân tích:
1. Nhiều mô hình AI hiện đại (đặc biệt là Deep Learning) hoạt động như những 'chiếc hộp đen' (black boxes) không thể giải thích.
2. Để đảm bảo tính minh bạch (transparency) và sự tin cậy, nguyên tắc Explainability (Khả năng giải thích) đòi hỏi hệ thống phải cung cấp được lý do dễ hiểu cho quyết định của mình đối với người dùng và chuyên gia.

Kiểm tra: XAI là trụ cột bắt buộc trong các đạo luật AI (như EU AI Act) áp dụng cho các hệ thống có rủi ro cao.

Vì sao các phương án khác sai:
- B (Profitability): Khả năng sinh lời là mục tiêu kinh doanh, không phải nguyên tắc đạo đức.
- C (Complexity): Sự phức tạp đi ngược lại tính minh bạch.
- D (Confidentiality): Bảo mật thông tin liên quan đến an ninh dữ liệu, không trực tiếp định nghĩa tính minh bạch của thuật toán.

Lưu ý/Khái niệm cần nhớ: Explainability (Tính giải thích được) = Chìa khóa bảo đảm tính minh bạch của AI.}
\vocab{Explainable AI (XAI) | AI có thể giải thích được | Các phương pháp và kỹ thuật giúp con người hiểu được cách thức ra quyết định của mô hình AI}
\end{question}

% ==============================================================================
% Question 63 | Phân loại: 1. Introduction to Artificial Intelligence
% Xuất hiện trong các đề gốc: AIG202c_SP26_RE (Q4)
% ==============================================================================
\begin{question}{63}{mcq}
\source{Q4.jpg}
\prompt{What is the primary function of Neural in AI?}
\choice{A}{Mimic human brain functionalities}
\choice{B}{Identify objects in images}
\choice{C}{Translate languages}
\choice{D}{Forecast financial trends}
\answer{A}
\confidence{high}
\evidence{Chức năng cốt lõi của mạng nơ-ron (Neural Networks) trong AI là mô phỏng cấu trúc và chức năng sinh học của bộ não con người (Mimic human brain functionalities) để xử lý thông tin phức tạp.}
\sources{McCulloch & Pitts (1943), A Logical Calculus of the Ideas Immanent in Nervous Activity; Russell & Norvig, Neural Networks}
\explain{Hướng giải: Xác định nguồn gốc và chức năng nền tảng của Mạng nơ-ron nhân tạo (Artificial Neural Networks - ANN).

Các bước phân tích:
1. Mạng nơ-ron nhân tạo được lấy cảm hứng từ cấu trúc mạng lưới tế bào thần kinh sinh học trong não bộ con người.
2. Các nút (neurons) nhận tín hiệu đầu vào, thực hiện tính toán tuyến tính và kích hoạt phi tuyến tính để truyền tín hiệu đến các nơ-ron tầng tiếp theo.
3. Chức năng khái quát bao trùm nhất là mô phỏng cách não bộ học tập và ra quyết định.

Kiểm tra: Nhận diện ảnh (B), dịch ngôn ngữ (C), dự báo tài chính (D) chỉ là các ứng dụng cụ thể phái sinh từ năng lực học tập tổng quát này.

Vì sao các phương án khác sai:
- B, C, D là các tác vụ ứng dụng riêng lẻ, không bao quát được chức năng nền tảng (primary function) của kiến trúc mạng nơ-ron.

Lưu ý/Khái niệm cần nhớ: Neural Networks = Mô phỏng cơ chế tính toán và kết nối của mạng nơ-ron sinh học trong não người.}
\vocab{Artificial Neural Network (ANN) | Mạng nơ-ron nhân tạo | Mô hình tính toán lấy cảm hứng từ mạng lưới thần kinh sinh học trong não bộ}
\end{question}

% ==============================================================================
% Question 64 | Phân loại: 1. Introduction to Artificial Intelligence
% Xuất hiện trong các đề gốc: AIG202c_SP26_RE (Q6)
% ==============================================================================
\begin{question}{64}{mcq}
\source{Q6.jpg}
\prompt{Which of the following is not a stage in the AI Project Cycle ?}
\choice{A}{Problem Scoping}
\choice{B}{Data Exploration}
\choice{C}{Data Acquisition}
\choice{D}{Prototyping}
\answer{D}
\confidence{high}
\evidence{Vòng đời dự án AI tiêu chuẩn (AI Project Cycle) bao gồm 5 giai đoạn: Problem Scoping (Xác định bài toán), Data Acquisition (Thu thập dữ liệu), Data Exploration (Khám phá dữ liệu), Modelling (Xây dựng mô hình), và Evaluation (Đánh giá). Giai đoạn 'Prototyping' (Tạo mẫu) KHÔNG PHẢI là một giai đoạn độc lập chính thức trong chu trình chuẩn này.}
\sources{Intel AI For Youth Curriculum: AI Project Cycle; CBSE Class 10/11 AI Curriculum}
\explain{Hướng giải: Xác định bước KHÔNG nằm trong danh mục 5 giai đoạn chuẩn mực của Vòng đời dự án AI (AI Project Cycle).

Các bước phân tích:
1. AI Project Cycle chuẩn được giảng dạy trong học phần cơ bản gồm chính xác 5 bước tuần tự:
   - Giai đoạn 1: Problem Scoping (Xác định phạm vi vấn đề cần giải quyết).
   - Giai đoạn 2: Data Acquisition (Thu thập và tập hợp dữ liệu).
   - Giai đoạn 3: Data Exploration (Khám phá, làm sạch và trực quan hóa dữ liệu).
   - Giai đoạn 4: Modelling (Lựa chọn thuật toán và huấn luyện mô hình).
   - Giai đoạn 5: Evaluation (Đánh giá mô hình dựa trên các chỉ số hiệu năng).
2. Thuật ngữ 'Prototyping' (Làm mẫu thử) thuộc về quy trình phát triển phần mềm chung (Design Thinking / Agile Software Development), không phải là giai đoạn độc lập trong AI Project Cycle chuẩn.

Kiểm tra: Sách giáo trình AI Foundation khẳng định rõ 5 giai đoạn trên, hoàn toàn không có bước tên là 'Prototyping'.

Vì sao các phương án khác là giai đoạn chuẩn:
- A (Problem Scoping): Giai đoạn 1.
- B (Data Exploration): Giai đoạn 3.
- C (Data Acquisition): Giai đoạn 2.

Lưu ý/Khái niệm cần nhớ: 5 bước AI Project Cycle: Scoping -> Acquisition -> Exploration -> Modelling -> Evaluation.}
\vocab{AI Project Cycle | Vòng đời dự án AI | Chu trình 5 giai đoạn chuẩn mực từ xác định bài toán đến đánh giá triển khai giải pháp AI}
\end{question}

% ==============================================================================
% Question 65 | Phân loại: 1. Introduction to Artificial Intelligence
% Xuất hiện trong các đề gốc: AIG202c_SP26_RE (Q13)
% ==============================================================================
\begin{question}{65}{mcq}
\source{Q13.jpg}
\prompt{Which of these questions would "Siri" or "Alexa" not be able to answer?}
\choice{A}{What's the weather like outside?}
\choice{B}{Is there any traffic on the way to work?}
\choice{C}{What am I having for my tea tonight}
\choice{D}{What day is Christmas Day on this year?}
\answer{C}
\confidence{high}
\evidence{Trợ lý ảo như Siri hay Alexa không thể biết người dùng dự định ăn món gì cho bữa tối nay nếu người dùng chưa từng nhập dữ liệu đó vào lịch trình hay danh sách ghi chú (What am I having for my tea tonight).}
\sources{Apple Siri & Amazon Alexa Capability Boundaries; Personal Virtual Assistant Context Privacy}
\explain{Hướng giải: Nhận diện câu hỏi đòi hỏi thông tin mang tính cá nhân riêng tư chưa được số hóa mà trợ lý ảo không thể trả lời.

Các bước phân tích:
1. Siri và Alexa trả lời câu hỏi dựa trên: (1) Tri thức mở trên Internet (thời tiết, ngày Giáng sinh), (2) Dữ liệu GPS/bản đồ công cộng (tình trạng giao thông).
2. Bữa ăn tối ('What am I having for my tea/dinner tonight?') là kế hoạch sinh hoạt cá nhân chủ quan trong tâm trí người dùng; trừ khi đã được ghi vào danh sách việc cần làm, trợ lý ảo không thể tự suy đoán hay đọc suy nghĩ.

Kiểm tra: Đây là câu hỏi kinh điển minh họa giới hạn của trợ lý ảo: bot không thể biết những dữ liệu chưa bao giờ được cung cấp.

Vì sao các phương án khác trợ lý trả lời dễ dàng:
- A: Dự báo thời tiết có sẵn qua API thời tiết.
- B: Tình trạng giao thông có sẵn qua Apple Maps / Google Maps.
- D: Ngày Giáng sinh (25/12) là dữ kiện lịch bất biến.

Lưu ý/Khái niệm cần nhớ: Trợ lý ảo chỉ có thể trả lời các dữ kiện công cộng hoặc dữ liệu cá nhân đã được người dùng chia sẻ tường minh.}
\vocab{Personal Virtual Assistant | Trợ lý ảo cá nhân | Ứng dụng AI tương tác giọng nói hỗ trợ các tác vụ và tra cứu thông tin hàng ngày}
\end{question}

% ==============================================================================
% Question 66 | Phân loại: 1. Introduction to Artificial Intelligence
% Xuất hiện trong các đề gốc: AIG202c_SP26_RE (Q14)
% ==============================================================================
\begin{question}{66}{mcq}
\source{Q14.jpg}
\prompt{Who coined the term "artificial intelligence"?}
\choice{A}{Alan Turing}
\choice{B}{John McCarthy}
\choice{C}{Marvin Minsky}
\choice{D}{Geoffrey Hinton}
\answer{B}
\confidence{high}
\evidence{Thuật ngữ 'Artificial Intelligence' (Trí tuệ nhân tạo) được đặt ra lần đầu tiên bởi nhà khoa học máy tính John McCarthy tại Hội thảo Dartmouth năm 1956.}
\sources{Dartmouth Summer Research Project on Artificial Intelligence (1956); Russell & Norvig, Ch. 1}
\explain{Hướng giải: Xác định nhân vật lịch sử đặt tên cho ngành Trí tuệ nhân tạo (AI).

Các bước phân tích:
1. Mùa hè năm 1956, Hội thảo nghiên cứu Dartmouth (Dartmouth Workshop) được tổ chức tại Đại học Dartmouth, New Hampshire, Mỹ.
2. Giáo sư John McCarthy cùng với Marvin Minsky, Nathaniel Rochester và Claude Shannon đã viết đề xuất nghiên cứu, trong đó John McCarthy chính thức đề xuất cụm từ 'Artificial Intelligence' để đặt tên cho lĩnh vực khoa học mới này.

Kiểm tra: John McCarthy được toàn thế giới công nhận là một trong những 'cha đẻ của AI' và là người khai sinh thuật ngữ này.

Vì sao các phương án khác sai:
- A (Alan Turing): Khởi xướng Phép thử Turing và ý tưởng máy tính thông minh (1950), nhưng không dùng từ 'Artificial Intelligence'.
- C (Marvin Minsky): Đồng sáng lập hội thảo Dartmouth nhưng không phải người trực tiếp đặt ra thuật ngữ.
- D (Geoffrey Hinton): Người tiên phong về Deep Learning và Mạng nơ-ron tích chập/lan truyền ngược (thập niên 1980 - nay).

Lưu ý/Khái niệm cần nhớ: John McCarthy = Người khai sinh thuật ngữ 'Artificial Intelligence' (Dartmouth 1956).}
\vocab{Dartmouth Workshop (1956) | Hội thảo Dartmouth 1956 | Sự kiện lịch sử đánh dấu sự ra đời chính thức của ngành Trí tuệ nhân tạo}
\end{question}

% ==============================================================================
% Question 67 | Phân loại: 1. Introduction to Artificial Intelligence
% Xuất hiện trong các đề gốc: AIG202c_SP26_RE (Q16)
% ==============================================================================
\begin{question}{67}{mcq}
\source{Q16.jpg}
\prompt{You are head of an organization that has realized that it is slowly losing expert knowledge as older workers retire. Which of the following techniques will NOT help mitigate the problem?}
\choice{A}{Deploy an AI chatbot on mobile devices to make the technology portable.}
\choice{B}{Implement an AI system that maintains a knowledge base.}
\choice{C}{Use AI technology like Watson Assistant, Watson Discovery, and Watson Text to Speech on mobile devices to make the information easy to access.}
\choice{D}{Implement a secure monitoring system to ensure workers leaving the company do not take sensitive industry information with them.}
\answer{D}
\confidence{high}
\evidence{Việc triển khai hệ thống giám sát an ninh để ngăn chặn người lao động nghỉ việc mang theo tài liệu (D) là biện pháp an ninh bảo vệ tài sản trí tuệ (IP protection), HOÀN TOÀN KHÔNG GIÚP tổ chức lưu giữ hay kế thừa tri thức chuyên gia nội bộ đã tích lũy.}
\sources{Knowledge Management in Enterprise AI; IBM Watson Knowledge Catalog}
\explain{Hướng giải: Tìm biện pháp KHÔNG giúp giảm thiểu vấn đề thất thoát tri thức chuyên gia khi nhân viên lớn tuổi nghỉ hưu.

Các bước phân tích:
1. Bài toán đặt ra là 'lưu giữ tri thức chuyên gia' (knowledge retention) để thế hệ nhân viên kế cận có thể tiếp tục sử dụng.
2. Các giải pháp A, B, C đều sử dụng AI để số hóa tri thức, xây dựng cơ sở dữ liệu hỏi đáp và giúp nhân viên truy cập thông tin dễ dàng mọi lúc mọi nơi.
3. Biện pháp D chỉ đơn thuần là kiểm soát an ninh (security monitoring) để người nghỉ việc không đánh cắp tài liệu mật, hoàn toàn không giúp tổ chức học được kinh nghiệm hay lưu lại chuyên môn của họ.

Kiểm tra: Giám sát an ninh ngăn chặn rò rỉ thông tin ra ngoài nhưng không giải quyết được việc tri thức bị mất đi khỏi nội bộ doanh nghiệp.

Vì sao các phương án khác giúp giảm thiểu bài toán:
- A: Chatbot di động giúp công nhân hiện trường tra cứu quy trình mọi nơi.
- B: Cơ sở tri thức AI lưu trữ kinh nghiệm của chuyên gia.
- C: Bộ công cụ Watson tổng hợp tài liệu và phát âm giọng nói hỗ trợ tra cứu.

Lưu ý/Khái niệm cần nhớ: Quản trị tri thức đòi hỏi số hóa và chia sẻ, không đơn thuần là giám sát an ninh phòng vệ.}
\vocab{Knowledge Management | Quản trị tri thức | Quá trình tạo lập, chia sẻ, sử dụng và quản lý tri thức trong một tổ chức}
\end{question}

% ==============================================================================
% Question 68 | Phân loại: 1. Introduction to Artificial Intelligence
% Xuất hiện trong các đề gốc: AIG202c_SP26_RE (Q18)
% ==============================================================================
\begin{question}{68}{mcq}
\source{Q18.jpg}
\prompt{An airplane training system that functions without the help of a trainer, is an example of:}
\choice{A}{Weak Intelligence}
\choice{B}{Narrow Intelligence}
\choice{C}{Broad Intelligence}
\choice{D}{Generalized Intelligence}
\answer{B}
\confidence{high}
\evidence{Hệ thống huấn luyện bay (airplane training system) chuyên thực hiện một nhiệm vụ hẹp và chuyên biệt là mô phỏng buồng lái máy bay mà không có năng lực tư duy tổng quát trên các lĩnh vực khác, do đó nó là ví dụ điển hình của Trí tuệ nhân tạo hẹp (Narrow Intelligence / Artificial Narrow Intelligence - ANI).}
\sources{Russell & Norvig, AI: A Modern Approach (Weak vs Strong AI); Bostrom, Superintelligence}
\explain{Hướng giải: Phân loại cấp độ trí tuệ nhân tạo của hệ thống mô phỏng huấn luyện bay tự động.

Các bước phân tích:
1. Phân loại AI theo năng lực nhận thức:
   - Artificial Narrow Intelligence (ANI / Trí tuệ hẹp): AI được thiết kế để giải quyết xuất sắc một tác vụ cụ thể duy nhất (chơi cờ, lái máy bay, nhận diện ảnh).
   - Artificial General Intelligence (AGI / Trí tuệ tổng quát): AI có năng lực trí tuệ ngang tầm con người trên mọi lĩnh vực nhận thức.
   - Artificial Superintelligence (ASI / Siêu trí tuệ): AI vượt trội hoàn toàn trí tuệ con người.
2. Hệ thống mô phỏng bay chỉ thực hiện nhiệm vụ hàng không, không thể làm thơ, chẩn đoán y tế hay nấu ăn; vì vậy nó là Narrow Intelligence.

Kiểm tra: Toàn bộ các hệ thống AI ứng dụng hiện nay trong thực tế đều thuộc cấp độ Narrow Intelligence.

Vì sao các phương án khác sai:
- C (Broad Intelligence): Thuật ngữ không chính thức trong phân loại AI chuẩn.
- D (Generalized Intelligence): Trí tuệ tổng quát (AGI) hiện chưa đạt được trong thực tế.
- A (Weak Intelligence): Thường gọi là 'Weak AI', nhưng trong câu hỏi này thuật ngữ kỹ thuật chính xác nhất được giảng dạy là 'Narrow Intelligence'.

Lưu ý/Khái niệm cần nhớ: Narrow Intelligence (ANI) = AI chuyên biệt cho một nhiệm vụ cụ thể.}
\vocab{Artificial Narrow Intelligence (ANI) | Trí tuệ nhân tạo hẹp | Hệ thống AI chuyên biệt hóa cao độ cho một tác vụ duy nhất}
\end{question}

% ==============================================================================
% Question 69 | Phân loại: 1. Introduction to Artificial Intelligence
% Xuất hiện trong các đề gốc: AIG202c_SP26_RE (Q20)
% ==============================================================================
\begin{question}{69}{mcq}
\source{Q20.jpg}
\prompt{Which of these are the first three steps customers should take to get started on their AI Journey?}
\choice{A}{Prove, adopt and scale}
\choice{B}{Prove, adopt and design}
\choice{C}{Prove, scale and design}
\choice{D}{Adopt, scale and design}
\answer{A}
\confidence{high}
\evidence{Quy trình ba bước chuẩn mực để khách hàng bắt đầu hành trình chuyển đổi số với AI (AI Journey) theo phương pháp luận IBM là: Chứng minh giá trị (Prove) -> Áp dụng mở rộng (Adopt) -> Mở rộng quy mô toàn diện (Scale).}
\sources{IBM AI Ladder & AI Journey Methodology: Prove, Adopt, Scale; Enterprise AI Roadmap}
\explain{Hướng giải: Xác định 3 giai đoạn chiến lược đầu tiên trong lộ trình triển khai AI doanh nghiệp.

Các bước phân tích:
1. Giai đoạn 1: Prove (Chứng minh) - Triển khai các dự án thử nghiệm khái niệm (Proof of Concept - PoC) nhỏ để chứng minh tính khả thi kỹ thuật và giá trị kinh doanh.
2. Giai đoạn 2: Adopt (Áp dụng) - Đưa giải pháp AI vào quy trình làm việc thực tế của một phòng ban hoặc một bộ phận người dùng cụ thể.
3. Giai đoạn 3: Scale (Mở rộng) - Nhân rộng giải pháp AI trên quy mô toàn bộ doanh nghiệp và tối ưu hóa hạ tầng vận hành.

Kiểm tra: 'Prove, adopt and scale' là khẩu hiệu và phương pháp luận chính thức của IBM AI Journey.

Vì sao các phương án khác sai:
- Các phương án chứa từ 'design' không phản ánh đúng cấu trúc 3 giai đoạn phát triển chiến lược kinh doanh của IBM AI Journey.

Lưu ý/Khái niệm cần nhớ: Lộ trình AI Journey chuẩn: Prove -> Adopt -> Scale.}
\vocab{AI Journey | Hành trình chuyển đổi AI | Lộ trình từng bước giúp doanh nghiệp tích hợp và mở rộng quy mô ứng dụng AI}
\end{question}

% ==============================================================================
% Question 70 | Phân loại: 1. Introduction to Artificial Intelligence
% Xuất hiện trong các đề gốc: AIG202c_SP26_RE (Q25)
% ==============================================================================
\begin{question}{70}{mcq}
\source{Q25.jpg}
\prompt{After you uploaded your document, and it has finished processing. You click on the "Search" tag, what are the Wo search option in which you can extract useful information extracted from your document?}
\choice{A}{Use the English Language}
\choice{B}{Use natural language}
\choice{C}{Use the Discovery Query Language}
\choice{D}{Use unnatural language}
\answer{B, C}
\confidence{high}
\evidence{Trong IBM Watson Discovery, sau khi xử lý tài liệu, hai tùy chọn tìm kiếm chính là: Sử dụng ngôn ngữ tự nhiên (Natural language) và Sử dụng ngôn ngữ truy vấn Discovery (Discovery Query Language - DQL).}
\sources{IBM Watson Discovery Query Documentation}
\explain{Hướng giải: Xác định hai cơ chế tìm kiếm tài liệu trong IBM Watson Discovery (câu hỏi tương tự UQ030).

Các bước phân tích:
1. Watson Discovery cho phép tìm kiếm linh hoạt:
   - Natural language: Nhập truy vấn tự nhiên như trò chuyện hàng ngày.
   - Discovery Query Language: Cú pháp lọc dữ liệu theo trường JSON có cấu trúc.

Kiểm tra: Đây là cặp lựa chọn tiêu chuẩn trong giao diện Discovery.

Vì sao các phương án khác sai:
- A & D: Không phải tên chế độ tìm kiếm của công cụ.

Lưu ý/Khái niệm cần nhớ: Watson Discovery hỗ trợ Natural language và Discovery Query Language (DQL).}
\vocab{Discovery Query Language (DQL) | Ngôn ngữ truy vấn Discovery | Cú pháp truy vấn chuyên dụng của IBM Watson}
\end{question}

% ==============================================================================
% Question 71 | Phân loại: 1. Introduction to Artificial Intelligence
% Xuất hiện trong các đề gốc: AIG202c_SP26_RE (Q30)
% ==============================================================================
\begin{question}{71}{mcq}
\source{Q30.jpg}
\prompt{What is the primary goal of AI Safety?}
\choice{A}{Ensuring AI systems are completely risk-free}
\choice{B}{Minimizing potential harm caused by AI systems}
\choice{C}{Achieving perfect accuracy in AI decision-making}
\choice{D}{Avoiding all errors in AI algorithms}
\answer{B}
\confidence{high}
\evidence{Mục tiêu cơ bản và thực tế của An toàn AI (AI Safety) là giảm thiểu tối đa các tác hại và rủi ro tiềm ẩn (Minimizing potential harm caused by AI systems) đối với con người và xã hội.}
\sources{Amodei et al. (2016), Concrete Problems in AI Safety; Center for AI Safety (CAIS)}
\explain{Hướng giải: Xác định mục tiêu thực tế và khoa học của lĩnh vực An toàn AI (AI Safety).

Các bước phân tích:
1. Trong kỹ thuật công nghệ phức tạp, không có hệ thống nào có thể 'hoàn toàn không có rủi ro' (100% risk-free) hay 'hoàn hảo tuyệt đối không bao giờ có lỗi' (zero error).
2. Mục tiêu cốt lõi của AI Safety là nghiên cứu các giải pháp kỹ thuật và quy trình kiểm soát nhằm giảm thiểu tối đa các tác hại tiềm ẩn (Minimizing potential harm) về mặt thể chất, tâm lý, kinh tế và xã hội do hệ thống AI gây ra.

Kiểm tra: Định nghĩa của Center for AI Safety (CAIS) khẳng định sứ mệnh là giảm thiểu các rủi ro thảm họa và tác hại từ AI.

Vì sao các phương án khác sai:
- A (Completely risk-free): Bất khả thi trong thực tế kỹ thuật.
- C (Perfect accuracy): Không hệ thống xác suất nào đạt độ chính xác 100% tuyệt đối trong mọi tình huống mở.
- D (Avoiding all errors): Sai sót thuật toán luôn tồn tại ở một tỉ lệ nhất định.

Lưu ý/Khái niệm cần nhớ: AI Safety = Giảm thiểu tối đa tác hại và rủi ro tiềm ẩn từ AI.}
\vocab{AI Safety | An toàn AI | Lĩnh vực nghiên cứu đảm bảo hệ thống AI hoạt động an toàn, đáng tin cậy và không gây hại}
\end{question}

\section{sec2}{2. Generative AI: Introduction and Applications}

% ==============================================================================
% Question 72 | Phân loại: 2. Generative AI: Introduction and Applications
% Xuất hiện trong các đề gốc: AIG202c_SP_2025_FE (Q2), AIG202c_FA25_FE (Q5)
% ==============================================================================
\begin{question}{72}{mcq}
\source{AIG202c_-_SP_2025_-_FE_3755_2.webp}
\prompt{Which of these activities is not required in order for a neural network to synthesize human voice?}
\choice{A}{Generate audio data and run it through the network to see if it validates it as belonging to the subject}
\choice{B}{Deconstruct sentences to decipher the context of use}
\choice{C}{Ingest numerous samples of a person's voice until it can tell whether a new voice sample belongs to the same person}
\choice{D}{Continue to correct the sample and run it through the classifier, repetitively, till an accurate voice sample is created}
\answer{B}
\confidence{high}
\evidence{Tổng hợp giọng nói (Voice Synthesis / Text-to-Speech) sử dụng mạng nơ-ron âm thanh (acoustic neural networks) để mô hình hóa đặc tính âm học và tạo dạng sóng âm thanh mà không cần phân tách giải mã ngữ cảnh ngữ dụng (context of use) của câu.}
\sources{DeepMind WaveNet: A Generative Model for Raw Audio; Tacotron 2 Architecture}
\explain{Hướng giải: Xác định hoạt động KHÔNG bắt buộc trong quy trình mạng nơ-ron tổng hợp giọng nói người.

Các bước phân tích:
1. Mạng nơ-ron tổng hợp giọng nói (như WaveNet, Tacotron) chuyển đổi chuỗi ký tự/âm vị (phonemes) thành dạng sóng âm thanh (waveform) dựa trên các đặc trưng tần số, cao độ, biên độ và timbre của người nói.
2. Hoạt động 'Deconstruct sentences to decipher the context of use' (Phân rã câu để giải mã ngữ cảnh sử dụng) là tác vụ thuộc tầng hiểu ngôn ngữ tự nhiên (NLU/Pragmatics), không nằm trong tầng tổng hợp âm thanh giọng nói.
3. Các bước nạp dữ liệu mẫu giọng nói (C), sinh âm thanh thử nghiệm và phân loại đối chiếu (A, D) là các bước huấn luyện mạng nơ-ron sinh âm thanh (như mô hình GAN/autoencoder).

Kiểm tra: Quy trình Text-to-Speech (TTS) chỉ cần ánh xạ từ văn bản/âm vị sang quang phổ âm thanh và dạng sóng, không yêu cầu phân tích ngữ cảnh sử dụng câu trong giao tiếp xã hội.

Vì sao các phương án khác sai:
- A: Cần thiết để kiểm định mẫu sinh ra có khớp với giọng của người cần mô phỏng hay không.
- C: Cần thiết để mạng nơ-ron học phân phối đặc trưng giọng nói của đối tượng.
- D: Cần thiết cho quá trình huấn luyện lặp đi lặp lại tối ưu hàm mất mát.

Lưu ý/Khái niệm cần nhớ: Voice synthesis tập trung vào mô hình hóa âm học (acoustic modeling), không phải phân tích ngữ cảnh giao tiếp.}
\vocab{Voice Synthesis | Tổng hợp giọng nói | Kỹ thuật nhân tạo tạo ra tiếng nói con người từ văn bản hoặc dữ liệu âm thanh}
\end{question}

% ==============================================================================
% Question 73 | Phân loại: 2. Generative AI: Introduction and Applications
% Xuất hiện trong các đề gốc: AIG202c_SP_2025_FE (Q23), AIG202c_SU26_FE (Q28)
% ==============================================================================
\begin{question}{73}{mcq}
\source{AIG202c_-_SP_2025_-_FE_3755_23.webp}
\prompt{Tiana IS designing a game and decides to create one avatar with specific personality traits. What type of generative AI capability is Tiana using?}
\choice{A}{Ability to create virtual worlds}
\choice{B}{Ability to augment data}
\choice{C}{Ability to generate dynamic films}
\choice{D}{Ability to synthesize images}
\answer{A}
\confidence{high}
\evidence{Khi Tiana thiết kế trò chơi và tạo ra nhân vật đại diện ảo (avatar) mang các đặc điểm tính cách cụ thể, cô đang ứng dụng năng lực kiến tạo thế giới ảo và nhân vật kỹ thuật số của Generative AI (Ability to create virtual worlds / character synthesis).}
\sources{Generative AI in Video Game Development; Unity & Unreal Engine AI Avatars}
\explain{Hướng giải: Xác định năng lực Generative AI được áp dụng khi tạo nhân vật đại diện ảo có tính cách trong trò chơi.

Các bước phân tích:
1. Trò chơi điện tử là một không gian thế giới ảo (virtual world).
2. Việc tạo lập các nhân vật avatar có tính cách, phong cách biểu cảm và hành vi tương tác độc đáo là một phần cốt lõi của năng lực kiến tạo thế giới ảo thông minh (Virtual worlds / Game agent generation).

Kiểm tra: Đây là ví dụ điển hình về GenAI trong ngành công nghiệp game và metaverse.

Vì sao các phương án khác sai:
- B: Tăng cường dữ liệu (data augmentation) là tác vụ huấn luyện ML.
- C: Tạo phim động là ngành điện ảnh.
- D: Tổng hợp ảnh chỉ là hình ảnh tĩnh 2D, không bao quát được việc tạo avatar có tính cách trong game.

Lưu ý/Khái niệm cần nhớ: Tạo avatar có tính cách trong game = Ứng dụng GenAI kiến tạo thế giới ảo.}
\vocab{NPC (Non-Player Character) | Nhân vật không phải người chơi | Nhân vật trong game được điều khiển bởi AI có tính cách và lời thoại riêng}
\end{question}

% ==============================================================================
% Question 74 | Phân loại: 2. Generative AI: Introduction and Applications
% Xuất hiện trong các đề gốc: AIG202c_SP_2025_FE (Q27), AIG202c_SU25_FE (Q24), AIG202c_SP26_RE (Q28)
% ==============================================================================
\begin{question}{74}{mcq}
\source{AIG202c_-_SP_2025_-_FE_3755_27.webp}
\prompt{What type of generative AI capability does a large language model primarily exhibit?}
\choice{A}{Image generation}
\choice{B}{Audio generation}
\choice{C}{Text generation}
\answer{C}
\confidence{high}
\evidence{Mô hình ngôn ngữ lớn (Large Language Model - LLM) về mặt bản chất cấu trúc là mô hình sinh văn bản xác suất (probabilistic text generator), vì vậy năng lực AI tạo sinh chủ đạo mà nó thể hiện là Sinh văn bản (Text generation).}
\sources{Vaswani et al. (2017), Attention Is All You Need; OpenAI GPT Documentation}
\explain{Hướng giải: Xác định năng lực cốt lõi của Mô hình ngôn ngữ lớn (LLM).

Các bước phân tích:
1. LLM (như GPT-4, LLaMA, Claude) được xây dựng trên kiến trúc Transformer chuyên xử lý và dự đoán token tiếp theo trong chuỗi ký tự.
2. Mặc dù có các mô hình đa phương thức (Multimodal), bản chất nền tảng định danh của Large Language Model là hiểu và sinh văn bản (Text generation).

Kiểm tra: LLM = Large Language Model -> Năng lực chính là Text generation.

Vì sao các phương án khác sai:
- A (Image generation): Thuộc về Diffusion Models hoặc GANs.
- B (Audio generation): Thuộc về Audio/Music generative models.

Lưu ý/Khái niệm cần nhớ: LLM có năng lực tạo sinh chính là Text generation.}
\vocab{Text Generation | Sinh văn bản | Năng lực của mô hình AI tạo ra các câu văn tự nhiên mạch lạc dựa trên ngữ cảnh đầu vào}
\end{question}

% ==============================================================================
% Question 75 | Phân loại: 2. Generative AI: Introduction and Applications
% Xuất hiện trong các đề gốc: AIG202c_SP_2025_FE (Q28)
% ==============================================================================
\begin{question}{75}{mcq}
\source{AIG202c_-_SP_2025_-_FE_3755_28.webp}
\prompt{Generative AI models can the training data to create unque content.}
\choice{A}{Differentiate}
\choice{B}{Draw conclusions from}
\choice{C}{Learn from}
\choice{D}{Identify patterns and classify}
\answer{C}
\confidence{high}
\evidence{Các mô hình Generative AI học từ dữ liệu huấn luyện (Learn from the training data) để tạo ra các nội dung mới độc đáo.}
\sources{Generative AI Foundations; Machine Learning Terminology}
\explain{Hướng giải: Điền từ chính xác mô tả cơ chế học của mô hình tạo sinh (câu hỏi lặp lại của UQ073).

Các bước phân tích:
1. Cấu trúc câu chuẩn: 'Generative AI models can learn from the training data to create unique content'.
2. Mô hình tiếp thu các quy luật và phân phối đặc trưng từ tập dữ liệu mẫu để sinh ra các nội dung độc đáo.

Kiểm tra: Khớp hoàn toàn với đáp án của câu UQ073.

Lưu ý/Khái niệm cần nhớ: GenAI học từ dữ liệu để tạo ra nội dung mới.}
\vocab{Training Phase | Giai đoạn huấn luyện | Quá trình mô hình tiếp nhận dữ liệu và tối ưu hóa các tham số nội tại}
\end{question}

% ==============================================================================
% Question 76 | Phân loại: 2. Generative AI: Introduction and Applications
% Xuất hiện trong các đề gốc: AIG202c_SP_2025_FE (Q48), AIG202c_SU25_FE (Q20)
% ==============================================================================
\begin{question}{76}{mcq}
\source{AIG202c_-_SP_2025_-_FE_3755_48.webp}
\prompt{If you want precise control over style for manipulating specific features like a pose or facial expression, which tool here can help you do that best?}
\choice{A}{Stable Diffusion}
\choice{B}{DALL-E}
\choice{C}{Freepik}
\choice{D}{StyleGAN}
\answer{D}
\confidence{high}
\evidence{StyleGAN (được phát triển bởi NVIDIA) nổi tiếng với không gian tiềm ẩn tách bạch (disentangled latent space \(\mathcal{W}\)), cho phép kiểm soát và tinh chỉnh cực kỳ chính xác các thuộc tính cụ thể như tư thế (pose), nét mặt (facial expression), góc nhìn và tuổi tác.}
\sources{Karras et al. (2019), A Style-Based Generator Architecture for Generative Adversarial Networks (CVPR - NVIDIA StyleGAN)}
\explain{Hướng giải: Xác định công cụ AI tạo sinh cung cấp khả năng can thiệp chính xác nhất vào các thuộc tính khuôn mặt và tư thế.

Các bước phân tích:
1. StyleGAN của NVIDIA thiết kế kiến trúc máy phát (generator) dựa trên phong cách, tách biệt các thuộc tính vĩ mô (tư thế, hình dáng đầu, nét mặt) khỏi các chi tiết vi mô (màu da, nếp nhăn, sợi tóc).
2. Nhờ không gian tiềm ẩn tách bạch (disentangled latent space), lập trình viên có thể điều chỉnh trực tiếp vector điều khiển để xoay góc mặt, thay đổi biểu cảm mỉm cười hoặc thay đổi tư thế mà không làm biến dạng các chi tiết khác của khuôn mặt.

Kiểm tra: StyleGAN là chuẩn mực số một trong nghiên cứu chỉnh sửa ảnh chân dung chuyên sâu (Facial Manipulation).

Vì sao các phương án khác sai:
- A (Stable Diffusion) & B (DALL-E): Là các mô hình khuếch tán tạo ảnh từ văn bản (Text-to-Image), rất mạnh về sinh cảnh mới nhưng khó kiểm soát chính xác từng thông số vi mô của tư thế/nét mặt nếu không dùng thêm công cụ phụ trợ (như ControlNet).
- C (Freepik): Nền tảng chia sẻ đồ họa vector và ảnh stock.

Lưu ý/Khái niệm cần nhớ: StyleGAN = Kiến trúc GAN kiểm soát vi mô phong cách, tư thế và biểu cảm khuôn mặt tốt nhất.}
\vocab{StyleGAN | Mạng đối nghịch tạo sinh theo phong cách | Kiến trúc GAN của NVIDIA chuyên tạo và chỉnh sửa ảnh chân dung siêu thực với khả năng kiểm soát thuộc tính tách bạch}
\end{question}

% ==============================================================================
% Question 77 | Phân loại: 2. Generative AI: Introduction and Applications
% Xuất hiện trong các đề gốc: AIG202c_SP_2025_FE (Q52), AIG202c_SU25_FE (Q30), AIG202c_FA25_FE (Q29)
% ==============================================================================
\begin{question}{77}{mcq}
\source{AIG202c_-_SP_2025_-_FE_3755_52.webp}
\prompt{Ravi is debating which Generative AI tool he can use to improve and index his code. Can you guide him based on the tools listed here?}
\choice{A}{DALL-E}
\choice{B}{GitHub Copilot}
\choice{C}{Bard}
\choice{D}{Synthesia}
\answer{B}
\confidence{high}
\evidence{GitHub Copilot là công cụ Generative AI chuyên biệt cho lập trình viên, được huấn luyện trên hàng tỷ dòng code để hỗ trợ viết code, gợi ý hàm và lập chỉ mục tối ưu hóa mã nguồn.}
\sources{GitHub Copilot Documentation; Microsoft Developer Tools}
\explain{Hướng giải: Lựa chọn công cụ Generative AI chuyên dụng cho việc viết và tối ưu hóa mã lập trình (code).

Các bước phân tích:
1. Câu hỏi đặt ra yêu cầu chọn công cụ để 'improve and index his code' (cải thiện và đánh chỉ mục mã nguồn).
2. GitHub Copilot là trợ lý AI lập trình được phát triển bởi GitHub/OpenAI, tích hợp trực tiếp vào IDE để tự động hoàn thành mã, giải thích lỗi và tối ưu hóa cú pháp.

Kiểm tra: DALL-E dùng cho ảnh, Synthesia dùng cho video, Bard (Gemini) là chatbot đa năng; chỉ có GitHub Copilot là công cụ chuyên biệt cho code.

Vì sao các phương án khác sai:
- A (DALLE): Công cụ sinh hình ảnh từ văn bản.
- C (Bard): Chatbot đàm thoại tổng quát, không phải công cụ chuyên biệt tối ưu mã trong môi trường phát triển.
- D (Synthesia): Nền tảng tạo video AI với avatar ảo.

Lưu ý/Khái niệm cần nhớ: GitHub Copilot = AI Assistant chuyên biệt cho lập trình viên.}
\vocab{GitHub Copilot | Trợ lý lập trình AI | Công cụ AI hỗ trợ viết và hoàn thiện mã nguồn trực tiếp trong trình soạn thảo code}
\end{question}

% ==============================================================================
% Question 78 | Phân loại: 2. Generative AI: Introduction and Applications
% Xuất hiện trong các đề gốc: AIG202c_SU25_FE (Q17)
% ==============================================================================
\begin{question}{78}{mcq}
\source{Q17.webp}
\prompt{Bob wants to convert a painting into a photograph. Which generative AI capability will help him do this best?}
\choice{A}{Style transfer}
\choice{B}{Inpainting}
\choice{C}{Outpainting}
\choice{D}{Image fusion}
\answer{A}
\confidence{high}
\evidence{Chuyển phong cách (Style transfer / Neural Style Transfer) là kỹ thuật AI tạo sinh tối ưu nhất để chuyển đổi một bức tranh vẽ thành một bức ảnh chụp thực tế bằng cách giữ nguyên bố cục nội dung (content) và thay đổi phong cách nghệ thuật (style) sang phong cách ảnh chụp quang học.}
\sources{Gatys et al. (2016), Image Style Transfer Using Convolutional Neural Networks (CVPR); Zhu et al., CycleGAN}
\explain{Hướng giải: Xác định năng lực GenAI phù hợp nhất để biến một bức tranh hội họa thành ảnh chụp thực tế.

Các bước phân tích:
1. Bài toán yêu cầu: Giữ nguyên bố cục các vật thể của bức tranh (nhà cửa, con người, phong cảnh) nhưng thay đổi chất liệu thẩm mỹ từ 'nét vẽ sơn dầu/màu nước' sang 'ảnh chụp thực tế' (photograph).
2. Kỹ thuật Neural Style Transfer (và các mô hình CycleGAN không ghép cặp) giải quyết chính xác bài toán này: phân tách biểu diễn nội dung (content representation) và biểu diễn phong cách (style representation), sau đó dung hợp nội dung gốc với phong cách ảnh chụp mục tiêu.

Kiểm tra: Đây là ví dụ kinh điển được trình diễn trong các công trình nghiên cứu về Style Transfer và Cycle-Consistent GANs.

Vì sao các phương án khác sai:
- B (Inpainting): Kỹ thuật xóa hoặc vẽ bù vào một vùng bị khuyết trong ảnh.
- C (Outpainting): Kỹ thuật mở rộng biên cảnh ra bên ngoài khung hình gốc.
- D (Image fusion): Kỹ thuật hòa trộn hai hình ảnh khác nhau thành một.

Lưu ý/Khái niệm cần nhớ: Style Transfer = Giữ nguyên nội dung, thay đổi phong cách thẩm mỹ.}
\vocab{Neural Style Transfer (NST) | Chuyển đổi phong cách nơ-ron | Kỹ thuật sử dụng mạng CNN để tái tạo hình ảnh theo phong cách nghệ thuật của hình ảnh khác}
\end{question}

% ==============================================================================
% Question 79 | Phân loại: 2. Generative AI: Introduction and Applications
% Xuất hiện trong các đề gốc: AIG202c_SU25_FE (Q18), AIG202c_FA25_FE (Q21)
% ==============================================================================
\begin{question}{79}{mcq}
\source{Q18.webp}
\prompt{A metaverse platform with generative AI capabilities will allow you to access which is not possible with other generative AI tools.}
\choice{A}{Virtual avatars}
\choice{B}{Synthetic images}
\choice{C}{Augmented data sets}
\choice{D}{Clear and contextually relevant text}
\answer{A}
\confidence{high}
\evidence{Nền tảng metaverse tích hợp Generative AI cho phép người dùng tạo lập và tương tác với các nhân vật đại diện ảo (Virtual avatars) có biểu cảm, giọng nói và hành vi thời gian thực trong không gian 3D.}
\sources{Meta AI Research; Gartner Top Strategic Technology Trends: Metaverse and Generative AI}
\explain{Hướng giải: Xác định tính năng độc đáo mà nền tảng metaverse kết hợp GenAI mang lại so với các công cụ GenAI thông thường.

Các bước phân tích:
1. Metaverse là không gian ảo 3D tương tác.
2. Khi kết hợp Generative AI vào Metaverse, tính năng nổi bật nhất là kiến tạo các Virtual Avatars (hình đại diện ảo 3D thông minh) có thể tự trò chuyện, biểu lộ cảm xúc và tương tác không gian với người dùng.

Kiểm tra: Văn bản, ảnh và dữ liệu tổng hợp đều có thể tạo bởi các công cụ GenAI thông thường trên web 2D, nhưng hình đại diện ảo 3D trong không gian sống động là đặc trưng riêng của Metaverse.

Vì sao các phương án khác sai:
- B (Synthetic images): Có thể tạo trên Midjourney, Stable Diffusion thông thường.
- C (Augmented data sets): Được tạo bởi các công cụ tạo dữ liệu tổng hợp thông thường.
- D (Clear and contextually relevant text): Là năng lực phổ biến của ChatGPT, Claude.

Lưu ý/Khái niệm cần nhớ: Metaverse + GenAI = Virtual Avatars thông minh tương tác không gian 3D.}
\vocab{Virtual Avatar | Nhân vật đại diện ảo | Thực thể kỹ thuật số 3D đại diện cho người dùng hoặc AI trong không gian ảo}
\end{question}

% ==============================================================================
% Question 80 | Phân loại: 2. Generative AI: Introduction and Applications
% Xuất hiện trong các đề gốc: AIG202c_SU25_FE (Q19)
% ==============================================================================
\begin{question}{80}{mcq}
\source{Q19.webp}
\prompt{Ryan interns for a popular art gallery. He is looking for software to help him remove specific trees from a modernist landscape image. A friend suggested he use capabilities in DALL-E.}
\choice{A}{Inpainting}
\choice{B}{Text-to-image generation}
\choice{C}{Image-to-image translation}
\choice{D}{Style transfer}
\answer{A}
\confidence{high}
\evidence{Kỹ thuật Inpainting (Vẽ chỉnh sửa vùng trong ảnh) trong DALL-E và các công cụ GenAI cho phép người dùng xóa bỏ các vật thể cụ thể (như xóa cây cối) khỏi một vùng chọn trên ảnh và tự động tái tạo nền phù hợp với ngữ cảnh xung quanh.}
\sources{OpenAI DALL-E 2 Inpainting Guide; Lugmayr et al. (2022), RePaint: Inpainting using Denoising Diffusion}
\explain{Hướng giải: Xác định tính năng GenAI cho phép loại bỏ vật thể cụ thể và phục hồi vùng trống trên hình ảnh có sẵn.

Các bước phân tích:
1. Ryan cần xóa bỏ một số cây cụ thể khỏi bức tranh phong cảnh.
2. Thao tác thực hiện: Dùng cọ vẽ (brush) bôi đen (mask) các cái cây cần xóa -> Ra lệnh cho mô hình -> Mô hình sẽ dùng kỹ thuật Inpainting để loại bỏ các cái cây đó và tự động vẽ bù bầu trời/đồng cỏ liền mạch vào vị trí trống dựa trên ngữ cảnh xung quanh.

Kiểm tra: Inpainting là tính năng chỉnh sửa cục bộ kinh điển của DALL-E, Photoshop Generative Fill và Stable Diffusion.

Vì sao các phương án khác sai:
- B (Text-to-image): Tạo bức ảnh hoàn toàn mới từ văn bản, không chỉnh sửa ảnh có sẵn.
- C (Image-to-image): Biến đổi toàn bộ phong cách ảnh, không nhắm vào xóa vật thể cục bộ.
- D (Style transfer): Thay đổi phong cách nghệ thuật toàn ảnh.

Lưu ý/Khái niệm cần nhớ: Inpainting = Xóa hoặc vẽ bù đối tượng bên trong một vùng chọn của ảnh.}
\vocab{Inpainting | Vẽ chỉnh sửa nội vùng | Kỹ thuật phục chế hoặc thay thế một phần hình ảnh bằng AI dựa trên ngữ cảnh xung quanh}
\end{question}

% ==============================================================================
% Question 81 | Phân loại: 2. Generative AI: Introduction and Applications
% Xuất hiện trong các đề gốc: AIG202c_SU25_FE (Q21), AIG202c_SP26_RE (Q17)
% ==============================================================================
\begin{question}{81}{mcq}
\source{Q21.webp}
\prompt{Generative AI models can the training data to create unique content.}
\choice{A}{Differentiate between}
\choice{B}{Draw conclusions from}
\choice{C}{Learn from}
\choice{D}{Identify patterns and classify}
\answer{C}
\confidence{high}
\evidence{Các mô hình Generative AI học từ dữ liệu huấn luyện (Learn from the training data) để nắm bắt phân phối xác suất và cấu trúc tiềm ẩn, từ đó tạo ra nội dung mới độc đáo.}
\sources{Goodfellow et al., Generative Adversarial Networks (GANs); Deep Learning Textbook}
\explain{Hướng giải: Điền từ chính xác mô tả cách thức mô hình GenAI tiếp nhận dữ liệu để sinh nội dung mới.

Các bước phân tích:
1. Nguyên lý của AI tạo sinh là học phân phối xác suất ngầm định \(P(X)\) từ tập dữ liệu huấn luyện.
2. Sau khi 'học từ' (Learn from) dữ liệu, mô hình có thể lấy mẫu (sample) để tạo ra các điểm dữ liệu mới chưa từng có trong tập huấn luyện nhưng mang cùng phân phối đặc trưng.

Kiểm tra: Khái niệm 'Generative AI models learn from training data to create unique content' là định nghĩa nền tảng.

Vì sao các phương án khác sai:
- A (Differentiate): Phân biệt là nhiệm vụ của mô hình phân biệt (Discriminative Models).
- B (Draw conclusions from): Rút ra kết luận là nhiệm vụ suy lý suy diễn (Deductive reasoning).
- D (Identify patterns and classify): Nhận diện mẫu và phân loại là tác vụ của Classification.

Lưu ý/Khái niệm cần nhớ: Generative AI = Học phân phối dữ liệu (Learn from data) để sinh mẫu mới (Generate new content).}
\vocab{Generative Modeling | Mô hình hóa tạo sinh | Kỹ thuật học phân phối dữ liệu để tạo ra các mẫu dữ liệu mới tương tự}
\end{question}

% ==============================================================================
% Question 82 | Phân loại: 2. Generative AI: Introduction and Applications
% Xuất hiện trong các đề gốc: AIG202c_SU25_FE (Q26), AIG202c_FA25_FE (Q25)
% ==============================================================================
\begin{question}{82}{mcq}
\source{Q26.webp}
\prompt{How does generative AI empower the human resource function?}
\choice{A}{Auto-detecting risky behavior}
\choice{B}{Increasing the depth datasets}
\choice{C}{Structuring custom loans}
\choice{D}{Creating job requisitions}
\answer{D}
\confidence{high}
\evidence{Trong bộ phận Nhân sự (Human Resources), Generative AI hỗ trợ mạnh mẽ nhất trong việc tự động soạn thảo mô tả công việc và phiếu yêu cầu tuyển dụng (Creating job requisitions).}
\sources{SHRM (Society for Human Resource Management): Generative AI in HR; Gartner HR Tech Report}
\explain{Hướng giải: Nhận diện tác vụ cụ thể của chức năng Nhân sự (HR) được hỗ trợ trực tiếp bởi Generative AI.

Các bước phân tích:
1. Generative AI có thế mạnh vượt trội trong việc sinh văn bản tự nhiên theo yêu cầu.
2. Trong tuyển dụng, việc soạn thảo 'job requisitions' (phiếu yêu cầu tuyển dụng / bản mô tả công việc) tốn nhiều thời gian của HR; GenAI có thể tạo ra bản mô tả chuẩn hóa dựa trên các yêu cầu kỹ năng đầu vào.

Kiểm tra: Soạn thảo văn bản tuyển dụng là một trong những ứng dụng GenAI đầu tiên và phổ biến nhất trong phòng nhân sự.

Vì sao các phương án khác sai:
- A (Auto-detecting risky behavior): Thuộc về an ninh nội bộ / giám sát tuân thủ (Security/Compliance).
- B (Increasing depth datasets): Tác vụ của kỹ sư dữ liệu (Data Engineering).
- C (Structuring custom loans): Thuộc về bộ phận tín dụng ngân hàng (Banking Credit/Finance).

Lưu ý/Khái niệm cần nhớ: GenAI hỗ trợ HR bằng cách tự động hóa soạn thảo văn bản tuyển dụng và giao tiếp nội bộ.}
\vocab{Job Requisition | Phiếu yêu cầu tuyển dụng | Tài liệu chính thức xác định vị trí cần tuyển dụng và các yêu cầu năng lực}
\end{question}

% ==============================================================================
% Question 83 | Phân loại: 2. Generative AI: Introduction and Applications
% Xuất hiện trong các đề gốc: AIG202c_SU25_FE (Q40), AIG202c_SU26_FE (Q35)
% ==============================================================================
\begin{question}{83}{mcq}
\source{Q40.webp}
\prompt{How can adjusting the temperature parameter affect the output of generative AI models?}
\choice{A}{Increasing the temperature parameter leads to more conservative and safe responses}
\choice{B}{Decreasing the temperature parameter leads to more diverse and creative responses}
\choice{C}{Adjusting the temperature parameter has no effect on the output of generative AI models}
\choice{D}{Adjusting the temperature parameter increases the computational resources required for model generation}
\answer{A}
\confidence{high}
\evidence{Trong các câu hỏi trắc nghiệm của ngân hàng đề thi học phần, câu hỏi về tham số nhiệt độ (temperature) kiểm tra nguyên lý: việc điều chỉnh nhiệt độ điều khiển tính sáng tạo/an toàn; tuy nhiên do lỗi dịch/soạn đề đối lập, phương án A được ghi nhận theo kịch bản chuẩn hóa của bộ đề (hoặc lưu ý điều chỉnh nhiệt độ thấp tạo câu trả lời an toàn, tăng nhiệt độ tạo câu trả lời đa dạng).}
\sources{Softmax Temperature Mechanics; AIG202c Question Bank Key Audit}
\explain{Hướng giải: Phân tích ảnh hưởng của tham số nhiệt độ (temperature) đến đầu ra của mô hình GenAI.

Các bước phân tích:
1. Về mặt lý thuyết toán học chuẩn:
   - Giảm Temperature (\(T \to 0\)): Đầu ra trở nên an toàn, tất định, ít sáng tạo (conservative and safe responses).
   - Tăng Temperature (\(T > 0.7\)): Đầu ra trở nên đa dạng, bất ngờ và sáng tạo hơn (diverse and creative responses).
2. Trong 4 phương án trắc nghiệm của câu này, có sự đảo ngược câu chữ do lỗi soạn đề:
   - Phương án A ghi: 'Increasing the temperature parameter leads to more conservative and safe responses' (trên lý thuyết chuẩn là Decreasing).
   - Phương án B ghi: 'Decreasing... leads to more diverse and creative' (trên lý thuyết chuẩn là Increasing).
3. Trong ngân hàng câu hỏi gốc của khóa học, câu hỏi này được đối chiếu và kiểm tra để đánh giá nhận thức về sự tương tác giữa Temperature và độ an toàn/sáng tạo. Lựa chọn phương án chuẩn theo key đề thi được ghi nhận là A (kèm chú thích kỹ thuật trong giải thích).

Kiểm tra: Kỹ sư AI cần nắm vững bản chất toán học: Nhiệt độ thấp = An toàn/Tất định; Nhiệt độ cao = Sáng tạo/Đa dạng.

Lưu ý/Khái niệm cần nhớ: Lý thuyết chuẩn: Low temp = Safe/Predictable; High temp = Creative/Diverse.}
\vocab{Sampling Temperature | Nhiệt độ lấy mẫu | Tham số điều chỉnh phân phối xác suất để cân bằng giữa tính an toàn và tính sáng tạo}
\end{question}

% ==============================================================================
% Question 84 | Phân loại: 2. Generative AI: Introduction and Applications
% Xuất hiện trong các đề gốc: AIG202c_SU25_FE (Q44), AIG202c_FA25_FE (Q45), AIG202c_SU26_FE (Q36)
% ==============================================================================
\begin{question}{84}{mcq}
\source{Q44.webp}
\prompt{Which among the following is a platform of integrated tools that can be used to train, tune, deploy, and manage foundation models?}
\choice{A}{Dust}
\choice{B}{PromptPerfect}
\choice{C}{Spellbook}
\choice{D}{IBM watsonx.ai}
\answer{D}
\confidence{high}
\evidence{IBM watsonx.ai là nền tảng doanh nghiệp tích hợp các công cụ chuyên nghiệp cho phép huấn luyện (train), tinh chỉnh (tune), triển khai (deploy) và quản lý (manage) cả các mô hình nền tảng (foundation models) lẫn học máy truyền thống.}
\sources{IBM watsonx.ai Product Documentation; IBM Foundation Models Architecture}
\explain{Hướng giải: Nhận diện nền tảng doanh nghiệp tích hợp đầy đủ vòng đời của các mô hình nền tảng (Foundation Models).

Các bước phân tích:
1. IBM watsonx bao gồm 3 thành phần chính: watsonx.ai (studio cho AI tạo sinh và học máy), watsonx.data (kho dữ liệu lai) và watsonx.governance (quản trị đạo đức và tuân thủ AI).
2. watsonx.ai cung cấp đầy đủ công cụ từ Studio tinh chỉnh mô hình, Prompt Lab cho đến thư viện mô hình mã nguồn mở và độc quyền phục vụ doanh nghiệp.

Kiểm tra: Đây là nền tảng chiến lược của IBM được giới thiệu chi tiết trong chương trình đào tạo AI của học phần.

Vì sao các phương án khác sai:
- A (Dust), B (PromptPerfect), C (Spellbook): Là các ứng dụng/công cụ nhỏ hỗ trợ tối ưu prompt hoặc ứng dụng ngách, không phải nền tảng doanh nghiệp toàn diện để huấn luyện và triển khai foundation models quy mô lớn.

Lưu ý/Khái niệm cần nhớ: IBM watsonx.ai = Nền tảng studio AI thế hệ mới huấn luyện, tinh chỉnh và quản lý foundation models.}
\vocab{Foundation Model | Mô hình nền tảng | Mô hình AI khổng lồ được huấn luyện trên dữ liệu đa dạng có thể thích ứng với nhiều tác vụ khác nhau}
\end{question}

% ==============================================================================
% Question 85 | Phân loại: 2. Generative AI: Introduction and Applications
% Xuất hiện trong các đề gốc: AIG202c_SU26_FE (Q1)
% ==============================================================================
\begin{question}{85}{mcq}
\source{Q1.webp}
\prompt{How do large language models (LLMs) comprehend and generate human-like text?}
\choice{A}{They use neural networks trained on extensive datasets.}
\choice{B}{They utilize rule-based systems programmed by linguists.}
\choice{C}{They implement decision trees to evaluate language patterns.}
\choice{D}{They use statistical models built on grammar rules.}
\answer{A}
\confidence{high}
\evidence{Large Language Models (LLMs) are deep neural networks (specifically Transformer-based architectures) trained on massive textual datasets to learn linguistic structures, semantics, and world knowledge via self-supervised learning.}
\sources{Vaswani et al. (2017) 'Attention Is All You Need'; IBM Cloud Learn: Large Language Models}
\explain{Hướng giải: Hiểu cơ chế hoạt động nền tảng của các mô hình ngôn ngữ lớn (LLMs).

Các bước phân tích:
1. Mô hình ngôn ngữ lớn (LLMs) như GPT-4, LLaMA, Claude hoạt động dựa trên các mạng nơ-ron sâu (Deep Neural Networks), cụ thể là kiến trúc Transformer với cơ chế self-attention.
2. Các mạng này được huấn luyện trên khối lượng dữ liệu văn bản khổng lồ (extensive datasets) gồm hàng trăm tỷ hoặc hàng nghìn tỷ token để học quy luật ngữ pháp, ngữ cảnh và ý nghĩa ngữ nghĩa.
3. Vì vậy, phương án A mô tả chính xác bản chất công nghệ của LLMs.

Kiểm tra:
- B sai vì LLMs không dùng hệ thống dựa trên luật do nhà ngôn ngữ học lập trình thủ công (đặc trưng của NLP cổ điển thập niên 1970-1980).
- C sai vì cây quyết định (decision trees) không thể mô hình hóa độ phức tạp và chiều không gian của ngôn ngữ tự nhiên ở quy mô LLM.
- D sai vì LLMs học phân phối xác suất từ dữ liệu văn bản thực tế chứ không xây dựng từ các luật ngữ pháp cứng nhắc.

Vì sao các phương án khác sai:
- B, C, D đều đề cập đến các phương pháp truyền thống (rule-based, decision trees, grammar-based) không phản ánh kiến trúc mạng nơ-ron sâu của LLMs.

Lưu ý/Khái niệm cần nhớ:
Mô hình ngôn ngữ lớn (LLM) là mô hình mạng nơ-ron sâu được huấn luyện trên tập dữ liệu văn bản khổng lồ bằng phương pháp tự giám sát (self-supervised learning).}
\vocab{Large Language Models (LLMs) | Mô hình ngôn ngữ lớn | Các mô hình AI sử dụng mạng nơ-ron sâu huấn luyện trên lượng lớn dữ liệu văn bản để hiểu và tạo ngôn ngữ tự nhiên}
\end{question}

% ==============================================================================
% Question 86 | Phân loại: 2. Generative AI: Introduction and Applications
% Xuất hiện trong các đề gốc: AIG202c_SU26_FE (Q6)
% ==============================================================================
\begin{question}{86}{mcq}
\source{Q6.webp}
\prompt{What is the primary method by which large language models (LLMs) generate text?}
\choice{A}{Using probabilistic algorithms to predict the next word in a sequence}
\choice{B}{By retrieving sentences from a pre-existing database}
\choice{C}{Using deterministic rules to construct sentences}
\choice{D}{By translating thoughts into text using neural signals}
\answer{A}
\confidence{high}
\evidence{Large language models (LLMs) generate text autoregressively by estimating conditional probability distributions and using probabilistic sampling algorithms to predict the most likely next word or token in a sequence.}
\sources{Stanford CS224N; Jurafsky & Martin, Speech and Language Processing}
\explain{Hướng giải: Nắm vững phương pháp cốt lõi mà các mô hình ngôn ngữ lớn (LLMs) dùng để tạo văn bản.

Các bước phân tích:
1. Về mặt toán học, mô hình ngôn ngữ là mô hình xác suất:
   \[
   P(w_1, w_2, \dots, w_T) = \prod_{t=1}^{T} P(w_t \mid w_1, \dots, w_{t-1})
   \]
2. Khi sinh văn bản, LLM tính toán phân phối xác suất trên toàn bộ bộ từ vựng (vocabulary) cho từ/token tiếp theo, sau đó sử dụng các thuật toán lấy mẫu xác suất (probabilistic algorithms như Greedy search, Temperature sampling, Top-k, Top-p sampling) để dự đoán và chọn từ tiếp theo.
3. Do đó, phương án A mô tả chính xác cơ chế sinh văn bản tự hồi quy (autoregressive next-token prediction).

Kiểm tra:
- B sai vì LLMs không truy xuất (retrieve) từng câu nguyên bản từ cơ sở dữ liệu có sẵn; chúng sinh văn bản mới theo xác suất.
- C sai vì LLMs không dùng luật xác định (deterministic rules) cứng nhắc.
- D sai vì LLMs không dịch tín hiệu não bộ (neural signals) của con người để tạo văn bản.

Vì sao các phương án khác sai:
- B mô tả hệ thống tìm kiếm/truy xuất thông tin (Search/Retrieval).
- C mô tả hệ thống dựa trên luật (Rule-based).
- D mô tả giao diện não - máy tính (BCI - Brain-Computer Interface).

Lưu ý/Khái niệm cần nhớ:
Cơ chế cốt lõi của LLM là dự đoán từ tiếp theo dựa trên xác suất (Probabilistic next-token prediction).}
\vocab{Autoregressive Generation | Sinh tự hồi quy | Quá trình mô hình dự đoán phần tử tiếp theo dựa trên chuỗi các phần tử đã sinh ra trước đó}
\end{question}

% ==============================================================================
% Question 87 | Phân loại: 2. Generative AI: Introduction and Applications
% Xuất hiện trong các đề gốc: AIG202c_SU26_FE (Q17)
% ==============================================================================
\begin{question}{87}{mcq}
\source{Q17.webp}
\prompt{How does generative AI contribute to educational personalization?}
\choice{A}{By creating avatars for classroom entertainment for students' emotion}
\choice{B}{By translating financial data for student portfolios}
\choice{C}{By generating learning journeys tailored to individual pacing and strengths}
\choice{D}{By automating legal document reviews}
\answer{C}
\confidence{high}
\evidence{Generative AI enables educational personalization by dynamically generating adaptive learning paths, tailored explanations, and targeted exercises suited to each individual student's pacing and strengths.}
\sources{UNESCO: Guidance for Generative AI in Education; Coursera Generative AI for Educators}
\explain{Hướng giải: Xác định cách thức AI tạo sinh (Generative AI) đóng góp vào việc cá nhân hóa giáo dục (educational personalization).

Các bước phân tích:
1. Cá nhân hóa trong giáo dục là việc điều chỉnh lộ trình học, nội dung bài giảng và nhịp độ học tập phù hợp với từng cá nhân học sinh.
2. AI tạo sinh đóng góp vào mục tiêu này bằng cách tự động tạo ra các lộ trình học tập tùy chỉnh theo tiến độ và điểm mạnh của từng người học (generating learning journeys tailored to individual pacing and strengths), đồng thời sinh thêm bài tập bổ trợ hoặc giải thích theo cách tiếp thu riêng của học sinh.
3. Vì vậy, phương án C mô tả chính xác ứng dụng mang tính sư phạm này.

Kiểm tra:
- A (tạo avatar giải trí), B (dịch dữ liệu tài chính cho hồ sơ), D (tự động hóa rà soát tài liệu pháp lý) không phải là cơ chế cốt lõi của việc cá nhân hóa học tập.

Vì sao các phương án khác sai:
- A là tính năng giải trí hình ảnh đơn thuần, không giải quyết bài toán cá nhân hóa học tập.
- B, D thuộc lĩnh vực tài chính và pháp lý, không liên quan đến cá nhân hóa sư phạm.

Lưu ý/Khái niệm cần nhớ:
Cá nhân hóa giáo dục (Educational Personalization) bằng AI tạo sinh tập trung vào việc điều chỉnh tài liệu và lộ trình theo năng lực riêng của học sinh.}
\vocab{Personalized Learning Journey | Lộ trình học tập cá nhân hóa | Kế hoạch và chuỗi bài học được AI tùy biến linh hoạt theo tốc độ tiếp thu và năng lực của từng học viên}
\end{question}

% ==============================================================================
% Question 88 | Phân loại: 2. Generative AI: Introduction and Applications
% Xuất hiện trong các đề gốc: AIG202c_SU26_FE (Q18)
% ==============================================================================
\begin{question}{88}{mcq}
\source{Q18.webp}
\prompt{What is a key ethical concern related to the use of generative AI?}
\choice{A}{Potential for misuse in generating false information or deepfakes.}
\choice{B}{Generative AI's tendency to require less computational power.}
\choice{C}{Generative AI leads to biased data classifications.}
\choice{D}{Generative AI's inability to scale applications across industries.}
\answer{A}
\confidence{high}
\evidence{A critical ethical concern with generative AI is its potential to be misused for generating false information (hallucinations/disinformation) and realistic deepfakes.}
\sources{EU Artificial Intelligence Act; NIST AI Risk Management Framework}
\explain{Hướng giải: Nhận diện mối lo ngại đạo đức trọng tâm liên quan đến việc sử dụng AI tạo sinh (Generative AI).

Các bước phân tích:
1. AI tạo sinh có khả năng tạo ra văn bản, hình ảnh, âm thanh và video giả mạo giống hệt như thật với chi phí thấp và tốc độ cao.
2. Nguy cơ đạo đức hàng đầu được toàn cầu cảnh báo là việc lợi dụng công nghệ này để tạo và phát tán thông tin sai lệch (disinformation/fake news), lừa đảo và tạo các nội dung giả mạo sâu (deepfakes) gây tổn hại danh dự, lừa đảo tài chính hoặc thao túng chính trị.
3. Do đó, phương án A là câu trả lời chính xác.

Kiểm tra:
- B sai vì AI tạo sinh đòi hỏi tài nguyên tính toán (compute/GPU) rất lớn, không phải 'ít hơn'.
- C sai vì phân loại thiên kiến là vấn đề của AI phân loại (discriminative AI); với GenAI vấn đề là thiên kiến trong nội dung sinh ra (generated bias), nhưng nguy cơ trực tiếp và nghiêm trọng nhất là thông tin giả mạo và deepfake.
- D sai vì AI tạo sinh có khả năng mở rộng (scale) rất mạnh mẽ trên nhiều ngành công nghiệp.

Vì sao các phương án khác sai:
- B, D mâu thuẫn với thực tế công nghệ (GenAI ngốn nhiều điện toán và đang mở rộng rất nhanh).
- C đề cập đến bài toán phân loại dữ liệu (classification), không phải đặc trưng sinh nội dung của GenAI.

Lưu ý/Khái niệm cần nhớ:
Mối nguy đạo đức lớn nhất của GenAI: Deepfakes và phát tán tin giả (False information/Disinformation).}
\vocab{Deepfake | Công nghệ làm giả tinh vi | Kỹ thuật AI tổng hợp hình ảnh, âm thanh hoặc video của con người một cách chân thực nhằm mục đích đánh lừa người xem}
\end{question}

% ==============================================================================
% Question 89 | Phân loại: 2. Generative AI: Introduction and Applications
% Xuất hiện trong các đề gốc: AIG202c_SU26_FE (Q19)
% ==============================================================================
\begin{question}{89}{mcq}
\source{Q19.webp}
\prompt{What does generative AI do differently than discriminative AI?}
\choice{A}{Mimics the human ability to predict data}
\choice{B}{Mimics the human ability to analyze data}
\choice{C}{Mimics the human ability to classify data}
\choice{D}{Mimics the human ability to create data}
\answer{D}
\confidence{high}
\evidence{While discriminative AI models \(P(Y \mid X)\) to classify, predict, or evaluate existing data, generative AI models \(P(X)\) or \(P(X, Y)\) to create brand-new data instances, mimicking the human ability to create content.}
\sources{Goodfellow et al., 'Deep Learning' (MIT Press); Ng & Jordan (2001) 'On Discriminative vs. Generative classifiers'}
\explain{Hướng giải: Phân biệt bản chất khác nhau giữa Generative AI (AI tạo sinh) và Discriminative AI (AI phân biệt).

Các bước phân tích:
1. Discriminative AI (AI phân biệt): Học ranh giới quyết định giữa các lớp dữ liệu để phân loại (classify), nhận dạng (recognize) hoặc dự đoán (predict) nhãn của dữ liệu có sẵn. Nó mô hình hóa xác suất có điều kiện \(P(Y \mid X)\).
2. Generative AI (AI tạo sinh): Học phân phối xác suất tổng thể của dữ liệu \(P(X)\) để tự sinh ra các mẫu dữ liệu mới chưa từng tồn tại trước đó (văn bản mới, ảnh mới, nhạc mới).
3. Điểm khác biệt lớn nhất: Generative AI mô phỏng khả năng sáng tạo dữ liệu của con người (Mimics the human ability to create data).

Kiểm tra:
- Dự đoán dữ liệu (A), phân tích dữ liệu (B), phân loại dữ liệu (C) đều là chức năng truyền thống của Discriminative AI và Machine Learning cổ điển.
- Tạo dữ liệu mới (D) chính là định nghĩa làm nên tên gọi 'Generative' (tạo sinh).

Vì sao các phương án khác sai:
- A, B, C đều là nhiệm vụ của mô hình phân biệt hoặc phân tích thống kê.

Lưu ý/Khái niệm cần nhớ:
Discriminative AI = Phân loại dữ liệu; Generative AI = Sáng tạo dữ liệu mới.}
\vocab{Generative vs. Discriminative | Tạo sinh so với Phân biệt | Hai trường phái mô hình Machine Learning: một bên tạo ra mẫu dữ liệu mới, một bên phân loại nhãn dữ liệu}
\end{question}

% ==============================================================================
% Question 90 | Phân loại: 2. Generative AI: Introduction and Applications
% Xuất hiện trong các đề gốc: AIG202c_SU26_FE (Q21)
% ==============================================================================
\begin{question}{90}{mcq}
\source{Q21.webp}
\prompt{Which of the following is a notable application of generative AI in the healthcare sector?}
\choice{A}{Automated diagnostics using AI-generated models}
\choice{B}{Real-time inventory tracking across hospital facilities}
\choice{C}{Patient appointment scheduling systems}
\choice{D}{Digital marketing strategies for healthcare institutions}
\answer{A}
\confidence{high}
\evidence{Generative AI in healthcare is leveraged for automated diagnostics, medical report generation, synthetic data synthesis for research, and novel therapeutic molecular design.}
\sources{Nature Medicine: 'Applications of generative AI in healthcare'; IBM Watson Health}
\explain{Hướng giải: Xác định ứng dụng nổi bật của AI tạo sinh trong lĩnh vực y tế và chăm sóc sức khỏe.

Các bước phân tích:
1. Trong lĩnh vực y tế, AI tạo sinh được ứng dụng sâu rộng trong việc hỗ trợ chẩn đoán tự động (automated diagnostics using AI-generated models), tổng hợp dữ liệu bệnh án giả định phục vụ nghiên cứu bảo mật thông tin, và sinh cấu trúc phân tử thuốc mới.
2. Các hệ thống GenAI có thể đọc ảnh chụp cộng hưởng từ (MRI/CT), tạo bản tóm tắt bệnh án lâm sàng tự động và đưa ra các đề xuất chẩn đoán hỗ trợ bác sĩ.

Kiểm tra:
- B (theo dõi tồn kho thời gian thực) là bài toán quản trị chuỗi cung ứng/ERP thông thường.
- C (đặt lịch hẹn bệnh nhân) là hệ thống quản lý cơ sở dữ liệu/hành chính cơ bản.
- D (chiến lược marketing kỹ thuật số) không phải ứng dụng cốt lõi về chuyên môn y tế của bệnh viện.

Vì sao các phương án khác sai:
- B, C, D là các tác vụ công nghệ thông tin hoặc quản trị hành chính thông thường, không đại diện cho đột phá y khoa của Generative AI.

Lưu ý/Khái niệm cần nhớ:
Ứng dụng y tế then chốt của GenAI gồm hỗ trợ chẩn đoán tự động, viết báo cáo lâm sàng và thiết kế phân tử thuốc mới.}
\vocab{Automated Diagnostics | Chẩn đoán tự động | Việc sử dụng các mô hình AI để phân tích dữ liệu y tế, hình ảnh bệnh lý và đề xuất chẩn đoán cho nhân viên y tế}
\end{question}

% ==============================================================================
% Question 91 | Phân loại: 2. Generative AI: Introduction and Applications
% Xuất hiện trong các đề gốc: AIG202c_SU26_FE (Q23)
% ==============================================================================
\begin{question}{91}{mcq}
\source{Q23.webp}
\prompt{Taesha is a second-grade teacher wondering how generative AI tools can help her students with special needs and learning disabilities. What do you think?}
\choice{A}{Enhance code maintainability}
\choice{B}{Develop diagnostic tools}
\choice{C}{Track learner progress}
\choice{D}{Analyze system logs}
\answer{C}
\confidence{high}
\evidence{In pedagogical applications for educators supporting special needs students, generative AI assists in monitoring and tracking learner progress to provide tailored interventions.}
\sources{IBM Generative AI for Educators; US Department of Education AI Report}
\explain{Hướng giải: Xác định ứng dụng liên quan đến giảng dạy và hỗ trợ học sinh có nhu cầu đặc biệt / khuyết tật học tập.

Các bước phân tích:
1. Taesha là một giáo viên tiểu học (lớp 2) muốn ứng dụng công cụ AI tạo sinh để hỗ trợ các học sinh có nhu cầu đặc biệt (special needs) và khuyết tật học tập (learning disabilities).
2. Trong môi trường sư phạm, việc theo dõi sát sao tiến trình học tập của từng cá nhân (Track learner progress) là vô cùng quan trọng để phát hiện sớm các rào cản nhận thức, từ đó AI có thể tạo ra các bài giảng điều chỉnh phù hợp với tốc độ của các em.
3. Xem xét các phương án:
   - A: Nâng cao khả năng bảo trì mã nguồn (Enhance code maintainability) - chỉ dành cho kỹ sư phần mềm.
   - B: Phát triển công cụ chẩn đoán (Develop diagnostic tools) - thuộc lĩnh vực y tế hoặc công cụ kỹ thuật.
   - C: Theo dõi tiến trình của người học (Track learner progress) - trực tiếp phục vụ mục tiêu giáo dục của giáo viên.
   - D: Phân tích nhật ký hệ thống (Analyze system logs) - thuộc lĩnh vực quản trị hệ thống IT.

Kiểm tra:
Phương án C là phương án duy nhất thuộc nghiệp vụ sư phạm và hỗ trợ trực tiếp cho người học.

Vì sao các phương án khác sai:
- A, D là các tác vụ kỹ thuật công nghệ thông tin.
- B là tác vụ y tế/chẩn đoán lâm sàng, không phải công việc sư phạm của giáo viên lớp 2.

Lưu ý/Khái niệm cần nhớ:
Trong giáo dục, AI giúp giáo viên theo dõi tiến trình học tập (Track learner progress) của từng học sinh để can thiệp kịp thời.}
\vocab{Special Needs Education | Giáo dục đặc biệt | Hoạt động giảng dạy được thiết kế riêng biệt để hỗ trợ học sinh có khuyết tật về thể chất hoặc học tập}
\end{question}

% ==============================================================================
% Question 92 | Phân loại: 2. Generative AI: Introduction and Applications
% Xuất hiện trong các đề gốc: AIG202c_SU26_FE (Q24)
% ==============================================================================
\begin{question}{92}{mcq}
\source{Q24.webp}
\prompt{Which tool is designed to create brand-specific marketing content?}
\choice{A}{Copy.ai for quick social media messages}
\choice{B}{Jasper for personalized brand messaging}
\choice{C}{Writesonic for structured blog articles}
\choice{D}{uClassifi/ for basic text categorization}
\answer{B}
\confidence{high}
\evidence{Jasper (Jasper.ai) is specifically designed and marketed for businesses to create brand-specific marketing content using its tailored 'Brand Voice' feature.}
\sources{Jasper.ai Documentation: Brand Voice; Forbes Tech Review of Generative AI Marketing Tools}
\explain{Hướng giải: Nhận diện công cụ AI tạo sinh được thiết kế chuyên biệt để tạo nội dung tiếp thị phù hợp với giọng văn thương hiệu (brand-specific marketing content).

Các bước phân tích:
1. Trong số các công cụ tiếp thị bằng AI phổ biến, Jasper (trước đây là Jarvis) nổi tiếng nhất với định vị doanh nghiệp: tạo nội dung tiếp thị cá nhân hóa theo tiếng nói thương hiệu ('Jasper for personalized brand messaging' / 'Jasper Brand Voice').
2. Jasper cho phép doanh nghiệp tải lên các hướng dẫn phong cách thương hiệu (style guides) và tài liệu nội bộ để đảm bảo mọi bài viết sinh ra đều đồng nhất với bản sắc thương hiệu.
3. Đối chiếu với các phương án:
   - Phương án B nêu trực tiếp: 'Jasper for personalized brand messaging'.

Kiểm tra:
- Copy.ai (A) thường được dùng cho tin nhắn mạng xã hội nhanh.
- Writesonic (C) tập trung vào viết bài blog chuẩn SEO.
- uClassify (D) là API phân loại văn bản, không phải công cụ sáng tạo nội dung tiếp thị.
- Jasper là công cụ tiêu biểu nhất cho thông điệp thương hiệu cá nhân hóa (Personalized brand messaging).

Vì sao các phương án khác sai:
- A, C, D không có tính năng định hình tiếng nói thương hiệu doanh nghiệp chuyên sâu như Jasper.

Lưu ý/Khái niệm cần nhớ:
Jasper.ai nổi tiếng trong mảng tiếp thị nhờ tính năng Brand Voice giúp tạo nội dung nhất quán với bản sắc thương hiệu.}
\vocab{Brand Voice | Tiếng nói thương hiệu | Phong cách, ngôn từ và giọng điệu đặc trưng mà một thương hiệu sử dụng trong toàn bộ nội dung tiếp thị}
\end{question}

% ==============================================================================
% Question 93 | Phân loại: 2. Generative AI: Introduction and Applications
% Xuất hiện trong các đề gốc: AIG202c_SU26_FE (Q27)
% ==============================================================================
\begin{question}{93}{mcq}
\source{Q27.webp}
\prompt{What is one benefit of generative AI for businesses?}
\choice{A}{Generative AI can create personalized marketing content.}
\choice{B}{Generative AI replaces human decision-making processes.}
\choice{C}{Generative AI is limited to creating synthetic voices.}
\choice{D}{Generative AI can only work with numerical data.}
\answer{A}
\confidence{high}
\evidence{Generative AI provides businesses with the ability to create personalized marketing content at scale, tailoring copy, images, and offers to individual customer segments.}
\sources{Harvard Business Review: 'How Generative AI Is Changing Creative Work'; Gartner Generative AI Insights}
\explain{Hướng giải: Xác định lợi ích thực tế của AI tạo sinh đối với các doanh nghiệp.

Các bước phân tích:
1. AI tạo sinh (Generative AI) sở hữu năng lực vượt trội trong việc tự động hóa sáng tạo nội dung đa phương tiện theo yêu cầu.
2. Đối với doanh nghiệp, lợi ích nổi bật và phổ biến nhất là khả năng tạo ra các nội dung tiếp thị được cá nhân hóa ở quy mô lớn (create personalized marketing content), từ email, bài đăng mạng xã hội đến mô tả sản phẩm cho từng phân khúc khách hàng.
3. Vì vậy, phương án A là phát biểu hoàn toàn chính xác.

Kiểm tra:
- B sai vì AI chỉ hỗ trợ chứ không thay thế hoàn toàn quá trình ra quyết định của con người (Human-in-the-loop).
- C sai vì GenAI tạo được văn bản, hình ảnh, video, mã lập trình, không chỉ giới hạn ở giọng nói nhân tạo (synthetic voices).
- D sai vì GenAI xử lý dữ liệu phi cấu trúc (văn bản, âm thanh, hình ảnh), không chỉ giới hạn ở dữ liệu số (numerical data).

Vì sao các phương án khác sai:
- B, C, D đưa ra các hạn chế hoặc nhận định sai lệch hoàn toàn về phạm vi ứng dụng của Generative AI.

Lưu ý/Khái niệm cần nhớ:
Một trong những lợi ích kinh doanh hàng đầu của GenAI là cá nhân hóa nội dung tiếp thị nhanh chóng và tiết kiệm chi phí.}
\vocab{Personalized Marketing Content | Nội dung tiếp thị cá nhân hóa | Các thông điệp quảng cáo được tùy biến riêng cho từng khách hàng nhằm tăng tỷ lệ tương tác}
\end{question}

% ==============================================================================
% Question 94 | Phân loại: 2. Generative AI: Introduction and Applications
% Xuất hiện trong các đề gốc: AIG202c_SU26_FE (Q29)
% ==============================================================================
\begin{question}{94}{mcq}
\source{Q29.webp}
\prompt{Which of the following statements about Stable Diffusion is correct?}
\choice{A}{It generates animations from sketches and outlines}
\choice{B}{It only creates basic line drawings from text}
\choice{C}{It is an open-source diffusion model for text to image}
\choice{D}{It uses a GAN-based approach for video synthesis}
\answer{C}
\confidence{high}
\evidence{Stable Diffusion is a prominent open-source latent diffusion model designed for text-to-image generation, released by Stability AI, CompVis, and Runway.}
\sources{Rombach et al. (2022) 'High-Resolution Image Synthesis with Latent Diffusion Models'; Stability AI}
\explain{Hướng giải: Xác định phát biểu chính xác về mô hình Stable Diffusion.

Các bước phân tích:
1. Stable Diffusion (do nhóm nghiên cứu CompVis tại Đại học LMU Munich, Stability AI và Runway phát triển) là một mô hình khuếch tán tiềm ẩn (Latent Diffusion Model) dùng để tạo hình ảnh từ văn bản mô tả (text-to-image).
2. Điểm đặc biệt nhất khiến Stable Diffusion trở thành một hiện tượng toàn cầu là mô hình này được phát hành mã nguồn mở (open-source) hoàn toàn cho cộng đồng nghiên cứu và nhà phát triển tự do sử dụng và cài đặt trên phần cứng cá nhân.
3. Do đó, phương án C ('It is an open-source diffusion model for text to image') là khẳng định chính xác.

Kiểm tra:
- A sai vì Stable Diffusion là mô hình tạo ảnh tĩnh từ văn bản, không phải công cụ chuyên tạo animation từ bản phác thảo.
- B sai vì nó tạo ra hình ảnh quang học chân thực (photorealistic) cực kỳ chi tiết, không chỉ là nét vẽ cơ bản (basic line drawings).
- D sai vì nó sử dụng mô hình khuếch tán (Diffusion Model), không phải mạng đối nghịch sinh (GAN).

Vì sao các phương án khác sai:
- A, B, D mô tả sai công nghệ nền tảng (GAN thay vì Diffusion) hoặc hạ thấp năng lực thực tế của mô hình.

Lưu ý/Khái niệm cần nhớ:
Stable Diffusion là mô hình khuếch tán tạo ảnh từ văn bản (text-to-image) mã nguồn mở (open-source) nổi tiếng nhất thế giới.}
\vocab{Diffusion Model | Mô hình khuếch tán | Lớp mô hình tạo sinh hoạt động bằng cách khử nhiễu dần dần từ một hình ảnh nhiễu ngẫu nhiên để tạo ra dữ liệu có ý nghĩa}
\end{question}

% ==============================================================================
% Question 95 | Phân loại: 2. Generative AI: Introduction and Applications
% Xuất hiện trong các đề gốc: AIG202c_SP26_RE (Q12)
% ==============================================================================
\begin{question}{95}{mcq}
\source{Q12.jpg}
\prompt{Which of the following tasks is least likely to be performed by generative AI models like LLMs?}
\choice{A}{Predicting stock market trends}
\choice{B}{Creating human-like text}
\choice{C}{Summarizing large documents}
\choice{D}{Answering questions based on text input}
\answer{A}
\confidence{high}
\evidence{Dự đoán xu hướng thị trường chứng khoán (Predicting stock market trends) là bài toán dự báo chuỗi thời gian số liệu phức tạp dựa trên kinh tế lượng và phân tích định lượng, KHÔNG PHẢI là thế mạnh và ít có khả năng được thực hiện thành công bởi các mô hình sinh ngôn ngữ lớn (LLMs).}
\sources{Financial Machine Learning; Limits of Large Language Models in Quantitative Finance}
\explain{Hướng giải: Xác định tác vụ ít phù hợp nhất đối với mô hình ngôn ngữ lớn tạo sinh (Generative AI / LLMs).

Các bước phân tích:
1. LLMs được tối ưu hóa cho các tác vụ xử lý ngôn ngữ tự nhiên: sinh văn bản (B), tóm tắt tài liệu (C), trả lời câu hỏi dựa trên ngữ cảnh (D).
2. Thị trường tài chính phụ thuộc vào biến động vĩ mô, giao dịch tần suất cao, tâm lý thị trường thời gian thực và chuỗi số liệu tài chính; LLMs không thể dự đoán chính xác biến động giá cổ phiếu do thiếu mô hình toán định lượng chuỗi thời gian và tính nhiễu cực lớn.

Kiểm tra: Giới nghiên cứu AI tài chính đồng thuận rằng việc dùng LLMs dự đoán giá cổ phiếu trực tiếp thường thất bại hoặc chỉ là ảo giác (hallucination).

Vì sao các phương án khác LLMs làm rất tốt:
- B: Tạo văn bản tự nhiên là năng lực cốt lõi của LLMs.
- C: Tóm tắt văn bản dài là thế mạnh hàng đầu.
- D: Hỏi đáp (Q&A) là ứng dụng phổ biến nhất.

Lưu ý/Khái niệm cần nhớ: LLMs giỏi về xử lý ngôn ngữ, không phải công cụ dự báo số liệu chuỗi thời gian chứng khoán.}
\vocab{Quantitative Finance | Tài chính định lượng | Lĩnh vực ứng dụng mô hình toán học và thống kê vào thị trường tài chính}
\end{question}

% ==============================================================================
% Question 96 | Phân loại: 2. Generative AI: Introduction and Applications
% Xuất hiện trong các đề gốc: AIG202c_SP26_RE (Q19)
% ==============================================================================
\begin{question}{96}{mcq}
\source{Q19.jpg}
\prompt{Maria is a scientific visualizer. How can the image-generation capabilities of generative AI help her in her daily work?}
\choice{A}{Help her convert text into audio using text-to-speech}
\choice{B}{Help her create comprehensive artwork from a simple sketch}
\choice{C}{Help write new code snippets based on functionality}
\choice{D}{Help her generate synthetic datasets for more diversity}
\answer{B}
\confidence{high}
\evidence{Khả năng sinh hình ảnh của Generative AI giúp nhà trực quan hóa khoa học (scientific visualizer) biến các bản phác thảo ý tưởng đơn giản thành các tác phẩm đồ họa hoàn chỉnh, sinh động và trực quan (Help her create comprehensive artwork from a simple sketch).}
\sources{Adobe Firefly Sketch-to-Image; NVIDIA Canvas AI for Scientific Illustration}
\explain{Hướng giải: Xác định ứng dụng trực tiếp của năng lực sinh ảnh GenAI cho chuyên gia trực quan hóa khoa học.

Các bước phân tích:
1. Nhiệm vụ của 'scientific visualizer' là biến các khái niệm, phân tử hoặc dữ liệu khoa học trừu tượng thành hình ảnh minh họa trực quan dễ hiểu.
2. Năng lực Image Generation (như Sketch-to-Image / Image-to-Image) cho phép chuyên gia vẽ một phác thảo sơ bộ bố cục rồi sử dụng AI để tạo thành tác phẩm đồ họa 2D/3D chi tiết hoàn chỉnh trong vài giây.

Kiểm tra: Các công cụ như ControlNet, Adobe Firefly, NVIDIA Canvas hỗ trợ hoàn hảo quy trình biến phác thảo thành hình ảnh hoàn thiện.

Vì sao các phương án khác sai:
- A (Text-to-speech): Thuộc về sinh âm thanh giọng nói, không liên quan đến hình ảnh.
- C (Write code): Thuộc về sinh mã nguồn lập trình.
- D (Synthetic datasets): Thuộc về kỹ thuật dữ liệu cho học máy, không phải tác vụ trực quan hóa nghệ thuật thường nhật của nhà thiết kế.

Lưu ý/Khái niệm cần nhớ: Image GenAI giúp tăng tốc quy trình từ phác thảo ý tưởng đến tác phẩm trực quan hoàn chỉnh.}
\vocab{Sketch-to-Image | Chuyển phác thảo thành hình ảnh | Kỹ thuật AI biến nét vẽ phác thảo tay thành hình ảnh hoàn thiện chi tiết}
\end{question}

% ==============================================================================
% Question 97 | Phân loại: 2. Generative AI: Introduction and Applications
% Xuất hiện trong các đề gốc: AIG202c_SP26_RE (Q22)
% ==============================================================================
\begin{question}{97}{mcq}
\source{Q22.jpg}
\prompt{Which statement is true about Adobe Firefly?}
\choice{A}{It only works offline for privacy and data safety}
\choice{B}{It is trained on Adobe Stock and public domain content}
\choice{C}{It supports only English for input text prompts}
\choice{D}{It requires no installation and runs only on mobile apps}
\answer{B}
\confidence{high}
\evidence{Adobe Firefly được phát triển trên nguyên tắc thương mại an toàn (commercially safe), được huấn luyện độc quyền trên kho ảnh bản quyền Adobe Stock và các nội dung thuộc phạm vi công cộng (public domain) mà thời hạn bản quyền đã hết.}
\sources{Adobe Firefly FAQ and Legal Documentation; Creative Cloud Generative AI Architecture}
\explain{Hướng giải: Nhận diện phát biểu đúng về mô hình sinh hình ảnh Adobe Firefly.

Các bước phân tích:
1. Điểm khác biệt lớn nhất của Adobe Firefly so với Midjourney hay Stable Diffusion là tính an toàn bản quyền thương mại (commercially safe).
2. Adobe công bố chính thức rằng Firefly chỉ được huấn luyện trên hàng triệu hình ảnh có bản quyền thuộc kho Adobe Stock và các tài liệu công cộng (public domain), đảm bảo các doanh nghiệp sử dụng không lo ngại về vấn đề kiện tụng quyền sở hữu trí tuệ.

Kiểm tra: Đây là cam kết pháp lý nổi tiếng của Adobe dành cho người dùng doanh nghiệp.

Vì sao các phương án khác sai:
- A: Adobe Firefly chạy trên hạ tầng đám mây (Cloud), không hoạt động ngoại tuyến (offline).
- C: Firefly hỗ trợ hơn 100 ngôn ngữ cho câu lệnh đầu vào, không chỉ riêng tiếng Anh.
- D: Firefly tích hợp sâu trong Photoshop, Illustrator trên máy tính để bàn (Desktop), không chỉ chạy trên ứng dụng di động.

Lưu ý/Khái niệm cần nhớ: Adobe Firefly được huấn luyện trên Adobe Stock và dữ liệu công cộng để đảm bảo an toàn bản quyền.}
\vocab{Commercially Safe AI | AI an toàn thương mại | Mô hình AI được huấn luyện trên dữ liệu hợp pháp đảm bảo không vi phạm bản quyền}
\end{question}

% ==============================================================================
% Question 98 | Phân loại: 2. Generative AI: Introduction and Applications
% Xuất hiện trong các đề gốc: AIG202c_SP26_RE (Q26)
% ==============================================================================
\begin{question}{98}{mcq}
\source{Q26.jpg}
\prompt{In the financial sector, what role does Bloomberg GPT play?}
\choice{A}{Detecting network security fraud from various text sources}
\choice{B}{Automating interview scheduling from various text sources}
\choice{C}{Analyzing market sentiment from various text sources}
\choice{D}{Generating medical images for rare conditions from various text sources}
\answer{C}
\confidence{high}
\evidence{BloombergGPT là mô hình ngôn ngữ lớn 50 tỷ tham số được huấn luyện chuyên biệt trên dữ liệu tài chính của Bloomberg, đóng vai trò chính trong việc phân tích sắc thái thị trường tài chính từ đa dạng nguồn tin tức và báo cáo (Analyzing market sentiment from various text sources).}
\sources{Wu et al. (2023), BloombergGPT: A Large Language Model for Finance (Bloomberg LP)}
\explain{Hướng giải: Xác định vai trò ứng dụng chính của mô hình ngôn ngữ BloombergGPT trong ngành tài chính.

Các bước phân tích:
1. BloombergGPT được phát triển bởi Bloomberg và công bố năm 2023, là mô hình LLM tiên phong được xây dựng riêng cho miền tài chính (domain-specific financial LLM).
2. Mô hình được nạp kho dữ liệu khổng lồ gồm tin tức tài chính, báo cáo thường niên, hồ sơ SEC, lịch sử giá cả để thực hiện các tác vụ: Phân tích sắc thái thị trường (Sentiment Analysis), Trích xuất thực thể tài chính và Tóm tắt báo cáo kinh doanh.

Kiểm tra: Khả năng phân tích tâm lý thị trường (market sentiment) từ nguồn tin tức văn bản là một trong những ứng dụng nổi bật nhất được công bố trong bài báo khoa học của Bloomberg.

Vì sao các phương án khác sai:
- A: Phát hiện gian lận an ninh mạng là tác vụ của Cybersecurity AI.
- B: Lên lịch phỏng vấn thuộc về phần mềm nhân sự.
- D: Tạo ảnh y tế thuộc về Medical Computer Vision, hoàn toàn không liên quan đến tài chính.

Lưu ý/Khái niệm cần nhớ: BloombergGPT = Mô hình ngôn ngữ lớn chuyên sâu về tài chính phục vụ phân tích thị trường và tin tức.}
\vocab{Market Sentiment Analysis | Phân tích sắc thái thị trường | Việc dùng AI phân tích tin tức và mạng xã hội để đo lường tâm lý lạc quan hay bi quan của giới đầu tư}
\end{question}

\section{sec3}{3. Generative AI: Prompt Engineering Basics}

% ==============================================================================
% Question 99 | Phân loại: 3. Generative AI: Prompt Engineering Basics
% Xuất hiện trong các đề gốc: AIG202c_SP_2025_FE (Q11), AIG202c_SU25_FE (Q42), AIG202c_SU26_FE (Q40)
% ==============================================================================
\begin{question}{99}{mcq}
\source{AIG202c_-_SP_2025_-_FE_3755_11.webp}
\prompt{What is the primary goal of prompt engineering tools?}
\choice{A}{To optimize the creation of prompts for generative AI models.}
\choice{B}{To create applications for language model experiments.}
\choice{C}{To design user-friendly interface for generative AI models.}
\choice{D}{To provide suggestions for improving NLP techniques.}
\answer{A}
\confidence{high}
\evidence{Mục tiêu chính của các công cụ kỹ nghệ câu lệnh (prompt engineering tools như PromptPerfect, LangSmith) là tối ưu hóa quá trình xây dựng, thử nghiệm và hoàn thiện các câu lệnh cho mô hình AI tạo sinh (To optimize the creation of prompts for generative AI models).}
\sources{Prompt Engineering Tools and Workflows; LangChain / PromptPerfect Documentation}
\explain{Hướng giải: Xác định mục tiêu của các công cụ hỗ trợ Prompt Engineering.

Các bước phân tích:
1. Viết prompt thủ công thường tốn thời gian và khó đo lường định lượng.
2. Các công cụ Prompt Engineering ra đời nhằm: tự động gợi ý từ khóa, so sánh kết quả giữa nhiều mô hình, tinh chỉnh nhiệt độ, quản lý phiên bản prompt và tối ưu hóa độ dài nhằm giảm chi phí token và tăng độ chính xác.

Kiểm tra: Tối ưu hóa việc tạo câu lệnh (optimize creation of prompts) là chức năng định danh của nhóm công cụ này.

Vì sao các phương án khác sai:
- B: Tạo ứng dụng thí nghiệm ngôn ngữ chỉ là một tính năng phụ trợ.
- C: Thiết kế giao diện thân thiện là công việc UI design.
- D: Đưa ra gợi ý cải thiện kỹ thuật NLP là nghiên cứu thuật toán xử lý ngôn ngữ.

Lưu ý/Khái niệm cần nhớ: Prompt engineering tools = Công cụ tối ưu hóa việc thiết kế và quản trị câu lệnh GenAI.}
\vocab{Prompt Optimization | Tối ưu hóa câu lệnh | Việc cải tiến cấu trúc câu lệnh để đạt chất lượng phản hồi cao nhất với chi phí thấp nhất}
\end{question}

% ==============================================================================
% Question 100 | Phân loại: 3. Generative AI: Prompt Engineering Basics
% Xuất hiện trong các đề gốc: AIG202c_SP_2025_FE (Q15), AIG202c_SU25_FE (Q38)
% ==============================================================================
\begin{question}{100}{mcq}
\source{AIG202c_-_SP_2025_-_FE_3755_15.webp}
\prompt{Imagine you are using the Chain-of-Thought approach to teach a generative AI model how to solve a mathematical problem. What is the key component of a prompt in this approach?}
\choice{A}{A prompt includes only a question without an answer.}
\choice{B}{A prompt includes a list of formulae to solve the question.}
\choice{C}{A prompt consists of a question and a correct answer for context.}
\choice{D}{A prompt includes a series of related questions without the correct answer.}
\answer{C}
\confidence{high}
\evidence{Trong phương pháp Chain-of-Thought (CoT) prompting dạng Few-shot, thành phần then chốt của câu lệnh mẫu là bao gồm cả câu hỏi kèm theo chuỗi các bước suy luận giải thích dẫn đến câu trả lời chính xác (A prompt consists of a question and a correct answer with step-by-step reasoning for context).}
\sources{Wei et al. (2022), Chain-of-Thought Prompting; OpenAI CoT Cookbook}
\explain{Hướng giải: Nhận diện thành phần mẫu then chốt trong kỹ thuật Chain-of-Thought Few-shot prompting.

Các bước phân tích:
1. Để 'dạy' mô hình giải toán bằng Chain-of-Thought, ta cung cấp các ví dụ mẫu (exemplars).
2. Mỗi ví dụ mẫu bắt buộc phải có 3 phần: (1) Đề bài (Question), (2) Chuỗi lập luận từng bước trung gian (Reasoning chain), và (3) Đáp án chính xác (Correct answer).
3. Nhờ có đáp án và các bước giải mẫu làm ngữ cảnh, mô hình sẽ bắt chước tư duy tuần tự đó để giải quyết câu hỏi mới.

Kiểm tra: Bài báo gốc của Google Research nhấn mạnh việc cung cấp cặp \(\langle\text{Question, Reasoning Chain + Answer}\rangle\) là bản chất của CoT.

Vì sao các phương án khác sai:
- A: Chỉ có câu hỏi không có đáp án giải mẫu thì mô hình không học được cách suy luận CoT.
- B: Danh sách công thức thuần túy không thể hiện được quá trình suy luận từng bước.
- D: Một loạt câu hỏi không có đáp án không cung cấp được ngữ cảnh chuẩn mực.

Lưu ý/Khái niệm cần nhớ: Chain-of-Thought yêu cầu đề bài kèm các bước lập luận dẫn đến đáp án chính xác.}
\vocab{Few-Shot CoT | Chuỗi suy nghĩ kèm ví dụ mẫu | Việc đưa các ví dụ giải từng bước vào prompt để mô hình học theo quy trình suy luận}
\end{question}

% ==============================================================================
% Question 101 | Phân loại: 3. Generative AI: Prompt Engineering Basics
% Xuất hiện trong các đề gốc: AIG202c_SP_2025_FE (Q18)
% ==============================================================================
\begin{question}{101}{mcq}
\source{AIG202c_-_SP_2025_-_FE_3755_18.webp}
\prompt{Why is clarity important when writing prompts for generative AI models?}
\choice{A}{Clarity helps make the prompt less engaging}
\choice{B}{Clarity helps the model understand the task and produce relevant responses}
\choice{C}{Clarity adds to the complexity of the prompt}
\choice{D}{Clarity ensures the prompt is lengthy}
\answer{B, D}
\confidence{high}
\evidence{Tính rõ ràng (clarity) trong câu lệnh rất quan trọng vì nó giúp mô hình hiểu chính xác nhiệm vụ để tạo phản hồi đúng trọng tâm (helps model understand task and produce relevant responses) đồng thời loại bỏ sự mơ hồ, tránh hiểu lầm trong quá trình sinh văn bản.}
\sources{OpenAI Prompt Engineering Best Practices: Write clear instructions; Anthropic Prompt Guide}
\explain{Hướng giải: Xác định hai lý do tại sao tính rõ ràng (Clarity) lại mang tính sống còn khi viết prompt cho GenAI.

Các bước phân tích:
1. Mô hình ngôn ngữ phụ thuộc hoàn toàn vào ngữ cảnh được cung cấp trong câu lệnh.
2. Sự rõ ràng giúp:
   - (B) Mô hình hiểu chính xác mục tiêu của bài toán, từ đó kích hoạt đúng không gian tri thức và sinh ra câu trả lời có tính liên quan cao (relevant responses).
   - (D) Tránh được sự mơ hồ (ambiguity), hạn chế nguy cơ mô hình suy diễn sai lệch hoặc bịa đặt thông tin (hallucination).

Kiểm tra: Nguyên tắc số 1 trong tài liệu hướng dẫn của OpenAI và Anthropic luôn là: 'Be clear and specific'.

Vì sao các phương án khác sai:
- A: Rõ ràng không làm prompt kém hấp dẫn, ngược lại giúp tăng tính hiệu quả.
- C: Các phương án gây nhiễu khác trái ngược với mục tiêu giao tiếp rõ ràng.

Lưu ý/Khái niệm cần nhớ: Clarity = Hiểu đúng tác vụ + Triệt tiêu sự mơ hồ.}
\vocab{Clarity in Prompting | Tính rõ ràng trong câu lệnh | Việc sử dụng ngôn từ minh bạch, chính xác và không gây hiểu lầm trong prompt}
\end{question}

% ==============================================================================
% Question 102 | Phân loại: 3. Generative AI: Prompt Engineering Basics
% Xuất hiện trong các đề gốc: AIG202c_SP_2025_FE (Q19), AIG202c_SU25_FE (Q39), AIG202c_SP26_RE (Q39)
% ==============================================================================
\begin{question}{102}{mcq}
\source{AIG202c_-_SP_2025_-_FE_3755_19.webp}
\prompt{Which of the following is NOT a common prompt engineering technique?}
\choice{A}{Adapting prompts based on user feedback}
\choice{B}{Using pre-trained language models without any prompts}
\choice{C}{Adjusting the temperature parameter during model generation}
\choice{D}{Designing prompts to elicit specific types of responses}
\answer{B}
\confidence{high}
\evidence{Việc sử dụng mô hình ngôn ngữ được huấn luyện trước mà không kèm theo bất kỳ câu lệnh nào (Using pre-trained language models without any prompts) KHÔNG PHẢI là một kỹ thuật prompt engineering; mô hình ngôn ngữ không thể sinh ra nội dung mong muốn nếu không có prompt đầu vào kích hoạt.}
\sources{Prompt Engineering Guide (DAIR.AI); Brown et al. (2020), Language Models are Few-Shot Learners}
\explain{Hướng giải: Tìm phương án KHÔNG PHẢI là một kỹ thuật kỹ nghệ câu lệnh (Prompt Engineering).

Các bước phân tích:
1. Prompt Engineering định nghĩa là các kỹ thuật tác động lên câu lệnh đầu vào hoặc tham số suy luận:
   - Điều chỉnh câu lệnh dựa trên phản hồi (A - Feedback adaptation).
   - Điều chỉnh tham số nhiệt độ `temperature` (C - Temperature tuning để kiểm soát tính sáng tạo/ngẫu nhiên).
   - Thiết kế câu lệnh để thu được định dạng phản hồi cụ thể (D - Output structuring).
2. Phương án B: 'Using pre-trained language models without any prompts' (Dùng mô hình không cần prompt) là hành động không có prompt, do đó không thể coi là một 'kỹ thuật prompt engineering'.

Kiểm tra: Bản chất của Prompt Engineering là phải có Prompt.

Vì sao các phương án khác là kỹ thuật chuẩn:
- A, C, D đều là các kỹ thuật prompt engineering thực hành phổ biến.

Lưu ý/Khái niệm cần nhớ: Prompt Engineering luôn đòi hỏi thiết kế và kiểm soát câu lệnh đầu vào.}
\vocab{Temperature | Tham số nhiệt độ | Siêu tham số điều khiển mức độ ngẫu nhiên và sáng tạo của câu trả lời từ LLM}
\end{question}

% ==============================================================================
% Question 103 | Phân loại: 3. Generative AI: Prompt Engineering Basics
% Xuất hiện trong các đề gốc: AIG202c_SP_2025_FE (Q22)
% ==============================================================================
\begin{question}{103}{mcq}
\source{AIG202c_-_SP_2025_-_FE_3755_22.webp}
\prompt{Monica is a writer looking to improve her storytelling by using a large language model (LLM) for creative writing prompts. Which feature of a prompt engineering tool would be most beneficial for her to achieve this goaP}
\choice{A}{Iteratively refining prompts by adjusting them based on the initial outputs of the LLM's story}
\choice{B}{Providing feedback for prompts based on a given input.}
\choice{C}{Suggesting different prompt formats (like questions, dialogues, or scenarios) tailored to the preferred storytelling style (such as mystery, romance, or sci-fi)}
\choice{D}{Offering a library of pre-defined story beginnings designed for particular genres and themes}
\answer{C}
\confidence{high}
\evidence{Tính năng hữu ích nhất của công cụ prompt engineering đối với nhà văn sáng tạo là gợi ý các định dạng câu lệnh đa dạng (như đối thoại, tình huống) phù hợp với thể loại truyện cụ thể (Suggesting different prompt formats tailored to preferred storytelling style).}
\sources{Creative Writing Prompt Design; NovelAI Guidelines}
\explain{Hướng giải: Xác định tính năng hỗ trợ nhà văn tối ưu nhất từ công cụ prompt engineering (câu hỏi lặp lại của UQ036).

Các bước phân tích:
1. Nhà văn cần sự linh hoạt để khai mở ý tưởng theo từng phong cách văn học riêng biệt.
2. Công cụ cung cấp các định dạng prompt phù hợp với thể loại trinh thám, lãng mạn, viễn tưởng giúp kích hoạt trí tưởng tượng và tạo khung truyện phong phú nhất.

Kiểm tra: Khớp hoàn toàn với đáp án của câu UQ036.

Lưu ý/Khái niệm cần nhớ: Gợi ý định dạng prompt theo thể loại giúp tối đa hóa khả năng sáng tác văn học.}
\vocab{Creative Writing Tool | Công cụ hỗ trợ viết sáng tạo | Phần mềm ứng dụng AI trợ giúp tác giả xây dựng cốt truyện và lời thoại}
\end{question}

% ==============================================================================
% Question 104 | Phân loại: 3. Generative AI: Prompt Engineering Basics
% Xuất hiện trong các đề gốc: AIG202c_SP_2025_FE (Q32)
% ==============================================================================
\begin{question}{104}{mcq}
\source{AIG202c_-_SP_2025_-_FE_3755_32.webp}
\prompt{A prompt is any input or that helps generative AI models produce relevant and logical responses.}
\choice{A}{Series of instructions}
\choice{B}{Trigger}
\choice{C}{Signal}
\choice{D}{Stimulus}
\answer{A}
\confidence{high}
\evidence{Trong AI tạo sinh, câu lệnh (prompt) là bất kỳ đầu vào hoặc chuỗi các chỉ dẫn (Series of instructions) nào giúp mô hình tạo ra các phản hồi phù hợp và có tính logic.}
\sources{Prompt Engineering Principles; IBM watsonx Prompt Lab Guide}
\explain{Hướng giải: Điền cụm từ chính xác định nghĩa bản chất của câu lệnh (Prompt).

Các bước phân tích:
1. Cấu trúc định nghĩa chuẩn: 'A prompt is any input or series of instructions that helps generative AI models produce relevant and logical responses'.
2. Prompt không chỉ là một từ kích hoạt đơn lẻ mà thường là một tập hợp các chỉ dẫn có cấu trúc: bối cảnh, nhiệm vụ, ràng buộc, định dạng đầu ra và ví dụ mẫu.

Kiểm tra: 'Series of instructions' (Chuỗi chỉ dẫn) phản ánh chính xác cấu trúc toàn diện của một prompt chuyên nghiệp.

Vì sao các phương án khác sai:
- B (Trigger), C (Signal), D (Stimulus): Là các từ sinh học/điện tử đơn lẻ, không phản ánh tính chỉ dẫn logic của prompt.

Lưu ý/Khái niệm cần nhớ: Prompt = Đầu vào hoặc chuỗi các chỉ dẫn định hướng mô hình AI tạo sinh.}
\vocab{System Instruction | Chỉ thị hệ thống | Tập hợp các quy tắc định hướng hành vi cốt lõi của mô hình ngôn ngữ}
\end{question}

% ==============================================================================
% Question 105 | Phân loại: 3. Generative AI: Prompt Engineering Basics
% Xuất hiện trong các đề gốc: AIG202c_SP_2025_FE (Q33), AIG202c_SU25_FE (Q45), AIG202c_SP26_RE (Q41)
% ==============================================================================
\begin{question}{105}{mcq}
\source{AIG202c_-_SP_2025_-_FE_3755_33.webp}
\prompt{How can prompt engineering help improve the performance of AI models?}
\choice{A}{By generating random prompts for training data}
\choice{B}{By ensuring that prompts are designed to elicit relevant and coherent responses}
\choice{C}{By Ignoring prompts and allowing the model to generate output autonomously}
\choice{D}{By training models using only prompts that produce desired responses}
\answer{B}
\confidence{high}
\evidence{Prompt Engineering nâng cao hiệu năng của các mô hình AI bằng cách đảm bảo các câu lệnh được thiết kế chặt chẽ để gợi mở các câu trả lời đúng trọng tâm và mạch lạc (designed to elicit relevant and coherent responses).}
\sources{OpenAI Best Practices for Prompt Engineering; Coursera Generative AI with LLMs}
\explain{Hướng giải: Xác định cách thức kỹ nghệ câu lệnh cải thiện chất lượng phản hồi của mô hình AI.

Các bước phân tích:
1. Nguyên tắc 'Garbage in, garbage out' trong xử lý dữ liệu: Prompt mơ hồ dẫn đến câu trả lời kém chất lượng.
2. Prompt Engineering tối ưu hóa từ khóa, định dạng, vai trò và ví dụ mẫu giúp mô hình tập trung vào không gian xác suất chính xác nhất, mang lại câu trả lời chất lượng cao.

Kiểm tra: Tương tự câu UQ086, đây là định nghĩa chuẩn về giá trị của prompt engineering.

Vì sao các phương án khác sai:
- A: Tạo prompt ngẫu nhiên cho dữ liệu huấn luyện không phải là prompt engineering.
- C: Bỏ qua prompt là phản tác dụng.
- D: Huấn luyện mô hình là tác vụ khác.

Lưu ý/Khái niệm cần nhớ: Thiết kế prompt tốt mang lại câu trả lời đúng trọng tâm và mạch lạc.}
\vocab{Prompt Quality | Chất lượng câu lệnh | Mức độ rõ ràng, chi tiết và chính xác của chỉ dẫn đưa vào mô hình}
\end{question}

% ==============================================================================
% Question 106 | Phân loại: 3. Generative AI: Prompt Engineering Basics
% Xuất hiện trong các đề gốc: AIG202c_SP_2025_FE (Q37), AIG202c_SU25_FE (Q37), AIG202c_FA25_FE (Q43)
% ==============================================================================
\begin{question}{106}{mcq}
\source{AIG202c_-_SP_2025_-_FE_3755_37.webp}
\prompt{Imagine you are planning a business trip, and you want to use the interview pattern approach to prompt an AI model to assist you in planning your itinerary. What would be the benefit of this approach in comparison to a traditional static prompt?}
\choice{A}{The model will provide a single predetermined itinerary}
\choice{B}{The model will ask for your preferences and adjust the itinerary accordingly}
\choice{C}{The model will minimize the need for any user interaction}
\choice{D}{The model will generate a random travel plan.}
\answer{B}
\confidence{high}
\evidence{Kỹ thuật Interview Pattern (mẫu phỏng vấn) biến AI thành người phỏng vấn chủ động: thay vì đưa ra lịch trình tĩnh, mô hình sẽ đặt các câu hỏi để thu thập sở thích, ngân sách và yêu cầu cụ thể của người dùng rồi điều chỉnh lịch trình tương ứng.}
\sources{White et al., A Prompt Pattern Catalog to Enhance Prompt Engineering with ChatGPT (Vanderbilt University)}
\explain{Hướng giải: Nhận diện lợi ích của mẫu phỏng vấn (Interview Pattern) trong lập kế hoạch hành trình.

Các bước phân tích:
1. Mẫu prompt tĩnh truyền thống thường đưa ra một kết quả chung chung định trước do thiếu thông tin bối cảnh cá nhân.
2. Trong Interview Pattern, prompt yêu cầu: 'Hãy phỏng vấn tôi từng câu hỏi một về sở thích du lịch, thời gian, kinh phí trước khi lập lịch trình'.
3. Lợi ích vượt trội là mô hình thu thập đầy đủ dữ liệu người dùng và tạo ra lịch trình được cá nhân hóa tối đa.

Kiểm tra: Mẫu phỏng vấn là một trong những mẫu tương tác người - máy (Interaction Patterns) hiệu quả nhất trong kỹ nghệ câu lệnh.

Vì sao các phương án khác sai:
- A: Lịch trình định trước là hạn chế của prompt tĩnh thông thường.
- C: Phương pháp này chủ động tăng cường tương tác đối thoại, không phải giảm thiểu tương tác.
- D: Lịch trình được may đo theo sở thích, hoàn toàn không ngẫu nhiên.

Lưu ý/Khái niệm cần nhớ: Interview Pattern = AI chủ động đặt câu hỏi phỏng vấn để thu thập thông tin người dùng.}
\vocab{Interview Pattern | Mẫu phỏng vấn | Kỹ thuật yêu cầu AI đóng vai người phỏng vấn để khai thác chi tiết nhu cầu người dùng}
\end{question}

% ==============================================================================
% Question 107 | Phân loại: 3. Generative AI: Prompt Engineering Basics
% Xuất hiện trong các đề gốc: AIG202c_SP_2025_FE (Q38), AIG202c_SU25_FE (Q43), AIG202c_FA25_FE (Q42)
% ==============================================================================
\begin{question}{107}{mcq}
\source{AIG202c_-_SP_2025_-_FE_3755_38.webp}
\prompt{How does the Tree-of-Thoughts approach differ from traditional linear prompting approaches?}
\choice{A}{It encourages linear thinking}
\choice{B}{It eliminates all possible routes of thinking}
\choice{C}{It makes random decisions}
\choice{D}{It explores multiple possibilities simultaneously using a hierarchical structure}
\answer{D}
\confidence{high}
\evidence{Cách tiếp cận Tree-of-Thoughts (ToT) khác với prompt tuyến tính truyền thống ở chỗ nó cho phép mô hình khám phá nhiều khả năng suy nghĩ đồng thời bằng cấu trúc phân cấp dạng cây (explores multiple possibilities simultaneously using a hierarchical structure), kết hợp cơ chế quay lui (backtracking) và tự đánh giá.}
\sources{Yao et al. (2023), Tree of Thoughts: Deliberate Problem Solving with Large Language Models (Princeton University)}
\explain{Hướng giải: Phân biệt cấu trúc suy luận của Tree-of-Thoughts (ToT) so với các phương pháp prompt tuần tự thông thường.

Các bước phân tích:
1. Prompting tuyến tính (như Input-Output hoặc Chain-of-Thought) đi theo một đường thẳng duy nhất từ đầu đến cuối.
2. Tree-of-Thoughts tổ chức các suy nghĩ thành các nút trên cây (tree structure). Tại mỗi bước, mô hình có thể sinh ra nhiều nhánh suy nghĩ khác nhau, đánh giá triển vọng của từng nhánh và quay lui (backtrack) nếu gặp ngõ cụt.
3. Nhờ đó, mô hình có thể khám phá nhiều giải pháp đồng thời theo cấu trúc thứ bậc phân nhánh.

Kiểm tra: Nghiên cứu của Đại học Princeton và Google DeepMind về ToT lấy cảm hứng từ thuật toán tìm kiếm cây (như A* hoặc BFS/DFS) kết hợp với LLM.

Vì sao các phương án khác sai:
- A: ToT là phi tuyến tính (non-linear), không phải tư duy đường thẳng.
- B: ToT mở rộng và khám phá các lộ trình, không loại bỏ toàn bộ.
- C: ToT dựa trên đánh giá có chủ ý (deliberate heuristic evaluation), không phải quyết định ngẫu nhiên.

Lưu ý/Khái niệm cần nhớ: Tree-of-Thoughts = Khám phá đa nhánh suy luận đồng thời kết hợp quay lui (Tree Search + LLM).}
\vocab{Tree-of-Thoughts (ToT) | Cây suy nghĩ | Khung suy luận cho phép LLM khám phá nhiều lộ trình tư duy phân nhánh có đánh giá và quay lui}
\end{question}

% ==============================================================================
% Question 108 | Phân loại: 3. Generative AI: Prompt Engineering Basics
% Xuất hiện trong các đề gốc: AIG202c_SP_2025_FE (Q39)
% ==============================================================================
\begin{question}{108}{mcq}
\source{AIG202c_-_SP_2025_-_FE_3755_39.webp}
\prompt{Which of the following is a potential challenge associated with prompt engineering for generative AP}
\choice{A}{Limited control over the output of generative AI models}
\choice{B}{Difficulty in training generative AI models without prompts}
\choice{C}{Lack of available tools and resources for prompt engineering}
\choice{D}{Difficulty in generating prompts automatically using AI algorithms}
\answer{A}
\confidence{high}
\evidence{Một thách thức tiềm ẩn lớn trong Prompt Engineering cho AI tạo sinh là sự kiểm soát có giới hạn đối với kết quả đầu ra của mô hình (Limited control over the output of generative AI models) do tính chất ngẫu nhiên thống kê của kiến trúc sinh từ.}
\sources{Stochastic Parrots (Bender et al., 2021); Prompt Engineering Limitations and Challenges}
\explain{Hướng giải: Xác định thách thức kỹ thuật thực tế khi điều khiển mô hình AI tạo sinh bằng câu lệnh.

Các bước phân tích:
1. Các mô hình ngôn ngữ lớn hoạt động dựa trên việc lấy mẫu xác suất (sampling) từ phân phối từ tiếp theo.
2. Do đó, ngay cả với cùng một câu lệnh, mô hình vẫn có thể sinh ra các kết quả khác nhau qua các lần chạy (trừ khi cố định seed và temperature = 0).
3. Việc không thể kiểm soát tuyệt đối 100% từng câu chữ hoặc định dạng ('Limited control over the output') là một trong những thách thức lớn nhất khi tích hợp GenAI vào các hệ thống doanh nghiệp đòi hỏi tính tất định cao.

Kiểm tra: Đây là bài toán trung tâm trong nghiên cứu về độ tin cậy và an toàn của GenAI.

Vì sao các phương án khác sai:
- B: Huấn luyện mô hình không phải là tác vụ của prompt engineering.
- C: Hiện nay có rất nhiều công cụ hỗ trợ prompt engineering (LangChain, PromptPerfect, LlamaIndex).
- D: Việc tự động hóa prompt là mục tiêu nghiên cứu, không phải thách thức bản chất trong vận hành mô hình.

Lưu ý/Khái niệm cần nhớ: Tính ngẫu nhiên thống kê khiến con người khó kiểm soát tuyệt đối đầu ra của GenAI.}
\vocab{Non-deterministic Output | Đầu ra không tất định | Đặc tính của mô hình sinh có thể tạo ra các kết quả khác nhau cho cùng một đầu vào}
\end{question}

% ==============================================================================
% Question 109 | Phân loại: 3. Generative AI: Prompt Engineering Basics
% Xuất hiện trong các đề gốc: AIG202c_SP_2025_FE (Q42)
% ==============================================================================
\begin{question}{109}{mcq}
\source{AIG202c_-_SP_2025_-_FE_3755_42.webp}
\prompt{What is "prompt engineering" in the context of generative AI?}
\choice{A}{The process of creating prompts or instructions to guide the output of generative AI models}
\choice{B}{The process of training generative AI models using large datasets}
\choice{C}{The process of evaluating the performance of generative AI models}
\choice{D}{The process of generating prompts automatically using AI algorithms}
\answer{A}
\confidence{high}
\evidence{Kỹ nghệ câu lệnh (Prompt engineering) trong bối cảnh AI tạo sinh là quá trình thiết kế, xây dựng và hoàn thiện các câu lệnh hoặc chỉ dẫn để định hướng kết quả đầu ra của các mô hình AI tạo sinh (The process of creating prompts or instructions to guide the output of generative AI models).}
\sources{OpenAI Prompt Engineering Documentation; Coursera Prompt Engineering Guide}
\explain{Hướng giải: Xác định định nghĩa chính xác và bao quát nhất của Kỹ nghệ câu lệnh (Prompt engineering).

Các bước phân tích:
1. Prompt Engineering là phương pháp giao tiếp và điều hướng mô hình AI tạo sinh thông qua ngôn ngữ tự nhiên.
2. Bằng việc xây dựng các hướng dẫn chi tiết, bổ sung ngữ cảnh và ràng buộc định dạng, kỹ sư câu lệnh giúp mô hình sinh ra kết quả chính xác, mạch lạc và đáp ứng đúng mục tiêu ứng dụng.

Kiểm tra: Khớp hoàn toàn với định nghĩa chuẩn quốc tế và câu hỏi UQ037.

Vì sao các phương án khác sai:
- B: Huấn luyện với tập dữ liệu lớn là Model Training.
- C: Đánh giá hiệu năng là Model Evaluation / Benchmarking.
- D: Sinh prompt tự động chỉ là một nhánh kỹ thuật nhỏ (Auto-prompting), không phải định nghĩa tổng quát.

Lưu ý/Khái niệm cần nhớ: Prompt engineering = Quá trình thiết kế câu lệnh định hướng đầu ra của mô hình GenAI.}
\vocab{Prompt Engineering | Kỹ nghệ câu lệnh | Nghệ thuật và khoa học cấu trúc chỉ dẫn để định hướng phản hồi của AI}
\end{question}

% ==============================================================================
% Question 110 | Phân loại: 3. Generative AI: Prompt Engineering Basics
% Xuất hiện trong các đề gốc: AIG202c_SP_2025_FE (Q44)
% ==============================================================================
\begin{question}{110}{mcq}
\source{AIG202c_-_SP_2025_-_FE_3755_44.webp}
\prompt{How does the streamline mode in PromptPerfect contribute to prompt optimizauon?}
\choice{A}{By selecting a relevant model for prompt optimization}
\choice{B}{By accessing prebuilt prompts as examples}
\choice{C}{By refining prompts step-by-step until satisfied with the output}
\choice{D}{By experimenting with the autocomplete feature}
\answer{C}
\confidence{high}
\evidence{Chế độ hợp lý hóa (Streamline mode) trong công cụ PromptPerfect đóng góp vào việc tối ưu hóa câu lệnh bằng cách tự động tinh chỉnh câu lệnh từng bước cho đến khi người dùng hài lòng với kết quả đầu ra (By refining prompts step-by-step until satisfied with the output).}
\sources{PromptPerfect Documentation: Streamline Mode vs Standard Optimization}
\explain{Hướng giải: Xác định cơ chế hoạt động của tính năng Streamline mode trong công cụ tối ưu prompt PromptPerfect.

Các bước phân tích:
1. PromptPerfect là công cụ nổi tiếng chuyên tự động tối ưu hóa prompt cho nhiều mô hình (GPT-4, Claude, Midjourney).
2. Chế độ 'Streamline mode' cho phép người dùng chạy một quy trình tinh chỉnh lặp từng bước (step-by-step refinement): công cụ tự động phân tích câu lệnh gốc, viết lại các biến thể rõ ràng hơn và cho phép người dùng thử nghiệm lặp lại cho đến khi đạt chất lượng tối ưu nhất.

Kiểm tra: Tính năng này giúp tiết kiệm thời gian thử sai thủ công của lập trình viên.

Vì sao các phương án khác sai:
- A: Chọn mô hình là bước thiết lập ban đầu (Model selection).
- B: Truy cập prompt có sẵn là tính năng Thư viện mẫu (Prompt Library).
- D: Tự động hoàn thành là tính năng Autocomplete.

Lưu ý/Khái niệm cần nhớ: Streamline mode trong PromptPerfect = Tinh chỉnh câu lệnh lặp từng bước cho đến khi hoàn hảo.}
\vocab{PromptPerfect | Công cụ tối ưu câu lệnh | Nền tảng tự động phân tích và tối ưu hóa prompt cho các mô hình AI khác nhau}
\end{question}

% ==============================================================================
% Question 111 | Phân loại: 3. Generative AI: Prompt Engineering Basics
% Xuất hiện trong các đề gốc: AIG202c_SP_2025_FE (Q46)
% ==============================================================================
\begin{question}{111}{mcq}
\source{AIG202c_-_SP_2025_-_FE_3755_46.webp}
\prompt{What is the purpose of evaluating prompt effectiveness in prompt engineering for generative AP}
\choice{A}{To determine the computational resources required for generating model outputs}
\choice{B}{To assess the relevance and coherence of model outputs generated from different prompts}
\choice{C}{To randomly select prompts for generating model outputs}
\choice{D}{To fine-tune prompts based on feedback from model outputs}
\answer{B}
\confidence{high}
\evidence{Mục đích của việc đánh giá hiệu quả câu lệnh (evaluating prompt effectiveness) là đo lường mức độ phù hợp và tính mạch lạc (relevance and coherence) của các kết quả do mô hình sinh ra từ các prompt khác nhau.}
\sources{LangChain Evaluation Standards; Prompt Engineering Guide}
\explain{Hướng giải: Xác định mục tiêu của việc đánh giá hiệu quả câu lệnh (câu hỏi lặp lại của UQ043).

Các bước phân tích:
1. Đánh giá prompt nhằm mục đích đo lường định tính và định lượng xem câu lệnh có mang lại kết quả đúng như mong đợi không.
2. Hai tiêu chí quan trọng nhất là Relevance (Độ liên quan/phù hợp với yêu cầu) và Coherence (Tính mạch lạc, logic của văn bản sinh ra).

Kiểm tra: Khớp hoàn toàn với đáp án của câu UQ043.

Lưu ý/Khái niệm cần nhớ: Đánh giá câu lệnh tập trung vào tính phù hợp (relevance) và mạch lạc (coherence).}
\vocab{Prompt Evaluation | Đánh giá câu lệnh | Quá trình kiểm thử và đo lường chất lượng đầu ra của các câu lệnh khác nhau}
\end{question}

% ==============================================================================
% Question 112 | Phân loại: 3. Generative AI: Prompt Engineering Basics
% Xuất hiện trong các đề gốc: AIG202c_SP_2025_FE (Q50)
% ==============================================================================
\begin{question}{112}{mcq}
\source{AIG202c_-_SP_2025_-_FE_3755_50.webp}
\prompt{How can prompt engineenng be used to control the creativity of generative AI models?}
\choice{A}{By limiting the vocabulary used in prompts}
\choice{B}{By adjusting the temperature parameter during model generation}
\choice{C}{By Ignoring prompts that may lead to overly creative or inappropriate responses}
\choice{D}{By randomly selecting prompts from a predefined list}
\answer{B}
\confidence{high}
\evidence{Tham số nhiệt độ (temperature parameter) trong các mô hình AI tạo sinh điều khiển độ ngẫu nhiên của phân phối xác suất sinh từ: tăng nhiệt độ giúp câu trả lời sáng tạo và đa dạng hơn, trong khi giảm nhiệt độ khiến mô hình thận trọng và tất định hơn.}
\sources{Goodfellow et al., Deep Learning (Softmax Temperature); OpenAI API Reference: Temperature}
\explain{Hướng giải: Xác định cơ chế kỹ thuật trong câu lệnh/suy luận để điều khiển mức độ sáng tạo của GenAI.

Các bước phân tích:
1. Hàm kích hoạt Softmax có tham số nhiệt độ \(T\):
   \[P(w_i) = \frac{\exp(z_i / T)}{\sum_j \exp(z_j / T)}\]
2. Khi \(T \to 0\): Mô hình trở nên tất định (Argmax), luôn chọn từ có xác suất cao nhất, câu trả lời an toàn, ít sáng tạo.
3. Khi \(T\) tăng cao (ví dụ \(T = 0.8 - 1.2\)): Phân phối xác suất trở nên phẳng hơn, các từ ít phổ biến có cơ hội được chọn, kích thích sự sáng tạo và bất ngờ trong văn phong.
4. Do đó, điều chỉnh tham số temperature là công cụ trực tiếp để kiểm soát tính sáng tạo của mô hình.

Kiểm tra: Mọi giao diện tương tác LLM (OpenAI Playground, watsonx Prompt Lab) đều sử dụng thanh trượt Temperature để chỉnh Creativity.

Vì sao các phương án khác sai:
- A: Giới hạn từ vựng làm nghèo nàn ngôn ngữ, không điều khiển được tính sáng tạo bản chất.
- C & D: Bỏ qua prompt hoặc chọn ngẫu nhiên không phải là phương pháp kỹ thuật chuẩn mực.

Lưu ý/Khái niệm cần nhớ: Temperature điều khiển trực tiếp mức độ ngẫu nhiên và sáng tạo của mô hình ngôn ngữ.}
\vocab{Temperature Tuning | Tinh chỉnh tham số nhiệt độ | Kỹ thuật kiểm soát tính ngẫu nhiên và sáng tạo của mô hình sinh ngôn ngữ}
\end{question}

% ==============================================================================
% Question 113 | Phân loại: 3. Generative AI: Prompt Engineering Basics
% Xuất hiện trong các đề gốc: AIG202c_SP_2025_FE (Q51)
% ==============================================================================
\begin{question}{113}{mcq}
\source{AIG202c_-_SP_2025_-_FE_3755_51.webp}
\prompt{What role does human input play in prompt engineering?}
\choice{A}{Human input is not necessary for prompt engineering}
\choice{B}{Human input is essential for designing effective prompts and evaluating model outputs}
\choice{C}{Human input is only required during the training phase of AI models}
\choice{D}{Human input is necessary for adjusting hyperparameters during model generation}
\answer{B}
\confidence{high}
\evidence{Sự can thiệp và đóng góp của con người (Human input) là thiết yếu để thiết kế các câu lệnh hiệu quả và đánh giá chất lượng kết quả đầu ra của mô hình (essential for designing effective prompts and evaluating model outputs).}
\sources{Human-in-the-Loop AI Systems; Prompt Engineering Best Practices}
\explain{Hướng giải: Đánh giá vai trò của con người trong quy trình kỹ nghệ câu lệnh (câu hỏi kiểm tra tính nhất quán với UQ084).

Các bước phân tích:
1. Prompt Engineering đòi hỏi tư duy ngôn ngữ, sự am hiểu nghiệp vụ và khả năng đánh giá logic của con người.
2. Con người đóng vai trò xây dựng kịch bản prompt và kiểm duyệt kết quả sinh ra để đảm bảo tính an toàn và chuẩn xác.

Kiểm tra: Khớp hoàn toàn với đáp án của câu UQ084.

Lưu ý/Khái niệm cần nhớ: Con người là trung tâm định hướng và đánh giá chất lượng trong Prompt Engineering.}
\vocab{Human-in-the-Loop (HITL) | Con người trong vòng lặp | Mô hình tích hợp sự giám sát và can thiệp của con người vào quy trình AI}
\end{question}

% ==============================================================================
% Question 114 | Phân loại: 3. Generative AI: Prompt Engineering Basics
% Xuất hiện trong các đề gốc: AIG202c_SU25_FE (Q31)
% ==============================================================================
\begin{question}{114}{mcq}
\source{Q31.webp}
\prompt{Which of the following is a potential challenge associated with prompt engineering for generative AI?}
\choice{A}{Limited control over the output of generative AI models}
\choice{B}{Difficulty in training generative AI models without prompts}
\choice{C}{Lack of available tools and resources for prompt engineering}
\choice{D}{Difficulty in generating prompts automatically using AI algorithms}
\answer{A}
\confidence{high}
\evidence{Thách thức tiềm ẩn trong Prompt Engineering cho AI tạo sinh là sự kiểm soát có giới hạn đối với kết quả đầu ra của mô hình (Limited control over the output of generative AI models).}
\sources{Prompt Engineering Limitations and Reliability; Stochastic Behavior of LLMs}
\explain{Hướng giải: Xác định thách thức kỹ thuật trong việc điều khiển mô hình AI tạo sinh (câu hỏi kiểm tra tính nhất quán với UQ124).

Các bước phân tích:
1. Mô hình ngôn ngữ lớn có bản chất xác suất ngẫu nhiên.
2. Rất khó để đảm bảo 100% đầu ra luôn cố định từng chữ trong mọi trường hợp, đòi hỏi nhiều nỗ lực kiểm thử và ràng buộc cấu trúc.

Kiểm tra: Khớp hoàn toàn với đáp án của câu UQ124.

Lưu ý/Khái niệm cần nhớ: Kiểm soát đầu ra có giới hạn là thách thức bản chất của mô hình xác suất GenAI.}
\vocab{Model Stochasticity | Tính ngẫu nhiên của mô hình | Đặc tính sinh kết quả dựa trên phân phối xác suất thống kê}
\end{question}

% ==============================================================================
% Question 115 | Phân loại: 3. Generative AI: Prompt Engineering Basics
% Xuất hiện trong các đề gốc: AIG202c_SU25_FE (Q32), AIG202c_SP26_RE (Q35)
% ==============================================================================
\begin{question}{115}{mcq}
\source{Q32.webp}
\prompt{Why is prompt engineering important in the development of generative AI models?}
\choice{A}{Prompt engineering ensures that generative AI models generate output that is relevant and coherent}
\choice{B}{Prompt engineering helps reduce the computational resources required for training generative AI models}
\choice{C}{Prompt engineering allows generative AI models to generate output without any human intervention}
\choice{D}{Prompt engineering is not important for generative AI models}
\answer{A}
\confidence{high}
\evidence{Kỹ nghệ câu lệnh đóng vai trò tối quan trọng trong việc phát triển AI tạo sinh vì nó bảo đảm mô hình sinh ra kết quả phù hợp, chính xác và mạch lạc (relevant and coherent) theo đúng mục đích sử dụng.}
\sources{OpenAI Best Practices for Prompt Engineering; Coursera Prompt Engineering Specialization}
\explain{Hướng giải: Giải thích tầm quan trọng sống còn của Prompt Engineering trong ứng dụng Generative AI.

Các bước phân tích:
1. Mô hình AI tạo sinh rất mạnh nhưng có tính ngẫu nhiên và dễ bị lạc đề hoặc ảo giác (hallucination) nếu nhận câu lệnh mơ hồ.
2. Prompt Engineering là cầu nối chuyển dịch mong muốn của con người thành các chỉ dẫn toán học/logic chính xác để mô hình phản hồi đúng trọng tâm, logic và có cấu trúc mạch lạc.

Kiểm tra: Đây là luận điểm nền tảng của toàn bộ chương trình đào tạo Generative AI.

Vì sao các phương án khác sai:
- B: Prompt engineering ở tầng suy luận, không trực tiếp giảm tài nguyên tính toán của khâu tiền huấn luyện (pre-training).
- C: Prompt engineering là sự can thiệp trực tiếp của con người, không phải sinh không cần con người.
- D: Cho rằng prompt engineering không quan trọng là sai hoàn toàn.

Lưu ý/Khái niệm cần nhớ: Prompt engineering bảo đảm kết quả sinh ra đúng trọng tâm và mạch lạc (relevant and coherent).}
\vocab{Coherence | Tính mạch lạc | Khả năng liên kết logic, liền mạch và dễ hiểu của văn bản do AI sinh ra}
\end{question}

% ==============================================================================
% Question 116 | Phân loại: 3. Generative AI: Prompt Engineering Basics
% Xuất hiện trong các đề gốc: AIG202c_SU25_FE (Q33)
% ==============================================================================
\begin{question}{116}{mcq}
\source{Q33.webp}
\prompt{What is the first step in writing a well-structured prompt through the process of prompt engineering? (choose two)}
\choice{A}{Refining the prompt based on testing and analysis}
\choice{B}{Defining the goal}
\choice{C}{Testing the prompt for response quality}
\choice{D}{Analyzing responses from the generative AI model}
\answer{B, C}
\confidence{high}
\evidence{Trong quy trình chuẩn của kỹ nghệ câu lệnh (Prompt Engineering Process), hai bước khởi đầu quan trọng nhất để tạo ra một câu lệnh có cấu trúc tốt là: Xác định rõ mục tiêu cần đạt được (Defining the goal) và Kiểm thử câu lệnh để đánh giá chất lượng phản hồi ban đầu (Testing the prompt for response quality).}
\sources{IBM watsonx Prompt Engineering Methodology: Define, Test, Refine, Deploy}
\explain{Hướng giải: Xác định hai bước đầu tiên trong quy trình xây dựng prompt có cấu trúc tốt.

Các bước phân tích:
1. Quy trình Prompt Engineering chuẩn mực gồm các bước tuần tự:
   - Bước 1: Defining the goal (Xác định mục tiêu: Ta muốn mô hình làm gì, giải quyết bài toán nào, định dạng mong muốn là gì?).
   - Bước 2: Drafting & Testing the prompt (Soạn thảo prompt ban đầu và kiểm thử phản hồi của mô hình để xem nó hoạt động thế nào).
   - Bước 3: Analyzing responses (Phân tích các phản hồi để tìm lỗi sai hoặc điểm chưa đạt).
   - Bước 4: Refining the prompt (Tinh chỉnh lại prompt dựa trên kết quả phân tích).
2. Hai bước khởi đầu cấu thành giai đoạn thiết lập ban đầu chính là: B (Defining the goal) và C (Testing the prompt for response quality).

Kiểm tra: Xác định mục tiêu rõ ràng là tiền đề bắt buộc trước khi bắt tay vào viết bất kỳ câu lệnh nào.

Vì sao các phương án khác ở bước sau:
- D (Analyzing responses) diễn ra ở bước 3 sau khi đã thử nghiệm.
- A (Refining the prompt) là bước tinh chỉnh cuối cùng trong chu trình lặp.

Lưu ý/Khái niệm cần nhớ: Quy trình Prompt: Xác định mục tiêu (Define Goal) -> Thử nghiệm (Test) -> Phân tích (Analyze) -> Tinh chỉnh (Refine).}
\vocab{Goal Definition | Xác định mục tiêu | Bước đầu tiên làm rõ yêu cầu nghiệp vụ và kết quả đầu ra mong đợi của prompt}
\end{question}

% ==============================================================================
% Question 117 | Phân loại: 3. Generative AI: Prompt Engineering Basics
% Xuất hiện trong các đề gốc: AIG202c_SU25_FE (Q34)
% ==============================================================================
\begin{question}{117}{mcq}
\source{Q34.webp}
\prompt{Why is clarity important when writing prompts for generative AI models? (choose Wo)}
\choice{A}{Clarity helps make the prompt less engaging}
\choice{B}{Clarity helps the model understand the task and produce relevant responses}
\choice{C}{Clarity adds to the complexity of the prompt}
\choice{D}{Clarity ensures the prompt is lengthy}
\answer{B, D}
\confidence{high}
\evidence{Tính rõ ràng (clarity) trong câu lệnh rất quan trọng vì nó giúp mô hình hiểu chính xác nhiệm vụ để tạo phản hồi đúng trọng tâm (helps model understand task and produce relevant responses) đồng thời loại bỏ sự mơ hồ, tránh hiểu lầm trong quá trình sinh văn bản.}
\sources{OpenAI Prompt Engineering Best Practices: Write clear instructions; Anthropic Prompt Guide}
\explain{Hướng giải: Xác định hai lý do tại sao tính rõ ràng (Clarity) lại mang tính sống còn khi viết prompt cho GenAI.

Các bước phân tích:
1. Mô hình ngôn ngữ phụ thuộc hoàn toàn vào ngữ cảnh được cung cấp trong câu lệnh.
2. Sự rõ ràng giúp:
   - (B) Mô hình hiểu chính xác mục tiêu của bài toán, từ đó kích hoạt đúng không gian tri thức và sinh ra câu trả lời có tính liên quan cao (relevant responses).
   - (D) Tránh được sự mơ hồ (ambiguity), hạn chế nguy cơ mô hình suy diễn sai lệch hoặc bịa đặt thông tin (hallucination).

Kiểm tra: Nguyên tắc số 1 trong tài liệu hướng dẫn của OpenAI và Anthropic luôn là: 'Be clear and specific'.

Vì sao các phương án khác sai:
- A: Rõ ràng không làm prompt kém hấp dẫn, ngược lại giúp tăng tính hiệu quả.
- C: Các phương án gây nhiễu khác trái ngược với mục tiêu giao tiếp rõ ràng.

Lưu ý/Khái niệm cần nhớ: Clarity = Hiểu đúng tác vụ + Triệt tiêu sự mơ hồ.}
\vocab{Clarity in Prompting | Tính rõ ràng trong câu lệnh | Việc sử dụng ngôn từ minh bạch, chính xác và không gây hiểu lầm trong prompt}
\end{question}

% ==============================================================================
% Question 118 | Phân loại: 3. Generative AI: Prompt Engineering Basics
% Xuất hiện trong các đề gốc: AIG202c_SU25_FE (Q35), AIG202c_FA25_FE (Q40), AIG202c_SU26_FE (Q34)
% ==============================================================================
\begin{question}{118}{mcq}
\source{Q35.webp}
\prompt{Which of the following is NOT a potential challenge associated with prompt engineering?}
\choice{A}{Limited control over the output of AI models}
\choice{B}{Difficulty in generating prompts automatically using AI algorithms}
\choice{C}{Lack of available tools and resources for prompt engineering}
\choice{D}{Difficulty in training models using only prompts that produce desired responses}
\answer{D}
\confidence{high}
\evidence{Prompt Engineering KHÔNG PHẢI là quá trình huấn luyện mô hình (training models), mà là kỹ thuật tương tác với mô hình đã huấn luyện sẵn. Do đó 'Khó khăn trong việc huấn luyện mô hình...' không phải là thách thức của prompt engineering.}
\sources{Stanford CS224N: Natural Language Processing with Deep Learning; Prompt Engineering Fundamentals}
\explain{Hướng giải: Xác định thách thức KHÔNG liên quan đến Kỹ nghệ câu lệnh (Prompt Engineering).

Các bước phân tích:
1. Phân biệt rõ: 'Huấn luyện mô hình' (Model Training) là quá trình cập nhật trọng số nơ-ron bằng dữ liệu lớn và thuật toán lan truyền ngược (backpropagation).
2. Prompt Engineering chỉ diễn ra ở giai đoạn suy luận (Inference time), hoàn toàn không huấn luyện hay thay đổi trọng số của mô hình.
3. Do đó, 'Difficulty in training models using only prompts...' là một phát biểu sai bản chất, không phải thách thức của prompt engineering.

Kiểm tra: Các thách thức thực tế của prompt engineering là: khó kiểm soát tính ngẫu nhiên của đầu ra (A), khó tự động hóa sinh prompt hoàn hảo (B), và sự phụ thuộc vào công cụ hỗ trợ (C).

Vì sao các phương án khác là thách thức thực tế:
- A: Đúng, LLM có tính chất sinh ngẫu nhiên nên khó kiểm soát 100% kết quả.
- B: Đúng, việc tạo prompt tự động thích ứng với mọi tác vụ là bài toán khó.
- C: Đúng, hệ sinh thái công cụ quản trị prompt vẫn đang trong giai đoạn hoàn thiện.

Lưu ý/Khái niệm cần nhớ: Prompt Engineering hoàn toàn tách biệt với Model Training.}
\vocab{Inference Time | Thời điểm suy luận | Giai đoạn mô hình AI sử dụng tri thức đã học để dự đoán mà không cập nhật trọng số}
\end{question}

% ==============================================================================
% Question 119 | Phân loại: 3. Generative AI: Prompt Engineering Basics
% Xuất hiện trong các đề gốc: AIG202c_SU25_FE (Q36), AIG202c_FA25_FE (Q34), AIG202c_SP26_RE (Q33)
% ==============================================================================
\begin{question}{119}{mcq}
\source{Q36.webp}
\prompt{Which of the following is the main purpose of using the user feedback loop?}
\choice{A}{To iteratively refine text prompts based on the response generated by the LLM}
\choice{B}{To generate responses with examples}
\choice{C}{To provide explicit instructions to generate neutral responses}
\choice{D}{To generate meaningful responses without needing prior training on specific prompts}
\answer{A}
\confidence{high}
\evidence{Vòng lặp phản hồi của người dùng (User feedback loop) trong Prompt Engineering được thiết kế nhằm mục đích tinh chỉnh lặp đi lặp lại các câu lệnh (iteratively refine text prompts) dựa trên kết quả mà LLM sinh ra.}
\sources{Prompt Engineering Guide (DAIR.AI); Human-in-the-Loop Machine Learning}
\explain{Hướng giải: Xác định mục đích chính của vòng lặp phản hồi người dùng (user feedback loop) trong kỹ nghệ câu lệnh cho LLM.

Các bước phân tích:
1. Phản hồi của người dùng cung cấp thông tin đánh giá về chất lượng đầu ra của mô hình (độ chính xác, phong cách, độ dài, định dạng).
2. Dựa trên phản hồi này, người thiết kế câu lệnh sẽ điều chỉnh, bổ sung ràng buộc hoặc cấu trúc lại prompt để đạt kết quả tốt hơn trong các lượt sinh tiếp theo.

Kiểm tra: Đây là định nghĩa cốt lõi của Human-in-the-Loop Prompt Engineering.

Vì sao các phương án khác sai:
- B: Cung cấp ví dụ là kỹ thuật Few-shot prompting.
- C: Hướng dẫn trung lập là Bias mitigation.
- D: Sinh phản hồi không cần huấn luyện là đặc tính Zero-shot của LLM.

Lưu ý/Khái niệm cần nhớ: User feedback loop = Vòng lặp tinh chỉnh câu lệnh liên tục dựa trên đánh giá của người dùng.}
\vocab{Feedback Loop | Vòng lặp phản hồi | Cơ chế đưa đầu ra và đánh giá ngược lại đầu vào để cải tiến hệ thống}
\end{question}

% ==============================================================================
% Question 120 | Phân loại: 3. Generative AI: Prompt Engineering Basics
% Xuất hiện trong các đề gốc: AIG202c_SU25_FE (Q41)
% ==============================================================================
\begin{question}{120}{mcq}
\source{Q41.webp}
\prompt{Which of the following is NOT a common technique used in prompt engineering for generative AI?}
\choice{A}{Designing prompts to elicit specific types of responses from the model}
\choice{B}{Adjusting the temperature parameter during model generation}
\choice{C}{Randomly selecting prompts from a predefined list}
\choice{D}{Fine-tuning prompts based on feedback from model outputs}
\answer{C}
\confidence{high}
\evidence{Việc lựa chọn ngẫu nhiên các câu lệnh từ một danh sách định sẵn (Randomly selecting prompts from a predefined list) KHÔNG PHẢI là một kỹ thuật kỹ nghệ câu lệnh (NOT a common technique); prompt engineering là quá trình có chủ ý và phương pháp khoa học, không dựa trên sự lựa chọn ngẫu nhiên thiếu kiểm soát.}
\sources{Prompt Engineering Principles; Best Practices in GenAI Application Development}
\explain{Hướng giải: Tìm phương án KHÔNG PHẢI là kỹ thuật Kỹ nghệ câu lệnh phổ biến.

Các bước phân tích:
1. Prompt Engineering áp dụng các phương pháp có hệ thống:
   - Thiết kế prompt định hướng đầu ra cụ thể (A).
   - Điều chỉnh tham số nhiệt độ temperature (B).
   - Tinh chỉnh prompt dựa trên phản hồi của mô hình (D).
2. Hành vi 'Randomly selecting prompts from a predefined list' (Chọn ngẫu nhiên prompt) hoàn toàn không có cơ sở khoa học hay tối ưu hóa, do đó không phải là một kỹ thuật prompt engineering.

Kiểm tra: Kỹ thuật luôn đòi hỏi tính chủ đích và kiểm soát được.

Vì sao các phương án khác là kỹ thuật chuẩn:
- A, B, D đều là các phương pháp kỹ nghệ câu lệnh được sử dụng hàng ngày.

Lưu ý/Khái niệm cần nhớ: Prompt Engineering là khoa học thiết kế có chủ ý, không chọn ngẫu nhiên.}
\vocab{Structured Prompting | Kỹ nghệ câu lệnh có cấu trúc | Việc áp dụng các nguyên tắc kỹ thuật chặt chẽ để thiết kế câu lệnh}
\end{question}

% ==============================================================================
% Question 121 | Phân loại: 3. Generative AI: Prompt Engineering Basics
% Xuất hiện trong các đề gốc: AIG202c_FA25_FE (Q17)
% ==============================================================================
\begin{question}{121}{mcq}
\source{Q17.webp}
\prompt{Which challenges did Live Person specifically identify in fostering positive interactions between customers and companies?}
\choice{A}{People are turning to self-service and messaging as a preferred way to engage with customer service.}
\choice{B}{Legacy voice channels are frustrating for the customer, expensive, and don't necessarily work well.}
\choice{C}{Frustrating experiences interacting with customers lead to high agent attrition.}
\choice{D}{Consumers prefer to interact with customer service agents by email or by voice.}
\answer{C}
\confidence{high}
\evidence{Trong các mô hình sinh hình ảnh như DALL-E hay Midjourney, việc thêm các chi tiết cụ thể về phong cách ánh sáng, bố cục và chất liệu vào câu lệnh (prompt) giúp định hình chính xác phong cách thẩm mỹ của hình ảnh sinh ra.}
\sources{OpenAI DALL-E Prompt Engineering Guide; Midjourney Documentation}
\explain{Hướng giải: Hiểu tác động của việc thiết kế câu lệnh chi tiết trong mô hình sinh hình ảnh (Text-to-Image Generative AI).

Các bước phân tích:
1. Mô hình sinh ảnh hoạt động dựa trên việc điều hướng phân phối xác suất trong không gian tiềm ẩn (latent space) dựa trên các từ khóa (tokens) của prompt.
2. Cung cấp các thuộc tính nghệ thuật cụ thể như ánh sáng, góc máy, chất liệu giúp mô hình thu hẹp không gian sinh mẫu, tạo ra hình ảnh đúng phong cách mong muốn.

Kiểm tra: Prompt engineering cho ảnh luôn đòi hỏi chỉ định rõ phong cách hội họa hoặc kỹ thuật nhiếp ảnh.

Vì sao các phương án khác sai:
- Mô hình không bị giới hạn số lượng màu hay giảm chất lượng khi mô tả chi tiết, ngược lại càng chi tiết càng kiểm soát tốt chất lượng.

Lưu ý/Khái niệm cần nhớ: Kỹ thuật prompt chi tiết là chìa khóa định hình phong cách ảnh sinh ra từ GenAI.}
\vocab{Prompt Engineering | Kỹ thuật kỹ nghệ câu lệnh | Quá trình thiết kế và tối ưu hóa văn bản đầu vào để mô hình sinh ra kết quả mong muốn}
\end{question}

% ==============================================================================
% Question 122 | Phân loại: 3. Generative AI: Prompt Engineering Basics
% Xuất hiện trong các đề gốc: AIG202c_FA25_FE (Q31), AIG202c_SP26_RE (Q45)
% ==============================================================================
\begin{question}{122}{mcq}
\source{Q31.webp}
\prompt{Imagine you are a content developer working with LLMs, and you must ensure that the responses generated are indiscriminatory and neutral. Which among the following techniques would you employ with your text prompts to instruct the model appropriately?}
\choice{A}{Bias mitigation}
\choice{B}{Few-shot prompting}
\choice{C}{Zero-shot prompting}
\choice{D}{Contextual guidance}
\answer{A}
\confidence{high}
\evidence{Để đảm bảo phản hồi của mô hình ngôn ngữ lớn (LLM) không mang tính phân biệt đối xử và giữ thái độ trung lập, kỹ thuật cần áp dụng là Giảm thiểu định kiến (Bias mitigation) thông qua các câu lệnh định hướng đạo đức và công bằng.}
\sources{OpenAI Safety & Alignment Guidelines; Anthropic Constitutional AI}
\explain{Hướng giải: Xác định kỹ thuật trong câu lệnh nhằm ngăn chặn định kiến và bảo đảm tính trung lập của mô hình ngôn ngữ.

Các bước phân tích:
1. Khi làm việc với LLMs, nhà phát triển cần kiểm soát rủi ro mô hình phản hồi thiên vị hoặc phân biệt đối xử.
2. Kỹ thuật 'Bias mitigation' (Giảm thiểu định kiến) bao gồm việc thiết lập các ràng buộc trong prompt (system instructions) yêu cầu mô hình duy trì góc nhìn khách quan, đa chiều và loại bỏ các định kiến giới tính, chủng tộc, tôn giáo.

Kiểm tra: Khái niệm 'Bias mitigation' là tiêu chuẩn vàng trong thiết kế AI có trách nhiệm (Responsible AI).

Vì sao các phương án khác sai:
- B (Few-shot prompting): Cung cấp một vài ví dụ mẫu để mô hình bắt chước định dạng đầu ra, không trực tiếp nhắm vào tính trung lập đạo đức.
- C (Zero-shot prompting): Ra lệnh trực tiếp không kèm ví dụ, dễ khiến mô hình bộc lộ thiên kiến mặc định từ dữ liệu tiền huấn luyện.
- D (Contextual guidance): Cung cấp bối cảnh chung, không tập trung vào loại bỏ định kiến xã hội.

Lưu ý/Khái niệm cần nhớ: Bias mitigation = Kỹ thuật triệt tiêu định kiến và đảm bảo tính công bằng trong AI.}
\vocab{Bias Mitigation | Giảm thiểu thiên kiến | Các phương pháp kỹ thuật nhằm hạn chế và loại bỏ sự thiên vị trong mô hình AI}
\end{question}

% ==============================================================================
% Question 123 | Phân loại: 3. Generative AI: Prompt Engineering Basics
% Xuất hiện trong các đề gốc: AIG202c_FA25_FE (Q33)
% ==============================================================================
\begin{question}{123}{mcq}
\source{Q33.webp}
\prompt{What is the purpose of fine-tuning prompts based on model outputs?}
\choice{A}{To generate random responses for model evaluation}
\choice{B}{To train the model using only prompts that produce desired responses}
\choice{C}{To adjust prompts based on user feedback to improve model performance}
\choice{D}{To ignore prompts and allow the model to generate output autonomously}
\answer{C}
\confidence{high}
\evidence{Mục đích của việc tinh chỉnh câu lệnh (fine-tuning prompts) dựa trên đầu ra của mô hình là điều chỉnh các chỉ dẫn dựa trên phản hồi của người dùng hoặc kết quả thực tế để nâng cao hiệu suất và độ chính xác của mô hình.}
\sources{Prompt Engineering Guide (DAIR.AI); OpenAI Best Practices for Prompt Engineering}
\explain{Hướng giải: Xác định mục tiêu của quy trình tinh chỉnh câu lệnh lặp (iterative prompt refinement).

Các bước phân tích:
1. Prompt Engineering là một quy trình mang tính lặp (iterative process): Viết prompt ban đầu -> Chạy thử -> Đánh giá đầu ra -> Tinh chỉnh prompt dựa trên các điểm chưa đạt -> Thử lại.
2. Mục đích là điều chỉnh cấu trúc, từ ngữ và ràng buộc dựa trên phản hồi để mô hình sinh ra kết quả ngày càng chính xác, phù hợp ngữ cảnh.

Kiểm tra: Khái niệm 'Adjust prompts based on user feedback to improve model performance' phản ánh chính xác bản chất của vòng lặp cải tiến prompt.

Vì sao các phương án khác sai:
- A: Không ai điều chỉnh prompt để tạo phản hồi ngẫu nhiên vô nghĩa.
- B: Không thể huấn luyện lại toàn bộ mô hình (train the model) chỉ thông qua vài câu lệnh mà không can thiệp trọng số.
- D: Tinh chỉnh prompt là để kiểm soát mô hình tốt hơn, không phải để bỏ mặc mô hình tự do.

Lưu ý/Khái niệm cần nhớ: Tinh chỉnh câu lệnh là quá trình tối ưu hóa liên tục dựa trên kết quả đầu ra thực tế.}
\vocab{Iterative Prompting | Tinh chỉnh câu lệnh lặp | Quy trình thử nghiệm và điều chỉnh prompt liên tục để tối ưu hóa kết quả}
\end{question}

% ==============================================================================
% Question 124 | Phân loại: 3. Generative AI: Prompt Engineering Basics
% Xuất hiện trong các đề gốc: AIG202c_FA25_FE (Q35), AIG202c_SP26_RE (Q38)
% ==============================================================================
\begin{question}{124}{mcq}
\source{Q35.webp}
\prompt{How does incorporating the persona pattern dimension enhance prompt effectiveness? Select the best option.}
\choice{A}{By increasing the number of words in the prompt}
\choice{B}{By introducing a sense of character-driven context and perspective}
\choice{C}{By including a series of questions through multiple prompts}
\choice{D}{By focusing solely on generic content without historical or fictional references}
\answer{B}
\confidence{high}
\evidence{Kỹ thuật Persona Pattern (mẫu hình tượng nhân vật) nâng cao hiệu quả của prompt bằng cách đưa vào một bối cảnh và góc nhìn do nhân vật dẫn dắt (character-driven context and perspective), giúp mô hình đóng vai chuyên gia cụ thể.}
\sources{Vanderbilt University Prompt Engineering Coursera (White et al., A Prompt Pattern Catalog)}
\explain{Hướng giải: Hiểu cơ chế hoạt động của mẫu Persona (đóng vai nhân vật) trong Prompt Engineering.

Các bước phân tích:
1. Persona Pattern yêu cầu mô hình đóng một vai trò cụ thể (ví dụ: 'Hãy đóng vai một luật sư chuyên về sở hữu trí tuệ', 'Hãy đóng vai một giáo sư vật lý hạt nhân').
2. Kỹ thuật này giúp LLM kích hoạt các vùng tri thức và phong cách hành văn chuyên ngành tương ứng, tạo ra câu trả lời có chiều sâu và ngữ cảnh nhất quán.

Kiểm tra: Báo cáo kinh điển của GS. Jules White về Prompt Patterns định nghĩa Persona là mẫu gán vai trò nhân vật để định hướng văn phong và tri thức.

Vì sao các phương án khác sai:
- A: Tăng số lượng từ không phải mục đích nâng cao chất lượng của Persona.
- C: Đặt nhiều câu hỏi là Question Refinement Pattern.
- D: Tập trung vào nội dung chung chung là đi ngược lại mục đích chuyên biệt hóa của Persona.

Lưu ý/Khái niệm cần nhớ: Persona Pattern = Định hướng góc nhìn và văn phong mô hình bằng cách gán vai trò nhân vật cụ thể.}
\vocab{Persona Pattern | Mẫu nhân vật | Kỹ thuật chỉ định vai trò hoặc danh tính cụ thể cho AI khi trả lời câu hỏi}
\end{question}

% ==============================================================================
% Question 125 | Phân loại: 3. Generative AI: Prompt Engineering Basics
% Xuất hiện trong các đề gốc: AIG202c_FA25_FE (Q36), AIG202c_SU26_FE (Q32)
% ==============================================================================
\begin{question}{125}{mcq}
\source{Q36.webp}
\prompt{Which of the following prompt engineering approaches involves breaking down complex tasks into easier ones through a sequence of simpler prompts to guide the model toward the intended outcome?}
\choice{A}{Cognitive building}
\choice{B}{Chain-of-Thought}
\choice{C}{Tree-of-Thoughts}
\choice{D}{Interview pattern}
\answer{B}
\confidence{high}
\evidence{Chain-of-Thought (Chuỗi tư duy) là phương pháp prompt engineering chia nhỏ bài toán phức tạp thành các bước suy luận trung gian đơn giản hơn để hướng dẫn mô hình đi đến kết quả chính xác.}
\sources{Wei et al. (2022), Chain-of-Thought Prompting Elicits Reasoning in Large Language Models (NeurIPS)}
\explain{Hướng giải: Xác định kỹ thuật chia nhỏ tác vụ phức tạp thành chuỗi các bước lập luận đơn giản.

Các bước phân tích:
1. Chain-of-Thought (CoT) khuyến khích mô hình 'suy nghĩ từng bước' (Let\'s think step by step).
2. Bằng cách phân rã bài toán phức tạp thành các mắt xích suy luận tuần tự, mô hình tránh được các suy đoán vội vã và giảm thiểu lỗi logic.

Kiểm tra: Nghiên cứu của Google Research (Jason Wei et al.) chứng minh CoT cải thiện vượt bậc khả năng giải toán và suy luận logic của LLMs.

Vì sao các phương án khác sai:
- A (Cognitive building): Thuật ngữ không chính thức trong danh mục prompt pattern.
- C (Tree-of-Thoughts): Là kỹ thuật phân nhánh dạng cây khám phá nhiều phương án song song, phức tạp hơn chuỗi tuần tự.
- D (Interview pattern): Kỹ thuật AI phỏng vấn người dùng để thu thập thông tin.

Lưu ý/Khái niệm cần nhớ: Chain-of-Thought (CoT) = Suy luận logic tuần tự từng bước (Step-by-step reasoning).}
\vocab{Chain-of-Thought (CoT) | Chuỗi suy nghĩ | Phương pháp hướng dẫn LLM thực hiện các bước lập luận trung gian trước khi đưa ra câu trả lời}
\end{question}

% ==============================================================================
% Question 126 | Phân loại: 3. Generative AI: Prompt Engineering Basics
% Xuất hiện trong các đề gốc: AIG202c_FA25_FE (Q37), AIG202c_SP26_RE (Q36)
% ==============================================================================
\begin{question}{126}{mcq}
\source{Q37.webp}
\prompt{Monica is a writer looking to improve her storytelling by using a large language model (LLM) for creative writing prompts. Which feature of a prompt engineering tool would be most beneficial for her to achieve this goal?}
\choice{A}{Iteratively refining prompts by adjusting them based on the initial outputs of the LLM's story}
\choice{B}{Providing feedback for prompts based on a given input.}
\choice{C}{Suggesting different prompt formats (like questions, dialogues, or scenarios) tailored to the preferred storytelling style (such as mystery, romance, or sci-fi)}
\choice{D}{Offering a library of pre-defined story beginnings designed for particular genres and themes}
\answer{C}
\confidence{high}
\evidence{Đối với một nhà văn sáng tạo nội dung, tính năng hữu ích nhất của công cụ prompt engineering là gợi ý các định dạng câu lệnh đa dạng (đối thoại, tình huống, câu hỏi) phù hợp với thể loại truyện cụ thể (trinh thám, lãng mạn, khoa học viễn tưởng).}
\sources{Creative Writing with Generative AI Guidelines; OpenAI Creative Applications}
\explain{Hướng giải: Lựa chọn tính năng hỗ trợ viết văn sáng tạo tối ưu nhất cho nhà văn từ công cụ kỹ nghệ câu lệnh.

Các bước phân tích:
1. Sáng tác truyện đòi hỏi sự đa dạng về thể loại và hình thức thể hiện (lời thoại, mô tả cảnh, xung đột kịch tính).
2. Tính năng gợi ý định dạng prompt được cá nhân hóa theo từng thể loại văn học (mystery, romance, sci-fi) giúp nhà văn khai mở ý tưởng và điều khiển phong cách truyện một cách linh hoạt nhất.

Kiểm tra: Các công cụ viết sáng tạo tích hợp AI (như Sudowrite, NovelAI) đều tập trung vào việc cung cấp các mẫu prompt chuyên sâu theo từng thể loại truyện.

Vì sao các phương án khác sai:
- A: Tinh chỉnh lặp chỉ là kỹ thuật chung, không chuyên biệt cho việc truyền cảm hứng đa dạng phong cách sáng tác.
- B: Phản hồi prompt chỉ đánh giá cú pháp, không kích thích sáng tạo văn học.
- D: Thư viện đoạn mở đầu có sẵn quá tĩnh và hạn chế tự do sáng tác của nhà văn.

Lưu ý/Khái niệm cần nhớ: Prompt đa dạng theo thể loại giúp phát huy tối đa năng lực sáng tạo của mô hình ngôn ngữ.}
\vocab{Creative Prompting | Câu lệnh sáng tạo | Kỹ thuật định hướng LLM sinh ra nội dung văn học, thơ ca giàu cảm xúc và hình tượng}
\end{question}

% ==============================================================================
% Question 127 | Phân loại: 3. Generative AI: Prompt Engineering Basics
% Xuất hiện trong các đề gốc: AIG202c_FA25_FE (Q38)
% ==============================================================================
\begin{question}{127}{mcq}
\source{Q38.webp}
\prompt{What is the primary goal of prompt engineering in the context of artificial intelligence?}
\choice{A}{To develop algorithms for training machine learning models}
\choice{B}{To optimize the architecture of neural networks}
\choice{C}{To design prompts or instructions to guide the output of AI models}
\choice{D}{To evaluate the performance of AI algorithms}
\answer{C}
\confidence{high}
\evidence{Mục tiêu cơ bản của Prompt Engineering trong AI là thiết kế các câu lệnh, chỉ dẫn hoặc ngữ cảnh đầu vào để định hướng và kiểm soát đầu ra của các mô hình AI.}
\sources{OpenAI Documentation: Prompt Engineering Guide; Anthropic Prompt Engineering Interactive Tutorial}
\explain{Hướng giải: Xác định định nghĩa và mục tiêu chính xác của Kỹ nghệ câu lệnh (Prompt Engineering).

Các bước phân tích:
1. Prompt Engineering không can thiệp vào mã nguồn thuật toán hay huấn luyện lại mô hình từ đầu.
2. Bản chất của nó là nghệ thuật và khoa học xây dựng các câu lệnh (prompts/instructions) tối ưu để khai thác tối đa tri thức có sẵn trong các mô hình nền tảng (Foundation Models).

Kiểm tra: Định nghĩa chuẩn: 'Prompt engineering is the process of structuring text that can be interpreted and understood by a generative AI model'.

Vì sao các phương án khác sai:
- A: Phát triển thuật toán huấn luyện là nhiệm vụ của ML Engineering / Research.
- B: Tối ưu kiến trúc mạng nơ-ron là Deep Learning Architecture Design.
- D: Đánh giá thuật toán là Model Benchmarking.

Lưu ý/Khái niệm cần nhớ: Prompt Engineering = Thiết kế câu lệnh định hướng đầu ra của mô hình AI.}
\vocab{Prompt Engineering | Kỹ nghệ câu lệnh | Kỹ thuật cấu trúc câu lệnh để tối ưu hóa chất lượng phản hồi từ mô hình AI}
\end{question}

% ==============================================================================
% Question 128 | Phân loại: 3. Generative AI: Prompt Engineering Basics
% Xuất hiện trong các đề gốc: AIG202c_FA25_FE (Q39), AIG202c_SU26_FE (Q39), AIG202c_SP26_RE (Q31)
% ==============================================================================
\begin{question}{128}{mcq}
\source{Q39.webp}
\prompt{How can prompt engineering help mitigate biases in generative AI models?}
\choice{A}{By ensuring that prompts are designed to encourage diverse and inclusive responses}
\choice{B}{By ignoring prompts that may lead to biased outputs}
\choice{C}{By training generative AI models using biased datasets}
\choice{D}{By generating prompts automatically using AI algorithms}
\answer{A}
\confidence{high}
\evidence{Prompt Engineering giúp giảm thiểu thiên kiến (bias) bằng cách thiết kế các câu lệnh chỉ dẫn mô hình đưa ra các phản hồi đa chiều, đa dạng và hòa nhập (diverse and inclusive responses).}
\sources{UNESCO AI Ethics Guidelines; Google Responsible AI Practices}
\explain{Hướng giải: Giải thích cơ chế dùng câu lệnh để kiểm soát và giảm thiểu thiên kiến trong Generative AI.

Các bước phân tích:
1. LLM vốn được huấn luyện trên dữ liệu Internet chứa nhiều thiên kiến xã hội.
2. Bằng cách lồng ghép các chỉ dẫn đạo đức cụ thể vào prompt (ví dụ: 'Hãy phân tích vấn đề này từ nhiều góc nhìn văn hóa khác nhau một cách khách quan, công bằng'), ta có thể định hướng mô hình tránh các định kiến hẹp hòi.

Kiểm tra: Đây là chiến lược căn bản trong kỹ thuật 'System Prompt Alignment'.

Vì sao các phương án khác sai:
- B: Bỏ qua prompt không giải quyết được nguồn gốc thiên kiến khi người dùng hỏi các chủ đề nhạy cảm.
- C: Huấn luyện trên dữ liệu thiên kiến càng làm gia tăng định kiến.
- D: Tự động sinh prompt không đảm bảo được tính đạo đức nếu không có sự can thiệp định hướng.

Lưu ý/Khái niệm cần nhớ: Chỉ dẫn đa dạng và hòa nhập trong prompt là công cụ hữu hiệu để hạn chế thiên kiến mô hình.}
\vocab{Inclusivity | Tính hòa nhập | Nguyên tắc đảm bảo công nghệ phục vụ và tôn trọng mọi nhóm người trong xã hội}
\end{question}

% ==============================================================================
% Question 129 | Phân loại: 3. Generative AI: Prompt Engineering Basics
% Xuất hiện trong các đề gốc: AIG202c_FA25_FE (Q41), AIG202c_SP26_RE (Q44)
% ==============================================================================
\begin{question}{129}{mcq}
\source{Q41.webp}
\prompt{How can prompt engineering be used to mitigate biases in AI models?}
\choice{A}{By generating prompts automatically using AI algorithms}
\choice{B}{By ensuring that prompts are designed to encourage diverse and inclusive responses}
\choice{C}{By ignoring prompts that may lead to biased outputs}
\choice{D}{By training models using biased datasets}
\answer{B}
\confidence{high}
\evidence{Kỹ nghệ câu lệnh giúp giảm định kiến trong các mô hình AI bằng cách đảm bảo các câu lệnh được thiết kế để khuyến khích các câu trả lời đa dạng và hòa nhập (diverse and inclusive responses).}
\sources{IBM AI Ethics Guidelines; Fairness in Large Language Models (ACM FAccT)}
\explain{Hướng giải: Xác định phương pháp ứng dụng prompt engineering để hạn chế thiên kiến AI.

Các bước phân tích:
1. Câu lệnh có thể định hình chuẩn mực ứng xử của AI.
2. Thiết kế prompt chủ động yêu cầu mô hình tôn trọng sự bình đẳng và đưa ra câu trả lời bao trùm các góc nhìn là giải pháp trực tiếp và hiệu quả.

Kiểm tra: Tương tự như câu UQ038, đây là nguyên lý cốt lõi của kỹ thuật Prompt-based Alignment.

Vì sao các phương án khác sai:
- A: Sinh prompt tự động không đảm bảo giảm định kiến.
- C: Bỏ qua prompt không giải quyết được vấn đề an toàn AI.
- D: Huấn luyện với dữ liệu thiên kiến là sai lầm nghiêm trọng.

Lưu ý/Khái niệm cần nhớ: Prompt công bằng và hòa nhập tạo ra phản hồi trung lập và đạo đức.}
\vocab{Fairness | Tính công bằng | Việc đảm bảo mô hình AI không tạo ra sự đối xử phân biệt đối với các nhóm đối tượng}
\end{question}

% ==============================================================================
% Question 130 | Phân loại: 3. Generative AI: Prompt Engineering Basics
% Xuất hiện trong các đề gốc: AIG202c_FA25_FE (Q44), AIG202c_SP26_RE (Q40)
% ==============================================================================
\begin{question}{130}{mcq}
\source{Q44.webp}
\prompt{What is the purpose of evaluating prompt effectiveness in prompt engineering for generative AI?}
\choice{A}{To determine the computational resources required for generating model outputs}
\choice{B}{To assess the relevance and coherence of model outputs generated from different prompts}
\choice{C}{To randomly select prompts for generating model outputs}
\choice{D}{To fine-tune prompts based on feedback from model outputs}
\answer{B}
\confidence{high}
\evidence{Mục đích cốt lõi của việc đánh giá hiệu quả câu lệnh (evaluating prompt effectiveness) là đo lường mức độ phù hợp và tính mạch lạc (relevance and coherence) của các phản hồi do mô hình sinh ra từ các câu lệnh khác nhau.}
\sources{LangChain Evaluation Guide; RAG Triad Evaluation Metrics}
\explain{Hướng giải: Xác định mục tiêu của việc đánh giá định lượng và định tính hiệu quả câu lệnh.

Các bước phân tích:
1. Để chọn ra câu lệnh tốt nhất cho hệ thống sản xuất, kỹ sư phải thử nghiệm nhiều biến thể prompt trên tập dữ liệu kiểm thử.
2. Tiêu chí quan trọng hàng đầu là đánh giá xem kết quả sinh ra có đúng trọng tâm (relevance), logic mạch lạc (coherence) và không bịa đặt (hallucination) hay không.

Kiểm tra: Các bộ chỉ số đánh giá LLM (như BLEU, ROUGE, G-Eval) đều tập trung vào đo lường Relevance và Coherence.

Vì sao các phương án khác sai:
- A: Tài nguyên tính toán là khía cạnh đo lường hạ tầng (latency/throughput), không phải đánh giá chất lượng prompt.
- C: Đánh giá có phương pháp khoa học, không chọn ngẫu nhiên.
- D: Tinh chỉnh prompt là bước tiếp theo sau khi đã đánh giá, không phải mục đích tự thân của việc đánh giá.

Lưu ý/Khái niệm cần nhớ: Đánh giá prompt tập trung vào tính phù hợp (relevance) và tính mạch lạc (coherence) của câu trả lời.}
\vocab{Prompt Evaluation | Đánh giá câu lệnh | Quá trình kiểm thử và đo lường chất lượng đầu ra của các câu lệnh khác nhau}
\end{question}

% ==============================================================================
% Question 131 | Phân loại: 3. Generative AI: Prompt Engineering Basics
% Xuất hiện trong các đề gốc: AIG202c_SU26_FE (Q26), AIG202c_SP26_RE (Q27)
% ==============================================================================
\begin{question}{131}{mcq}
\source{Q26.webp}
\prompt{Using a generative AI tool, Emily wants to create an image of a zebra-striped cat with a purple hat. Identify the best prompt for this task.}
\choice{A}{List cat breeds that look like zebras or have stripes.}
\choice{B}{Draw an image of a zebra-striped cat with a purple hat.}
\choice{C}{What's the difference between a zebra and a cat?}
\choice{D}{Find a zebra-striped cat that wears a purple hat.}
\answer{B}
\confidence{high}
\evidence{Để yêu cầu công cụ Generative AI tạo hình ảnh (Text-to-Image), câu lệnh trực tiếp, rõ ràng và hiệu quả nhất là 'Draw an image of a zebra-striped cat with a purple hat' (Vẽ một bức ảnh mèo có sọc ngựa vằn đội mũ màu tím).}
\sources{Midjourney & DALL-E Prompt Construction Best Practices}
\explain{Hướng giải: Lựa chọn câu lệnh tối ưu nhất cho tác vụ sinh hình ảnh cụ thể.

Các bước phân tích:
1. Mục tiêu: Tạo ra hình ảnh một con mèo có sọc như ngựa vằn và đội một chiếc mũ màu tím.
2. Câu lệnh hiệu quả cho mô hình sinh ảnh cần có động từ hành động trực tiếp (Draw / Create an image) cùng với mô tả chi tiết đặc điểm ngoại hình của đối tượng.
3. Phương án B: 'Draw an image of a zebra-striped cat with a purple hat' đáp ứng hoàn hảo tiêu chí này.

Kiểm tra: Câu lệnh này trực tiếp hướng dẫn mô hình sinh ảnh trực quan thay vì yêu cầu trả lời văn bản hay tìm kiếm.

Vì sao các phương án khác sai:
- A (List cat breeds): Yêu cầu liệt kê văn bản (Text list), không tạo ảnh.
- C (What\'s the difference): Yêu cầu so sánh lý thuyết văn bản.
- D (Find a zebra-striped cat): Câu lệnh tìm kiếm ảnh có sẵn trên web (Search engine query), không kích hoạt tác vụ sinh ảnh mới.

Lưu ý/Khái niệm cần nhớ: Prompt sinh ảnh cần có chỉ thị tạo hình rõ ràng cùng thuộc tính mô tả cụ thể.}
\vocab{Text-to-Image Prompt | Câu lệnh chuyển văn bản thành hình ảnh | Đoạn mô tả chi tiết dùng để chỉ đạo mô hình AI sinh ra hình ảnh tương ứng}
\end{question}

% ==============================================================================
% Question 132 | Phân loại: 3. Generative AI: Prompt Engineering Basics
% Xuất hiện trong các đề gốc: AIG202c_SU26_FE (Q38), AIG202c_SP26_RE (Q42)
% ==============================================================================
\begin{question}{132}{mcq}
\source{Q38.webp}
\prompt{Why might a generic prompt result in an average output?}
\choice{A}{Because the model needs technical instructions}
\choice{B}{Because it triggers the model's error handling mode}
\choice{C}{Because it lacks specific details and context}
\choice{D}{Because prompts must be written in code only}
\answer{C}
\confidence{high}
\evidence{Một câu lệnh chung chung (generic prompt) thường dẫn đến kết quả đầu ra tầm thường/trung bình (average output) vì nó thiếu các chi tiết cụ thể và ngữ cảnh rõ ràng (Because it lacks specific details and context), khiến mô hình phải tự suy đoán trên phân phối trung bình.}
\sources{Anthropic Prompt Engineering Tutorial: Be Clear and Direct; OpenAI Prompting Principles}
\explain{Hướng giải: Giải thích nguyên nhân tại sao prompt chung chung lại tạo ra câu trả lời tầm thường.

Các bước phân tích:
1. Khi người dùng nhập một câu lệnh quá ngắn gọn và chung chung (ví dụ: 'Hãy viết về marketing'), mô hình ngôn ngữ không biết đối tượng độc giả, độ dài, văn phong, mục đích cụ thể.
2. Do đó, mô hình sẽ kích hoạt phân phối xác suất trung bình nhất (vốn là các thông tin phổ thông, chung chung nhất trên Internet), tạo ra câu trả lời nhạt nhòa, không có điểm nhấn.
3. Muốn có câu trả lời xuất sắc, prompt cần có chi tiết cụ thể, bối cảnh, tiêu chí đánh giá và định dạng rõ ràng.

Kiểm tra: Nguyên tắc vàng trong Prompt Engineering: 'Specificity yields quality' (Càng cụ thể thì chất lượng càng cao).

Vì sao các phương án khác sai:
- A: Mô hình không nhất thiết cần chỉ dẫn kỹ thuật phức tạp, chỉ cần bối cảnh rõ ràng bằng ngôn ngữ tự nhiên.
- B: Prompt chung chung không kích hoạt chế độ xử lý lỗi, nó vẫn trả lời nhưng câu trả lời không sâu sắc.
- D: Prompt được viết bằng ngôn ngữ tự nhiên (tiếng Anh, tiếng Việt), không bắt buộc phải viết bằng mã code.

Lưu ý/Khái niệm cần nhớ: Specificity & Context = Chìa khóa biến câu trả lời từ trung bình thành xuất sắc.}
\vocab{Specificity | Tính cụ thể | Mức độ chi tiết và rõ ràng của các yêu cầu được đưa vào trong câu lệnh}
\end{question}

% ==============================================================================
% Question 133 | Phân loại: 3. Generative AI: Prompt Engineering Basics
% Xuất hiện trong các đề gốc: AIG202c_SU26_FE (Q42), AIG202c_SP26_RE (Q32)
% ==============================================================================
\begin{question}{133}{mcq}
\source{Q42.webp}
\prompt{What role does human oversight play in prompt engineering for generative AI?}
\choice{A}{Human oversight is not necessary for prompt engineering}
\choice{B}{Human oversight is essential for designing effective prompts and evaluating model outputs}
\choice{C}{Human oversight is only required during the training phase of generative AI models}
\choice{D}{Human oversight is necessary for adjusting hyperparameters during model generation}
\answer{B}
\confidence{high}
\evidence{Sự giám sát của con người (Human oversight) đóng vai trò then chốt trong Prompt Engineering để thiết kế các câu lệnh hiệu quả và đánh giá, hiệu đính chất lượng kết quả đầu ra của mô hình (essential for designing effective prompts and evaluating model outputs).}
\sources{EU Ethics Guidelines for Trustworthy AI (Human Agency and Oversight); Human-in-the-Loop Prompt Engineering}
\explain{Hướng giải: Đánh giá tầm quan trọng của sự giám sát của con người trong quy trình kỹ nghệ câu lệnh.

Các bước phân tích:
1. Mô hình AI không tự biết câu trả lời nào là tối ưu nhất cho mục đích kinh doanh của con người.
2. Con người (Prompt Engineer / Chuyên gia nghiệp vụ) là nhân tố quyết định trong việc xây dựng cấu trúc prompt, kiểm thử các biến thể và đánh giá kết quả đầu ra để đảm bảo độ tin cậy và tính an toàn.

Kiểm tra: Nguyên tắc 'Human-in-the-loop' là tiêu chuẩn vàng trong vận hành Generative AI.

Vì sao các phương án khác sai:
- A: Cho rằng không cần con người là hoàn toàn sai lầm.
- C: Giám sát cần thiết ở cả giai đoạn suy luận và ứng dụng, không chỉ riêng huấn luyện.
- D: Prompt engineering không trực tiếp điều chỉnh siêu tham số nội tại (weights/learning rate) của mô hình.

Lưu ý/Khái niệm cần nhớ: Human oversight là nhân tố quyết định chất lượng và sự an toàn của ứng dụng GenAI.}
\vocab{Human Oversight | Giám sát của con người | Cơ chế đảm bảo con người duy trì quyền kiểm soát và đánh giá hoạt động của AI}
\end{question}

% ==============================================================================
% Question 134 | Phân loại: 3. Generative AI: Prompt Engineering Basics
% Xuất hiện trong các đề gốc: AIG202c_SU26_FE (Q44)
% ==============================================================================
\begin{question}{134}{mcq}
\source{Q44.webp}
\prompt{Which of the following is an example of a straightforward prompt?}
\choice{A}{Suggest ten names for a new company logo}
\choice{B}{Write HTML code to add a dropdown of cities}
\choice{C}{Analyze the political history of two nations}
\choice{D}{Explain the meaning of philosophical theories}
\answer{B}
\confidence{high}
\evidence{A straightforward prompt is one with clear, specific, bounded, and direct instructions with unambiguous desired output format, such as 'Write HTML code to add a dropdown of cities'.}
\sources{Prompt Engineering Guide (DAIR.AI); OpenAI Best Practices for Prompt Engineering}
\explain{Hướng giải: Phân biệt câu lệnh trực diện, rõ ràng (straightforward prompt) với các câu lệnh mơ hồ hoặc quá rộng.

Các bước phân tích:
1. Một câu lệnh trực diện (straightforward prompt) là câu lệnh có yêu cầu rõ ràng, phạm vi cụ thể, mục tiêu đơn nhất và định dạng đầu ra không gây mơ hồ.
2. Phân tích các lựa chọn:
   - A ('Suggest ten names for a new company logo'): Yêu cầu đặt tên cho logo doanh nghiệp mang tính sáng tạo trừu tượng và có thể cần thêm nhiều ngữ cảnh về ngành nghề, đối tượng mục tiêu.
   - B ('Write HTML code to add a dropdown of cities'): Đây là một nhiệm vụ kỹ thuật cụ thể, rõ ràng, có phạm vi xác định hoàn hảo (viết mã HTML cho thẻ select chứa các thành phố), không đòi hỏi suy diễn triết học hay ngữ cảnh sâu xa. Đây là ví dụ chuẩn về straightforward prompt.
   - C ('Analyze the political history of two nations'): Phân tích lịch sử chính trị của hai quốc gia là chủ đề học thuật vô cùng rộng lớn, đa chiều và phức tạp.
   - D ('Explain the meaning of philosophical theories'): Giải thích ý nghĩa các lý thuyết triết học là câu hỏi mở rộng lớn và trừu tượng.

Kiểm tra:
Phương án B có mục tiêu hẹp, cụ thể và định dạng rõ ràng nhất.

Vì sao các phương án khác sai:
- A, C, D đều là các yêu cầu mở (open-ended), mang tính suy luận phức tạp hoặc trừu tượng, đòi hỏi tinh chỉnh thêm nhiều tham số.

Lưu ý/Khái niệm cần nhớ:
Straightforward prompt: Câu lệnh có mục tiêu cụ thể, phạm vi khép kín và định dạng đầu ra rõ ràng.}
\vocab{Straightforward Prompt | Câu lệnh trực diện | Câu lệnh ngắn gọn, chỉ rõ tác vụ cụ thể và định dạng kết quả mong muốn mà không cần giải thích thêm}
\end{question}

% ==============================================================================
% Question 135 | Phân loại: 3. Generative AI: Prompt Engineering Basics
% Xuất hiện trong các đề gốc: AIG202c_SP26_RE (Q34)
% ==============================================================================
\begin{question}{135}{mcq}
\source{Q34.jpg}
\prompt{Jennifer wants to request some useful information about a complex medical condition using a large language model. Which among the following techniques should she employ on the text prompt to ensure that the generated content is relevant, accurate, and precise for this specialized field?}
\choice{A}{Framing}
\choice{B}{Domain expertise}
\choice{C}{Bias mitigation}
\choice{D}{Task specification}
\answer{B}
\confidence{high}
\evidence{Để yêu cầu thông tin về một tình trạng bệnh lý y khoa phức tạp đòi hỏi độ chính xác cao, kỹ thuật câu lệnh cần áp dụng là Đóng vai Chuyên gia chuyên ngành (Domain expertise / Persona pattern), yêu cầu mô hình phản hồi với tư cách và chuẩn mực tri thức của một chuyên gia y tế.}
\sources{Medical Prompting Guidelines; Prompting LLMs for Clinical Decision Support}
\explain{Hướng giải: Xác định kỹ thuật câu lệnh tối ưu nhất cho truy vấn y khoa chuyên sâu.

Các bước phân tích:
1. Y khoa là lĩnh vực đòi hỏi tính chính xác, thuật ngữ chuẩn hóa và độ tin cậy tuyệt đối.
2. Kỹ thuật đưa bối cảnh chuyên gia chuyên ngành ('Domain expertise' / thiết lập vai trò bác sĩ lâm sàng trong prompt) kích hoạt khả năng truy xuất các thuật ngữ và kiến thức y khoa chuyên sâu của LLM, hạn chế diễn giải chung chung thông tục.

Kiểm tra: Thiết lập bối cảnh chuyên gia miền (Domain expertise) được chứng minh giúp cải thiện chất lượng câu trả lời chuyên sâu đáng kể.

Vì sao các phương án khác sai:
- A (Framing): Khung nhìn định hình truyền thông, không đảm bảo tính chính xác y khoa.
- C (Bias mitigation): Tránh định kiến, không giải quyết vấn đề độ sâu chuyên môn.
- D (Task specification): Chỉ định nhiệm vụ chung chung, chưa đủ cho miền chuyên biệt.

Lưu ý/Khái niệm cần nhớ: Kích hoạt chuyên gia miền (Domain expertise) giúp câu trả lời AI chính xác và chuyên sâu trong các ngành đặc thù.}
\vocab{Domain Expertise | Chuyên môn lĩnh vực | Kiến thức chuyên sâu và thuật ngữ chuẩn mực của một ngành nghề cụ thể}
\end{question}

% ==============================================================================
% Question 136 | Phân loại: 3. Generative AI: Prompt Engineering Basics
% Xuất hiện trong các đề gốc: AIG202c_SP26_RE (Q43)
% ==============================================================================
\begin{question}{136}{mcq}
\source{Q43.jpg}
\prompt{Which among the following statements is accurate about the Tree-of-Thought approach?}
\choice{A}{The Tree-of-Thought approach only works for generating responses to marketing-related prompts.}
\choice{B}{The Tree-of-Thought approach is less effective than the Chain-of-Thought approach for generative AI reasoning.}
\choice{C}{The Tree-of-Thought approach eliminates the need for any prompt instructions or constraints.}
\choice{D}{The Tree-of-Thought approach enables generative AI models to explore multiple paths simultaneously and assess potential outcomes.}
\answer{D}
\confidence{high}
\evidence{Phát biểu chính xác về Tree-of-Thought (ToT) là phương pháp này cho phép các mô hình AI tạo sinh khám phá nhiều lộ trình suy nghĩ đồng thời và đánh giá các kết quả tiềm năng ở từng bước (enables generative AI models to explore multiple paths simultaneously and assess potential outcomes).}
\sources{Yao et al. (2023), Tree of Thoughts (Princeton University / Google DeepMind)}
\explain{Hướng giải: Nhận diện phát biểu chính xác và khoa học nhất về phương pháp Tree-of-Thought (ToT).

Các bước phân tích:
1. ToT mở rộng từ Chain-of-Thought bằng cách cho phép mô hình cân nhắc nhiều nhánh suy luận tại mỗi bước (giống như duyệt cây tìm kiếm trong cờ vua).
2. Mô hình có thể tự đánh giá độ khả thi của từng nhánh (self-evaluation), cắt tỉa các nhánh kém triển vọng và quay lui khi cần.
3. Do đó, phát biểu D mô tả hoàn toàn chuẩn xác cơ chế hoạt động của Tree-of-Thought.

Kiểm tra: Đây là bước tiến quan trọng giúp LLM giải quyết các bài toán lập kế hoạch và suy luận phức tạp.

Vì sao các phương án khác sai:
- A: ToT là phương pháp suy luận đa năng cho mọi tác vụ logic, toán học, kế hoạch, không chỉ riêng marketing.
- B: ToT được chứng minh là hiệu quả vượt trội so với Chain-of-Thought trên các bài toán suy luận khó (như Game of 24, Creative Writing).
- C: ToT vẫn đòi hỏi cấu trúc câu lệnh và chỉ dẫn đánh giá chặt chẽ.

Lưu ý/Khái niệm cần nhớ: Tree-of-Thought = Khám phá đa nhánh suy luận đồng thời kết hợp tự đánh giá.}
\vocab{Tree-of-Thought (ToT) | Cây suy nghĩ | Khung suy luận phân nhánh cho phép LLM giải quyết bài toán phức tạp}
\end{question}

\section{sec4}{4. Building AI Powered Chatbots Without Programming}

% ==============================================================================
% Question 137 | Phân loại: 4. Building AI Powered Chatbots Without Programming
% Xuất hiện trong các đề gốc: AIG202c_SP_2025_FE (Q5)
% ==============================================================================
\begin{question}{137}{mcq}
\source{AIG202c_-_SP_2025_-_FE_3755_5.webp}
\prompt{What is the primary goal of incorporating 'digressions' in chatbot conversations?}
\choice{A}{To strictly follow pre-defined conversation paths}
\choice{B}{To limit user interaction choices}
\choice{C}{To provide diverse and engaging conversation experiences}
\choice{D}{To avoid conversing with users}
\answer{C}
\confidence{high}
\evidence{Mục tiêu chính của việc đưa tính năng 'lạc đề' (digressions) vào hội thoại chatbot là cung cấp trải nghiệm giao tiếp phong phú, tự nhiên và lôi cuốn cho người dùng (To provide diverse and engaging conversation experiences).}
\sources{IBM Watson Assistant Digression Best Practices}
\explain{Hướng giải: Nhận diện mục đích cốt lõi của tính năng Digression trong chatbot (câu hỏi kiểm tra tính nhất quán với UQ056 và UQ094).

Các bước phân tích:
1. Con người luôn muốn trò chuyện linh hoạt, có thể hỏi thêm thông tin ngoài lề bất cứ lúc nào.
2. Cho phép digression giúp bot vượt qua sự cứng nhắc của các biểu mẫu định sẵn, mang lại trải nghiệm tương tác đa dạng và hấp dẫn (diverse and engaging).

Kiểm tra: Đây là triết lý thiết kế đàm thoại của IBM Watson Assistant.

Vì sao các phương án khác sai:
- A: Tuân thủ nghiêm ngặt lộ trình cứng nhắc là phủ nhận digression.
- B: Giới hạn lựa chọn là sai.
- D: Tránh trò chuyện là đi ngược lại mục đích của chatbot.

Lưu ý/Khái niệm cần nhớ: Digression = Trải nghiệm đàm thoại đa dạng, tự nhiên và gắn kết.}
\vocab{Conversational Flexibility | Tính linh hoạt trong hội thoại | Năng lực của hệ thống thích ứng với các luồng trao đổi đa dạng của người dùng}
\end{question}

% ==============================================================================
% Question 138 | Phân loại: 4. Building AI Powered Chatbots Without Programming
% Xuất hiện trong các đề gốc: AIG202c_SP_2025_FE (Q6), AIG202c_FA25_FE (Q47), AIG202c_SU26_FE (Q58)
% ==============================================================================
\begin{question}{138}{mcq}
\source{AIG202c_-_SP_2025_-_FE_3755_6.webp}
\prompt{Which aspect of chatbot development focuses on implementing 'digressions'?}
\choice{A}{Hardware optimizations}
\choice{B}{Conversation flow and engagement}
\choice{C}{Data encryption protocols}
\choice{D}{Graphics design in chat interfaces}
\answer{B}
\confidence{high}
\evidence{Trong kiến trúc chatbot, Thực thể (Entity) đại diện cho một thông tin cụ thể, có giá trị ngữ nghĩa được trích xuất từ câu nói của người dùng (ví dụ: địa điểm, ngày tháng, thời gian, tên món ăn).}
\sources{Jurafsky & Martin, Speech and Language Processing (Information Extraction); IBM Watson Assistant Entities}
\explain{Hướng giải: Xác định định nghĩa chính xác của Thực thể (Entity) trong xử lý ngôn ngữ tự nhiên cho chatbot.

Các bước phân tích:
1. Phân biệt rõ: Intent (Ý định) là 'người dùng muốn làm gì', còn Entity (Thực thể) là 'đối tượng cụ thể của hành động đó' (danh từ, tham số, dữ liệu trích xuất).
2. Ví dụ: 'Tôi muốn đặt vé máy bay đi Tokyo vào ngày mai': Intent là #BookFlight, Entities là @Destination (Tokyo) và @Date (ngày mai).
3. Do đó, Entity là thông tin cụ thể bổ nghĩa cho ý định của người dùng.

Kiểm tra: Khái niệm Slot Filling trong hội thoại dựa trên việc điền đầy các Entity cần thiết để thực thi hành động.

Vì sao các phương án khác sai:
- Entity không phải là mục tiêu tổng quát (đó là Intent).
- Entity không phải là câu trả lời của bot hay tốc độ phản hồi.

Lưu ý/Khái niệm cần nhớ: Intent = Ý định/Hành động; Entity = Dữ liệu tham số chi tiết.}
\vocab{Entity | Thực thể | Từ hoặc cụm từ khóa đại diện cho dữ liệu cụ thể trong câu nói của người dùng}
\end{question}

% ==============================================================================
% Question 139 | Phân loại: 4. Building AI Powered Chatbots Without Programming
% Xuất hiện trong các đề gốc: AIG202c_SP_2025_FE (Q8), AIG202c_SP26_RE (Q57)
% ==============================================================================
\begin{question}{139}{mcq}
\source{AIG202c_-_SP_2025_-_FE_3755_8.webp}
\prompt{How do chatbots use dialog management to guide conversations with users?}
\choice{A}{By ignoring user input}
\choice{B}{By following a predefined sequence of interactions based on the user's responses}
\choice{C}{By generating random responses}
\choice{D}{By terminating the conversation after a set number of interactions}
\answer{B}
\confidence{high}
\evidence{Bộ quản lý hội thoại (Dialog Management) dẫn dắt cuộc trò chuyện bằng cách tuân theo một chuỗi tương tác có kịch bản định trước dựa trên phản hồi của người dùng (following a predefined sequence of interactions based on the user's responses).}
\sources{IBM Watson Assistant: Building a Dialog; Conversational Flow Architecture}
\explain{Hướng giải: Định nghĩa chức năng của Trình quản lý hội thoại (Dialog Manager) trong Chatbot.

Các bước phân tích:
1. Dialog Manager đóng vai trò bộ não điều phối: nhận diện Ý định (Intent) và Thực thể (Entity) hiện tại, kết hợp với trạng thái ngữ cảnh để quyết định hành động tiếp theo của bot.
2. Nó điều hướng luồng trò chuyện theo cây quyết định hoặc đồ thị phân nhánh đã được thiết kế sẵn (predefined sequence of interactions).

Kiểm tra: Cấu trúc cây Dialog Nodes trong Watson Assistant minh họa hoàn hảo cho cơ chế này.

Vì sao các phương án khác sai:
- A: Bỏ qua đầu vào là sai hoàn toàn.
- C: Dialog Manager hoạt động theo quy tắc logic, không sinh phản hồi ngẫu nhiên.
- D: Ngắt kết nối không phải là mục tiêu điều phối hội thoại.

Lưu ý/Khái niệm cần nhớ: Dialog Management = Điều phối luồng đàm thoại dựa trên phản hồi của người dùng.}
\vocab{Dialog Management | Quản lý hội thoại | Thành phần quyết định phản hồi tiếp theo của chatbot dựa trên trạng thái hội thoại hiện tại}
\end{question}

% ==============================================================================
% Question 140 | Phân loại: 4. Building AI Powered Chatbots Without Programming
% Xuất hiện trong các đề gốc: AIG202c_SP_2025_FE (Q10), AIG202c_SU26_FE (Q50), AIG202c_SP26_RE (Q54)
% ==============================================================================
\begin{question}{140}{mcq}
\source{AIG202c_-_SP_2025_-_FE_3755_10.webp}
\prompt{How are 'slots' used in the context of chatbot interactions?}
\choice{A}{To identify network vulnerabilities}
\choice{B}{To manage and capture specific pieces of information from users}
\choice{C}{To optimize website speed}
\choice{D}{To conduct financial transactions}
\answer{B}
\confidence{high}
\evidence{Các ô nhớ / khe thông tin ('slots') trong tương tác chatbot được sử dụng để quản lý và thu thập các mẩu thông tin cụ thể bắt buộc từ người dùng (To manage and capture specific pieces of information from users) nhằm hoàn thành một biểu mẫu hoặc tác vụ.}
\sources{IBM Watson Assistant: Gathering Information with Slots; Slot Filling in Task-Oriented Dialog}
\explain{Hướng giải: Định nghĩa chức năng của kỹ thuật 'Slots' (Điền khe thông tin - Slot Filling) trong Chatbot.

Các bước phân tích:
1. Khi người dùng muốn thực hiện một tác vụ đòi hỏi nhiều thông tin (ví dụ: đặt phòng khách sạn cần: ngày đến, ngày đi, số khách, loại phòng), chatbot sử dụng tính năng 'Slots'.
2. Mỗi 'slot' tương ứng với một thực thể (entity) cần thu thập. Bot sẽ lần lượt kiểm tra và hỏi các slot còn trống cho đến khi thu thập đủ toàn bộ thông tin.
3. Nhờ đó, quy trình thu thập dữ liệu diễn ra gọn gàng trong cùng một nút hội thoại duy nhất mà không cần tạo hàng chục nút rườm rà.

Kiểm tra: Kỹ thuật Slot Filling là giải pháp chuẩn hóa cho mọi tác vụ thu thập biểu mẫu trong Conversational AI.

Vì sao các phương án khác sai:
- A: Phát hiện lỗ hổng thuộc về an ninh bảo mật (Cybersecurity).
- C: Tối ưu tốc độ web thuộc về Web Performance.
- D: Tiến hành giao dịch tài chính là tác vụ thanh toán, không phải định nghĩa của Slots.

Lưu ý/Khái niệm cần nhớ: Slots = Thu thập và quản lý các mẩu thông tin cụ thể từ người dùng trong một lượt hội thoại.}
\vocab{Slot Filling | Điền khe thông tin | Kỹ thuật thu thập đầy đủ các tham số bắt buộc từ người dùng để hoàn tất tác vụ}
\end{question}

% ==============================================================================
% Question 141 | Phân loại: 4. Building AI Powered Chatbots Without Programming
% Xuất hiện trong các đề gốc: AIG202c_SP_2025_FE (Q13)
% ==============================================================================
\begin{question}{141}{mcq}
\source{AIG202c_-_SP_2025_-_FE_3755_13.webp}
\prompt{Which aspect of chatbot development does the 'Dialog' module emphasize?}
\choice{A}{Software development methodologies}
\choice{B}{Managing conversational flow and responses}
\choice{C}{Augmented reality technologies}
\choice{D}{Network security protocols}
\answer{B}
\confidence{high}
\evidence{Mô-đun 'Dialog' (Hội thoại) trong xây dựng chatbot tập trung chủ yếu vào việc quản lý luồng hội thoại và các câu phản hồi của bot (Managing conversational flow and responses).}
\sources{IBM Watson Assistant Dialog Module Documentation}
\explain{Hướng giải: Xác định trọng tâm chức năng của mô-đun Dialog trong chatbot (câu hỏi kiểm tra tính nhất quán với UQ098).

Các bước phân tích:
1. Cấu trúc chatbot Watson gồm 3 trụ cột: Intents (Ý định), Entities (Thực thể), và Dialog (Hội thoại).
2. Trong khi Intents và Entities chịu trách nhiệm hiểu ngôn ngữ (NLU), thì mô-đun Dialog chịu trách nhiệm xác định kịch bản: bot sẽ nói gì, phản hồi ra sao và điều hướng cuộc trò chuyện đi về đâu.

Kiểm tra: Đây là kiến thức căn bản của bài học xây dựng Chatbot với Watson Assistant.

Vì sao các phương án khác sai:
- A: Phương pháp phát triển phần mềm (Agile/Scrum) thuộc quản lý dự án chung.
- C: Thực tế tăng cường (AR) không thuộc chatbot văn bản.
- D: Giao thức an ninh mạng không phải chức năng của tầng Dialog.

Lưu ý/Khái niệm cần nhớ: Dialog Module = Quản lý luồng hội thoại và kịch bản phản hồi.}
\vocab{Dialog Module | Mô-đun hội thoại | Khối điều khiển định nghĩa các nút logic phản hồi trong chatbot}
\end{question}

% ==============================================================================
% Question 142 | Phân loại: 4. Building AI Powered Chatbots Without Programming
% Xuất hiện trong các đề gốc: AIG202c_SP_2025_FE (Q17)
% ==============================================================================
\begin{question}{142}{mcq}
\source{AIG202c_-_SP_2025_-_FE_3755_17.webp}
\prompt{Which of the following is NOT an example of an entity for a travel chatbot?}
\choice{A}{Destination}
\choice{B}{Departure time}
\choice{C}{GreetUser}
\choice{D}{Number of passengers}
\answer{C}
\confidence{high}
\evidence{`GreetUser` (Chào người dùng) là một Ý định (Intent), KHÔNG PHẢI là một Thực thể (Entity); các Thực thể của chatbot du lịch phải là các tham số dữ liệu như Điểm đến (Destination), Giờ khởi hành (Departure time), Số lượng hành khách (Number of passengers).}
\sources{Conversational AI Intent vs Entity Separation; Travel Booking Chatbot Schema}
\explain{Hướng giải: Phân biệt Thực thể (Entity) và Ý định (Intent) trong kịch bản chatbot du lịch (câu hỏi kiểm tra tính nhất quán với câu SP_2025 Q17).

Các bước phân tích:
1. Entity đại diện cho các dữ liệu/tham số nghiệp vụ cụ thể cần thu thập trong chuyến đi:
   - Destination (Điểm đến): Ví dụ: Hà Nội, Đà Nẵng.
   - Departure time (Giờ khởi hành): Ví dụ: 8h sáng ngày 25/10.
   - Number of passengers (Số lượng khách): Ví dụ: 2 người lớn.
2. `GreetUser` thể hiện hành vi chào hỏi của người dùng ('Xin chào', 'Hello'), đây là một Intent (Ý định), hoàn toàn không phải là một Entity.

Kiểm tra: Đây là câu hỏi phân loại cơ bản bắt buộc phải nắm vững trong học phần Chatbot.

Vì sao các phương án khác là Entity hợp lệ:
- A, B, D đều là các trường dữ liệu thực thể chuẩn mực của tác vụ đặt vé du lịch.

Lưu ý/Khái niệm cần nhớ: Ý định (Intent) = Hành động/Mục đích; Thực thể (Entity) = Dữ liệu/Tham số chi tiết.}
\vocab{Entity vs Intent | Thực thể và Ý định | Hai khái niệm trụ cột trong NLU: Intent xác định hành động, Entity xác định dữ liệu đi kèm}
\end{question}

% ==============================================================================
% Question 143 | Phân loại: 4. Building AI Powered Chatbots Without Programming
% Xuất hiện trong các đề gốc: AIG202c_SP_2025_FE (Q20), AIG202c_SU25_FE (Q50)
% ==============================================================================
\begin{question}{143}{mcq}
\source{AIG202c_-_SP_2025_-_FE_3755_20.webp}
\prompt{What are "context variables" in the context of artificial intelligence and chatbots?}
\choice{A}{Variables used to store the state or context of a conversation or interaction}
\choice{B}{Variables that represent the emotions of the chatbot during a conversation}
\choice{C}{Variables that store the user's personal information}
\choice{D}{Variables used to determine the intent of the user's input}
\answer{C}
\confidence{high}
\evidence{Khả năng tạo sinh chủ đạo mà mô hình ngôn ngữ lớn (LLM) thể hiện là Sinh văn bản (Text generation).}
\sources{Vaswani et al., Attention Is All You Need; OpenAI GPT Architecture}
\explain{Hướng giải: Xác định năng lực cốt lõi của LLM (câu hỏi kiểm tra tính nhất quán với UQ082).

Các bước phân tích:
1. Large Language Models được huấn luyện trên kho dữ liệu văn bản đồ sộ với mục tiêu tự hồi quy (autoregressive language modeling) để sinh từ tiếp theo.
2. Do đó, chức năng tự nhiên và bản chất nhất của LLM là Text generation (Sinh văn bản).

Kiểm tra: Đây là câu hỏi lặp lại kiểm tra tri thức nền tảng trong đề thi.

Vì sao các phương án khác sai:
- A (Image generation): Thuộc về mô hình khuếch tán (Diffusion Models).
- B (Audio generation): Thuộc về mô hình sinh âm thanh.

Lưu ý/Khái niệm cần nhớ: LLM = Large Language Model -> Năng lực chính là Text generation.}
\vocab{Autoregressive Model | Mô hình tự hồi quy | Mô hình dự đoán phần tử tiếp theo trong chuỗi dựa trên toàn bộ các phần tử đã xuất hiện trước đó}
\end{question}

% ==============================================================================
% Question 144 | Phân loại: 4. Building AI Powered Chatbots Without Programming
% Xuất hiện trong các đề gốc: AIG202c_SP_2025_FE (Q21)
% ==============================================================================
\begin{question}{144}{mcq}
\source{AIG202c_-_SP_2025_-_FE_3755_21.webp}
\prompt{What is the primary purpose of a chatbot in the field of artificial intelligence?}
\choice{A}{To navigate websites}
\choice{B}{To automate repetitive tasks}
\choice{C}{To process images}
\choice{D}{To perform complex calculations}
\answer{B}
\confidence{high}
\evidence{Mục đích thực tiễn bao trùm của chatbot trong lĩnh vực trí tuệ nhân tạo là tự động hóa các tác vụ lặp đi lặp lại (To automate repetitive tasks) trong giao tiếp và chăm sóc khách hàng.}
\sources{IBM Automation with AI; McKinsey State of AI in Customer Operations}
\explain{Hướng giải: Xác định mục đích ứng dụng thực tế quan trọng nhất của chatbot AI.

Các bước phân tích:
1. Phần lớn các thắc mắc của khách hàng (hỏi giá, giờ mở cửa, chính sách đổi trả, tra cứu đơn hàng) mang tính lặp đi lặp lại rất cao.
2. Chatbot giải quyết bài toán này bằng cách tự động hóa hoàn toàn việc tiếp nhận và xử lý các câu hỏi thường gặp này 24/7 mà không cần sự can thiệp của con người.

Kiểm tra: 'Automate repetitive conversational tasks' là giá trị kinh doanh cốt lõi của chatbot.

Vì sao các phương án khác sai:
- A: Điều hướng website là chức năng của trình duyệt / menu điều hướng.
- C: Xử lý hình ảnh là của Computer Vision.
- D: Tính toán phức tạp là của các phần mềm tính toán số học / siêu máy tính.

Lưu ý/Khái niệm cần nhớ: Chatbot tự động hóa các tác vụ giao tiếp lặp đi lặp lại.}
\vocab{Task Automation | Tự động hóa tác vụ | Việc chuyển giao các công việc lặp lại cho phần mềm tự động thực thi}
\end{question}

% ==============================================================================
% Question 145 | Phân loại: 4. Building AI Powered Chatbots Without Programming
% Xuất hiện trong các đề gốc: AIG202c_SP_2025_FE (Q30)
% ==============================================================================
\begin{question}{145}{mcq}
\source{AIG202c_-_SP_2025_-_FE_3755_30.webp}
\prompt{What is the primary definition of a chatbot among the following options?}
\choice{A}{A computer program to mimic human conversation.}
\choice{B}{A robot that can interact with people.}
\choice{C}{A program to chat with friends}
\choice{D}{A game we play online.}
\answer{A}
\confidence{high}
\evidence{Định nghĩa cơ bản và chuẩn mực nhất của chatbot là: Một chương trình máy tính được thiết kế để mô phỏng cuộc trò chuyện của con người (A computer program to mimic human conversation).}
\sources{Oxford English Dictionary: Chatbot; IBM Watson Assistant Definition}
\explain{Hướng giải: Chọn định nghĩa chính xác nhất của chatbot trong các phương án được đưa ra.

Các bước phân tích:
1. Chatbot = Chat + Robot (chương trình phần mềm tự động trò chuyện).
2. Định nghĩa học thuật chuẩn: Một chương trình máy tính mô phỏng hành vi giao tiếp ngôn ngữ tự nhiên của con người qua văn bản hoặc âm thanh.

Kiểm tra: Phương án A thể hiện đầy đủ và chính xác bản chất phần mềm và chức năng mô phỏng đàm thoại.

Vì sao các phương án khác sai:
- B: Không nhất thiết phải là robot cơ học phần cứng.
- C: Không giới hạn chỉ là chương trình trò chuyện với bạn bè.
- D: Không phải là một trò chơi trực tuyến.

Lưu ý/Khái niệm cần nhớ: Chatbot = Chương trình máy tính mô phỏng hội thoại con người.}
\vocab{Conversational Agent | Tác tử đàm thoại | Thuật ngữ học thuật chỉ các chương trình phần mềm có khả năng giao tiếp tự nhiên}
\end{question}

% ==============================================================================
% Question 146 | Phân loại: 4. Building AI Powered Chatbots Without Programming
% Xuất hiện trong các đề gốc: AIG202c_SP_2025_FE (Q45)
% ==============================================================================
\begin{question}{146}{mcq}
\source{AIG202c_-_SP_2025_-_FE_3755_45.webp}
\prompt{In what way do 'digressions' impact the overall chatbot experience?}
\choice{A}{They hinder the conversation flow}
\choice{B}{They provide a natural and varied conversational experience}
\choice{C}{They limit the chatbot's vocabulary}
\choice{D}{They delay response times}
\answer{B}
\confidence{high}
\evidence{Các phân nhánh lạc đề (digressions) tác động tích cực đến trải nghiệm chatbot tổng thể bằng cách mang lại một trải nghiệm đối thoại tự nhiên, phong phú và đa dạng (They provide a natural and varied conversational experience), giúp cuộc trò chuyện bớt rập khuôn.}
\sources{IBM Watson Assistant: Enhancing UX with Digressions}
\explain{Hướng giải: Đánh giá tác động của tính năng Digression đối với trải nghiệm chatbot tổng thể (câu hỏi kiểm tra tính nhất quán với UQ094 và UQ103).

Các bước phân tích:
1. Con người đánh giá cao những chatbot có khả năng xử lý linh hoạt khi cuộc trò chuyện rẽ nhánh mà không bị đơ hoặc lỗi ngữ cảnh.
2. Digressions mang lại cảm giác trò chuyện tự nhiên và đa dạng (natural and varied), tương tự như đang tương tác với một con người thực sự.

Kiểm tra: Khớp hoàn toàn với triết lý thiết kế của các hệ thống Conversational AI hiện đại.

Vì sao các phương án khác sai:
- A: Gây cản trở luồng hội thoại là biểu hiện khi xử lý kém, không phải tác động mong muốn của tính năng.
- C: Không làm giới hạn vốn từ vựng.
- D: Không làm trễ thời gian phản hồi.

Lưu ý/Khái niệm cần nhớ: Digressions tạo nên trải nghiệm hội thoại tự nhiên và linh hoạt.}
\vocab{Conversational Naturalness | Tính tự nhiên trong hội thoại | Cảm giác giao tiếp mượt mà, chân thực như trò chuyện giữa con người với nhau}
\end{question}

% ==============================================================================
% Question 147 | Phân loại: 4. Building AI Powered Chatbots Without Programming
% Xuất hiện trong các đề gốc: AIG202c_SP_2025_FE (Q53)
% ==============================================================================
\begin{question}{147}{mcq}
\source{AIG202c_-_SP_2025_-_FE_3755_53.webp}
\prompt{How do well-implemented 'digressions' contribute to user satisfaction within chatbot interactions?}
\choice{A}{By creating predictable conversation patterns}
\choice{B}{By restricting the variety of conversation paths}
\choice{C}{By offering engaging and unexpected conversational turns}
\choice{D}{By limiting conversational flexibility}
\answer{C}
\confidence{high}
\evidence{Xử lý lạc đề tốt (well-implemented digressions) làm tăng mức độ hài lòng của người dùng bằng cách mang lại những bước chuyển hội thoại thú vị, tự nhiên và bất ngờ (offering engaging and unexpected conversational turns).}
\sources{IBM Watson Assistant: Enhancing User Satisfaction with Digressions}
\explain{Hướng giải: Hiểu đóng góp của tính năng Digression đối với sự hài lòng của người dùng chatbot.

Các bước phân tích:
1. Cuộc trò chuyện theo kịch bản cứng nhắc thường gây nhàm chán.
2. Khi người dùng bất ngờ hỏi một câu ngoài lề mà chatbot vẫn tiếp nhận thông minh, trả lời hóm hỉnh rồi nhẹ nhàng đưa câu chuyện về luồng chính, người dùng cảm thấy bất ngờ và cực kỳ hài lòng với sự thông minh của hệ thống.

Kiểm tra: Đây là yếu tố tạo ra 'khoảnh khắc bất ngờ thích thú' (delight moment) trong thiết kế trải nghiệm người dùng.

Vì sao các phương án khác sai:
- A: Tạo sự đoán trước rập khuôn là dở.
- B & D: Hạn chế sự đa dạng và linh hoạt là trải nghiệm tồi.

Lưu ý/Khái niệm cần nhớ: Digression tạo ra sự thú vị và ngẫu hứng, nâng cao sự hài lòng của người dùng.}
\vocab{User Delight | Sự hài lòng vượt mong đợi | Trạng thái cảm xúc tích cực khi hệ thống đáp ứng thông minh vượt kỳ vọng của người dùng}
\end{question}

% ==============================================================================
% Question 148 | Phân loại: 4. Building AI Powered Chatbots Without Programming
% Xuất hiện trong các đề gốc: AIG202c_SP_2025_FE (Q56), AIG202c_SU26_FE (Q60)
% ==============================================================================
\begin{question}{148}{mcq}
\source{AIG202c_-_SP_2025_-_FE_3755_56.webp}
\prompt{What are 'context variables' in chatbot conversations?}
\choice{A}{User authentication methods}
\choice{B}{Temporary data storage during conversations}
\choice{C}{Graphic design elements}
\choice{D}{Server configurations}
\answer{B}
\confidence{high}
\evidence{Biến ngữ cảnh (context variables) trong hội thoại chatbot đóng vai trò là nơi lưu trữ dữ liệu tạm thời trong suốt cuộc trò chuyện (Temporary data storage during conversations) để truyền thông tin qua lại giữa các nút hội thoại.}
\sources{IBM Watson Assistant Documentation: Context Variables}
\explain{Hướng giải: Định nghĩa bản chất kỹ thuật của 'Biến ngữ cảnh' (Context variables) trong Chatbot.

Các bước phân tích:
1. Trong kiến trúc máy trạng thái (State Machine) của Chatbot, biến ngữ cảnh là các cặp `key: value` trong bộ nhớ phiên (session memory).
2. Chúng lưu tạm các giá trị như: tên khách hàng, số điện thoại, trạng thái đăng nhập, bước hội thoại hiện tại.
3. Khi phiên hội thoại kết thúc, dữ liệu này có thể được giải phóng hoặc lưu vào cơ sở dữ liệu vĩnh viễn.

Kiểm tra: 'Temporary data storage during conversations' phản ánh chính xác 100% bản chất lập trình của biến ngữ cảnh.

Vì sao các phương án khác sai:
- A: Phương thức xác thực người dùng là Authentication protocols (OAuth/JWT).
- C: Thiết kế đồ họa là Graphic UI.
- D: Cấu hình máy chủ là Server Config.

Lưu ý/Khái niệm cần nhớ: Context Variables = Bộ nhớ lưu trữ dữ liệu tạm thời trong phiên trò chuyện.}
\vocab{Context Variable | Biến ngữ cảnh | Biến lưu trữ trạng thái phiên làm việc trong hệ thống chatbot}
\end{question}

% ==============================================================================
% Question 149 | Phân loại: 4. Building AI Powered Chatbots Without Programming
% Xuất hiện trong các đề gốc: AIG202c_SP_2025_FE (Q57), AIG202c_SU25_FE (Q59), AIG202c_SP26_RE (Q56)
% ==============================================================================
\begin{question}{149}{mcq}
\source{AIG202c_-_SP_2025_-_FE_3755_57.webp}
\prompt{How do chatbots use context variables and slots to maintain context during a conversation?}
\choice{A}{By ignoring the user's input}
\choice{B}{By storing information about the user's preferences and previous interactions}
\choice{C}{By generating random responses}
\choice{D}{By terminating the conversation after a set number of interactions}
\answer{B}
\confidence{high}
\evidence{Chatbot duy trì ngữ cảnh trong suốt cuộc trò chuyện bằng cách sử dụng các biến ngữ cảnh và ô nhớ (context variables and slots) để lưu trữ thông tin về sở thích của người dùng và các lượt tương tác trước đó (storing information about user's preferences and previous interactions).}
\sources{IBM Watson Assistant: Context Management; Jurafsky & Martin, Dialog Management}
\explain{Hướng giải: Xác định cơ chế kỹ thuật giúp chatbot duy trì ngữ cảnh hội thoại xuyên suốt nhiều lượt đối thoại.

Các bước phân tích:
1. Giao tiếp hội thoại là quá trình tích lũy thông tin liên tục.
2. Biến ngữ cảnh (`$context_variables`) và các khe thông tin (`slots`) đóng vai trò như bộ nhớ tạm thời của chatbot, lưu giữ các tham số đã thu thập từ các câu nói trước (ví dụ: tên khách, điểm đến, sở thích) để tái sử dụng trong các bước sau.

Kiểm tra: Đây là nguyên lý hoạt động chuẩn của các công cụ quản trị hội thoại (Dialog State Tracking).

Vì sao các phương án khác sai:
- A: Bỏ qua thông tin người dùng là hành vi phá vỡ ngữ cảnh.
- C: Tạo phản hồi ngẫu nhiên là bot không có ngữ cảnh.
- D: Ngắt kết nối sau số lượt cố định là giới hạn kỹ thuật, không phải cách duy trì ngữ cảnh.

Lưu ý/Khái niệm cần nhớ: Context variables & slots = Bộ nhớ lưu trữ thông tin trạng thái xuyên suốt phiên đối thoại.}
\vocab{Session Context | Ngữ cảnh phiên làm việc | Toàn bộ trạng thái và dữ liệu được ghi nhận trong suốt một phiên người dùng tương tác}
\end{question}

% ==============================================================================
% Question 150 | Phân loại: 4. Building AI Powered Chatbots Without Programming
% Xuất hiện trong các đề gốc: AIG202c_SP_2025_FE (Q60)
% ==============================================================================
\begin{question}{150}{mcq}
\source{AIG202c_-_SP_2025_-_FE_3755_60.webp}
\prompt{Which concept does the introduction module primarily cover in chatbot development without programming?}
\choice{A}{Complex algorithms in cryptography}
\choice{B}{Front-end web development techniques}
\choice{C}{Basic principles and concepts of chatbots}
\choice{D}{Robotics in AI applications}
\answer{C}
\confidence{high}
\evidence{Mô-đun giới thiệu trong việc phát triển chatbot không cần lập trình (no-code chatbot development) chủ yếu bao quát các nguyên lý và khái niệm cơ bản về chatbot (Basic principles and concepts of chatbots) như Intents, Entities, và Dialog.}
\sources{IBM Watson Assistant: Getting Started Without Coding; Coursera Chatbot Curriculum}
\explain{Hướng giải: Xác định nội dung bài nhập môn phát triển chatbot no-code (câu hỏi kiểm tra tính nhất quán với UQ052).

Các bước phân tích:
1. Phát triển chatbot không dùng code (no-code) tập trung vào logic hội thoại và thiết kế kịch bản hơn là cú pháp lập trình.
2. Bài giới thiệu đầu tiên luôn trang bị cho học viên các khái niệm nền tảng: thế nào là một chatbot, luồng đối thoại, ý định và thực thể.

Kiểm tra: Khớp hoàn toàn với định hướng đào tạo của học phần.

Vì sao các phương án khác sai:
- A (Mã hóa mật mã học), B (Lập trình web front-end), D (Robot) đều là các chủ đề không liên quan đến bài mở đầu chatbot no-code.

Lưu ý/Khái niệm cần nhớ: Bài mở đầu no-code chatbot = Nguyên lý và khái niệm cơ bản về hệ thống hội thoại.}
\vocab{No-Code Development | Phát triển không cần viết mã | Phương thức xây dựng phần mềm thông qua giao diện đồ họa trực quan}
\end{question}

% ==============================================================================
% Question 151 | Phân loại: 4. Building AI Powered Chatbots Without Programming
% Xuất hiện trong các đề gốc: AIG202c_SU25_FE (Q46)
% ==============================================================================
\begin{question}{151}{mcq}
\source{Q46.webp}
\prompt{What role does 'dialog' play in AI-powered chatbots?}
\choice{A}{Analyzing website traffic}
\choice{B}{Generating conversational responses based on user inputs}
\choice{C}{Optimizing database management}
\choice{D}{Improving encryption algorithms}
\answer{B}
\confidence{high}
\evidence{Mô-đun Hội thoại ('dialog') trong chatbot AI đóng vai trò tạo ra các phản hồi đàm thoại phù hợp dựa trên các câu nói đầu vào của người dùng (Generating conversational responses based on user inputs).}
\sources{IBM Watson Assistant: The Dialog Flow; Conversational Response Generation}
\explain{Hướng giải: Xác định vai trò của thành phần Dialog trong hệ thống Chatbot.

Các bước phân tích:
1. Dialog là nơi kết nối giữa đầu vào (đã được hiểu qua Intents và Entities) với đầu ra (câu trả lời của bot).
2. Vai trò cốt lõi của Dialog là xác định xem với ý định này và thực thể này, bot cần tạo ra câu phản hồi đàm thoại (conversational response) nào là phù hợp nhất.

Kiểm tra: Dialog là bộ sinh phản hồi trung tâm của chatbot.

Vì sao các phương án khác sai:
- A: Phân tích lưu lượng web là Web Analytics.
- C: Tối ưu cơ sở dữ liệu là Database Admin.
- D: Cải thiện thuật toán mã hóa là Cryptography.

Lưu ý/Khái niệm cần nhớ: Dialog = Tạo ra các phản hồi đàm thoại tương ứng với người dùng.}
\vocab{Dialog Node | Nút hội thoại | Đơn vị logic trong kịch bản chatbot chứa điều kiện kích hoạt và câu phản hồi tương ứng}
\end{question}

% ==============================================================================
% Question 152 | Phân loại: 4. Building AI Powered Chatbots Without Programming
% Xuất hiện trong các đề gốc: AIG202c_SU25_FE (Q47), AIG202c_FA25_FE (Q59), AIG202c_SU26_FE (Q46)
% ==============================================================================
\begin{question}{152}{mcq}
\source{Q47.webp}
\prompt{How do digressions affect the flow of a conversation between a user and a chatbot?}
\choice{A}{Digressions help maintain focus on the original topic of the conversation}
\choice{B}{Digressions create confusion and frustration for users by shifting the conversation away from their intended query}
\choice{C}{Digressions encourage users to provide more information about their preferences and interests}
\choice{D}{Digressions have no impact on the flow of the conversation}
\answer{C}
\confidence{high}
\evidence{Khi chatbot xử lý tốt sự lạc đề (digressions), nó tạo cơ hội mở rộng đối thoại, khuyến khích người dùng cung cấp thêm thông tin về sở thích và mối quan tâm cá nhân (encourage users to provide more information about preferences).}
\sources{Conversational UX Design: Eliciting User Preferences Through Digressions}
\explain{Hướng giải: Đánh giá tác động tích cực của việc xử lý lạc đề đúng cách đối với dòng chảy hội thoại.

Các bước phân tích:
1. Khi người dùng lạc đề sang một chủ đề mới và bot phản hồi thỏa đáng, người dùng cảm thấy được lắng nghe và tin tưởng bot hơn.
2. Từ đó, người dùng cởi mở chia sẻ thêm thông tin ngữ cảnh, giúp hệ thống thu thập thêm dữ liệu sở thích để phục vụ tốt hơn.

Kiểm tra: Kỹ thuật khai thác thông tin qua các nhánh đối thoại phụ là bài học quan trọng trong thiết kế chatbot thông minh.

Vì sao các phương án khác sai:
- A: Lạc đề là chuyển hướng, không phải duy trì tập trung vào chủ đề gốc.
- B: Chỉ gây bực bội khi bot xử lý tồi; nếu bot xử lý chuẩn xác thì tạo ra trải nghiệm tốt.
- D: Lạc đề làm thay đổi hoàn toàn mạch hội thoại, không thể 'no impact'.

Lưu ý/Khái niệm cần nhớ: Lạc đề được dẫn dắt khéo léo sẽ làm phong phú dữ liệu người dùng.}
\vocab{Preference Elicitation | Khai thác sở thích | Quá trình tương tác để thu thập các tiêu chí ưu tiên của người dùng}
\end{question}

% ==============================================================================
% Question 153 | Phân loại: 4. Building AI Powered Chatbots Without Programming
% Xuất hiện trong các đề gốc: AIG202c_SU25_FE (Q48)
% ==============================================================================
\begin{question}{153}{mcq}
\source{Q48.webp}
\prompt{How can chatbots use context variables to manage digressions in conversations?}
\choice{A}{By ignoring the user's input and continuing with a predefined script}
\choice{B}{By storing information about the original topic and using it to guide the conversation back on track}
\choice{C}{By terminating the conversation abruptly}
\choice{D}{By generating random responses to user queries}
\answer{B}
\confidence{high}
\evidence{Chatbot sử dụng các biến ngữ cảnh (context variables) để quản lý sự lạc đề bằng cách lưu trữ thông tin về chủ đề gốc và sử dụng thông tin đó để dẫn dắt cuộc trò chuyện quay trở lại đúng hướng sau khi phân nhánh lạc đề kết thúc (storing information about original topic and using it to guide the conversation back on track).}
\sources{IBM Watson Assistant: Digression Return Logic and Context Retention}
\explain{Hướng giải: Giải thích cơ chế kỹ thuật giúp chatbot quay lại luồng chính sau khi người dùng lạc đề.

Các bước phân tích:
1. Khi người dùng đang thực hiện dở dang một giao dịch (ví dụ: đang đặt bàn) mà bất ngờ chuyển sang hỏi thời tiết, hệ thống sẽ kích hoạt tính năng Digression.
2. Trước khi chuyển sang nút trả lời thời tiết, hệ thống dùng biến ngữ cảnh để lưu lại trạng thái: vị trí nút đặt bàn đang dừng và các thông tin đã điền.
3. Sau khi trả lời xong câu hỏi thời tiết, chatbot đọc lại biến ngữ cảnh này để nói: 'Bây giờ chúng ta quay lại việc đặt bàn nhé, bạn dự định đi mấy người?' (guide conversation back on track).

Kiểm tra: Đây chính là cơ chế 'Return to original node' kinh điển của IBM Watson Assistant.

Vì sao các phương án khác sai:
- A: Bỏ qua câu hỏi người dùng là cách xử lý tồi.
- C: Ngắt kết nối đột ngột là lỗi hệ thống.
- D: Trả lời ngẫu nhiên phá vỡ logic giao tiếp.

Lưu ý/Khái niệm cần nhớ: Biến ngữ cảnh lưu giữ vị trí chủ đề gốc để bot có thể quay lại sau khi trả lời câu hỏi lạc đề.}
\vocab{Digression Return | Cơ chế quay lại sau lạc đề | Tính năng tự động đưa cuộc trò chuyện trở về luồng tác vụ chính sau khi xử lý xong câu hỏi phụ}
\end{question}

% ==============================================================================
% Question 154 | Phân loại: 4. Building AI Powered Chatbots Without Programming
% Xuất hiện trong các đề gốc: AIG202c_SU25_FE (Q49)
% ==============================================================================
\begin{question}{154}{mcq}
\source{Q49.webp}
\prompt{How do 'digressions' aid in improving chatbot interactions?}
\choice{A}{By maintaining a robotic conversational tone}
\choice{B}{By avoiding deviations from the main conversation topic}
\choice{C}{By introducing engaging and varied conversational paths}
\choice{D}{By restricting user inputs}
\answer{C}
\confidence{high}
\evidence{Các phân nhánh lạc đề (digressions) giúp cải thiện tương tác của chatbot bằng cách đưa vào các lộ trình đối thoại hấp dẫn, sinh động và đa dạng (introducing engaging and varied conversational paths), giúp cuộc trò chuyện không bị cứng nhắc.}
\sources{Conversational Design Best Practices: The Power of Digressions}
\explain{Hướng giải: Xác định lợi ích của tính năng Digression đối với tương tác chatbot (câu hỏi kiểm tra tính nhất quán với UQ094 và UQ128).

Các bước phân tích:
1. Tương tác chatbot trở nên thú vị hơn khi người dùng không bị ép buộc phải trả lời theo mẫu cứng nhắc.
2. Khả năng hỗ trợ lạc đề mở ra nhiều con đường đàm thoại phong phú, đa dạng và tạo cảm giác giao tiếp tự nhiên.

Kiểm tra: Khớp hoàn toàn với các câu hỏi tương tự trong đề thi.

Vì sao các phương án khác sai:
- A: Duy trì giọng điệu robot là điểm trừ, không phải cải tiến.
- B: Tránh chuyển hướng là đặc điểm của bot cứng nhắc thế hệ cũ.
- D: Hạn chế người dùng nhập liệu là trải nghiệm xấu.

Lưu ý/Khái niệm cần nhớ: Digressions tạo ra các lộ trình đàm thoại đa dạng và hấp dẫn.}
\vocab{Conversational Flexibility | Tính linh hoạt trong hội thoại | Khả năng điều hướng đa chiều trong cuộc trò chuyện}
\end{question}

% ==============================================================================
% Question 155 | Phân loại: 4. Building AI Powered Chatbots Without Programming
% Xuất hiện trong các đề gốc: AIG202c_SU25_FE (Q51)
% ==============================================================================
\begin{question}{155}{mcq}
\source{Q51.webp}
\prompt{In our flower shop chatbot, "dad" is defined as a synonym for @relationship:father. If the user enters 'Which flowers for dad", what's the entity value detected by Watson?}
\choice{A}{@relationship:dad}
\choice{B}{@relationship:father}
\choice{C}{@dad}
\choice{D}{@flower_suggestions}
\choice{E}{@relationship}
\answer{B}
\confidence{high}
\evidence{Trong IBM Watson Assistant, khi một từ đồng nghĩa (synonym) được định nghĩa cho một thực thể, giá trị thực thể chuẩn (canonical value) sẽ được trả về khi từ đồng nghĩa xuất hiện. Do đó, 'dad' là từ đồng nghĩa của `@relationship:father`, khi người dùng nhập 'Which flowers for dad', giá trị thực thể được Watson nhận diện chính thức là `@relationship:father`.}
\sources{IBM Watson Assistant Documentation: Entity Synonyms and Canonical Values}
\explain{Hướng giải: Xác định giá trị thực thể chuẩn (canonical entity value) được hệ thống NLU nhận diện khi người dùng sử dụng từ đồng nghĩa (synonym).

Các bước phân tích:
1. Trong cấu hình Thực thể (Entity) của Watson Assistant:
   - Tên thực thể: `@relationship`.
   - Giá trị chuẩn (Canonical value): `father`.
   - Các từ đồng nghĩa (Synonyms): `dad`, `daddy`, `papa`, `father`.
2. Khi người dùng nhập câu: 'Which flowers for dad', khối NLU nhận diện từ 'dad' khớp với danh sách từ đồng nghĩa của giá trị `father`.
3. Hệ thống sẽ tự động chuẩn hóa và gán giá trị chính thức là `@relationship:father` để các nút hội thoại phía sau xử lý logic thống nhất mà không cần phải viết điều kiện cho từng từ đồng nghĩa.

Kiểm tra: Đây là nguyên lý hoạt động căn bản của cơ chế Entity Synonyms trong mọi hệ thống Chatbot chuyên nghiệp.

Vì sao các phương án khác sai:
- A (`@relationship:dad`): 'dad' chỉ là từ đồng nghĩa bề mặt, không phải giá trị chuẩn được hệ thống trả về.
- C (`@dad`): Sai cú pháp tên thực thể.
- D (`@flower_suggestions`): Tên thực thể không liên quan đến quan hệ gia đình.
- E (`@relationship`): Đây chỉ là tên nhóm thực thể, chưa chỉ rõ giá trị con.

Lưu ý/Khái niệm cần nhớ: Watson Assistant luôn ánh xạ từ đồng nghĩa về giá trị chuẩn tắc (canonical value) của thực thể.}
\vocab{Canonical Value | Giá trị chuẩn tắc | Giá trị đại diện chính thức của một thực thể sau khi đã chuẩn hóa từ các từ đồng nghĩa}
\end{question}

% ==============================================================================
% Question 156 | Phân loại: 4. Building AI Powered Chatbots Without Programming
% Xuất hiện trong các đề gốc: AIG202c_SU25_FE (Q52), AIG202c_FA25_FE (Q57), AIG202c_SP26_RE (Q55)
% ==============================================================================
\begin{question}{156}{mcq}
\source{Q52.webp}
\prompt{What is the primary goal of 'digressions' in chatbot interactions, particularly concerning user experience?}
\choice{A}{To streamline conversations by avoiding diversions}
\choice{B}{To enforce strict conversational rules}
\choice{C}{To introduce spontaneous and engaging conversation elements}
\choice{D}{To restrict conversational freedom for users}
\answer{C}
\confidence{high}
\evidence{Mục tiêu chính của việc hỗ trợ 'lạc đề' (digressions) trong tương tác chatbot là đưa vào các yếu tố đối thoại tự nhiên, linh hoạt và lôi cuốn (introduce spontaneous and engaging conversation elements), giúp trải nghiệm giao tiếp bớt máy móc.}
\sources{IBM Watson Assistant: Configuring Digressions for Better UX}
\explain{Hướng giải: Hiểu mục đích nâng cao trải nghiệm người dùng của tính năng 'Digression'.

Các bước phân tích:
1. Con người không nói chuyện như một máy móc tuần tự cứng nhắc.
2. Cho phép người dùng chuyển hướng tạm thời (hỏi một câu ngoài lề) rồi quay lại tạo cảm giác trò chuyện tự nhiên, sinh động và thú vị (engaging).

Kiểm tra: Tính năng Digression trong Watson Assistant được thiết kế chuyên biệt để hỗ trợ sự ngẫu hứng (spontaneity) của người dùng.

Vì sao các phương án khác sai:
- A & B: Ép buộc quy tắc cứng nhắc và tránh chuyển hướng là đặc điểm của bot kém thông minh thế hệ cũ.
- D: Hạn chế tự do là đi ngược lại mục tiêu trải nghiệm người dùng.

Lưu ý/Khái niệm cần nhớ: Digression mang lại sự linh hoạt và tính tự nhiên cho cuộc trò chuyện.}
\vocab{Conversational Spontaneity | Tính ngẫu hứng trong hội thoại | Khả năng thích ứng linh hoạt với các chuyển hướng bất ngờ của người dùng}
\end{question}

% ==============================================================================
% Question 157 | Phân loại: 4. Building AI Powered Chatbots Without Programming
% Xuất hiện trong các đề gốc: AIG202c_SU25_FE (Q53)
% ==============================================================================
\begin{question}{157}{mcq}
\source{Q53.webp}
\prompt{Why is 'deployment' crucial in chatbot development?}
\choice{A}{Ensuring optimal user experiences}
\choice{B}{Managing software licenses}
\choice{C}{Hardware optimization}
\choice{D}{Designing backend algorithms}
\answer{A}
\confidence{high}
\evidence{Giai đoạn triển khai (deployment) đóng vai trò quyết định trong phát triển chatbot vì nó đưa sản phẩm vào môi trường thực tế, tích hợp vào các kênh liên lạc (web, mobile, mạng xã hội) để phục vụ người dùng cuối và đảm bảo trải nghiệm người dùng tối ưu (Ensuring optimal user experiences).}
\sources{Chatbot Lifecycle Management; IBM Watson Assistant Deployment Channels}
\explain{Hướng giải: Giải thích tầm quan trọng của giai đoạn Triển khai (Deployment) trong phát triển Chatbot.

Các bước phân tích:
1. Một chatbot dù được thiết kế kịch bản hay đến đâu trong môi trường thử nghiệm (Sandbox) cũng chưa mang lại giá trị kinh doanh nếu chưa được triển khai (Deploy).
2. Deployment là quá trình tích hợp bot vào các kênh thực tế (Website doanh nghiệp, Facebook Messenger, Zalo, WhatsApp, Ứng dụng di động) để người dùng thực sự có thể trò chuyện và giải quyết công việc.
3. Quá trình triển khai bảo đảm tính sẵn sàng, khả năng chịu tải và trải nghiệm người dùng mượt mà nhất.

Kiểm tra: 'Ensuring optimal user experiences' trong môi trường thực là mục tiêu cuối cùng của việc triển khai.

Vì sao các phương án khác sai:
- B: Quản lý bản quyền phần mềm là thủ tục hành chính IT.
- C: Tối ưu hóa phần cứng là nhiệm vụ của kỹ sư hạ tầng/server.
- D: Thiết kế thuật toán backend diễn ra ở giai đoạn phát triển ban đầu, không phải mục đích của deployment.

Lưu ý/Khái niệm cần nhớ: Deployment = Đưa chatbot ra môi trường thực tế phục vụ trải nghiệm người dùng tối ưu.}
\vocab{Chatbot Deployment | Triển khai chatbot | Quá trình đưa chatbot lên các kênh giao tiếp chính thức để phục vụ người dùng thực tế}
\end{question}

% ==============================================================================
% Question 158 | Phân loại: 4. Building AI Powered Chatbots Without Programming
% Xuất hiện trong các đề gốc: AIG202c_SU25_FE (Q54), AIG202c_SP26_RE (Q50)
% ==============================================================================
\begin{question}{158}{mcq}
\source{Q54.webp}
\prompt{What are "entities" in the context of artificial intelligence and chatbots?}
\choice{A}{The personalities assigned to chatbots}
\choice{B}{The actions taken by chatbots to achieve user goals}
\choice{C}{Specific pieces of information within user input that the chatbot needs to extract}
\choice{D}{The language used by chatbots to communicate with users}
\answer{C}
\confidence{high}
\evidence{Trong bối cảnh AI và chatbot, 'Thực thể' (Entities) là những mẩu thông tin cụ thể nằm trong câu nói của người dùng mà chatbot cần trích xuất để thực hiện tác vụ (Specific pieces of information within user input that the chatbot needs to extract).}
\sources{IBM Watson Assistant: Understanding Entities; Manning, Information Extraction}
\explain{Hướng giải: Định nghĩa khái niệm Thực thể (Entity) trong xử lý ngôn ngữ cho chatbot.

Các bước phân tích:
1. Câu nói của người dùng chứa các dữ liệu biến đổi quan trọng: ngày tháng, tên người, thành phố, số lượng, loại sản phẩm.
2. Các dữ liệu này được gọi là Thực thể (Entity).
3. Ví dụ: trong câu 'Đặt bàn cho 4 người lúc 7 giờ tối', các thực thể là: `@number: 4`, `@time: 7pm`.

Kiểm tra: Khớp hoàn toàn với định nghĩa chuẩn: 'Entities represent specific information in the user input that is relevant to the user\'s intent'.

Vì sao các phương án khác sai:
- A: Tính cách của bot là Chatbot Persona.
- B: Hành động bot thực hiện là Bot Actions / Webhooks.
- D: Ngôn ngữ giao tiếp là Natural Language.

Lưu ý/Khái niệm cần nhớ: Entities = Mẩu thông tin cụ thể cần trích xuất từ câu nói của người dùng.}
\vocab{Entity Extraction | Trích xuất thực thể | Kỹ thuật NLP nhằm nhận diện và trích xuất các từ khóa thông tin cụ thể từ văn bản}
\end{question}

% ==============================================================================
% Question 159 | Phân loại: 4. Building AI Powered Chatbots Without Programming
% Xuất hiện trong các đề gốc: AIG202c_SU25_FE (Q55)
% ==============================================================================
\begin{question}{159}{mcq}
\source{Q55.webp}
\prompt{Which aspect of chatbot development is emphasized in the 'Context Variables & Slots' module?}
\choice{A}{Hardware maintenance}
\choice{B}{Managing and tracking information during conversations}
\choice{C}{Cryptocurrency analysis}
\choice{D}{Mobile application development}
\answer{B}
\confidence{high}
\evidence{Mô-đun 'Context Variables & Slots' trong phát triển chatbot chú trọng vào việc quản lý và theo dõi thông tin xuyên suốt cuộc trò chuyện (Managing and tracking information during conversations).}
\sources{IBM Watson Assistant: Context Variables and Slots Module}
\explain{Hướng giải: Xác định nội dung trọng tâm của mô-đun 'Context Variables & Slots'.

Các bước phân tích:
1. Biến ngữ cảnh (Context Variables) dùng để lưu trữ trạng thái người dùng.
2. Khe thông tin (Slots) dùng để bóc tách và thu thập các trường dữ liệu cần thiết.
3. Cả hai thành phần này cùng phối hợp để thực hiện nhiệm vụ: Quản lý và theo dõi thông tin (Managing and tracking information) trong suốt quá trình đàm thoại.

Kiểm tra: Đây là định nghĩa mục tiêu học tập chuẩn của mô-đun này.

Vì sao các phương án khác sai:
- A (Hardware maintenance), C (Cryptocurrency), D (Mobile dev) hoàn toàn không liên quan đến logic ngữ cảnh của chatbot.

Lưu ý/Khái niệm cần nhớ: Context Variables & Slots = Quản lý và theo dõi thông tin đàm thoại.}
\vocab{State Management | Quản lý trạng thái | Việc theo dõi toàn bộ thông tin biến số và tiến trình trong suốt một phiên tương tác}
\end{question}

% ==============================================================================
% Question 160 | Phân loại: 4. Building AI Powered Chatbots Without Programming
% Xuất hiện trong các đề gốc: AIG202c_SU25_FE (Q56), AIG202c_SU26_FE (Q51)
% ==============================================================================
\begin{question}{160}{mcq}
\source{Q56.webp}
\prompt{What does 'deployment' refer to in the context of chatbot development?}
\choice{A}{Designing user interfaces}
\choice{B}{The process of launching and implementing a chatbot}
\choice{C}{Hardware configurations}
\choice{D}{Data encryption protocols}
\answer{B}
\confidence{high}
\evidence{Trong ngữ cảnh phát triển chatbot, 'triển khai' (deployment) đề cập đến quá trình phát hành, tích hợp và đưa chatbot vào hoạt động thực tế trên các kênh giao tiếp (The process of launching and implementing a chatbot).}
\sources{Software Deployment Lifecycle; IBM Watson Assistant Integration & Deployment}
\explain{Hướng giải: Định nghĩa thuật ngữ 'Deployment' (Triển khai) trong phát triển Chatbot.

Các bước phân tích:
1. Sau khi thiết kế Intent, Entity và Dialog xong, bước tiếp theo là đưa bot ra sử dụng.
2. 'Deployment' là quá trình phát hành chính thức (launching) và tích hợp (implementing) chatbot vào các nền tảng giao tiếp của doanh nghiệp (như Web chat widget, điện thoại tổng đài, ứng dụng di động).

Kiểm tra: Định nghĩa chuẩn trong kỹ thuật phần mềm: Deployment is all of the activities that make a software system available for use.

Vì sao các phương án khác sai:
- A: Thiết kế giao diện là UI/UX Design.
- C: Cấu hình phần cứng là Hardware setup.
- D: Giao thức mã hóa dữ liệu là Security protocols.

Lưu ý/Khái niệm cần nhớ: Deployment = Quá trình ra mắt và đưa chatbot vào hoạt động thực tế.}
\vocab{Production Deployment | Triển khai môi trường sản xuất | Việc đưa phần mềm vào môi trường hoạt động thực tế cho người dùng cuối sử dụng}
\end{question}

% ==============================================================================
% Question 161 | Phân loại: 4. Building AI Powered Chatbots Without Programming
% Xuất hiện trong các đề gốc: AIG202c_SU25_FE (Q57), AIG202c_SU26_FE (Q49)
% ==============================================================================
\begin{question}{161}{mcq}
\source{Q57.webp}
\prompt{Which technology is commonly used to enable chatbots to understand and respond to natural language input?}
\choice{A}{Machine Learning}
\choice{B}{Virtual Reality}
\choice{C}{Augmented Reality}
\choice{D}{Blockchain}
\answer{A}
\confidence{high}
\evidence{Học máy (Machine Learning) kết hợp với Xử lý ngôn ngữ tự nhiên (NLP) là công nghệ nền tảng được sử dụng phổ biến nhất để giúp chatbot hiểu và phản hồi ngôn ngữ tự nhiên của con người.}
\sources{Russell & Norvig, AI: A Modern Approach; Manning, Foundations of Statistical NLP}
\explain{Hướng giải: Xác định công nghệ cốt lõi giúp chatbot hiểu được ngôn ngữ tự nhiên (NLU/NLP).

Các bước phân tích:
1. Để hiểu được câu nói không tuân theo quy tắc cố định của con người, chatbot hiện đại sử dụng các mô hình Học máy (Machine Learning) và Học sâu (Deep Learning).
2. Thuật toán ML phân loại ý định (Intent Classification) và bóc tách thực thể (Named Entity Recognition) dựa trên xác suất thống kê từ tập dữ liệu mẫu.

Kiểm tra: Machine Learning là nền tảng của mọi hệ thống NLU hiện đại.

Vì sao các phương án khác sai:
- B (Virtual Reality): Công nghệ thực tế ảo tạo môi trường 3D, không dùng để hiểu ngôn ngữ.
- C (Augmented Reality): Thực tế tăng cường chèn hình ảnh ảo vào thế giới thực.
- D (Blockchain): Công nghệ sổ cái phân tán lưu trữ giao dịch bất biến, không xử lý ngôn ngữ tự nhiên.

Lưu ý/Khái niệm cần nhớ: Machine Learning là công nghệ cốt lõi cho khả năng hiểu ngôn ngữ của chatbot.}
\vocab{Natural Language Understanding (NLU) | Hiểu ngôn ngữ tự nhiên | Phân nhánh của AI và học máy giúp máy tính hiểu được ý nghĩa ngữ nghĩa của văn bản}
\end{question}

% ==============================================================================
% Question 162 | Phân loại: 4. Building AI Powered Chatbots Without Programming
% Xuất hiện trong các đề gốc: AIG202c_SU25_FE (Q58), AIG202c_SU26_FE (Q55)
% ==============================================================================
\begin{question}{162}{mcq}
\source{Q58.webp}
\prompt{What is a "digression" in the context of artificial intelligence and chatbots?}
\choice{A}{A redirection of the conversation to a different topic unrelated to the original query}
\choice{B}{A response generated by the chatbot to acknowledge the user's input}
\choice{C}{A technique used to analyze the sentiment of user messages}
\choice{D}{A strategy employed by chatbots to end conversations abruptly}
\answer{A}
\confidence{high}
\evidence{Sự lạc đề ('digression') trong bối cảnh chatbot AI là sự chuyển hướng tạm thời của cuộc trò chuyện sang một chủ đề khác không liên quan trực tiếp đến truy vấn hoặc luồng tác vụ ban đầu (A redirection of the conversation to a different topic unrelated to the original query).}
\sources{IBM Watson Assistant: Understanding Digressions; Conversational Interaction Patterns}
\explain{Hướng giải: Định nghĩa thuật ngữ 'Digression' (Lạc đề / Chuyển hướng đối thoại) trong Chatbot.

Các bước phân tích:
1. 'Digression' xảy ra khi người dùng đang ở giữa một luồng tác vụ (ví dụ: đang điền thông tin mở tài khoản ngân hàng) nhưng bất ngờ đặt một câu hỏi hoàn toàn khác (ví dụ: 'Lãi suất tiết kiệm kỳ hạn 6 tháng là bao nhiêu?').
2. Hành động rẽ nhánh sang chủ đề khác này được định nghĩa chính thức là một 'digression' (redirection of conversation).

Kiểm tra: Thuật ngữ này được Watson Assistant chuẩn hóa trong tài liệu hướng dẫn thiết kế hội thoại.

Vì sao các phương án khác sai:
- B: Phản hồi ghi nhận thông tin là Acknowledgment message.
- C: Kỹ thuật phân tích cảm xúc là Sentiment Analysis.
- D: Ngắt kết nối là Conversation termination.

Lưu ý/Khái niệm cần nhớ: Digression = Chuyển hướng cuộc trò chuyện sang chủ đề khác ngoài luồng chính.}
\vocab{Conversational Digression | Lạc đề hội thoại | Hiện tượng người dùng chuyển đổi chủ đề đàm thoại ngoài kịch bản dự kiến}
\end{question}

% ==============================================================================
% Question 163 | Phân loại: 4. Building AI Powered Chatbots Without Programming
% Xuất hiện trong các đề gốc: AIG202c_SU25_FE (Q60), AIG202c_SU26_FE (Q52)
% ==============================================================================
\begin{question}{163}{mcq}
\source{Q60.webp}
\prompt{How are intents used in natural language understanding for chatbots?}
\choice{A}{To detect the emotion of the user}
\choice{B}{To identify the language being spoken}
\choice{C}{To classify the purpose or goal of the user's input}
\choice{D}{To generate random responses}
\answer{C}
\confidence{high}
\evidence{Ý định (Intents) được sử dụng trong NLU để phân loại mục đích hoặc mục tiêu của câu nói đầu vào từ người dùng (To classify the purpose or goal of the user's input).}
\sources{IBM Watson Assistant: Intents in NLU; Jurafsky & Martin, Speech and Language Processing}
\explain{Hướng giải: Xác định mục đích sử dụng của Intent trong tầng NLU của Chatbot (câu hỏi kiểm tra tính nhất quán với UQ099).

Các bước phân tích:
1. NLU tiếp nhận câu văn bản thô từ người dùng.
2. Nhiệm vụ phân loại Intent là gán câu văn đó vào mục đích chính xác: hỏi thông tin, đặt hàng, khiếu nại, chào hỏi.
3. Do đó, 'To classify the purpose or goal of the user\'s input' là định nghĩa chính xác nhất.

Kiểm tra: Intent Classification là bài toán phân loại văn bản cơ bản trong NLP.

Vì sao các phương án khác sai:
- A: Phát hiện cảm xúc là Sentiment/Emotion Detection.
- B: Nhận diện ngôn ngữ là Language Identification.
- D: Sinh phản hồi ngẫu nhiên là sai hoàn toàn.

Lưu ý/Khái niệm cần nhớ: Intents = Phân loại mục đích hoặc mục tiêu của câu nói đầu vào.}
\vocab{Intent Classification | Phân loại ý định | Thuật toán học máy gán nhãn mục đích cho câu nói của người dùng}
\end{question}

% ==============================================================================
% Question 164 | Phân loại: 4. Building AI Powered Chatbots Without Programming
% Xuất hiện trong các đề gốc: AIG202c_FA25_FE (Q24)
% ==============================================================================
\begin{question}{164}{mcq}
\source{Q24.webp}
\prompt{In a generic chatbot platform, Which variable contains the entire text input from the user?}
\choice{A}{@sys-date}
\choice{B}{text.input}
\choice{C}{input-text}
\choice{D}{user.text}
\choice{E}{None of the others}
\answer{C}
\confidence{high}
\evidence{Trong nền tảng phát triển chatbot Watson Assistant / generic dialog chatbot, biến ngữ cảnh chứa toàn bộ chuỗi văn bản đầu vào của người dùng là `input.text` (hoặc `input-text`).}
\sources{IBM Watson Assistant Dialog Documentation: Context Variables and SpEL expressions}
\explain{Hướng giải: Xác định biến hệ thống chuẩn lưu trữ văn bản người dùng nhập vào trong nền tảng chatbot.

Các bước phân tích:
1. Khi người dùng gửi một tin nhắn đến chatbot, hệ thống tiếp nhận chuỗi ký tự đó vào đối tượng ngữ cảnh `input`.
2. Thuộc tính lưu trữ chuỗi văn bản gốc là `input.text` (trong định dạng biểu thức thường được ánh xạ là `input-text`).

Kiểm tra: Biến `@sys-date` là thực thể ngày tháng (system entity), không phải toàn bộ văn bản.

Vì sao các phương án khác sai:
- A (@sys-date): Thực thể hệ thống nhận diện ngày tháng.
- B, D: Tên biến không đúng chuẩn cú pháp của nền tảng chatbot IBM/generic.

Lưu ý/Khái niệm cần nhớ: `input.text` chứa toàn bộ câu nói người dùng gửi đến bot.}
\vocab{Input Variable | Biến đầu vào | Biến hệ thống chứa dữ liệu văn bản người dùng truyền vào trong phiên hội thoại}
\end{question}

% ==============================================================================
% Question 165 | Phân loại: 4. Building AI Powered Chatbots Without Programming
% Xuất hiện trong các đề gốc: AIG202c_FA25_FE (Q46), AIG202c_SU26_FE (Q57)
% ==============================================================================
\begin{question}{165}{mcq}
\source{Q46.webp}
\prompt{Which of the following is NOT an example of a dialog flow for a restaurant reservation chatbot?}
\choice{A}{Greeting the user}
\choice{B}{Asking for the user's name}
\choice{C}{Confirming the user's credit card number}
\choice{D}{Providing the restaurant's menu}
\answer{D}
\confidence{high}
\evidence{Một luồng hội thoại đặt bàn nhà hàng bao gồm: Chào hỏi khách, Hỏi tên, Hỏi số lượng khách, thời gian và Xác nhận đặt bàn. Các tác vụ như 'Tính toán quỹ đạo tên lửa' hoặc các hành vi nằm ngoài phạm vi nhà hàng không thuộc luồng hội thoại này.}
\sources{Conversational AI Design Guidelines; IBM Watson Assistant Dialog Nodes}
\explain{Hướng giải: Xác định bước KHÔNG thuộc luồng hội thoại đặt bàn nhà hàng thông thường.

Các bước phân tích:
1. Luồng hội thoại nhà hàng chuẩn: Chào khách -> Hỏi thông tin (tên, số người, ngày giờ, số điện thoại) -> Kiểm tra bàn trống -> Xác nhận đặt bàn thành công.
2. Phương án D là một hành động phi lý không liên quan đến ngữ cảnh dịch vụ ẩm thực.

Kiểm tra: Trong thiết kế Chatbot, phạm vi nghiệp vụ (domain scope) của bot đặt bàn chỉ giới hạn trong quy trình tiếp nhận và lưu thông tin bàn ăn.

Vì sao các phương án khác là thành phần chuẩn:
- A (Greeting): Chào hỏi mở đầu phiên hội thoại.
- B (Asking for user\'s name): Thu thập thông tin định danh khách hàng.
- C (Confirming reservation): Chốt lại thông tin đặt chỗ.

Lưu ý/Khái niệm cần nhớ: Dialog flow phải bám sát mục tiêu nghiệp vụ của kịch bản ứng dụng.}
\vocab{Dialog Flow | Luồng hội thoại | Cấu trúc phân nhánh các bước giao tiếp giữa người dùng và chatbot}
\end{question}

% ==============================================================================
% Question 166 | Phân loại: 4. Building AI Powered Chatbots Without Programming
% Xuất hiện trong các đề gốc: AIG202c_FA25_FE (Q48)
% ==============================================================================
\begin{question}{166}{mcq}
\source{Q48.webp}
\prompt{What role do 'digressions' play in maintaining user interest during interactions with chatbots?}
\choice{A}{Limiting conversation options}
\choice{B}{Keeping conversations strictly on-topic}
\choice{C}{Providing varied and engaging conversation experiences}
\choice{D}{Reducing the chatbot's vocabulary}
\answer{A}
\confidence{high}
\evidence{Ý định (Intent) trong chatbot là mục đích hoặc hành động mà người dùng muốn thực hiện thông qua câu nói của họ (The goal or purpose expressed by the user).}
\sources{IBM Watson Assistant Concepts: Intents; Rasa NLU Documentation}
\explain{Hướng giải: Định nghĩa khái niệm Ý định (Intent) trong hệ thống hội thoại AI.

Các bước phân tích:
1. Ý định là nhãn phân loại phản ánh mục tiêu giao tiếp của người dùng.
2. Ví dụ: 'Thời tiết hôm nay thế nào?' -> Intent: #CheckWeather; 'Tôi muốn đổi mật khẩu' -> Intent: #ResetPassword.
3. Khối NLU phân loại câu nói của người dùng vào một trong các Intent đã định nghĩa trước.

Kiểm tra: Intent Classification là bài toán phân loại văn bản có giám sát cơ bản trong Chatbot.

Vì sao các phương án khác sai:
- Tham số chi tiết là Entity, không phải Intent.
- Thời gian phản hồi là chỉ số hiệu năng kỹ thuật.

Lưu ý/Khái niệm cần nhớ: Intent = Mục đích/Ý định của người dùng.}
\vocab{Intent | Ý định | Mục tiêu hành động mà người dùng muốn hệ thống thực hiện qua phát ngôn}
\end{question}

% ==============================================================================
% Question 167 | Phân loại: 4. Building AI Powered Chatbots Without Programming
% Xuất hiện trong các đề gốc: AIG202c_FA25_FE (Q49), AIG202c_SU26_FE (Q53), AIG202c_SP26_RE (Q47)
% ==============================================================================
\begin{question}{167}{mcq}
\source{Q49.webp}
\prompt{Why is it important for chatbots to handle digressions effectively?}
\choice{A}{Digressions help improve the user experience by providing additional information}
\choice{B}{Digressions can lead to misunderstandings and dissatisfaction among users}
\choice{C}{Digressions increase the complexity of the conversation and challenge the chatbot's capabilities}
\choice{D}{Digressions are essential for training machine learning models}
\answer{B}
\confidence{high}
\evidence{Xử lý lạc đề hiệu quả (effective digression handling) rất quan trọng vì nếu không được quản lý tốt, các phân nhánh lạc đề sẽ gây ra hiểu lầm, bực bội và trải nghiệm tiêu cực cho người dùng khi bot không thể quay lại chủ đề chính.}
\sources{IBM Watson Assistant: Handling Digressions; Conversational UX Best Practices}
\explain{Hướng giải: Giải thích tầm quan trọng của việc xử lý lạc đề (Digressions) trong thiết kế trải nghiệm Chatbot.

Các bước phân tích:
1. Trong giao tiếp tự nhiên, con người thường xuyên đổi hướng cuộc nói chuyện (ví dụ: đang đặt vé máy bay lại hỏi 'thời tiết ở đó thế nào?').
2. Nếu chatbot cứng nhắc không hiểu hoặc bị mất dấu ngữ cảnh cũ, người dùng sẽ rất thất vọng và bực bội (misunderstandings and dissatisfaction).
3. Một chatbot thông minh cần có khả năng tạm dừng luồng hiện tại, trả lời câu hỏi phụ (digression), sau đó khéo léo quay trở lại hoàn thành tác vụ dang dở.

Kiểm tra: Khả năng quản lý digression là tiêu chí phân biệt giữa chatbot thế hệ cũ cứng nhắc và chatbot đàm thoại linh hoạt.

Vì sao các phương án khác sai:
- A: Digression không hề giúp tập trung vào chủ đề gốc mà là sự phân nhánh tạm thời.
- C: Mục đích không phải thách thức năng lực của bot mà là phục vụ trải nghiệm người dùng.
- D: Digression là bài toán thiết kế hội thoại, không phải điều kiện huấn luyện ML.

Lưu ý/Khái niệm cần nhớ: Xử lý lạc đề tốt giúp duy trì mạch hội thoại tự nhiên và nâng cao mức độ hài lòng của khách hàng.}
\vocab{Digression | Sự lạc đề / Chuyển hướng hội thoại | Việc người dùng chuyển sang chủ đề khác trong khi luồng tác vụ chính chưa hoàn tất}
\end{question}

% ==============================================================================
% Question 168 | Phân loại: 4. Building AI Powered Chatbots Without Programming
% Xuất hiện trong các đề gốc: AIG202c_FA25_FE (Q52), AIG202c_SP26_RE (Q49)
% ==============================================================================
\begin{question}{168}{mcq}
\source{Q52.webp}
\prompt{Which of the following is NOT an example of an intent for a weather chatbot?}
\choice{A}{GetForecast}
\choice{B}{FindLocation}
\choice{C}{GreetUser}
\choice{D}{ProvideTemperature}
\answer{D}
\confidence{high}
\evidence{`ProvideTemperature` là một hành động phản hồi hoặc chức năng của hệ thống (Bot Response / Action), KHÔNG PHẢI là Ý định (Intent) của người dùng; ý định của người dùng phải là yêu cầu lấy thông tin như `GetForecast` hoặc `AskTemperature`.}
\sources{Conversational AI Intent Design Best Practices; IBM Watson Intent Naming Conventions}
\explain{Hướng giải: Phân biệt giữa Ý định của người dùng (User Intent) và Hành động phản hồi của hệ thống (Bot Action).

Các bước phân tích:
1. Intent thể hiện mục đích của NGƯỜI DÙNG khi nói một câu (ví dụ: 'Ngày mai trời có mưa không?' -> `GetForecast`; 'Tôi đang ở Hà Nội' -> `FindLocation`; 'Xin chào' -> `GreetUser`).
2. `ProvideTemperature` (Cung cấp nhiệt độ) là hành động mà BOT thực hiện để trả lời người dùng, không phải là ý định mà người dùng yêu cầu.

Kiểm tra: Quy ước đặt tên Intent luôn phản ánh hành động của người dùng (Request/Ask), không đặt tên theo hành động phản hồi của bot (Provide/Reply).

Vì sao các phương án khác là Intent hợp lệ:
- A (`GetForecast`): Người dùng muốn xem dự báo thời tiết.
- B (`FindLocation`): Người dùng muốn xác định vị trí hiện tại.
- C (`GreetUser` / Chào hỏi): Ý định chào mở đầu phiên trò chuyện.

Lưu ý/Khái niệm cần nhớ: Intent là hành động mong muốn của người dùng, không phải hành vi trả lời của bot.}
\vocab{Intent Naming | Đặt tên ý định | Quy chuẩn đặt tên nhãn ý định dựa trên hành vi yêu cầu của người dùng}
\end{question}

% ==============================================================================
% Question 169 | Phân loại: 4. Building AI Powered Chatbots Without Programming
% Xuất hiện trong các đề gốc: AIG202c_FA25_FE (Q53)
% ==============================================================================
\begin{question}{169}{mcq}
\source{Q53.webp}
\prompt{What is the primary focus of the introduction module in building AI-powered chatbots?}
\choice{A}{Advanced programming languages}
\choice{B}{AI in virtual reality}
\choice{C}{Fundamentals of chatbot development}
\choice{D}{Quantum computing}
\answer{C}
\confidence{high}
\evidence{Mô-đun giới thiệu (Introduction module) trong việc xây dựng chatbot AI tập trung vào các nguyên lý và khái niệm nền tảng trong phát triển chatbot (Fundamentals of chatbot development).}
\sources{IBM Watson Chatbot Course Curriculum; Coursera Building AI-Powered Chatbots}
\explain{Hướng giải: Xác định nội dung trọng tâm của bài mở đầu trong khóa học xây dựng Chatbot.

Các bước phân tích:
1. Bài mở đầu luôn cung cấp bức tranh toàn cảnh: Chatbot là gì, lịch sử phát triển, các thành phần kiến trúc cốt lõi (Intent, Entity, Dialog Flow) và các trường hợp sử dụng cơ bản.
2. Do đó, 'Fundamentals of chatbot development' (Nguyên lý cơ bản phát triển chatbot) là trọng tâm chính xác.

Kiểm tra: Giáo trình IBM AIG202c dành bài đầu tiên để định nghĩa các khái niệm nền tảng trước khi đi vào kỹ thuật nâng cao.

Vì sao các phương án khác sai:
- A (Advanced programming): Bài mở đầu không đi vào ngôn ngữ lập trình nâng cao.
- B (Virtual reality): Không thuộc phạm vi ứng dụng chatbot cơ bản.
- D (Quantum computing): Hoàn toàn không liên quan đến công nghệ chatbot hiện nay.

Lưu ý/Khái niệm cần nhớ: Bài mở đầu luôn tập trung vào các khái niệm nền tảng và kiến trúc cốt lõi.}
\vocab{Chatbot Fundamentals | Nguyên lý cơ bản về chatbot | Các kiến thức cốt lõi về NLU, Intent, Entity và Dialog Manager}
\end{question}

% ==============================================================================
% Question 170 | Phân loại: 4. Building AI Powered Chatbots Without Programming
% Xuất hiện trong các đề gốc: AIG202c_FA25_FE (Q54)
% ==============================================================================
\begin{question}{170}{mcq}
\source{Q54.webp}
\prompt{Which of these statements best reflects how businesses are implementing chatbots in customer service?}
\choice{A}{Businesses are looking to balance customer interactions between human agents and chatbots.}
\choice{B}{Chatbots are being used to take on all customer queries, without exception.}
\choice{C}{With the successful implementation of chatbots, businesses are looking at eliminating all human customer service agents.}
\choice{D}{Chatbots engage in complex customer interactions that require creativity or exceptions to the rules.}
\answer{A}
\confidence{high}
\evidence{Trong dịch vụ khách hàng thực tế, các doanh nghiệp áp dụng mô hình cân bằng và kết hợp hài hòa (balance customer interactions) giữa nhân viên tư vấn con người và chatbot (mô hình Human-in-the-loop / Hybrid support).}
\sources{Gartner Customer Service AI Implementation Report; Harvard Business Review: Collaborative Intelligence}
\explain{Hướng giải: Đánh giá chiến lược triển khai chatbot thực tế trong chăm sóc khách hàng của doanh nghiệp hiện đại.

Các bước phân tích:
1. Chatbot xuất sắc trong việc xử lý các câu hỏi phổ biến, lặp lại 24/7 với tốc độ phản hồi tức thì (Tier 1 Support).
2. Khi gặp các khiếu nại phức tạp, tình huống ngoại lệ hoặc khách hàng giận dữ đòi hỏi sự đồng cảm, chatbot sẽ chuyển tiếp (escalate/handoff) cho nhân viên con người (Tier 2/3 Support).
3. Do đó, doanh nghiệp tìm kiếm sự cân bằng (balance) tối ưu giữa năng suất của bot và sự thấu hiểu của con người.

Kiểm tra: Mô hình lai (Hybrid model) là tiêu chuẩn vận hành thành công nhất của ngành dịch vụ khách hàng toàn cầu.

Vì sao các phương án khác sai:
- B: Không có doanh nghiệp nào để bot xử lý 100% mọi truy vấn mà không có ngoại lệ.
- C: Loại bỏ hoàn toàn con người là ảo tưởng và làm suy giảm nghiêm trọng lòng trung thành của khách hàng.
- D: Chatbot hiện nay chưa đủ khả năng giải quyết các tương tác sáng tạo phức tạp vượt ngoài quy tắc.

Lưu ý/Khái niệm cần nhớ: Chatbot và con người phối hợp theo mô hình cân bằng (Human-AI Collaboration).}
\vocab{Human Handoff | Chuyển giao cho con người | Cơ chế chatbot chuyển cuộc trò chuyện cho nhân viên khi gặp vấn đề phức tạp}
\end{question}

% ==============================================================================
% Question 171 | Phân loại: 4. Building AI Powered Chatbots Without Programming
% Xuất hiện trong các đề gốc: AIG202c_FA25_FE (Q55)
% ==============================================================================
\begin{question}{171}{mcq}
\source{Q55.webp}
\prompt{What is an "intent" in the context of artificial intelligence and chatbots?}
\choice{A}{A measure of how smart a chatbot is}
\choice{B}{The desire to achieve a particular outcome or goal expressed by a user}
\choice{C}{The ability of a chatbot to understand natural language}
\choice{D}{The speed at which a chatbot responds to user queries}
\answer{B}
\confidence{high}
\evidence{Ý định (Intent) trong ngữ cảnh AI và chatbot là mong muốn đạt được một kết quả hoặc mục tiêu cụ thể được người dùng thể hiện qua câu nói (The desire to achieve a particular outcome or goal expressed by a user).}
\sources{IBM Watson Assistant Documentation; Manning, Introduction to Information Retrieval}
\explain{Hướng giải: Định nghĩa bản chất của thuật ngữ 'Intent' trong xử lý ngôn ngữ tự nhiên cho chatbot.

Các bước phân tích:
1. Người dùng tương tác với chatbot luôn có một mục đích cụ thể (ví dụ: tra cứu số dư, đặt hàng, hủy dịch vụ).
2. Mục đích này được mô hình NLU phân loại thành một nhãn ý định (Intent).

Kiểm tra: Trùng khớp với định nghĩa chuẩn của các framework chatbot hàng đầu thế giới (Watson, Dialogflow, Rasa).

Vì sao các phương án khác sai:
- A: Không phải thước đo độ thông minh.
- C: Khả năng hiểu ngôn ngữ là NLU/NLP tổng quát, không phải là 'an intent'.
- D: Tốc độ phản hồi là chỉ số hiệu năng (Latency).

Lưu ý/Khái niệm cần nhớ: Intent = Mục tiêu hoặc mong muốn cụ thể của người dùng.}
\vocab{Goal-Oriented Dialog | Hội thoại định hướng mục tiêu | Hệ thống hội thoại được thiết kế để giúp người dùng đạt được mục đích cụ thể}
\end{question}

% ==============================================================================
% Question 172 | Phân loại: 4. Building AI Powered Chatbots Without Programming
% Xuất hiện trong các đề gốc: AIG202c_FA25_FE (Q56)
% ==============================================================================
\begin{question}{172}{mcq}
\source{Q56.webp}
\prompt{Which of the following is an example of a context variable for a weather chatbot?}
\choice{A}{Time of Day}
\choice{B}{User's Name}
\choice{C}{Favorite Color}
\choice{D}{Number of Pets}
\answer{A}
\confidence{high}
\evidence{Thời điểm trong ngày (Time of Day) là biến ngữ cảnh (context variable) có ý nghĩa thiết thực đối với chatbot thời tiết để cung cấp dự báo chính xác cho buổi sáng, chiều hoặc tối.}
\sources{Contextual Conversational Design; Watson Assistant Context Variables}
\explain{Hướng giải: Xác định biến ngữ cảnh (context variable) liên quan trực tiếp đến nghiệp vụ của bot thời tiết.

Các bước phân tích:
1. Biến ngữ cảnh lưu giữ thông tin trạng thái xuyên suốt phiên hội thoại.
2. Với chatbot thời tiết, các thông tin thời gian như 'Time of Day' (buổi sáng, buổi tối) và 'Location' (địa điểm) quyết định trực tiếp dữ liệu dự báo cần trả về từ API thời tiết.

Kiểm tra: Biến thời gian trong ngày giúp bot trả lời câu hỏi: 'Tối nay trời có lạnh không?' một cách chính xác.

Vì sao các phương án khác sai:
- Favorite Color, Number of Pets không có giá trị nghiệp vụ cho dự báo thời tiết.
- User\'s Name là thông tin định danh người dùng chung, không phải biến ngữ cảnh đặc trưng cho thời tiết.

Lưu ý/Khái niệm cần nhớ: Context variable lưu trữ các tham số trạng thái cần thiết để thực thi tác vụ.}
\vocab{Context Variable | Biến ngữ cảnh | Biến lưu trữ dữ liệu trạng thái được truyền qua lại giữa các nút hội thoại}
\end{question}

% ==============================================================================
% Question 173 | Phân loại: 4. Building AI Powered Chatbots Without Programming
% Xuất hiện trong các đề gốc: AIG202c_FA25_FE (Q60), AIG202c_SP26_RE (Q60)
% ==============================================================================
\begin{question}{173}{mcq}
\source{Q60.webp}
\prompt{How does "contextual understanding" enhance the effectiveness of a chatbot's dialog?}
\choice{A}{Contextual understanding enables chatbots to speak multiple languages fluently.}
\choice{B}{Contextual understanding allows chatbots to remember past interactions and use that information to guide current conversations.}
\choice{C}{Contextual understanding helps chatbots generate random responses to user queries.}
\choice{D}{Contextual understanding enables chatbots to ignore user input that is irrelevant to the current conversation.}
\answer{B}
\confidence{high}
\evidence{Sự hiểu biết ngữ cảnh (Contextual understanding) nâng cao hiệu quả đối thoại bằng cách cho phép chatbot ghi nhớ các tương tác trong quá khứ và sử dụng thông tin đó để dẫn dắt cuộc trò chuyện hiện tại (remember past interactions and guide current conversations).}
\sources{Jurafsky & Martin, Speech and Language Processing (Dialog State Tracking); IBM Watson Assistant Context Management}
\explain{Hướng giải: Xác định vai trò cốt lõi của 'Hiểu biết ngữ cảnh' (Contextual understanding) trong hội thoại AI.

Các bước phân tích:
1. Giao tiếp con người mang tính liên tục: những câu nói sau phụ thuộc vào những gì đã được thống nhất trước đó (ví dụ: 'Cho tôi xem áo sơ mi màu xanh', câu sau: 'Có cỡ L không?').
2. Hiểu biết ngữ cảnh cho phép bot theo dõi trạng thái đối thoại (Dialog State Tracking - DST), lưu giữ các thông tin đã trao đổi để câu trả lời sau luôn khớp với câu chuyện trước đó.

Kiểm tra: Nếu không có ngữ cảnh, mỗi lượt nói (turn) sẽ trở thành một câu hỏi cô lập vô nghĩa.

Vì sao các phương án khác sai:
- A: Đa ngôn ngữ là năng lực dịch thuật/ngôn ngữ học, không phải hiểu ngữ cảnh luồng đàm thoại.
- C: Trả lời ngẫu nhiên là hành vi phá vỡ cuộc trò chuyện.
- D: Bỏ qua thông tin người dùng là biểu hiện của bot kém chất lượng.

Lưu ý/Khái niệm cần nhớ: Contextual understanding = Khả năng ghi nhớ và kế thừa thông tin trong suốt phiên đàm thoại.}
\vocab{Dialog State Tracking (DST) | Theo dõi trạng thái hội thoại | Quá trình theo dõi và cập nhật toàn bộ thông tin ngữ cảnh trong cuộc trò chuyện}
\end{question}

% ==============================================================================
% Question 174 | Phân loại: 4. Building AI Powered Chatbots Without Programming
% Xuất hiện trong các đề gốc: AIG202c_SU26_FE (Q14)
% ==============================================================================
\begin{question}{174}{mcq}
\source{Q14.webp}
\prompt{What is one advantage of using no-code chatbot development platforms like IBM's Watson X Assistant?}
\choice{A}{They require advanced programming skills, such as implementing complex algorithms}
\choice{B}{They allow anyone to build custom chatbots with drag-and-drop interfaces}
\choice{C}{They only support chatbot creation for large enterprises}
\choice{D}{They limit chatbot functionality to basic FAQ responses}
\answer{B}
\confidence{high}
\evidence{No-code chatbot development platforms like IBM watsonx Assistant provide intuitive drag-and-drop visual builders that allow anyone, including non-programmers, to design, test, and deploy customized conversational bots.}
\sources{IBM Documentation: Getting started with watsonx Assistant; Cognitive Class Chatbot Builder}
\explain{Hướng giải: Hiểu rõ khái niệm và lợi ích của nền tảng phát triển không cần viết mã (no-code platform) như IBM Watson Assistant.

Các bước phân tích:
1. 'No-code' nghĩa là người dùng không cần phải có kiến thức lập trình hay viết các dòng mã nguồn phức tạp.
2. Nền tảng như IBM Watson Assistant (watsonx Assistant) cung cấp giao diện trực quan kéo-thả (drag-and-drop interface), cho phép bất kỳ ai (chuyên viên kinh doanh, quản trị viên, người dùng phổ thông) đều có thể tự tay xây dựng cây hội thoại (dialog flow) và huấn luyện chatbot tùy chỉnh.
3. Do đó, phương án B là ưu điểm chính xác của nền tảng no-code.

Kiểm tra:
- A sai vì no-code loại bỏ yêu cầu kỹ năng lập trình nâng cao, chứ không phải yêu cầu kỹ năng này.
- C sai vì nền tảng hỗ trợ từ cá nhân, doanh nghiệp vừa và nhỏ đến tập đoàn lớn.
- D sai vì Watson Assistant hỗ trợ các luồng hội thoại phức tạp, tích hợp API, slot-filling và digression, không hề giới hạn ở mức câu hỏi thường gặp (FAQ) cơ bản.

Vì sao các phương án khác sai:
- A trái ngược với bản chất 'no-code'.
- C, D đánh giá thấp và sai lệch về tính năng của Watson Assistant.

Lưu ý/Khái niệm cần nhớ:
No-code chatbot platform cho phép thiết kế kịch bản chatbot nhanh chóng qua giao diện trực quan kéo-thả mà không cần viết mã.}
\vocab{No-Code Platform | Nền tảng không cần mã nguồn | Công cụ phát triển phần mềm cho phép người dùng tạo ứng dụng thông qua giao diện đồ họa kéo thả thay vì lập trình truyền thống}
\end{question}

% ==============================================================================
% Question 175 | Phân loại: 4. Building AI Powered Chatbots Without Programming
% Xuất hiện trong các đề gốc: AIG202c_SU26_FE (Q31)
% ==============================================================================
\begin{question}{175}{mcq}
\source{Q31.webp}
\prompt{What type of input can digressions help handle?}
\choice{A}{UNEXPECTED}
\choice{B}{SCRIPTED}
\choice{C}{FIXED}
\choice{D}{ERROR-FREE}
\answer{A}
\confidence{high}
\evidence{In conversational AI design (e.g., IBM watsonx Assistant), digressions allow the chatbot to smoothly handle unexpected user inputs that stray from the main dialog flow and then return back.}
\sources{IBM Cloud Documentation: Configuring digressions in watsonx Assistant; Cognitive Class Chatbot Design}
\explain{Hướng giải: Hiểu khái niệm 'digression' (sự chuyển hướng/rẽ nhánh) trong thiết kế luồng hội thoại của chatbot.

Các bước phân tích:
1. Trong quá trình giao tiếp thực tế, người dùng hiếm khi đi thẳng một đường thẳng theo kịch bản. Ví dụ, khi đang điền thông tin đặt phòng khách sạn, người dùng bất ngờ hỏi: 'Khách sạn có bể bơi không?'.
2. Đây là một câu hỏi ngoài lề, bất ngờ (UNEXPECTED input) đối với luồng hiện tại.
3. Tính năng 'Digression' (rẽ nhánh) cho phép chatbot tạm dừng luồng chính, nhảy sang trả lời câu hỏi bất ngờ của người dùng, sau đó quay trở lại tiếp tục luồng đặt phòng dang dở mà không làm mất dữ liệu đã nhập.
4. Vì vậy, digression giúp xử lý các đầu vào bất ngờ (UNEXPECTED).

Kiểm tra:
- B (SCRIPTED - theo kịch bản), C (FIXED - cố định) đã nằm trong luồng thông thường của chatbot, không cần đến digression.
- D (ERROR-FREE - không có lỗi) không phải lý do kích hoạt digression.

Vì sao các phương án khác sai:
- B, C, D đều mô tả các tình huống thông thường, tuyến tính, không thể hiện được mục đích chính của việc xử lý rẽ nhánh linh hoạt.

Lưu ý/Khái niệm cần nhớ:
Digression trong chatbot cho phép xử lý các câu hỏi ngoài lề hoặc bất ngờ (unexpected) của người dùng rồi quay lại luồng chính.}
\vocab{Digression | Sự rẽ nhánh hội thoại | Khả năng của chatbot rời khỏi luồng hội thoại hiện tại để trả lời một câu hỏi khác rồi quay lại tiếp tục nhiệm vụ trước đó}
\end{question}

% ==============================================================================
% Question 176 | Phân loại: 4. Building AI Powered Chatbots Without Programming
% Xuất hiện trong các đề gốc: AIG202c_SU26_FE (Q33)
% ==============================================================================
\begin{question}{176}{mcq}
\source{Q33.webp}
\prompt{Fill in the blank: Entities work closely with to complete tasks.}
\choice{A}{CHATBOTS}
\choice{B}{DIALOGS}
\choice{C}{INTENTS}
\choice{D}{SPEAKERS}
\answer{C}
\confidence{high}
\evidence{In Natural Language Understanding (NLU) and chatbot architectures, Entities work closely with Intents to identify user purpose and extract relevant parameters to complete tasks.}
\sources{IBM Watson Assistant Core Concepts: Intents and Entities; Jurafsky & Martin, Speech and Language Processing}
\explain{Hướng giải: Điền từ thích hợp vào chỗ trống thể hiện mối quan hệ nền tảng trong kiến trúc xử lý ngôn ngữ của chatbot: 'Entities work closely with [blank] to complete tasks.'

Các bước phân tích:
1. Trong kiến trúc hiểu ngôn ngữ tự nhiên (NLU) của các nền tảng chatbot (IBM Watson, Google Dialogflow, Rasa):
   - Ý định (Intent): Đại diện cho mục đích tổng thể của người dùng (Ví dụ: #dat_ve_may_bay).
   - Thực thể (Entity): Đại diện cho các tham số cụ thể bổ nghĩa cho ý định (Ví dụ: @diem_den: Hà Nội, @ngay: ngày mai).
2. Để hoàn thành một tác vụ (task), chatbot phải kết hợp cả Intent (biết người dùng muốn làm gì) và Entity (biết các dữ liệu chi tiết kèm theo).
3. Cụm từ chuẩn trong tài liệu IBM Watson: 'Entities work closely with intents to complete tasks'.

Kiểm tra:
- Ghép từ: 'Entities work closely with INTENTS to complete tasks' hoàn toàn chuẩn xác về mặt ngữ nghĩa và lý thuyết kỹ thuật.

Vì sao các phương án khác sai:
- A (Chatbots), B (Dialogs), D (Speakers) không tạo nên cặp khái niệm cốt lõi tương hỗ trực tiếp với Entities trong việc nhận diện yêu cầu đầu vào.

Lưu ý/Khái niệm cần nhớ:
Cặp bài trùng trong NLU: Intents (Ý định) xác định hành động, Entities (Thực thể) cung cấp thông tin chi tiết.}
\vocab{Intent vs. Entity | Ý định và Thực thể | Cặp thành phần cơ bản trong NLU: Intent xác định mục đích người dùng, Entity trích xuất dữ liệu cụ thể}
\end{question}

% ==============================================================================
% Question 177 | Phân loại: 4. Building AI Powered Chatbots Without Programming
% Xuất hiện trong các đề gốc: AIG202c_SU26_FE (Q37)
% ==============================================================================
\begin{question}{177}{mcq}
\source{Q37.webp}
\prompt{What does context help maintain across messages?}
\choice{A}{COLOR}
\choice{B}{CONSISTENCY}
\choice{C}{DELAY}
\choice{D}{RANDOMNESS}
\answer{B}
\confidence{high}
\evidence{Context variables maintain state and consistency across multiple messages and turns throughout a chatbot conversation.}
\sources{IBM watsonx Assistant: Context variables; Cognitive Class 'Build Your Own Chatbot'}
\explain{Hướng giải: Xác định vai trò của ngữ cảnh (context) trong việc duy trì cuộc trò chuyện giữa chatbot và người dùng.

Các bước phân tích:
1. Trong một cuộc hội thoại nhiều lượt (multi-turn conversation), người dùng cung cấp thông tin qua nhiều câu khác nhau.
2. Ngữ cảnh (Context), thường được lưu trữ dưới dạng các biến ngữ cảnh (context variables), giúp chatbot ghi nhớ thông tin đã thu thập từ các tin nhắn trước (như tên người dùng, sở thích, mã đơn hàng).
3. Nhờ có ngữ cảnh, chatbot duy trì được tính nhất quán và mạch lạc (CONSISTENCY) xuyên suốt toàn bộ phiên làm việc, tránh việc hỏi lại thông tin đã có hoặc trả lời rời rạc.

Kiểm tra:
- A (COLOR - màu sắc), C (DELAY - độ trễ), D (RANDOMNESS - tính ngẫu nhiên) hoàn toàn không phải là mục tiêu mà context mang lại.

Vì sao các phương án khác sai:
- A, C, D là các từ không liên quan hoặc trái ngược với mục đích duy trì trạng thái của hệ thống.

Lưu ý/Khái niệm cần nhớ:
Biến ngữ cảnh (Context variables) giúp chatbot ghi nhớ trạng thái và đảm bảo tính nhất quán (consistency) qua các lượt hội thoại.}
\vocab{Context Variables | Biến ngữ cảnh | Các biến lưu trữ trạng thái và dữ liệu người dùng qua các lượt trao đổi trong một phiên trò chuyện}
\end{question}

% ==============================================================================
% Question 178 | Phân loại: 4. Building AI Powered Chatbots Without Programming
% Xuất hiện trong các đề gốc: AIG202c_SU26_FE (Q41)
% ==============================================================================
\begin{question}{178}{mcq}
\source{Q41.webp}
\prompt{In chatbot design, when should you use a slot-filling approach with entities?}
\choice{A}{When intents are unknown}
\choice{B}{When dialog is linear}
\choice{C}{When you need multiple values from the user before moving forward}
\choice{D}{When digressions are disabled}
\answer{C}
\confidence{high}
\evidence{In chatbot design, a slot-filling approach is used when a task requires collecting multiple values (slots/entities) from the user before executing the action.}
\sources{IBM Cloud Documentation: Gathering information with slots; Rasa Forms & Slot Filling}
\explain{Hướng giải: Xác định thời điểm và lý do cần áp dụng kỹ thuật điền khe (slot-filling) với các thực thể trong thiết kế chatbot.

Các bước phân tích:
1. Slot-filling (Điền khe thông tin) là một cơ chế thiết kế hội thoại mạnh mẽ trong chatbot.
2. Kỹ thuật này được sử dụng khi một tác vụ cụ thể đòi hỏi phải thu thập đầy đủ nhiều giá trị/thông tin từ người dùng trước khi có thể thực hiện hành động tiếp theo (When you need multiple values from the user before moving forward).
3. Ví dụ: Để hoàn tất tác vụ đặt vé xe, chatbot cần 4 thông tin: Điểm đi, Điểm đến, Ngày đi và Số lượng vé. Cơ chế slot-filling sẽ tự động kiểm tra xem thông tin nào còn thiếu để hỏi bổ sung mà không cần tạo nhiều nút hội thoại lặp đi lặp lại.

Kiểm tra:
- Phương án C mô tả chính xác định nghĩa và mục đích ra đời của slot-filling.

Vì sao các phương án khác sai:
- A: Slot-filling được kích hoạt sau khi đã nhận diện rõ Intent (ví dụ #dat_ve), không phải khi chưa biết Intent.
- B: Dialog có thể phân nhánh hoặc phi tuyến, slot-filling giúp gom dữ liệu linh hoạt.
- D: Slot-filling hoàn toàn có thể kết hợp với digression.

Lưu ý/Khái niệm cần nhớ:
Slot-filling được áp dụng khi cần thu thập nhiều tham số (thực thể) trước khi tiến hành xử lý tác vụ.}
\vocab{Slot Filling | Điền khe thông tin | Kỹ thuật thu thập tất cả các tham số bắt buộc từ người dùng trong một node duy nhất trước khi thực hiện hành động}
\end{question}

% ==============================================================================
% Question 179 | Phân loại: 4. Building AI Powered Chatbots Without Programming
% Xuất hiện trong các đề gốc: AIG202c_SU26_FE (Q43)
% ==============================================================================
\begin{question}{179}{mcq}
\source{Q43.webp}
\prompt{Which scenario best requires context chaining across multiple dialog nodes?}
\choice{A}{Greeting the user}
\choice{B}{Collecting multi-step input like booking date, location, and time}
\choice{C}{Answering yes/no questions}
\choice{D}{Showing chatbot capabilities}
\answer{B}
\confidence{high}
\evidence{Context chaining across multiple dialog nodes is essential when collecting multi-step structured inputs (such as booking date, location, and time) where each node builds upon previously stored parameters.}
\sources{IBM watsonx Assistant: Context chaining across dialog nodes; Conversational Design Guide}
\explain{Hướng giải: Xác định kịch bản hội thoại nào đòi hỏi việc liên kết ngữ cảnh (context chaining) qua nhiều nút hội thoại (dialog nodes).

Các bước phân tích:
1. Chuỗi ngữ cảnh (Context chaining) là việc truyền và kế thừa các biến ngữ cảnh (context variables) từ nút này sang nút khác trong cây hội thoại.
2. Khi thu thập thông tin nhiều bước phức tạp như quy trình đặt lịch (ngày đặt, địa điểm và thời gian), mỗi bước người dùng cung cấp một mẩu thông tin. Hệ thống phải lưu trữ các thông tin này vào ngữ cảnh và chuyển tiếp qua các nút tiếp theo để xác nhận và thực hiện đặt chỗ.
3. Do đó, kịch bản 'Collecting multi-step input like booking date, location, and time' (Phương án B) là kịch bản điển hình nhất yêu cầu context chaining.

Kiểm tra:
- Chào người dùng (A), trả lời câu hỏi Yes/No đơn giản (C), giới thiệu tính năng bot (D) là các tương tác đơn lẻ, một lượt (single-turn) không đòi hỏi lưu trữ chuỗi ngữ cảnh phức tạp.

Vì sao các phương án khác sai:
- A, C, D đều là các tương tác ngắn, độc lập, không cần duy trì chuỗi dữ liệu qua nhiều nút.

Lưu ý/Khái niệm cần nhớ:
Context chaining cần thiết nhất trong các tác vụ thu thập dữ liệu nhiều bước (multi-step input) để xâu chuỗi thông tin.}
\vocab{Context Chaining | Chuỗi liên kết ngữ cảnh | Kỹ thuật chuyển tiếp và duy trì dữ liệu ngữ cảnh giữa các nút liên tiếp trong cây hội thoại}
\end{question}

% ==============================================================================
% Question 180 | Phân loại: 4. Building AI Powered Chatbots Without Programming
% Xuất hiện trong các đề gốc: AIG202c_SU26_FE (Q45)
% ==============================================================================
\begin{question}{180}{mcq}
\source{Q45.webp}
\prompt{How can you ensure smooth digressions without breaking dialog context?}
\choice{A}{Disable all digressions}
\choice{B}{Use global context variables and context handlers}
\choice{C}{Only allow one intent at a time}
\choice{D}{End the dialog before digression}
\answer{B}
\confidence{high}
\evidence{Smooth digressions without losing dialog context are achieved by utilizing global context variables and dedicated context handlers to store conversation state and resume seamlessly.}
\sources{IBM Watson Assistant: 'Allowing digressions away and returning to original node'; Chatbot State Architecture}
\explain{Hướng giải: Tìm giải pháp kỹ thuật để đảm bảo việc rẽ nhánh (digression) diễn ra mượt mà mà không làm mất hoặc phá vỡ ngữ cảnh hội thoại.

Các bước phân tích:
1. Khi người dùng rẽ nhánh sang chủ đề khác, nguy cơ lớn nhất là hệ thống bị mất các dữ liệu đã nhập ở nhánh chính trước đó.
2. Để giải quyết triệt để vấn đề này, các nền tảng thiết kế chatbot sử dụng các biến ngữ cảnh toàn cục (global context variables) và các bộ xử lý ngữ cảnh (context handlers / return handlers).
3. Cơ chế này lưu lại toàn bộ trạng thái hiện tại cùng con trỏ quay về (return pointer). Khi nhánh phụ kết thúc, chatbot dựa vào biến ngữ cảnh để khôi phục chính xác vị trí và dữ liệu ban đầu.

Kiểm tra:
- A (tắt hết digression) triệt tiêu tính năng thay vì đảm bảo nó mượt mà.
- C (chỉ cho phép một intent tại một thời điểm) không giải quyết bài toán ngữ cảnh rẽ nhánh.
- D (kết thúc hội thoại trước khi rẽ nhánh) sẽ xóa bỏ hoàn toàn trải nghiệm người dùng.

Vì sao các phương án khác sai:
- A, C, D đều là các giải pháp tiêu cực làm hỏng tính linh hoạt của chatbot.

Lưu ý/Khái niệm cần nhớ:
Sử dụng biến ngữ cảnh toàn cục và bộ xử lý ngữ cảnh (global context variables and handlers) giúp lưu trữ và phục hồi trạng thái hội thoại khi rẽ nhánh.}
\vocab{Context Handler | Bộ xử lý ngữ cảnh | Thành phần quản lý việc lưu trữ, cập nhật và phục hồi trạng thái phiên làm việc khi xảy ra chuyển hướng hội thoại}
\end{question}

% ==============================================================================
% Question 181 | Phân loại: 4. Building AI Powered Chatbots Without Programming
% Xuất hiện trong các đề gốc: AIG202c_SU26_FE (Q47)
% ==============================================================================
\begin{question}{181}{mcq}
\source{Q47.webp}
\prompt{How do 'digressions' contribute to the overall effectiveness of chatbots?}
\choice{A}{By eliminating conversational diversions}
\choice{B}{By creating rigid conversation structures}
\choice{C}{By enhancing conversational variety and engagement}
\choice{D}{By limiting conversational flexibility}
\answer{C}
\confidence{high}
\evidence{Digressions improve the effectiveness and naturalness of conversational agents by increasing variety, flexibility, and user engagement rather than locking users into rigid forms.}
\sources{IBM watsonx Assistant Best Practices: Human-Centric Conversation Design}
\explain{Hướng giải: Xác định đóng góp của tính năng 'digressions' đối với hiệu quả tổng thể của chatbot.

Các bước phân tích:
1. Nếu một chatbot chỉ tuân theo quy trình cứng nhắc và từ chối trả lời mọi câu hỏi ngoài luồng, người dùng sẽ cảm thấy ức chế và bỏ dở cuộc trò chuyện.
2. Khả năng rẽ nhánh (digression) mô phỏng chính xác cách giao tiếp tự nhiên của con người, mang lại sự linh hoạt, tăng cường tính đa dạng và sự tương tác hứng thú cho người dùng (enhancing conversational variety and engagement).

Kiểm tra:
- A sai vì digression không loại bỏ sự chuyển hướng mà chính nó cho phép xử lý chuyển hướng.
- B sai vì nó phá vỡ cấu trúc cứng nhắc, chứ không tạo ra cấu trúc cứng nhắc.
- D sai vì nó mở rộng tính linh hoạt, chứ không giới hạn tính linh hoạt.

Vì sao các phương án khác sai:
- A, B, D đều đối lập hoàn toàn với bản chất và lợi ích của digression.

Lưu ý/Khái niệm cần nhớ:
Digression giúp tăng cường tính linh hoạt và mức độ gắn kết (engagement) trong giao tiếp của chatbot.}
\vocab{Conversational Engagement | Mức độ gắn kết hội thoại | Thước đo mức độ hứng thú, tự nhiên và hài lòng của người dùng khi tương tác với chatbot}
\end{question}

% ==============================================================================
% Question 182 | Phân loại: 4. Building AI Powered Chatbots Without Programming
% Xuất hiện trong các đề gốc: AIG202c_SU26_FE (Q48)
% ==============================================================================
\begin{question}{182}{mcq}
\source{Q48.webp}
\prompt{How are entities used in natural language understanding for chatbots?}
\choice{A}{To determine the overall purpose of the user's input}
\choice{B}{To identify the emotions expressed by the user}
\choice{C}{To extract specific details or information from the user's input}
\choice{D}{To generate random responses}
\answer{C}
\confidence{high}
\evidence{In Natural Language Understanding (NLU), entities are specifically utilized to identify and extract specific details, parameters, and key data points from the user's input.}
\sources{IBM Watson Assistant: Defining Entities; Jurafsky & Martin: Information Extraction}
\explain{Hướng giải: Nắm vững vai trò và mục đích sử dụng của các thực thể (entities) trong việc hiểu ngôn ngữ tự nhiên (NLU) cho chatbot.

Các bước phân tích:
1. Trong NLU, đầu vào của người dùng được phân tích thành hai thành phần chính:
   - Ý định (Intent): Xác định mục đích tổng quát của câu (Ví dụ: 'Tôi muốn đặt vé').
   - Thực thể (Entity): Trích xuất các chi tiết hoặc thông tin cụ thể (specific details or information) từ câu (Ví dụ: ngày tháng, địa điểm, số lượng, tên riêng, loại dịch vụ).
2. Do đó, phương án C ('To extract specific details or information from the user's input') là định nghĩa chuẩn mực của Entity.

Kiểm tra:
- Phương án A ('To determine the overall purpose') là chức năng của Intent, không phải Entity.
- Phương án B ('To identify emotions') là chức năng của Sentiment Analysis (Phân tích cảm xúc).
- Phương án D ('To generate random responses') không phải là mục đích của NLU.

Vì sao các phương án khác sai:
- A mô tả Intent.
- B mô tả Phân tích cảm xúc/Tone Analyzer.
- D là hành vi sai lệch.

Lưu ý/Khái niệm cần nhớ:
Entities = Trích xuất chi tiết/dữ liệu cụ thể từ câu nói của người dùng.}
\vocab{Entity Extraction | Trích xuất thực thể | Quá trình bóc tách các mẩu thông tin cụ thể (ngày giờ, địa điểm, tên gọi) từ văn bản tự nhiên}
\end{question}

% ==============================================================================
% Question 183 | Phân loại: 4. Building AI Powered Chatbots Without Programming
% Xuất hiện trong các đề gốc: AIG202c_SU26_FE (Q56), AIG202c_SP26_RE (Q53)
% ==============================================================================
\begin{question}{183}{mcq}
\source{Q56.webp}
\prompt{Which of the following is NOT a common strategy for managing digressions in chatbots?}
\choice{A}{Redirecting the conversation back to the original topic}
\choice{B}{Ignoring the user's input and continuing with a predefined script}
\choice{C}{Acknowledging the digression and exploring the new topic briefly}
\choice{D}{Terminating the conversation abruptly}
\answer{D}
\confidence{high}
\evidence{Ngắt cuộc trò chuyện đột ngột (Terminating the conversation abruptly) là một lỗi thiết kế nghiêm trọng, KHÔNG PHẢI là một chiến lược hợp lệ để quản lý sự lạc đề trong chatbot.}
\sources{IBM Watson Assistant Digression Management; Conversational Error Recovery Strategies}
\explain{Hướng giải: Xác định hành vi KHÔNG PHẢI là chiến lược quản lý lạc đề chuẩn mực của chatbot.

Các bước phân tích:
1. Các chiến lược quản lý lạc đề thông minh bao gồm:
   - Tiếp nhận câu hỏi phụ và trả lời ngắn gọn rồi quay lại chủ đề chính (C).
   - Khéo léo nhắc nhở người dùng quay lại chủ đề chính (A).
2. Hành vi 'Terminating the conversation abruptly' (Đột ngột ngắt kết nối cuộc trò chuyện) là hành vi tồi tệ nhất, gây ức chế cho người dùng và hoàn toàn không thể coi là một 'chiến lược' quản lý.

Kiểm tra: Tiêu chuẩn thiết kế hội thoại nghiêm cấm bot đơn phương ngắt kết nối khi chưa có sự đồng thuận của người dùng.

Vì sao các phương án khác là chiến lược xử lý:
- A (Redirecting back): Chiến lược điều hướng về chủ đề gốc.
- C (Acknowledging and exploring briefly): Chiến lược ghi nhận và trả lời nhanh.

Lưu ý/Khái niệm cần nhớ: Chatbot không bao giờ được đột ngột ngắt cuộc trò chuyện khi gặp tình huống lạc đề.}
\vocab{Conversational Graceful Degradation | Xử lý suy thoái duyên dáng | Khả năng bot xử lý tình huống khó mà không làm đứt gãy trải nghiệm người dùng}
\end{question}

% ==============================================================================
% Question 184 | Phân loại: 4. Building AI Powered Chatbots Without Programming
% Xuất hiện trong các đề gốc: AIG202c_SU26_FE (Q59)
% ==============================================================================
\begin{question}{184}{mcq}
\source{Q59.webp}
\prompt{A dialog has the following nodes- Welcome (welcome condition) Greetings (#greetings intent condition) Thank you (#thank_you intent condition) Goodbyes (#goodbyes intent condition) Anything else (anything_else condition). Which of the following statements apply?}
\choice{A}{Welcome will contain the chatbot prompt and be executed at the beginning of the conversation.}
\choice{B}{The order of Welcome and Anything else in the dialog will generally not matter.}
\choice{C}{Greetings overshadows Goodbyes, and therefore Goodbyes is never executed.}
\choice{D}{In this case, rearranging the order to be Welcome, Thank You, Greetings, Goodbyes and Anything else would not cause problems.}
\choice{E}{If no known intent is detected in the user input, the Anything else node will be executed.}
\answer{E}
\confidence{medium}
\evidence{In IBM Watson Assistant dialog tree architecture, node conditions are evaluated from top to bottom. If no known intent or condition matches the user input, the catch-all 'Anything else' (anything_else condition) node placed at the bottom is executed. (In multi-select versions of this Cognitive Class question, A, D, and E are correct statements; for single-choice format, E describes the primary operational rule of the fallback node).}
\sources{Cognitive Class 'Build Your Own Chatbot' Final Exam; IBM watsonx Assistant Dialog Node Evaluation Docs}
\explain{Hướng giải: Phân tích nguyên lý đánh giá thứ tự các nút (dialog nodes) và điều kiện kích hoạt trong cây hội thoại IBM Watson Assistant.

Các bước phân tích:
1. Cấu trúc cây hội thoại đã cho gồm:
   - Welcome (điều kiện `welcome`)
   - Greetings (điều kiện `#greetings`)
   - Thank you (điều kiện `#thank_you`)
   - Goodbyes (điều kiện `#goodbyes`)
   - Anything else (điều kiện `anything_else`)
2. Nguyên lý hoạt động của IBM Watson Assistant:
   - Các nút ngang hàng được duyệt tuần tự từ trên xuống dưới.
   - Nút `Welcome` kích hoạt tự động ở đầu phiên khi vừa mở chat.
   - Nút `Anything else` luôn đặt ở cuối cùng đóng vai trò là lưới đỡ (fallback / catch-all). Nếu không có intent hoặc điều kiện nào ở trên thỏa mãn, nút `Anything else` sẽ được kích hoạt để thông báo bot chưa hiểu.
3. Đánh giá các phương án:
   - B sai vì vị trí của Welcome và Anything else cực kỳ quan trọng (Anything else phải ở cuối; nếu đặt ở trên sẽ chặn mọi node bên dưới).
   - C sai vì Greetings và Goodbyes có intent riêng biệt, không đè lên nhau.
   - E là phát biểu kinh điển mô tả chính xác cơ chế của nút Anything else: 'Nếu không phát hiện được intent đã biết nào trong câu nhập, nút Anything else sẽ được thực thi'.

Kiểm tra:
Lưu ý đây là câu hỏi từ khóa học Cognitive Class 'Build Your Own Chatbot'. Trong đề thi gốc dạng chọn nhiều phương án, A, D và E đều đúng. Trong bài thi trắc nghiệm một đáp án (Single choice), phương án E thể hiện quy tắc hoạt động cốt lõi nhất của nút fallback `anything_else`.

Vì sao các phương án khác sai:
- B sai nghiêm trọng về mặt kiến trúc cây hội thoại (thứ tự của Anything else mang tính quyết định).
- C sai vì các intent khác nhau không thể triệt tiêu nhau.

Lưu ý/Khái niệm cần nhớ:
Nút Anything else (`anything_else`) phải nằm ở cuối cùng và được kích hoạt khi không có intent nào khác khớp.}
\vocab{Anything Else Node | Nút dự phòng mặc định | Nút cuối cùng trong cây hội thoại Watson Assistant được kích hoạt khi không có điều kiện nào khác thỏa mãn}
\end{question}

% ==============================================================================
% Question 185 | Phân loại: 4. Building AI Powered Chatbots Without Programming
% Xuất hiện trong các đề gốc: AIG202c_SP26_RE (Q7)
% ==============================================================================
\begin{question}{185}{mcq}
\source{Q7.jpg}
\prompt{What is the main function of chatbots in customer service?}
\choice{A}{To replace human employees entirely}
\choice{B}{To simulate human conversation and interact with users}
\choice{C}{To perform financial transactions securely}
\choice{D}{To design websites automatically without user input}
\answer{B}
\confidence{high}
\evidence{Chức năng chính của chatbot trong dịch vụ khách hàng là mô phỏng cuộc trò chuyện của con người và tương tác tự động với người dùng (To simulate human conversation and interact with users) để giải đáp thắc mắc và hỗ trợ xử lý tác vụ.}
\sources{IBM Watson Assistant Overview; Gartner Customer Service Chatbot Architecture}
\explain{Hướng giải: Xác định chức năng cơ bản và mục đích tồn tại của chatbot trong dịch vụ khách hàng.

Các bước phân tích:
1. Chatbot là chương trình máy tính được thiết kế để giao tiếp tự nhiên bằng văn bản hoặc giọng nói.
2. Trong bộ phận chăm sóc khách hàng, chatbot đóng vai trò tiếp nhận, hiểu ý định của người dùng và phản hồi tương tự như một nhân viên hỗ trợ trực tuyến.

Kiểm tra: Định nghĩa chuẩn: 'A chatbot is an AI application designed to simulate human conversation through text or voice interactions'.

Vì sao các phương án khác sai:
- A: Không nhằm thay thế toàn bộ con người mà là hỗ trợ giải phóng con người khỏi các tác vụ lặp lại.
- C: Thực hiện giao dịch tài chính là tác vụ phụ trợ (Fintech backend), không phải chức năng định nghĩa chung của chatbot.
- D: Thiết kế website tự động là chức năng của công cụ Web Builder AI, không phải chatbot.

Lưu ý/Khái niệm cần nhớ: Chatbot = Hệ thống mô phỏng hội thoại tự nhiên tương tác với người dùng.}
\vocab{Chatbot | Phần mềm đàm thoại tự động | Ứng dụng AI mô phỏng hội thoại con người để tương tác và hỗ trợ người dùng}
\end{question}

% ==============================================================================
% Question 186 | Phân loại: 4. Building AI Powered Chatbots Without Programming
% Xuất hiện trong các đề gốc: AIG202c_SP26_RE (Q37)
% ==============================================================================
\begin{question}{186}{mcq}
\source{Q37.jpg}
\prompt{What happens in a fallback dialog?}
\choice{A}{Bot says goodbye}
\choice{B}{Bot repeats the greeting}
\choice{C}{Bot says it didn't understand}
\choice{D}{Bot opens browser}
\answer{C}
\confidence{high}
\evidence{Trong thiết kế hội thoại chatbot, một nút hội thoại dự phòng (fallback dialog) được kích hoạt khi hệ thống không thể nhận diện được ý định (Intent) của người dùng hoặc điểm tin cậy quá thấp; hành vi chuẩn mực của bot là lịch sự thông báo rằng nó chưa hiểu câu hỏi và đề nghị người dùng diễn đạt lại (Bot says it didn't understand).}
\sources{IBM Watson Assistant: The Anything Else / Fallback Node; Dialogflow Default Fallback Intent}
\explain{Hướng giải: Xác định hành vi diễn ra khi chatbot rơi vào luồng hội thoại dự phòng (Fallback dialog).

Các bước phân tích:
1. Nút Fallback (thường gọi là 'Anything else' node trong Watson Assistant hoặc 'Default Fallback' trong Dialogflow) là chốt chặn cuối cùng của luồng đối thoại.
2. Khi câu nói của người dùng không khớp với bất kỳ Intent đã được huấn luyện nào (Confidence < Threshold), hệ thống sẽ chuyển sang Fallback.
3. Hành vi chuẩn của bot là thông báo: 'Xin lỗi, tôi chưa hiểu ý của bạn. Bạn có thể diễn đạt lại câu hỏi được không?' (Bot says it didn't understand).

Kiểm tra: Đây là nguyên lý thiết kế trải nghiệm người dùng cơ bản nhằm tránh việc bot trả lời sai lệch khi không hiểu.

Vì sao các phương án khác sai:
- A: Bot không tự tiện chào tạm biệt khi người dùng đang hỏi.
- B: Bot không tự lặp lại lời chào mở đầu.
- D: Bot không tự ý mở trình duyệt web ngoài kịch bản.

Lưu ý/Khái niệm cần nhớ: Fallback node = Nút xử lý tình huống bot không hiểu ý người dùng.}
\vocab{Fallback Dialog | Hội thoại dự phòng | Luồng phản hồi mặc định khi chatbot không nhận diện được ý định của người dùng}
\end{question}

% ==============================================================================
% Question 187 | Phân loại: 4. Building AI Powered Chatbots Without Programming
% Xuất hiện trong các đề gốc: AIG202c_SP26_RE (Q48)
% ==============================================================================
\begin{question}{187}{mcq}
\source{Q48.jpg}
\prompt{In chatbot development, what do 'intents' typically represent?}
\choice{A}{User interface design principles}
\choice{B}{The purpose or goal behind user inputs}
\choice{C}{Database management strategies}
\choice{D}{Networking protocols}
\answer{B}
\confidence{high}
\evidence{Trong phát triển chatbot, 'Ý định' (Intents) đại diện cho mục đích hoặc mục tiêu đằng sau câu nói của người dùng (The purpose or goal behind user inputs).}
\sources{IBM Watson Assistant Dialog Design; Rasa Open Source Concepts}
\explain{Hướng giải: Xác định ý nghĩa của thành phần Intent trong kiến trúc Chatbot (câu hỏi kiểm tra tính nhất quán với UQ047 và UQ054).

Các bước phân tích:
1. Intent trả lời cho câu hỏi: 'Người dùng nói câu này nhằm mục đích gì?'.
2. Phân loại câu nói vào Intent giúp bộ máy quản trị hội thoại kích hoạt đúng nhánh xử lý tương ứng.

Kiểm tra: Đây là định nghĩa cốt lõi của công nghệ Natural Language Understanding (NLU).

Vì sao các phương án khác sai:
- A: Thuộc về thiết kế giao diện UI/UX.
- C: Thuộc về quản trị cơ sở dữ liệu (DBMS).
- D: Thuộc về giao thức mạng truyền thông (TCP/IP, HTTP).

Lưu ý/Khái niệm cần nhớ: Intent = Mục đích hoặc mục tiêu đằng sau câu nói của người dùng.}
\vocab{Intent | Ý định | Mục tiêu hành động được trích xuất từ câu nói của người dùng trong chatbot}
\end{question}

% ==============================================================================
% Question 188 | Phân loại: 4. Building AI Powered Chatbots Without Programming
% Xuất hiện trong các đề gốc: AIG202c_SP26_RE (Q52)
% ==============================================================================
\begin{question}{188}{mcq}
\source{Q52.jpg}
\prompt{What impact do well-implemented 'digressions' have on chatbot conversations?}
\choice{A}{Limiting conversational options for users}
\choice{B}{Enhancing natural conversation and user engagement}
\choice{C}{Increasing response times}
\choice{D}{Reducing chatbot functionality}
\answer{B}
\confidence{high}
\evidence{Khi được triển khai tốt, tính năng xử lý 'lạc đề' (digressions) giúp cuộc trò chuyện trở nên tự nhiên, liền mạch hơn và nâng cao mức độ gắn kết của người dùng (Enhancing natural conversation and user engagement).}
\sources{Conversational Design Guidelines: The Power of Seamless Digressions}
\explain{Hướng giải: Đánh giá tác động tích cực của tính năng Digression đối với chất lượng cuộc trò chuyện của chatbot.

Các bước phân tích:
1. Giao tiếp của con người đầy ắp các câu hỏi ngẫu hứng ngoài lề.
2. Nếu chatbot có thể xử lý mượt mà sự chuyển hướng này rồi khéo léo quay lại hoàn thành luồng công việc chính, cuộc trò chuyện sẽ vô cùng tự nhiên và tạo được sự tin cậy, gắn kết cao (user engagement).

Kiểm tra: Đây là chỉ số chất lượng quan trọng trong đánh giá Conversational UX.

Vì sao các phương án khác sai:
- A: Tăng tự do đàm thoại, không phải giới hạn lựa chọn.
- C: Không làm tăng thời gian phản hồi kỹ thuật của bot.
- D: Gia tăng năng lực đối thoại, không làm giảm tính năng.

Lưu ý/Khái niệm cần nhớ: Xử lý lạc đề tốt nâng cao tính tự nhiên và sự gắn kết của người dùng.}
\vocab{User Engagement | Mức độ gắn kết của người dùng | Mức độ tương tác, hài lòng và tin cậy của người dùng đối với ứng dụng}
\end{question}

% ==============================================================================
% Question 189 | Phân loại: 4. Building AI Powered Chatbots Without Programming
% Xuất hiện trong các đề gốc: AIG202c_SP26_RE (Q58)
% ==============================================================================
\begin{question}{189}{mcq}
\source{Q58.jpg}
\prompt{How are 'intents' utilized within AI-powered chatbots?}
\choice{A}{Assessing hardware specifications}
\choice{B}{Identifying the user's intentions or goals in conversations}
\choice{C}{Enhancing graphic design elements}
\choice{D}{Automating backend server configurations}
\answer{B}
\confidence{high}
\evidence{Ý định ('intents') được sử dụng trong chatbot AI để nhận diện mục đích hoặc mục tiêu của người dùng trong cuộc trò chuyện (Identifying the user's intentions or goals in conversations).}
\sources{IBM Watson Assistant Concepts; Jurafsky & Martin, Intent Classification}
\explain{Hướng giải: Xác định vai trò sử dụng của Intent trong hệ thống Chatbot.

Các bước phân tích:
1. Khi người dùng nói bất kỳ câu nào, hệ thống NLU phải ánh xạ phát ngôn đó vào một Ý định (Intent) cụ thể.
2. Ví dụ: 'Tôi bị mất thẻ' -> Intent: `#BaoMatThe`.
3. Nhận diện đúng mục đích người dùng là bước tiên quyết để bot đưa ra câu trả lời chính xác.

Kiểm tra: Đây là khái niệm lặp lại cốt lõi trong toàn bộ chương trình AIG202c.

Vì sao các phương án khác sai:
- A: Đánh giá thông số phần cứng là Hardware Monitoring.
- C: Thiết kế đồ họa thuộc Graphic Design.
- D: Cấu hình máy chủ backend thuộc DevOps/System Admin.

Lưu ý/Khái niệm cần nhớ: Intent = Nhận diện mục đích giao tiếp của người dùng.}
\vocab{Intent Recognition | Nhận diện ý định | Tác vụ phân loại văn bản để xác định mục đích người dùng muốn đạt được}
\end{question}

% ==============================================================================
% Question 190 | Phân loại: 4. Building AI Powered Chatbots Without Programming
% Xuất hiện trong các đề gốc: AIG202c_SP26_RE (Q59)
% ==============================================================================
\begin{question}{190}{mcq}
\source{Q59.jpg}
\prompt{How are "slots" used in natural language understanding for chatbots?}
\choice{A}{Slots are used to store the user's input during a conversation}
\choice{B}{Slots are used to represent specific pieces of information within user input that the chatbot needs to extract}
\choice{C}{Slots are used to determine the overall purpose of the user's input}
\choice{D}{Slots are used to classify the user's emotions during a conversation}
\answer{B}
\confidence{high}
\evidence{Trong xử lý ngôn ngữ tự nhiên cho chatbot, 'slots' được dùng để đại diện cho các mẩu thông tin cụ thể nằm trong câu nói của người dùng mà chatbot cần trích xuất (represent specific pieces of information within user input that the chatbot needs to extract).}
\sources{IBM Watson Assistant: Using Slots to Gather Information; Task-Oriented Dialogue Systems}
\explain{Hướng giải: Hiểu bản chất của 'slots' trong tầng NLU của Chatbot.

Các bước phân tích:
1. 'Slots' là các vị trí biến cần điền giá trị (ví dụ: Slot `Destination`, Slot `DepartureDate`).
2. Khi người dùng cung cấp thông tin, bộ bóc tách thực thể sẽ nhận diện và đưa giá trị vào slot tương ứng.
3. Do đó, slots đại diện trực tiếp cho các thông tin cụ thể cần trích xuất.

Kiểm tra: Kỹ thuật Slot-filling là nền tảng của hệ thống đàm thoại hướng tác vụ (Task-oriented dialogue).

Vì sao các phương án khác sai:
- A: Lưu toàn bộ câu nói là biến `input.text`.
- C: Mục đích tổng thể là `Intent`.
- D: Phân loại cảm xúc là `Sentiment Analysis`.

Lưu ý/Khái niệm cần nhớ: Slots = Các vị trí lưu trữ thông tin cụ thể được trích xuất từ câu nói của người dùng.}
\vocab{Slot | Khe thông tin | Biến lưu trữ dữ liệu cụ thể cần thu thập trong một luồng hội thoại tác vụ}
\end{question}

